%!TEX program = xelatex
\documentclass[11pt,a4paper]{article}
\usepackage[margin=2.2cm]{geometry}
\usepackage{amsmath,amssymb,amsthm,mathtools}
\usepackage{bm}
\usepackage{enumitem}
\usepackage{booktabs}
\usepackage{array}
\usepackage{tabularx}
\usepackage{graphicx}
\newcolumntype{Y}{>{\raggedright\arraybackslash}X}
\usepackage{tikz}
\usetikzlibrary{arrows.meta,calc,positioning}
\usepackage{xcolor}
\newcommand{\mserr}[1]{\textcolor{red}{#1}}
\newcommand{\errpanel}[1]{%
\begin{center}
\fcolorbox{red!75!black}{red!8}{%
\begin{minipage}{0.94\linewidth}
\setlength{\parskip}{0.45em plus 0.1em}
#1
\end{minipage}}
\end{center}
}
\usepackage[round,authoryear]{natbib}
\usepackage{hyperref}
\usepackage{fancyhdr}
\usepackage{xeCJK}

\hypersetup{
    colorlinks=true,
    linkcolor=blue!60!black,
    citecolor=blue!60!black,
    urlcolor=blue!60!black
}

\pagestyle{fancy}
\fancyhf{}
\rhead{\small Qiuzhen QE Probability \& Statistics --- Spring 2026}
\lhead{\small Official Qualifying Exam}
\rfoot{\small Page \thepage}

\DeclareMathOperator*{\argmax}{arg\,max}
\DeclareMathOperator*{\argmin}{arg\,min}
\DeclareMathOperator{\Var}{Var}
\DeclareMathOperator{\Cov}{Cov}
\DeclareMathOperator{\Corr}{Corr}
\DeclareMathOperator{\E}{\mathbb{E}}
\DeclareMathOperator{\Pbb}{\mathbb{P}}

\setlist[itemize]{topsep=3pt,itemsep=2pt,parsep=0pt}
\setlist[enumerate]{topsep=3pt,itemsep=2pt,parsep=0pt}

\newcommand{\examstatement}[1]{%
  \par\addvspace{0.85em}%
  \noindent{\bfseries #1}\par\nobreak\addvspace{0.28em}%
}

% Shared amsthm note/remark/theorem styles + cleveref
% Shared math extras for qe-review.
% Requires already-loaded: amsmath, amssymb.
%
% Classic pitfall: AMS blackboard-bold fonts have no digit glyphs, so
% \mathbb{1} renders as overlapping garbage. Map the digit 1 to dsfont's
% \mathds{1}; leave \mathbb{R}, \mathbb{P}, etc. on AMS.
%
% Usage (path relative to the main .tex being compiled):
%   notes:  % Shared math extras for qe-review.
% Requires already-loaded: amsmath, amssymb.
%
% Classic pitfall: AMS blackboard-bold fonts have no digit glyphs, so
% \mathbb{1} renders as overlapping garbage. Map the digit 1 to dsfont's
% \mathds{1}; leave \mathbb{R}, \mathbb{P}, etc. on AMS.
%
% Usage (path relative to the main .tex being compiled):
%   notes:  % Shared math extras for qe-review.
% Requires already-loaded: amsmath, amssymb.
%
% Classic pitfall: AMS blackboard-bold fonts have no digit glyphs, so
% \mathbb{1} renders as overlapping garbage. Map the digit 1 to dsfont's
% \mathds{1}; leave \mathbb{R}, \mathbb{P}, etc. on AMS.
%
% Usage (path relative to the main .tex being compiled):
%   notes:  \input{../../common/qe-math-extras.tex}
%   exams:  \input{../../../common/qe-math-extras.tex}

\usepackage{dsfont}
\usepackage{pdftexcmds}

\makeatletter
\let\qe@amsmmathbb\mathbb
\renewcommand{\mathbb}[1]{%
  \ifnum\pdf@strcmp{\unexpanded{#1}}{1}=\z@
    \mathds{1}%
  \else
    \qe@amsmmathbb{#1}%
  \fi
}
\makeatother

% Preferred indicator macros (either style is safe after the remap above).
\newcommand{\bbone}{\mathds{1}}
\newcommand{\ind}{\mathbf{1}}

%   exams:  % Shared math extras for qe-review.
% Requires already-loaded: amsmath, amssymb.
%
% Classic pitfall: AMS blackboard-bold fonts have no digit glyphs, so
% \mathbb{1} renders as overlapping garbage. Map the digit 1 to dsfont's
% \mathds{1}; leave \mathbb{R}, \mathbb{P}, etc. on AMS.
%
% Usage (path relative to the main .tex being compiled):
%   notes:  \input{../../common/qe-math-extras.tex}
%   exams:  \input{../../../common/qe-math-extras.tex}

\usepackage{dsfont}
\usepackage{pdftexcmds}

\makeatletter
\let\qe@amsmmathbb\mathbb
\renewcommand{\mathbb}[1]{%
  \ifnum\pdf@strcmp{\unexpanded{#1}}{1}=\z@
    \mathds{1}%
  \else
    \qe@amsmmathbb{#1}%
  \fi
}
\makeatother

% Preferred indicator macros (either style is safe after the remap above).
\newcommand{\bbone}{\mathds{1}}
\newcommand{\ind}{\mathbf{1}}


\usepackage{dsfont}
\usepackage{pdftexcmds}

\makeatletter
\let\qe@amsmmathbb\mathbb
\renewcommand{\mathbb}[1]{%
  \ifnum\pdf@strcmp{\unexpanded{#1}}{1}=\z@
    \mathds{1}%
  \else
    \qe@amsmmathbb{#1}%
  \fi
}
\makeatother

% Preferred indicator macros (either style is safe after the remap above).
\newcommand{\bbone}{\mathds{1}}
\newcommand{\ind}{\mathbf{1}}

%   exams:  % Shared math extras for qe-review.
% Requires already-loaded: amsmath, amssymb.
%
% Classic pitfall: AMS blackboard-bold fonts have no digit glyphs, so
% \mathbb{1} renders as overlapping garbage. Map the digit 1 to dsfont's
% \mathds{1}; leave \mathbb{R}, \mathbb{P}, etc. on AMS.
%
% Usage (path relative to the main .tex being compiled):
%   notes:  % Shared math extras for qe-review.
% Requires already-loaded: amsmath, amssymb.
%
% Classic pitfall: AMS blackboard-bold fonts have no digit glyphs, so
% \mathbb{1} renders as overlapping garbage. Map the digit 1 to dsfont's
% \mathds{1}; leave \mathbb{R}, \mathbb{P}, etc. on AMS.
%
% Usage (path relative to the main .tex being compiled):
%   notes:  \input{../../common/qe-math-extras.tex}
%   exams:  \input{../../../common/qe-math-extras.tex}

\usepackage{dsfont}
\usepackage{pdftexcmds}

\makeatletter
\let\qe@amsmmathbb\mathbb
\renewcommand{\mathbb}[1]{%
  \ifnum\pdf@strcmp{\unexpanded{#1}}{1}=\z@
    \mathds{1}%
  \else
    \qe@amsmmathbb{#1}%
  \fi
}
\makeatother

% Preferred indicator macros (either style is safe after the remap above).
\newcommand{\bbone}{\mathds{1}}
\newcommand{\ind}{\mathbf{1}}

%   exams:  % Shared math extras for qe-review.
% Requires already-loaded: amsmath, amssymb.
%
% Classic pitfall: AMS blackboard-bold fonts have no digit glyphs, so
% \mathbb{1} renders as overlapping garbage. Map the digit 1 to dsfont's
% \mathds{1}; leave \mathbb{R}, \mathbb{P}, etc. on AMS.
%
% Usage (path relative to the main .tex being compiled):
%   notes:  \input{../../common/qe-math-extras.tex}
%   exams:  \input{../../../common/qe-math-extras.tex}

\usepackage{dsfont}
\usepackage{pdftexcmds}

\makeatletter
\let\qe@amsmmathbb\mathbb
\renewcommand{\mathbb}[1]{%
  \ifnum\pdf@strcmp{\unexpanded{#1}}{1}=\z@
    \mathds{1}%
  \else
    \qe@amsmmathbb{#1}%
  \fi
}
\makeatother

% Preferred indicator macros (either style is safe after the remap above).
\newcommand{\bbone}{\mathds{1}}
\newcommand{\ind}{\mathbf{1}}


\usepackage{dsfont}
\usepackage{pdftexcmds}

\makeatletter
\let\qe@amsmmathbb\mathbb
\renewcommand{\mathbb}[1]{%
  \ifnum\pdf@strcmp{\unexpanded{#1}}{1}=\z@
    \mathds{1}%
  \else
    \qe@amsmmathbb{#1}%
  \fi
}
\makeatother

% Preferred indicator macros (either style is safe after the remap above).
\newcommand{\bbone}{\mathds{1}}
\newcommand{\ind}{\mathbf{1}}


\usepackage{dsfont}
\usepackage{pdftexcmds}

\makeatletter
\let\qe@amsmmathbb\mathbb
\renewcommand{\mathbb}[1]{%
  \ifnum\pdf@strcmp{\unexpanded{#1}}{1}=\z@
    \mathds{1}%
  \else
    \qe@amsmmathbb{#1}%
  \fi
}
\makeatother

% Preferred indicator macros (either style is safe after the remap above).
\newcommand{\bbone}{\mathds{1}}
\newcommand{\ind}{\mathbf{1}}

% Shared amsthm environments for qe-review (minimal).
% Requires already-loaded: amsthm.
% Safe to \input after \usepackage{amsmath,amssymb,amsthm,...}.
%
% Shared numbering uses aliascnt so cleveref keys off the environment
% name (Lemma, Proposition, Exercise, ...) rather than the parent
% counter (Theorem, Definition, Remark). Without aliascnt, \cref{lem:...}
% prints ``Theorem 9.13'' instead of ``Lemma 9.13''.

\usepackage{aliascnt}

\theoremstyle{plain}
\newtheorem{theorem}{Theorem}[section]

\newaliascnt{lemma}{theorem}
\newtheorem{lemma}[lemma]{Lemma}
\aliascntresetthe{lemma}

\newaliascnt{proposition}{theorem}
\newtheorem{proposition}[proposition]{Proposition}
\aliascntresetthe{proposition}

\newaliascnt{corollary}{theorem}
\newtheorem{corollary}[corollary]{Corollary}
\aliascntresetthe{corollary}

\newaliascnt{claim}{theorem}
\newtheorem{claim}[claim]{Claim}
\aliascntresetthe{claim}

\newaliascnt{fact}{theorem}
\newtheorem{fact}[fact]{Fact}
\aliascntresetthe{fact}

\theoremstyle{definition}
\newtheorem{definition}{Definition}[section]

\newaliascnt{exercise}{definition}
\newtheorem{exercise}[exercise]{Exercise}
\aliascntresetthe{exercise}

\newaliascnt{assumption}{definition}
\newtheorem{assumption}[assumption]{Assumption}
\aliascntresetthe{assumption}

\newaliascnt{notation}{definition}
\newtheorem{notation}[notation]{Notation}
\aliascntresetthe{notation}

\newaliascnt{construction}{definition}
\newtheorem{construction}[construction]{Construction}
\aliascntresetthe{construction}

% Example: document-wide numbering (not per section / chapter)
\newtheorem{example}{Example}

\theoremstyle{remark}
\newtheorem{remark}{Remark}[section]
\newtheorem*{remark*}{Remark}

\newaliascnt{heuristic}{remark}
\newtheorem{heuristic}[heuristic]{Heuristic}
\aliascntresetthe{heuristic}

\newaliascnt{convention}{remark}
\newtheorem{convention}[convention]{Convention}
\aliascntresetthe{convention}

\newaliascnt{warning}{remark}
\newtheorem{warning}[warning]{Warning}
\aliascntresetthe{warning}

\newaliascnt{proofsketch}{remark}
\newtheorem{proofsketch}[proofsketch]{Proof sketch}
\aliascntresetthe{proofsketch}

\newaliascnt{checkpoint}{remark}
\newtheorem{checkpoint}[checkpoint]{Checkpoint}
\aliascntresetthe{checkpoint}

% Note: document-wide numbering (not per section)
\newtheorem{note}{Note}
\newtheorem*{note*}{Note}

\newaliascnt{examnote}{note}
\newtheorem{examnote}[examnote]{Exam note}
\aliascntresetthe{examnote}

\newaliascnt{study}{note}
\newtheorem{study}[study]{Study note}
\aliascntresetthe{study}

\newaliascnt{summary}{note}
\newtheorem{summary}[summary]{Summary}
\aliascntresetthe{summary}

\newtheorem{examproblem}{Problem}[section]
\newtheorem*{examproblem*}{Problem}

% Bilingual aliases
\newenvironment{dingli}{\begin{theorem}}{\end{theorem}}
\newenvironment{yinli}{\begin{lemma}}{\end{lemma}}
\newenvironment{mingti}{\begin{proposition}}{\end{proposition}}
\newenvironment{tuilun}{\begin{corollary}}{\end{corollary}}
\newenvironment{dingyi}{\begin{definition}}{\end{definition}}
\newenvironment{li}{\begin{example}}{\end{example}}
\newenvironment{zhu}{\begin{remark}}{\end{remark}}
\newenvironment{biji}{\begin{note}}{\end{note}}
\newenvironment{kaoshi}{\begin{examnote}}{\end{examnote}}

\renewcommand{\qedsymbol}{\ensuremath{\square}}

% Shared cleveref setup for qe-review.
% Load AFTER: hyperref, and AFTER all \newtheorem / amsthm environments.
% Shared theorem counters are aliased in qe-amsthm-envs.tex (aliascnt),
% otherwise \cref on a lemma/proposition prints ``Theorem''.
%
% Usage (path relative to the main .tex being compiled):
%   notes:  % Shared cleveref setup for qe-review.
% Load AFTER: hyperref, and AFTER all \newtheorem / amsthm environments.
% Shared theorem counters are aliased in qe-amsthm-envs.tex (aliascnt),
% otherwise \cref on a lemma/proposition prints ``Theorem''.
%
% Usage (path relative to the main .tex being compiled):
%   notes:  % Shared cleveref setup for qe-review.
% Load AFTER: hyperref, and AFTER all \newtheorem / amsthm environments.
% Shared theorem counters are aliased in qe-amsthm-envs.tex (aliascnt),
% otherwise \cref on a lemma/proposition prints ``Theorem''.
%
% Usage (path relative to the main .tex being compiled):
%   notes:  \input{../../common/qe-cleveref.tex}
%   exams:  \input{../../../common/qe-cleveref.tex}
%
% Prefer \cref / \Cref over \ref. Use \Cref at sentence start.
% Unnumbered sectioning (e.g. \paragraph with default secnumdepth)
% produces an empty cleveref type; map that to "Paragraph".

\usepackage[nameinlink,noabbrev]{cleveref}

% Fallback for unnumbered \paragraph / \subsubsection labels
\crefname{}{Paragraph}{Paragraphs}
\Crefname{}{Paragraph}{Paragraphs}

% Numbered theorem-like (plain)
\crefname{theorem}{Theorem}{Theorems}
\Crefname{theorem}{Theorem}{Theorems}
\crefname{lemma}{Lemma}{Lemmas}
\Crefname{lemma}{Lemma}{Lemmas}
\crefname{proposition}{Proposition}{Propositions}
\Crefname{proposition}{Proposition}{Propositions}
\crefname{corollary}{Corollary}{Corollaries}
\Crefname{corollary}{Corollary}{Corollaries}
\crefname{claim}{Claim}{Claims}
\Crefname{claim}{Claim}{Claims}
\crefname{fact}{Fact}{Facts}
\Crefname{fact}{Fact}{Facts}

% Definition-like
\crefname{definition}{Definition}{Definitions}
\Crefname{definition}{Definition}{Definitions}
\crefname{example}{Example}{Examples}
\Crefname{example}{Example}{Examples}
\crefname{exercise}{Exercise}{Exercises}
\Crefname{exercise}{Exercise}{Exercises}
\crefname{assumption}{Assumption}{Assumptions}
\Crefname{assumption}{Assumption}{Assumptions}
\crefname{notation}{Notation}{Notations}
\Crefname{notation}{Notation}{Notations}
\crefname{construction}{Construction}{Constructions}
\Crefname{construction}{Construction}{Constructions}

% Remark-like
\crefname{remark}{Remark}{Remarks}
\Crefname{remark}{Remark}{Remarks}
\crefname{heuristic}{Heuristic}{Heuristics}
\Crefname{heuristic}{Heuristic}{Heuristics}
\crefname{convention}{Convention}{Conventions}
\Crefname{convention}{Convention}{Conventions}
\crefname{warning}{Warning}{Warnings}
\Crefname{warning}{Warning}{Warnings}
\crefname{proofsketch}{Proof sketch}{Proof sketches}
\Crefname{proofsketch}{Proof sketch}{Proof sketches}
\crefname{checkpoint}{Checkpoint}{Checkpoints}
\Crefname{checkpoint}{Checkpoint}{Checkpoints}

% Document-wide notes / exam problems
\crefname{note}{Note}{Notes}
\Crefname{note}{Note}{Notes}
\crefname{examnote}{Exam note}{Exam notes}
\Crefname{examnote}{Exam note}{Exam notes}
\crefname{study}{Study note}{Study notes}
\Crefname{study}{Study note}{Study notes}
\crefname{summary}{Summary}{Summaries}
\Crefname{summary}{Summary}{Summaries}
\crefname{examproblem}{Problem}{Problems}
\Crefname{examproblem}{Problem}{Problems}

% Sectioning / floats (explicit, noabbrev-friendly)
\crefname{section}{Section}{Sections}
\Crefname{section}{Section}{Sections}
\crefname{subsection}{Subsection}{Subsections}
\Crefname{subsection}{Subsection}{Subsections}
\crefname{subsubsection}{Subsubsection}{Subsubsections}
\Crefname{subsubsection}{Subsubsection}{Subsubsections}
\crefname{paragraph}{Paragraph}{Paragraphs}
\Crefname{paragraph}{Paragraph}{Paragraphs}
\crefname{equation}{Equation}{Equations}
\Crefname{equation}{Equation}{Equations}
\crefname{figure}{Figure}{Figures}
\Crefname{figure}{Figure}{Figures}
\crefname{table}{Table}{Tables}
\Crefname{table}{Table}{Tables}
\crefname{enumi}{Item}{Items}
\Crefname{enumi}{Item}{Items}

%   exams:  % Shared cleveref setup for qe-review.
% Load AFTER: hyperref, and AFTER all \newtheorem / amsthm environments.
% Shared theorem counters are aliased in qe-amsthm-envs.tex (aliascnt),
% otherwise \cref on a lemma/proposition prints ``Theorem''.
%
% Usage (path relative to the main .tex being compiled):
%   notes:  \input{../../common/qe-cleveref.tex}
%   exams:  \input{../../../common/qe-cleveref.tex}
%
% Prefer \cref / \Cref over \ref. Use \Cref at sentence start.
% Unnumbered sectioning (e.g. \paragraph with default secnumdepth)
% produces an empty cleveref type; map that to "Paragraph".

\usepackage[nameinlink,noabbrev]{cleveref}

% Fallback for unnumbered \paragraph / \subsubsection labels
\crefname{}{Paragraph}{Paragraphs}
\Crefname{}{Paragraph}{Paragraphs}

% Numbered theorem-like (plain)
\crefname{theorem}{Theorem}{Theorems}
\Crefname{theorem}{Theorem}{Theorems}
\crefname{lemma}{Lemma}{Lemmas}
\Crefname{lemma}{Lemma}{Lemmas}
\crefname{proposition}{Proposition}{Propositions}
\Crefname{proposition}{Proposition}{Propositions}
\crefname{corollary}{Corollary}{Corollaries}
\Crefname{corollary}{Corollary}{Corollaries}
\crefname{claim}{Claim}{Claims}
\Crefname{claim}{Claim}{Claims}
\crefname{fact}{Fact}{Facts}
\Crefname{fact}{Fact}{Facts}

% Definition-like
\crefname{definition}{Definition}{Definitions}
\Crefname{definition}{Definition}{Definitions}
\crefname{example}{Example}{Examples}
\Crefname{example}{Example}{Examples}
\crefname{exercise}{Exercise}{Exercises}
\Crefname{exercise}{Exercise}{Exercises}
\crefname{assumption}{Assumption}{Assumptions}
\Crefname{assumption}{Assumption}{Assumptions}
\crefname{notation}{Notation}{Notations}
\Crefname{notation}{Notation}{Notations}
\crefname{construction}{Construction}{Constructions}
\Crefname{construction}{Construction}{Constructions}

% Remark-like
\crefname{remark}{Remark}{Remarks}
\Crefname{remark}{Remark}{Remarks}
\crefname{heuristic}{Heuristic}{Heuristics}
\Crefname{heuristic}{Heuristic}{Heuristics}
\crefname{convention}{Convention}{Conventions}
\Crefname{convention}{Convention}{Conventions}
\crefname{warning}{Warning}{Warnings}
\Crefname{warning}{Warning}{Warnings}
\crefname{proofsketch}{Proof sketch}{Proof sketches}
\Crefname{proofsketch}{Proof sketch}{Proof sketches}
\crefname{checkpoint}{Checkpoint}{Checkpoints}
\Crefname{checkpoint}{Checkpoint}{Checkpoints}

% Document-wide notes / exam problems
\crefname{note}{Note}{Notes}
\Crefname{note}{Note}{Notes}
\crefname{examnote}{Exam note}{Exam notes}
\Crefname{examnote}{Exam note}{Exam notes}
\crefname{study}{Study note}{Study notes}
\Crefname{study}{Study note}{Study notes}
\crefname{summary}{Summary}{Summaries}
\Crefname{summary}{Summary}{Summaries}
\crefname{examproblem}{Problem}{Problems}
\Crefname{examproblem}{Problem}{Problems}

% Sectioning / floats (explicit, noabbrev-friendly)
\crefname{section}{Section}{Sections}
\Crefname{section}{Section}{Sections}
\crefname{subsection}{Subsection}{Subsections}
\Crefname{subsection}{Subsection}{Subsections}
\crefname{subsubsection}{Subsubsection}{Subsubsections}
\Crefname{subsubsection}{Subsubsection}{Subsubsections}
\crefname{paragraph}{Paragraph}{Paragraphs}
\Crefname{paragraph}{Paragraph}{Paragraphs}
\crefname{equation}{Equation}{Equations}
\Crefname{equation}{Equation}{Equations}
\crefname{figure}{Figure}{Figures}
\Crefname{figure}{Figure}{Figures}
\crefname{table}{Table}{Tables}
\Crefname{table}{Table}{Tables}
\crefname{enumi}{Item}{Items}
\Crefname{enumi}{Item}{Items}

%
% Prefer \cref / \Cref over \ref. Use \Cref at sentence start.
% Unnumbered sectioning (e.g. \paragraph with default secnumdepth)
% produces an empty cleveref type; map that to "Paragraph".

\usepackage[nameinlink,noabbrev]{cleveref}

% Fallback for unnumbered \paragraph / \subsubsection labels
\crefname{}{Paragraph}{Paragraphs}
\Crefname{}{Paragraph}{Paragraphs}

% Numbered theorem-like (plain)
\crefname{theorem}{Theorem}{Theorems}
\Crefname{theorem}{Theorem}{Theorems}
\crefname{lemma}{Lemma}{Lemmas}
\Crefname{lemma}{Lemma}{Lemmas}
\crefname{proposition}{Proposition}{Propositions}
\Crefname{proposition}{Proposition}{Propositions}
\crefname{corollary}{Corollary}{Corollaries}
\Crefname{corollary}{Corollary}{Corollaries}
\crefname{claim}{Claim}{Claims}
\Crefname{claim}{Claim}{Claims}
\crefname{fact}{Fact}{Facts}
\Crefname{fact}{Fact}{Facts}

% Definition-like
\crefname{definition}{Definition}{Definitions}
\Crefname{definition}{Definition}{Definitions}
\crefname{example}{Example}{Examples}
\Crefname{example}{Example}{Examples}
\crefname{exercise}{Exercise}{Exercises}
\Crefname{exercise}{Exercise}{Exercises}
\crefname{assumption}{Assumption}{Assumptions}
\Crefname{assumption}{Assumption}{Assumptions}
\crefname{notation}{Notation}{Notations}
\Crefname{notation}{Notation}{Notations}
\crefname{construction}{Construction}{Constructions}
\Crefname{construction}{Construction}{Constructions}

% Remark-like
\crefname{remark}{Remark}{Remarks}
\Crefname{remark}{Remark}{Remarks}
\crefname{heuristic}{Heuristic}{Heuristics}
\Crefname{heuristic}{Heuristic}{Heuristics}
\crefname{convention}{Convention}{Conventions}
\Crefname{convention}{Convention}{Conventions}
\crefname{warning}{Warning}{Warnings}
\Crefname{warning}{Warning}{Warnings}
\crefname{proofsketch}{Proof sketch}{Proof sketches}
\Crefname{proofsketch}{Proof sketch}{Proof sketches}
\crefname{checkpoint}{Checkpoint}{Checkpoints}
\Crefname{checkpoint}{Checkpoint}{Checkpoints}

% Document-wide notes / exam problems
\crefname{note}{Note}{Notes}
\Crefname{note}{Note}{Notes}
\crefname{examnote}{Exam note}{Exam notes}
\Crefname{examnote}{Exam note}{Exam notes}
\crefname{study}{Study note}{Study notes}
\Crefname{study}{Study note}{Study notes}
\crefname{summary}{Summary}{Summaries}
\Crefname{summary}{Summary}{Summaries}
\crefname{examproblem}{Problem}{Problems}
\Crefname{examproblem}{Problem}{Problems}

% Sectioning / floats (explicit, noabbrev-friendly)
\crefname{section}{Section}{Sections}
\Crefname{section}{Section}{Sections}
\crefname{subsection}{Subsection}{Subsections}
\Crefname{subsection}{Subsection}{Subsections}
\crefname{subsubsection}{Subsubsection}{Subsubsections}
\Crefname{subsubsection}{Subsubsection}{Subsubsections}
\crefname{paragraph}{Paragraph}{Paragraphs}
\Crefname{paragraph}{Paragraph}{Paragraphs}
\crefname{equation}{Equation}{Equations}
\Crefname{equation}{Equation}{Equations}
\crefname{figure}{Figure}{Figures}
\Crefname{figure}{Figure}{Figures}
\crefname{table}{Table}{Tables}
\Crefname{table}{Table}{Tables}
\crefname{enumi}{Item}{Items}
\Crefname{enumi}{Item}{Items}

%   exams:  % Shared cleveref setup for qe-review.
% Load AFTER: hyperref, and AFTER all \newtheorem / amsthm environments.
% Shared theorem counters are aliased in qe-amsthm-envs.tex (aliascnt),
% otherwise \cref on a lemma/proposition prints ``Theorem''.
%
% Usage (path relative to the main .tex being compiled):
%   notes:  % Shared cleveref setup for qe-review.
% Load AFTER: hyperref, and AFTER all \newtheorem / amsthm environments.
% Shared theorem counters are aliased in qe-amsthm-envs.tex (aliascnt),
% otherwise \cref on a lemma/proposition prints ``Theorem''.
%
% Usage (path relative to the main .tex being compiled):
%   notes:  \input{../../common/qe-cleveref.tex}
%   exams:  \input{../../../common/qe-cleveref.tex}
%
% Prefer \cref / \Cref over \ref. Use \Cref at sentence start.
% Unnumbered sectioning (e.g. \paragraph with default secnumdepth)
% produces an empty cleveref type; map that to "Paragraph".

\usepackage[nameinlink,noabbrev]{cleveref}

% Fallback for unnumbered \paragraph / \subsubsection labels
\crefname{}{Paragraph}{Paragraphs}
\Crefname{}{Paragraph}{Paragraphs}

% Numbered theorem-like (plain)
\crefname{theorem}{Theorem}{Theorems}
\Crefname{theorem}{Theorem}{Theorems}
\crefname{lemma}{Lemma}{Lemmas}
\Crefname{lemma}{Lemma}{Lemmas}
\crefname{proposition}{Proposition}{Propositions}
\Crefname{proposition}{Proposition}{Propositions}
\crefname{corollary}{Corollary}{Corollaries}
\Crefname{corollary}{Corollary}{Corollaries}
\crefname{claim}{Claim}{Claims}
\Crefname{claim}{Claim}{Claims}
\crefname{fact}{Fact}{Facts}
\Crefname{fact}{Fact}{Facts}

% Definition-like
\crefname{definition}{Definition}{Definitions}
\Crefname{definition}{Definition}{Definitions}
\crefname{example}{Example}{Examples}
\Crefname{example}{Example}{Examples}
\crefname{exercise}{Exercise}{Exercises}
\Crefname{exercise}{Exercise}{Exercises}
\crefname{assumption}{Assumption}{Assumptions}
\Crefname{assumption}{Assumption}{Assumptions}
\crefname{notation}{Notation}{Notations}
\Crefname{notation}{Notation}{Notations}
\crefname{construction}{Construction}{Constructions}
\Crefname{construction}{Construction}{Constructions}

% Remark-like
\crefname{remark}{Remark}{Remarks}
\Crefname{remark}{Remark}{Remarks}
\crefname{heuristic}{Heuristic}{Heuristics}
\Crefname{heuristic}{Heuristic}{Heuristics}
\crefname{convention}{Convention}{Conventions}
\Crefname{convention}{Convention}{Conventions}
\crefname{warning}{Warning}{Warnings}
\Crefname{warning}{Warning}{Warnings}
\crefname{proofsketch}{Proof sketch}{Proof sketches}
\Crefname{proofsketch}{Proof sketch}{Proof sketches}
\crefname{checkpoint}{Checkpoint}{Checkpoints}
\Crefname{checkpoint}{Checkpoint}{Checkpoints}

% Document-wide notes / exam problems
\crefname{note}{Note}{Notes}
\Crefname{note}{Note}{Notes}
\crefname{examnote}{Exam note}{Exam notes}
\Crefname{examnote}{Exam note}{Exam notes}
\crefname{study}{Study note}{Study notes}
\Crefname{study}{Study note}{Study notes}
\crefname{summary}{Summary}{Summaries}
\Crefname{summary}{Summary}{Summaries}
\crefname{examproblem}{Problem}{Problems}
\Crefname{examproblem}{Problem}{Problems}

% Sectioning / floats (explicit, noabbrev-friendly)
\crefname{section}{Section}{Sections}
\Crefname{section}{Section}{Sections}
\crefname{subsection}{Subsection}{Subsections}
\Crefname{subsection}{Subsection}{Subsections}
\crefname{subsubsection}{Subsubsection}{Subsubsections}
\Crefname{subsubsection}{Subsubsection}{Subsubsections}
\crefname{paragraph}{Paragraph}{Paragraphs}
\Crefname{paragraph}{Paragraph}{Paragraphs}
\crefname{equation}{Equation}{Equations}
\Crefname{equation}{Equation}{Equations}
\crefname{figure}{Figure}{Figures}
\Crefname{figure}{Figure}{Figures}
\crefname{table}{Table}{Tables}
\Crefname{table}{Table}{Tables}
\crefname{enumi}{Item}{Items}
\Crefname{enumi}{Item}{Items}

%   exams:  % Shared cleveref setup for qe-review.
% Load AFTER: hyperref, and AFTER all \newtheorem / amsthm environments.
% Shared theorem counters are aliased in qe-amsthm-envs.tex (aliascnt),
% otherwise \cref on a lemma/proposition prints ``Theorem''.
%
% Usage (path relative to the main .tex being compiled):
%   notes:  \input{../../common/qe-cleveref.tex}
%   exams:  \input{../../../common/qe-cleveref.tex}
%
% Prefer \cref / \Cref over \ref. Use \Cref at sentence start.
% Unnumbered sectioning (e.g. \paragraph with default secnumdepth)
% produces an empty cleveref type; map that to "Paragraph".

\usepackage[nameinlink,noabbrev]{cleveref}

% Fallback for unnumbered \paragraph / \subsubsection labels
\crefname{}{Paragraph}{Paragraphs}
\Crefname{}{Paragraph}{Paragraphs}

% Numbered theorem-like (plain)
\crefname{theorem}{Theorem}{Theorems}
\Crefname{theorem}{Theorem}{Theorems}
\crefname{lemma}{Lemma}{Lemmas}
\Crefname{lemma}{Lemma}{Lemmas}
\crefname{proposition}{Proposition}{Propositions}
\Crefname{proposition}{Proposition}{Propositions}
\crefname{corollary}{Corollary}{Corollaries}
\Crefname{corollary}{Corollary}{Corollaries}
\crefname{claim}{Claim}{Claims}
\Crefname{claim}{Claim}{Claims}
\crefname{fact}{Fact}{Facts}
\Crefname{fact}{Fact}{Facts}

% Definition-like
\crefname{definition}{Definition}{Definitions}
\Crefname{definition}{Definition}{Definitions}
\crefname{example}{Example}{Examples}
\Crefname{example}{Example}{Examples}
\crefname{exercise}{Exercise}{Exercises}
\Crefname{exercise}{Exercise}{Exercises}
\crefname{assumption}{Assumption}{Assumptions}
\Crefname{assumption}{Assumption}{Assumptions}
\crefname{notation}{Notation}{Notations}
\Crefname{notation}{Notation}{Notations}
\crefname{construction}{Construction}{Constructions}
\Crefname{construction}{Construction}{Constructions}

% Remark-like
\crefname{remark}{Remark}{Remarks}
\Crefname{remark}{Remark}{Remarks}
\crefname{heuristic}{Heuristic}{Heuristics}
\Crefname{heuristic}{Heuristic}{Heuristics}
\crefname{convention}{Convention}{Conventions}
\Crefname{convention}{Convention}{Conventions}
\crefname{warning}{Warning}{Warnings}
\Crefname{warning}{Warning}{Warnings}
\crefname{proofsketch}{Proof sketch}{Proof sketches}
\Crefname{proofsketch}{Proof sketch}{Proof sketches}
\crefname{checkpoint}{Checkpoint}{Checkpoints}
\Crefname{checkpoint}{Checkpoint}{Checkpoints}

% Document-wide notes / exam problems
\crefname{note}{Note}{Notes}
\Crefname{note}{Note}{Notes}
\crefname{examnote}{Exam note}{Exam notes}
\Crefname{examnote}{Exam note}{Exam notes}
\crefname{study}{Study note}{Study notes}
\Crefname{study}{Study note}{Study notes}
\crefname{summary}{Summary}{Summaries}
\Crefname{summary}{Summary}{Summaries}
\crefname{examproblem}{Problem}{Problems}
\Crefname{examproblem}{Problem}{Problems}

% Sectioning / floats (explicit, noabbrev-friendly)
\crefname{section}{Section}{Sections}
\Crefname{section}{Section}{Sections}
\crefname{subsection}{Subsection}{Subsections}
\Crefname{subsection}{Subsection}{Subsections}
\crefname{subsubsection}{Subsubsection}{Subsubsections}
\Crefname{subsubsection}{Subsubsection}{Subsubsections}
\crefname{paragraph}{Paragraph}{Paragraphs}
\Crefname{paragraph}{Paragraph}{Paragraphs}
\crefname{equation}{Equation}{Equations}
\Crefname{equation}{Equation}{Equations}
\crefname{figure}{Figure}{Figures}
\Crefname{figure}{Figure}{Figures}
\crefname{table}{Table}{Tables}
\Crefname{table}{Table}{Tables}
\crefname{enumi}{Item}{Items}
\Crefname{enumi}{Item}{Items}

%
% Prefer \cref / \Cref over \ref. Use \Cref at sentence start.
% Unnumbered sectioning (e.g. \paragraph with default secnumdepth)
% produces an empty cleveref type; map that to "Paragraph".

\usepackage[nameinlink,noabbrev]{cleveref}

% Fallback for unnumbered \paragraph / \subsubsection labels
\crefname{}{Paragraph}{Paragraphs}
\Crefname{}{Paragraph}{Paragraphs}

% Numbered theorem-like (plain)
\crefname{theorem}{Theorem}{Theorems}
\Crefname{theorem}{Theorem}{Theorems}
\crefname{lemma}{Lemma}{Lemmas}
\Crefname{lemma}{Lemma}{Lemmas}
\crefname{proposition}{Proposition}{Propositions}
\Crefname{proposition}{Proposition}{Propositions}
\crefname{corollary}{Corollary}{Corollaries}
\Crefname{corollary}{Corollary}{Corollaries}
\crefname{claim}{Claim}{Claims}
\Crefname{claim}{Claim}{Claims}
\crefname{fact}{Fact}{Facts}
\Crefname{fact}{Fact}{Facts}

% Definition-like
\crefname{definition}{Definition}{Definitions}
\Crefname{definition}{Definition}{Definitions}
\crefname{example}{Example}{Examples}
\Crefname{example}{Example}{Examples}
\crefname{exercise}{Exercise}{Exercises}
\Crefname{exercise}{Exercise}{Exercises}
\crefname{assumption}{Assumption}{Assumptions}
\Crefname{assumption}{Assumption}{Assumptions}
\crefname{notation}{Notation}{Notations}
\Crefname{notation}{Notation}{Notations}
\crefname{construction}{Construction}{Constructions}
\Crefname{construction}{Construction}{Constructions}

% Remark-like
\crefname{remark}{Remark}{Remarks}
\Crefname{remark}{Remark}{Remarks}
\crefname{heuristic}{Heuristic}{Heuristics}
\Crefname{heuristic}{Heuristic}{Heuristics}
\crefname{convention}{Convention}{Conventions}
\Crefname{convention}{Convention}{Conventions}
\crefname{warning}{Warning}{Warnings}
\Crefname{warning}{Warning}{Warnings}
\crefname{proofsketch}{Proof sketch}{Proof sketches}
\Crefname{proofsketch}{Proof sketch}{Proof sketches}
\crefname{checkpoint}{Checkpoint}{Checkpoints}
\Crefname{checkpoint}{Checkpoint}{Checkpoints}

% Document-wide notes / exam problems
\crefname{note}{Note}{Notes}
\Crefname{note}{Note}{Notes}
\crefname{examnote}{Exam note}{Exam notes}
\Crefname{examnote}{Exam note}{Exam notes}
\crefname{study}{Study note}{Study notes}
\Crefname{study}{Study note}{Study notes}
\crefname{summary}{Summary}{Summaries}
\Crefname{summary}{Summary}{Summaries}
\crefname{examproblem}{Problem}{Problems}
\Crefname{examproblem}{Problem}{Problems}

% Sectioning / floats (explicit, noabbrev-friendly)
\crefname{section}{Section}{Sections}
\Crefname{section}{Section}{Sections}
\crefname{subsection}{Subsection}{Subsections}
\Crefname{subsection}{Subsection}{Subsections}
\crefname{subsubsection}{Subsubsection}{Subsubsections}
\Crefname{subsubsection}{Subsubsection}{Subsubsections}
\crefname{paragraph}{Paragraph}{Paragraphs}
\Crefname{paragraph}{Paragraph}{Paragraphs}
\crefname{equation}{Equation}{Equations}
\Crefname{equation}{Equation}{Equations}
\crefname{figure}{Figure}{Figures}
\Crefname{figure}{Figure}{Figures}
\crefname{table}{Table}{Tables}
\Crefname{table}{Table}{Tables}
\crefname{enumi}{Item}{Items}
\Crefname{enumi}{Item}{Items}

%
% Prefer \cref / \Cref over \ref. Use \Cref at sentence start.
% Unnumbered sectioning (e.g. \paragraph with default secnumdepth)
% produces an empty cleveref type; map that to "Paragraph".

\usepackage[nameinlink,noabbrev]{cleveref}

% Fallback for unnumbered \paragraph / \subsubsection labels
\crefname{}{Paragraph}{Paragraphs}
\Crefname{}{Paragraph}{Paragraphs}

% Numbered theorem-like (plain)
\crefname{theorem}{Theorem}{Theorems}
\Crefname{theorem}{Theorem}{Theorems}
\crefname{lemma}{Lemma}{Lemmas}
\Crefname{lemma}{Lemma}{Lemmas}
\crefname{proposition}{Proposition}{Propositions}
\Crefname{proposition}{Proposition}{Propositions}
\crefname{corollary}{Corollary}{Corollaries}
\Crefname{corollary}{Corollary}{Corollaries}
\crefname{claim}{Claim}{Claims}
\Crefname{claim}{Claim}{Claims}
\crefname{fact}{Fact}{Facts}
\Crefname{fact}{Fact}{Facts}

% Definition-like
\crefname{definition}{Definition}{Definitions}
\Crefname{definition}{Definition}{Definitions}
\crefname{example}{Example}{Examples}
\Crefname{example}{Example}{Examples}
\crefname{exercise}{Exercise}{Exercises}
\Crefname{exercise}{Exercise}{Exercises}
\crefname{assumption}{Assumption}{Assumptions}
\Crefname{assumption}{Assumption}{Assumptions}
\crefname{notation}{Notation}{Notations}
\Crefname{notation}{Notation}{Notations}
\crefname{construction}{Construction}{Constructions}
\Crefname{construction}{Construction}{Constructions}

% Remark-like
\crefname{remark}{Remark}{Remarks}
\Crefname{remark}{Remark}{Remarks}
\crefname{heuristic}{Heuristic}{Heuristics}
\Crefname{heuristic}{Heuristic}{Heuristics}
\crefname{convention}{Convention}{Conventions}
\Crefname{convention}{Convention}{Conventions}
\crefname{warning}{Warning}{Warnings}
\Crefname{warning}{Warning}{Warnings}
\crefname{proofsketch}{Proof sketch}{Proof sketches}
\Crefname{proofsketch}{Proof sketch}{Proof sketches}
\crefname{checkpoint}{Checkpoint}{Checkpoints}
\Crefname{checkpoint}{Checkpoint}{Checkpoints}

% Document-wide notes / exam problems
\crefname{note}{Note}{Notes}
\Crefname{note}{Note}{Notes}
\crefname{examnote}{Exam note}{Exam notes}
\Crefname{examnote}{Exam note}{Exam notes}
\crefname{study}{Study note}{Study notes}
\Crefname{study}{Study note}{Study notes}
\crefname{summary}{Summary}{Summaries}
\Crefname{summary}{Summary}{Summaries}
\crefname{examproblem}{Problem}{Problems}
\Crefname{examproblem}{Problem}{Problems}

% Sectioning / floats (explicit, noabbrev-friendly)
\crefname{section}{Section}{Sections}
\Crefname{section}{Section}{Sections}
\crefname{subsection}{Subsection}{Subsections}
\Crefname{subsection}{Subsection}{Subsections}
\crefname{subsubsection}{Subsubsection}{Subsubsections}
\Crefname{subsubsection}{Subsubsection}{Subsubsections}
\crefname{paragraph}{Paragraph}{Paragraphs}
\Crefname{paragraph}{Paragraph}{Paragraphs}
\crefname{equation}{Equation}{Equations}
\Crefname{equation}{Equation}{Equations}
\crefname{figure}{Figure}{Figures}
\Crefname{figure}{Figure}{Figures}
\crefname{table}{Table}{Tables}
\Crefname{table}{Table}{Tables}
\crefname{enumi}{Item}{Items}
\Crefname{enumi}{Item}{Items}

% PDF sidebar outline + printed TOC for QE exam transcripts.
% Load in the preamble AFTER hyperref and AFTER \newcommand{\examstatement}.
\hypersetup{
  unicode=true,
  bookmarks=true,
  bookmarksopen=true,
  bookmarksopenlevel=3,
  bookmarksnumbered=false,
  bookmarksdepth=3
}
\IfFileExists{bookmark.sty}{%
  \usepackage{bookmark}%
  \bookmarksetup{open,openlevel=3,numbered=false,depth=3}%
}{}
% Unnumbered in print, but still written to the TOC / PDF outline
% down to \subsubsection. Starred headings create neither.
\setcounter{secnumdepth}{-1}
\setcounter{tocdepth}{3}
\makeatletter
\@ifundefined{examstatement}{}{%
  \let\qe@origexamstatement\examstatement
  \renewcommand{\examstatement}[1]{%
    \par\phantomsection
    \addcontentsline{toc}{subsection}{#1}%
    \qe@origexamstatement{#1}%
  }%
}
\makeatother


\title{\textbf{QUALIFYING EXAM: PROBABILITY \& STATISTICS}\\[0.5em]\Large SPRING 2026}
\date{}
\author{}

\begin{document}

\maketitle
\vspace{-3em}

\tableofcontents

\section{Instructions}
\begin{itemize}
    \item There are 11 problems in this exam (4 pages). You need to choose 8 of them to solve. If you select more than 8, only the first 8 that you have worked on will be graded. Note that 4 of the problems are worth 15 points each and the rest 10 points each.
    \item You must follow all the rules of exam taking. Misconducts will be subject to proper disciplinary actions by the Center.
    \item You must provide all necessary details for full credits. A final answer with no or little explanation/derivation, even if correct, receives a minimal credit.
    \item \( \mathbb{R} \) denotes the set of real numbers and \( \mathbb{N} = \{1, 2, 3, \dots\} \) denotes the set of positive integers. \( \xrightarrow{(d)} \) and \( \stackrel{(d)}{=} \) mean ``converges in distribution'' and ``equal in distribution'', respectively.
    \item The pmf of a \( \mathrm{Poisson}(\lambda) \) random variable \( X \) is given by \( P(X = x \mid \lambda) = \frac{e^{-\lambda} \lambda^x}{x!} \); \( x = 0, 1, \dots \); \( \lambda > 0 \).
\end{itemize}

\section{Statement of the problems}

\examstatement{Problem 1. [10 pts.]}
Let \( X = (X_1, \dots, X_n)^T \) be an \( \mathbb{R}^n \)-valued Gaussian random vector under a probability measure \( \mathbb{P} \), with mean vector \( \mu \in \mathbb{R}^n \) and covariance matrix \( V \in \mathbb{R}^{n \times n} \). Let \( \alpha \in \mathbb{R}^n \) and define a new probability measure \( \mathbb{Q} \) by

\[
\frac{\mathrm{d}\mathbb{Q}}{\mathrm{d}\mathbb{P}} = \frac{\exp(\alpha^T X)}{\mathbb{E}_{\mathbb{P}}[\exp(\alpha^T X)]}.
\]

Show that \( X \) is still Gaussian under \( \mathbb{Q} \) and identify the mean and covariance matrix of \( X \) under \( \mathbb{Q} \).

\examstatement{Problem 2. [15 pts.]}
Let \( \{X_n\}_{n \in \mathbb{N}} \) be a sequence of independent random variables with distribution:

\[
\begin{cases}
\mathbb{P}[X_n = n^2] = \mathbb{P}[X_n = -n^2] = \dfrac{1}{6n^2}, \\[0.6em]
\mathbb{P}[X_n = n] = \mathbb{P}[X_n = -n] = \dfrac{1}{6}, \\[0.6em]
\mathbb{P}[X_n = 0] = \dfrac{2}{3} - \dfrac{1}{3n^2}.
\end{cases}
\]

Let \( S_n = \sum_{j=1}^n X_j \). Find a sequence of positive real numbers \( \{b_n\}_{n \in \mathbb{N}} \) such that the law of \( S_n / b_n \) converges to standard normal distribution \( \mathcal{N}(0, 1) \).

\examstatement{Problem 3. [10 pts.]}
We have 100 noodles in a bowl. Each noodle has two free ends, so initially there are 200 free ends in total. We are blindfolded and repeatedly perform the following operation: choose uniformly at random two free ends among all available free ends and connect them together. We continue until there are no free ends left. At the end of this procedure, the noodles form a random collection of disjoint loops. Compute the expectation of the number of loops.

\examstatement{Problem 4. [10 pts.]}
Let \( \{B_t\}_{t \ge 0} \) be a one-dimensional standard Brownian motion started from the origin. Find and prove the limit (in distribution) of

\[
\left[ \int_0^{T^2} \exp(B_t)\, dt \right]^{1/T} \qquad \text{as } T \to \infty.
\]

Write your answer as a function of the standard normal distribution.

\examstatement{Problem 5. [10 pts.]}
A Galton--Watson branching process is a stochastic process \( \{Z_n\}_{n \in \mathbb{N} \cup \{0\}} \) which evolves according to the recurrence formula

\[
Z_0 := 1, \qquad Z_{n+1} := \sum_{j=1}^{Z_n} X(j, n),
\]

where \( \{X(j, n) : j \in \mathbb{N},\, n \in \mathbb{N} \cup \{0\}\} \) is a set of i.i.d.\ \( \mathbb{N} \cup \{0\} \)-valued random variables with mean \( m := \mathbb{E}[X(1, 1)] < \infty \). Let \( p_k := \mathbb{P}[X(1, 1) = k] \), \( k \in \mathbb{N} \cup \{0\} \), and we assume \( p_0 > 0 \), \( p_0 + p_1 < 1 \). Let \( f(s) := \mathbb{E}[s^{X(1,1)}] \) for \( s \in [0, 1] \) (with the convention \( 0^0 = 1 \)). Let

\[
\mathcal{F}_n := \sigma(Z_0, \dots, Z_n), \qquad n \in \mathbb{N} \cup \{0\}.
\]

\begin{enumerate}[label=(\alph*)]
    \item Compute

    \[
    \mathbb{E}[Z_{n+1} \mid \mathcal{F}_n],\quad
    \mathbb{E}[Z_n],\quad
    \mathbb{E}[s^{Z_{n+1}} \mid \mathcal{F}_n],\quad
    \mathbb{E}[s^{Z_n}], \qquad n \in \mathbb{N} \cup \{0\},\; s \in [0, 1].
    \]
    \item Determine the number of solutions of \( f(s) = s \) on \( [0, 1] \).
    \item Let \( \tau := \inf\{n \in \mathbb{N} : Z_n = 0\} \) (with \( \inf \emptyset := \infty \)). Call the event \( \{\tau < \infty\} \) extinction. Compute \( \mathbb{P}[\tau < \infty] \). Write your answer as a function of \( q \), where \( q \) is the smallest solution of \( f(s) = s \) on \( [0, 1] \).
\end{enumerate}

\examstatement{Problem 6. [15 pts.]}
In the above definition of Galton--Watson branching process, \( Z_n \) can be thought of as the number of descendants in the \( n \)-th generation, and \( X(j, n) \) can be thought of as the number of children of the \( j \)-th of these descendants. By tracking genealogical relationships, we obtain a tree \( T \) rooted at the single individual in generation 0 with a vertex for each individual in the progeny and an edge for each parent-child relationship. \( T \) is called a Galton--Watson tree.

A property of rooted trees is said to be inherited if this property holds for all finite trees, and whenever it holds for a tree, it also holds for all subtrees rooted at the children of the root.
\begin{enumerate}[label=(\alph*)]
    \item Prove that for a Galton--Watson tree \( T \), conditioned on non-extinction, an inherited property \( A \) has a probability of either 0 or 1. (Hint: One may use \( s \le f(s) \) implies \( s \in [0, q] \cup \{1\} \).)
    \item Let \( T \) be the Galton--Watson tree for an offspring distribution with mean \( m > 1 \). Perform percolation on \( T \) with density \( p \) (i.e., each edge of \( T \) is open with probability \( p \) and closed with probability \( 1 - p \), independently of all other edges). Let \( \mathcal{C}_0 \) be the maximal subtree of the root formed by open edges in \( T \). Define the critical probability

    \[
    p_c(T) := \sup\{p \in [0, 1] : \mathbb{P}_p[|\mathcal{C}_0| = \infty \mid T] = 0\},
    \]

    where \( |\mathcal{C}_0| \) is the number of vertices in \( \mathcal{C}_0 \). Conditioned on non-extinction of \( T \), compute \( p_c(T) \).
\end{enumerate}

\examstatement{Problem 7. [10 pts.]}
Let \( X_1, X_2, \dots, X_n \) be i.i.d.\ observations with pdf

\[
f(x; \theta) = \frac{1}{2\theta\sqrt{x}} e^{-\sqrt{x}/\theta} I_{\{x>0\}},
\]

where \( \theta > 0 \) is the unknown parameter.
\begin{enumerate}[label=(\alph*)]
    \item Show that the family \( f(x_1, \dots, x_n; \theta) \) (\( \theta > 0 \)) has a monotone likelihood ratio in a certain statistic \( T \).
    \item Find a size \( \alpha \) (\( 0 < \alpha < 1 \)) uniformly most powerful test for \( H_0 : \theta \ge 1 \) versus \( H_1 : \theta < 1/2 \). Make your rejection region expressed as explicitly as possible.
\end{enumerate}

\examstatement{Problem 8. [15 pts.]}
Let \( X_1, X_2, \dots, X_n \) be i.i.d.\ Bernoulli random variables with unknown success probability \( 0 < \theta < 1 \). Given \( 0 < p < 1 \), let \( \hat{\xi}_p \) be the sample quantile based on the \( n \) observations, i.e., \( \hat{\xi}_p = \inf\{x : F_n(x) \ge p\} \), where \( F_n \) denotes the empirical CDF based on the data.
\begin{enumerate}[label=(\alph*)]
    \item Does \( \hat{\xi}_p \) converge in probability to the population quantile \( \xi_p \) for each \( 0 < p < 1 \) as \( n \to \infty \)? Prove your answer.
    \item Is it possible that for a non-degenerate discrete distribution \( F \) we have \( \hat{\xi}_p \) converges in probability to \( \xi_p \) for all \( 0 < p < 1 \)? Prove your answer.
\end{enumerate}

\examstatement{Problem 9. [10 pts.]}
Suppose a single observation \( Y \mid \mu \sim N(\mu, 1) \), and the prior on \( \mu \) is \( \pi(\mu) = 0.5 \cdot I(\mu = -1) + 0.5 \cdot I(\mu = 1) \) where \( I(\cdot) \) is the indicator function. That is, under this prior, \( \mu \) is supported on \( \{-1, 1\} \) and takes either value with a probability 0.5. For estimating \( \mu \), let the action space be \( \{-1, 1\} \).
\begin{enumerate}[label=(\alph*)]
    \item Write down the posterior \( \pi(\mu \mid Y = y) \) and derive the Bayes rule \( \delta \) under the absolute error loss \( L(\mu, a) = |\mu - a| \).
    \item Is \( \delta \) above a minimax rule under the absolute error loss? Please justify your answer.
\end{enumerate}

\examstatement{Problem 10. [10 pts.]}
Let \( X_1, X_2, \dots, X_n \) be i.i.d.\ random variables following \( \mathrm{Poisson}(\sqrt{\lambda}) \) distribution where \( \lambda > 0 \).
\begin{enumerate}[label=(\alph*)]
    \item Find the BUE (UMVUE) \( \tilde{\lambda} \) for \( \lambda \).
    \item Let \( \hat{\lambda} \) denote the MLE for \( \lambda \) based on \( X_1, X_2, \dots, X_n \), show that \( \sqrt{n}(\tilde{\lambda} - \hat{\lambda}) \to 0 \) in probability as \( n \to \infty \).
\end{enumerate}

\examstatement{Problem 11. [15 pts.]}
Let \( Y = \sigma\bigl(\rho |U| + \sqrt{1 - \rho^2}\, V\bigr) \) where \( U, V \) are independent \( N(0, 1) \) random variables and \( \rho \in (-1, 1) \), \( \sigma \in (0, \infty) \) are unknown parameters.
\begin{enumerate}[label=(\alph*)]
    \item Show that the pdf of \( Y \) is given by

    \[
    f(y \mid \rho, \sigma)
    =
    \frac{a}{\sqrt{2\pi\sigma^2}}
    \exp\left(-\frac{y^2}{2\sigma^2}\right)
    \Phi\left( \frac{\rho y}{b\sigma\sqrt{1-\rho^2}} \right),
    \qquad -\infty < y < \infty
    \]

    for some positive constants \( a, b \). Write down the values of \( a \) and \( b \). In the above expression, \( \Phi(z) \) denotes the CDF for a standard normal random variable.
    \item Consider \( n \) observations \( Y_1, \dots, Y_n \), modeled as \( Y_i \stackrel{\mathrm{iid}}{\sim} f(y_i \mid \rho, \sigma) \), \( \rho \in (-1, 1) \), \( \sigma \in (0, \infty) \). Does a maximum likelihood estimate of \( (\rho, \sigma) \in (-1, 1) \times (0, \infty) \) always exist? Please justify your answer.
    \item For the same model, suppose we are interested in testing \( H_0 : \rho = 0 \) versus \( H_1 : \rho \neq 0 \). What is the value of \( c \) such that the test that rejects \( H_0 \) if

    \[
    \frac{\sqrt{n}\, |\overline{Y}|}{s_Y} > c
    \]

    has size \( \alpha = 0.05 \)? Identify \( c \) as a specific quantile of a named distribution. In the above expression, \( \overline{Y} \) is the sample mean and \( s_Y \) is the sample standard deviation.
\end{enumerate}

\noindent\rule{\linewidth}{0.4pt}

\section{Statement and solutions}

\subsection{Problem 1. [10 pts.]}
Let \( X = (X_1, \dots, X_n)^T \) be an \( \mathbb{R}^n \)-valued Gaussian random vector under a probability measure \( \mathbb{P} \), with mean vector \( \mu \in \mathbb{R}^n \) and covariance matrix \( V \in \mathbb{R}^{n \times n} \). Let \( \alpha \in \mathbb{R}^n \) and define a new probability measure \( \mathbb{Q} \) by

\[
\frac{\mathrm{d}\mathbb{Q}}{\mathrm{d}\mathbb{P}} = \frac{\exp(\alpha^T X)}{\mathbb{E}_{\mathbb{P}}[\exp(\alpha^T X)]}.
\]

Show that \( X \) is still Gaussian under \( \mathbb{Q} \) and identify the mean and covariance matrix of \( X \) under \( \mathbb{Q} \).

\setcounter{section}{1}

\subsubsection{Definitions used in Problem 1}

\begin{definition}[Signed measure]
\label{def:signed-measure}
Let \((\Omega, \mathcal{F})\) be a measurable space. A \emph{signed measure} on \((\Omega, \mathcal{F})\) is a map \(\nu \colon \mathcal{F} \to [-\infty, +\infty]\) such that \(\nu(\emptyset) = 0\), \(\nu\) attains at most one of the two values \(\pm\infty\), and \(\nu\) is countably additive: whenever \(\{A_k\}_{k \ge 1} \subset \mathcal{F}\) are pairwise disjoint,

\[
\nu\Bigl(\bigcup_{k=1}^{\infty} A_k\Bigr)
=
\sum_{k=1}^{\infty} \nu(A_k),
\]

where the series converges absolutely if the sum is finite. A (positive) measure is a signed measure taking values in \([0, +\infty]\); a probability measure is a positive measure with \(\nu(\Omega) = 1\). Every signed measure admits a Jordan decomposition \(\nu = \nu^{+} - \nu^{-}\) as a difference of mutually singular positive measures. In particular, the probability measures \(\mathbb{P}\) and \(\mathbb{Q}\) in Problem~1 are finite positive measures, hence signed measures.
\end{definition}

\begin{definition}[Radon--Nikodym derivative]
\label{def:radon-nikodym}
Let \(\mu\) be a \(\sigma\)-finite positive measure and \(\nu\) a \(\sigma\)-finite signed measure on \((\Omega, \mathcal{F})\). We say that \(\nu\) is \emph{absolutely continuous} with respect to \(\mu\), written \(\nu \ll \mu\), if \(\mu(A) = 0\) implies \(\nu(A) = 0\). In that case there exists a \(\mu\)-almost everywhere unique measurable function \(f \colon \Omega \to \mathbb{R}\), called the \emph{Radon--Nikodym derivative} of \(\nu\) with respect to \(\mu\) and written \(f = \frac{\mathrm{d}\nu}{\mathrm{d}\mu}\), such that

\[
\nu(A)
=
\int_{A} f\,\mathrm{d}\mu
\qquad
\text{for every } A \in \mathcal{F}.
\]

If \(\mathbb{P}\) and \(\mathbb{Q}\) are probability measures with \(\mathbb{Q} \ll \mathbb{P}\), then \(\frac{\mathrm{d}\mathbb{Q}}{\mathrm{d}\mathbb{P}} \ge 0\) holds \(\mathbb{P}\)-a.s., \(\E_{\mathbb{P}}\bigl[\frac{\mathrm{d}\mathbb{Q}}{\mathrm{d}\mathbb{P}}\bigr] = 1\), and

\[
\E_{\mathbb{Q}}[Y]
=
\E_{\mathbb{P}}\Bigl[Y \frac{\mathrm{d}\mathbb{Q}}{\mathrm{d}\mathbb{P}}\Bigr]
\]

for every nonnegative measurable random variable \(Y\), and for every \(\mathbb{Q}\)-integrable \(Y\). Conversely, any nonnegative random variable \(Z\) with \(\E_{\mathbb{P}}[Z] = 1\) defines a probability measure \(\mathbb{Q}\) by \(\frac{\mathrm{d}\mathbb{Q}}{\mathrm{d}\mathbb{P}} \coloneqq Z\). This is the construction used in Problem~1: the exponential tilt is already normalized so that its \(\mathbb{P}\)-expectation equals \(1\).
\end{definition}

\begin{definition}[\(\mathbb{R}^{n}\)-valued Gaussian random vector]
\label{def:gaussian-random-vector}
Let \(X = (X_{1}, \dots, X_{n})^{T}\) be an \(\mathbb{R}^{n}\)-valued random vector on a probability space \((\Omega, \mathcal{F}, \mathbb{P})\). We say that \(X\) is \emph{Gaussian} under \(\mathbb{P}\) (possibly degenerate) if for every \(a \in \mathbb{R}^{n}\) the real random variable \(a^{T} X\) is univariate Gaussian.\footnote{Univariate means real-valued, as opposed to \(\mathbb{R}^{n}\)-valued. A real random variable \(Y\) is univariate Gaussian if \(Y\sim\mathcal{N}(m,\sigma^{2})\) for some \(m\in\mathbb{R}\) and \(\sigma^{2}\ge 0\), equivalently \(\E[e^{itY}]=\exp(itm-\tfrac12\sigma^{2}t^{2})\) for all \(t\in\mathbb{R}\). Degenerate means \(\sigma=0\): then \(Y=m\) almost surely, i.e.\ the law is the point mass \(\delta_{m}\), also written \(\mathcal{N}(m,0)\). This case is allowed because \(a^{T}X\) can be constant (take \(a=0\), or any \(a\) in the kernel of a singular covariance).} Equivalently, there exist \(\bm{\mu} \in \mathbb{R}^{n}\) and a symmetric positive semidefinite matrix \(V \in \mathbb{R}^{n \times n}\) such that the characteristic function of \(X\) is

\[
\E_{\mathbb{P}}\bigl[e^{i t^{T} X}\bigr]
=
\exp\Bigl(i t^{T} \bm{\mu} - \frac{1}{2} t^{T} V t\Bigr),
\qquad
t \in \mathbb{R}^{n},
\]

or, equivalently, the moment generating function exists on all of \(\mathbb{R}^{n}\) and equals

\[
\E_{\mathbb{P}}\bigl[e^{t^{T} X}\bigr]
=
\exp\Bigl(t^{T} \bm{\mu} + \frac{1}{2} t^{T} V t\Bigr).
\]

Then \(\E_{\mathbb{P}}[X] = \bm{\mu}\) and \(\mathrm{cov}_{\mathbb{P}}(X) = V\), and we write \(X \sim \mathcal{N}(\bm{\mu}, V)\) under \(\mathbb{P}\). See \citep[\S3.3, \S3.9]{durrett2010}. If \(V\) is positive definite, \(X\) admits the Lebesgue density

\[
(2\pi)^{-n/2} (\det V)^{-1/2}
\exp\Bigl(
-\frac{1}{2} (x - \bm{\mu})^{T} V^{-1} (x - \bm{\mu})
\Bigr)
\]

on \(\mathbb{R}^{n}\); if \(V\) is only positive semidefinite, the law of \(X\) is supported on an affine subspace and need not be absolutely continuous with respect to Lebesgue measure on \(\mathbb{R}^{n}\).
\end{definition}

\subsubsection{Solution of Problem 1}

\paragraph{Answer.}

Under \(\mathbb{Q}\), the vector \(X\) is Gaussian with mean \(\bm{\mu}+V\alpha\) and the same covariance \(V\). Equivalently \(X\sim\mathcal{N}(\bm{\mu}+V\alpha,V)\) under \(\mathbb{Q}\).

\begin{proof}[Correct proof]
\textbf{Step 0 (the new measure is a probability).}
The random variable
\[
Z
\coloneqq
\frac{\exp(\alpha^{T}X)}{\E_{\mathbb{P}}[\exp(\alpha^{T}X)]}
\]
is nonnegative and satisfies \(\E_{\mathbb{P}}[Z]=1\). By \cref{def:radon-nikodym} it is a Radon--Nikodym derivative, and \(\mathbb{Q}\) is a probability measure with \(\mathbb{Q}\ll\mathbb{P}\) and
\[
\E_{\mathbb{Q}}[Y]
=
\E_{\mathbb{P}}[YZ]
\]
for every nonnegative measurable \(Y\) (and every \(\mathbb{Q}\)-integrable \(Y\)). The denominator is finite: \cref{def:gaussian-random-vector} supplies a moment generating function on all of \(\mathbb{R}^{n}\), so
\[
\E_{\mathbb{P}}\bigl[e^{\alpha^{T}X}\bigr]
=
\exp\Bigl(\alpha^{T}\bm{\mu}+\tfrac12\alpha^{T}V\alpha\Bigr)
\in(0,\infty).
\]

\textbf{Step 1 (moment generating function under \(\mathbb{Q}\)).}
Fix an arbitrary \(t\in\mathbb{R}^{n}\). Then
\[
\begin{aligned}
\E_{\mathbb{Q}}\bigl[e^{t^{T}X}\bigr]
&=
\E_{\mathbb{P}}\bigl[e^{t^{T}X}Z\bigr]
=
\frac{\E_{\mathbb{P}}\bigl[e^{(t+\alpha)^{T}X}\bigr]}{\E_{\mathbb{P}}\bigl[e^{\alpha^{T}X}\bigr]}.
\end{aligned}
\]
Apply the \(\mathbb{P}\)-moment generating function from \cref{def:gaussian-random-vector} at the two real vectors \(t+\alpha\) and \(\alpha\):

\[
\E_{\mathbb{P}}\bigl[e^{(t+\alpha)^{T}X}\bigr]
=
\exp\Bigl(
(t+\alpha)^{T}\bm{\mu}
+
\tfrac12(t+\alpha)^{T}V(t+\alpha)
\Bigr).
\]

Since \(V\) is symmetric,
\[
(t+\alpha)^{T}V(t+\alpha)
=
t^{T}Vt
+
2\alpha^{T}Vt
+
\alpha^{T}V\alpha.
\]
The ratio of the two exponentials is therefore
\[
\begin{aligned}
\E_{\mathbb{Q}}\bigl[e^{t^{T}X}\bigr]
&=
\exp\Bigl(
t^{T}\bm{\mu}
+
\alpha^{T}Vt
+
\tfrac12 t^{T}Vt
\Bigr)
=
\exp\Bigl(
t^{T}(\bm{\mu}+V\alpha)
+
\tfrac12 t^{T}Vt
\Bigr).
\end{aligned}
\]

\textbf{Step 2 (read off the law).}
The right-hand side is the moment generating function of \(\mathcal{N}(\bm{\mu}+V\alpha,V)\). Equivalently, replacing \(t\) by \(it\) (or quoting the characteristic-function form in \cref{def:gaussian-random-vector}),

\[
\E_{\mathbb{Q}}\bigl[e^{it^{T}X}\bigr]
=
\exp\Bigl(
it^{T}(\bm{\mu}+V\alpha)
-
\tfrac12 t^{T}Vt
\Bigr),
\qquad
t\in\mathbb{R}^{n}.
\]

Hence \(X\) is Gaussian under \(\mathbb{Q}\), with mean \(\bm{\mu}+V\alpha\) and covariance \(V\).
\end{proof}

\subsubsection{Manuscript proof idea}

The calculation below follows the handwritten notes. The opening identifications of the Gaussian characteristic function, the moment generating function, and \(\frac{\mathrm{d}\mathbb{Q}}{\mathrm{d}\mathbb{P}}\) are correct (they match \cref{def:gaussian-random-vector,def:radon-nikodym}). \mserr{Red} marks a false claim or a missing justification. The notes write \(d\) for the ambient dimension; this is the \(n\) of the exam.

The argument begins correctly:

\[
\E_{\mathbb{Q}}\bigl[e^{it^{T}X}\bigr]
=
\E_{\mathbb{P}}\Bigl[e^{it^{T}X}\frac{\mathrm{d}\mathbb{Q}}{\mathrm{d}\mathbb{P}}\Bigr]
=
\frac{\E_{\mathbb{P}}\bigl[e^{(it+\alpha)^{T}X}\bigr]}{\E_{\mathbb{P}}\bigl[e^{\alpha^{T}X}\bigr]}
=
\frac{\E_{\mathbb{P}}\bigl[e^{i(t-i\alpha)^{T}X}\bigr]}{\E_{\mathbb{P}}\bigl[e^{\alpha^{T}X}\bigr]}.
\]

This identity is correct. The last rewriting is the characteristic function of \(X\) under \(\mathbb{P}\), evaluated at the complex vector \(t-i\alpha\). \mserr{Gap: the formula in \cref{def:gaussian-random-vector} is stated for real frequencies. Extending it to \(t-i\alpha\in\mathbb{C}^{n}\) is legitimate because the Gaussian moment generating function exists on all of \(\mathbb{R}^{n}\) and is entire, but that justification was not written. The clean way to avoid the issue is to stay on the real axis and use the moment generating function, as in the correct proof.}

The notes compute the denominator by the same substitution, \(\alpha=i(-i\alpha)\), and write

\[
\E_{\mathbb{P}}\bigl[e^{\alpha^{T}X}\bigr]
=
\E_{\mathbb{P}}\bigl[e^{i(-i\alpha)^{T}X}\bigr]
=
\mserr{\exp\Bigl(i(-i\alpha)^{T}\bm{\mu}+\tfrac12(-i\alpha)^{T}V(-i\alpha)\Bigr)}
=
\mserr{\exp\bigl(-i\alpha^{T}\bm{\mu}-\tfrac12\alpha^{T}V\alpha\bigr)}.
\]

\mserr{False: the characteristic function uses \(-\frac12 t^{T}Vt\), not \(+\frac12 t^{T}Vt\). With the correct quadratic sign one gets \(\exp(\alpha^{T}\bm{\mu}+\frac12\alpha^{T}V\alpha)\), which is also the moment generating function already recorded at the top of the notes. Independently, \(i(-i\alpha)^{T}\bm{\mu}=\alpha^{T}\bm{\mu}\), not \(-i\alpha^{T}\bm{\mu}\).} (The same moment generating function was written correctly as ``eq.'' on the notes; it was not used.)

The numerator is then expanded as if the characteristic function had a positive quadratic term:

\[
\E_{\mathbb{P}}\bigl[e^{i(t-i\alpha)^{T}X}\bigr]
=
\mserr{\exp\Bigl(
i(t-i\alpha)^{T}\bm{\mu}
+
\tfrac12(t-i\alpha)^{T}V(t-i\alpha)
\Bigr)}.
\]

\mserr{False: this mixes the two formulae. Either use the characteristic function with a minus quadratic at the complex frequency \(u=t-i\alpha\), namely \(\exp\bigl(iu^{T}\bm{\mu}-\frac12 u^{T}Vu\bigr)\), or use the moment generating function at the real-plus-imaginary vector \(z=\alpha+it\), namely \(\exp\bigl(z^{T}\bm{\mu}+\frac12 z^{T}Vz\bigr)\). The displayed line takes the \(i(t-i\alpha)\) from the first formula and the \(+\frac12\) from the second.}

Dividing by the (already incorrect) denominator and expanding, the notes obtain

\[
\begin{aligned}
&\exp\bigl(
i\alpha^{T}\bm{\mu}
+
\tfrac12\alpha^{T}V\alpha
+
t^{T}\bm{\mu}
-
i\alpha^{T}\bm{\mu}
+
\tfrac12 t^{T}Vt
-
i\alpha^{T}Vt
-
it^{T}V\alpha
-
\tfrac12\alpha^{T}V\alpha
\bigr)
\\
&\qquad=
\mserr{\exp\bigl(
t^{T}\bm{\mu}
+
\tfrac12 t^{T}Vt
-
i\alpha^{T}Vt
-
it^{T}V\alpha
\bigr)}.
\end{aligned}
\]

The cancellation of \(\pm\frac12\alpha^{T}V\alpha\) and of \(\pm i\alpha^{T}\bm{\mu}\) is algebraically consistent with the two preceding false formulae, so it does not repair them. \mserr{The output is not a characteristic function: the linear term \(t^{T}\bm{\mu}\) is missing a factor of \(i\), and the quadratic term has the moment-generating sign \(+\frac12 t^{T}Vt\) instead of \(-\frac12 t^{T}Vt\).} After restoring the two signs and using symmetry of \(V\), the cross terms combine to \(it^{T}V\alpha\), and one would have reached
\[
\exp\Bigl(it^{T}(\bm{\mu}+V\alpha)-\tfrac12 t^{T}Vt\Bigr).
\]
\mserr{Gap: the notes stop at the last exponential and never identify the mean \(\bm{\mu}+V\alpha\) and covariance \(V\).}

\subsection{Problem 2. [15 pts.]}
Let \( \{X_n\}_{n \in \mathbb{N}} \) be a sequence of independent random variables with distribution:

\[
\begin{cases}
\mathbb{P}[X_n = n^2] = \mathbb{P}[X_n = -n^2] = \dfrac{1}{6n^2}, \\[0.6em]
\mathbb{P}[X_n = n] = \mathbb{P}[X_n = -n] = \dfrac{1}{6}, \\[0.6em]
\mathbb{P}[X_n = 0] = \dfrac{2}{3} - \dfrac{1}{3n^2}.
\end{cases}
\]

Let \( S_n = \sum_{j=1}^n X_j \). Find a sequence of positive real numbers \( \{b_n\}_{n \in \mathbb{N}} \) such that the law of \( S_n / b_n \) converges to standard normal distribution \( \mathcal{N}(0, 1) \).

\setcounter{section}{2}
\setcounter{definition}{0}

\subsubsection{Theorems used in Problem 2}

\begin{definition}[Characteristic function]
\label{def:chf}
If \(X\) is a real random variable, its \emph{characteristic function} is

\[
\varphi_{X}(t)
\coloneqq
\E\bigl[e^{itX}\bigr]
=
\E[\cos(tX)] + i\,\E[\sin(tX)],
\qquad
t \in \mathbb{R}.
\]

Always \(\varphi_{X}(0) = 1\), \(|\varphi_{X}(t)| \le 1\), and \(\varphi_{X}\) is uniformly continuous on \(\mathbb{R}\). If \(X_{1}, \dots, X_{n}\) are independent, then

\[
\varphi_{X_{1}+\cdots+X_{n}}(t)
=
\prod_{j=1}^{n} \varphi_{X_{j}}(t).
\]

The law of \(X\) is uniquely determined by \(\varphi_{X}\). If \(Y = aX + b\) with \(a, b \in \mathbb{R}\), then \(\varphi_{Y}(t) = e^{itb}\,\varphi_{X}(at)\). See \citep[Theorems~3.3.1 and~3.3.2]{durrett2010} and \citep[Chapter~6]{chung1968}.
\end{definition}

\begin{note}[Common characteristic functions]
\label{note:chf-table}
The uniqueness theorem turns the following dictionary into an identification tool for Lévy limits. Parameterizations: Gamma is shape-rate, \(\mathrm{Exp}(\lambda)=\mathrm{Gamma}(1,\lambda)\), \(\chi^{2}(k)=\mathrm{Gamma}(k/2,1/2)\), Geometric is the number of trials until the first success, and \(\mathrm{NB}(r,p)\) is the number of failures before \(r\) successes. Cauchy is included because its moment generating function does not exist, while the characteristic function does.

\begin{center}
\renewcommand{\arraystretch}{1.45}
\resizebox{\linewidth}{!}{%
\begin{tabular}{@{}lll@{}}
\toprule
Law & pmf / pdf & \(\varphi_{X}(t)\) \\
\midrule
\multicolumn{3}{@{}l}{\emph{Discrete}} \\
\addlinespace
\(X\sim\delta_{a}\) &
\(\displaystyle \Pbb(X=a)=1\) &
\(\displaystyle e^{ita}\) \\
\(X\sim\mathrm{Bernoulli}(p)\) &
\(\displaystyle \Pbb(X=1)=p,\ \Pbb(X=0)=1-p\) &
\(\displaystyle 1-p+p e^{it}\) \\
\(X\sim\mathrm{Bin}(n,p)\) &
\(\displaystyle \Pbb(X=k)=\binom{n}{k}p^{k}(1-p)^{n-k},\ k=0,\dots,n\) &
\(\displaystyle (1-p+p e^{it})^{n}\) \\
\(X\sim\mathrm{Geo}(p)\) &
\(\displaystyle \Pbb(X=k)=p(1-p)^{k-1},\ k=1,2,\dots\) &
\(\displaystyle \dfrac{p e^{it}}{1-(1-p)e^{it}}\) \\
\(X\sim\mathrm{NB}(r,p)\) &
\(\displaystyle \Pbb(X=k)=\binom{k+r-1}{k}p^{r}(1-p)^{k},\ k=0,1,\dots\) &
\(\displaystyle \Bigl(\dfrac{p}{1-(1-p)e^{it}}\Bigr)^{r}\) \\
\(X\sim\mathrm{Poisson}(\lambda)\) &
\(\displaystyle \Pbb(X=k)=\dfrac{e^{-\lambda}\lambda^{k}}{k!},\ k=0,1,\dots\) &
\(\displaystyle \exp\bigl(\lambda(e^{it}-1)\bigr)\) \\
\midrule
\multicolumn{3}{@{}l}{\emph{Continuous}} \\
\addlinespace
\(X\sim\mathrm{Unif}[a,b]\) &
\(\displaystyle f(x)=\dfrac{1}{b-a},\ x\in[a,b]\) &
\(\displaystyle \dfrac{e^{itb}-e^{ita}}{it(b-a)}\) (\(t\neq 0\); \(1\) at \(0\)) \\
\(X\sim\mathrm{Exp}(\lambda)\) &
\(\displaystyle f(x)=\lambda e^{-\lambda x},\ x>0\) &
\(\displaystyle \dfrac{\lambda}{\lambda-it}\) \\
\(X\sim\mathrm{Gamma}(\alpha,\beta)\) &
\(\displaystyle f(x)=\dfrac{\beta^{\alpha}}{\Gamma(\alpha)}\,x^{\alpha-1}e^{-\beta x},\ x>0\) &
\(\displaystyle \bigl(1-\dfrac{it}{\beta}\bigr)^{-\alpha}\) \\
\(X\sim\chi^{2}(k)\) &
\(\displaystyle f(x)=\dfrac{1}{2^{k/2}\Gamma(k/2)}\,x^{k/2-1}e^{-x/2},\ x>0\) &
\(\displaystyle (1-2it)^{-k/2}\) \\
\(X\sim\mathcal{N}(\mu,\sigma^{2})\) &
\(\displaystyle f(x)=\dfrac{1}{\sqrt{2\pi\sigma^{2}}}\exp\bigl(-\dfrac{(x-\mu)^{2}}{2\sigma^{2}}\bigr),\ x\in\mathbb{R}\) &
\(\displaystyle \exp\bigl(i\mu t-\tfrac{1}{2}\sigma^{2} t^{2}\bigr)\) \\
\(X\sim\mathrm{Cauchy}(\mu,\sigma)\) &
\(\displaystyle f(x)=\dfrac{1}{\pi\sigma\bigl(1+\bigl(\frac{x-\mu}{\sigma}\bigr)^{2}\bigr)},\ x\in\mathbb{R}\) &
\(\displaystyle \exp\bigl(i\mu t-\sigma|t|\bigr)\) \\
\bottomrule
\end{tabular}%
}
\end{center}

Beta, Student's \(t\), \(F\), and hypergeometric laws have no elementary closed-form characteristic function. The table entries themselves are the output of \cref{prop:chf-compute}; recovering a law from a closed form \(\varphi\) is \cref{thm:chf-inversion}, with the three worked cases \cref{ex:chf-poisson,ex:chf-dirac,ex:chf-gamma}. In Problem~2 the Gaussian identification is the last row with \(\mu=0\) and \(\sigma^{2}=1\): \(S_{n}/b_{n}\xrightarrow{(d)}\mathcal{N}(0,1)\) if and only if the characteristic functions tend to \(\exp(-t^{2}/2)\). See \cref{thm:levy-continuity} and \citep[Theorem~3.3.6]{durrett2010}.
\end{note}

The definition \(\varphi_{X}(t)=\E[e^{itX}]\) is already an algorithm: expand the expectation against the law of \(X\). That is the forward map (law \(\mapsto\varphi\)). The inverse map is inversion plus uniqueness: a characteristic function determines at most one probability measure, and the formulae in \cref{thm:chf-inversion} reconstruct it. In exam work one almost never evaluates the inversion integral; one computes a product of named transforms and reads the law off \cref{note:chf-table}. Inversion is why that reading is legal.

\begin{proposition}[Computing \(\varphi_{X}\) from the law]
\label{prop:chf-compute}
Let \(X\) be real with law \(\mu=\mathrm{Law}(X)\). Then, in all cases,

\[
\varphi_{X}(t)
=
\int_{\mathbb{R}} e^{itx}\,\mu(\mathrm{d}x),
\qquad
t\in\mathbb{R}.
\]

Specialisations that are the actual computations:

\begin{enumerate}[label=\textup{(\roman*)}]
    \item
    If \(\mu\) is discrete, supported on a countable set \(\{x_{k}\}\), then

\[
\varphi_{X}(t)
=
\sum_{k} e^{itx_{k}}\,\Pbb(X=x_{k}),
\]

the series being absolutely convergent.

    \item
    If \(\mu\) is absolutely continuous with density \(f\), then

\[
\varphi_{X}(t)
=
\int_{\mathbb{R}} e^{itx} f(x)\,\mathrm{d}x.
\]

    \item
    Independence and affine images, already in \cref{def:chf}: if \(X_{1},\dots,X_{n}\) are independent then \(\varphi_{X_{1}+\cdots+X_{n}}=\prod_{j}\varphi_{X_{j}}\); if \(Y=aX+b\) then \(\varphi_{Y}(t)=e^{itb}\varphi_{X}(at)\).
\end{enumerate}

The integral in (ii) is the Fourier transform of \(f\) (up to the sign convention \(e^{+itx}\)). It always exists as a bounded continuous function, because \(|e^{itx}f(x)|\le f(x)\). The moment generating function \(\E[e^{tX}]\) is the same integral with \(t\) real, and may fail to exist (Cauchy). See \citep[Theorems~3.3.1 and~3.3.2]{durrett2010} and \citep[Chapter~6]{chung1968}.
\end{proposition}

\begin{theorem}[Inversion formulae]
\label{thm:chf-inversion}
Let \(\mu\) be a probability measure on \(\mathbb{R}\) with characteristic function \(\varphi\).

\begin{enumerate}[label=\textup{(\roman*)}]
    \item
    \emph{L\'evy's inversion.} For \(a<b\),

\[
\mu(a,b)+\tfrac12\mu(\{a\})+\tfrac12\mu(\{b\})
=
\lim_{T\to\infty}
\frac{1}{2\pi}
\int_{-T}^{T}
\frac{e^{-ita}-e^{-itb}}{it}\,\varphi(t)\,\mathrm{d}t.
\]

The integrand is extended by continuity at \(t=0\), where it equals \(b-a\). Equivalently, \((e^{-ita}-e^{-itb})/(it)=\int_{a}^{b}e^{-itx}\,\mathrm{d}x\). In particular \(\varphi\) determines \(\mu\), which is \citep[Theorem~3.3.2]{durrett2010}.

    \item
    \emph{Density inversion.} If \(\int_{\mathbb{R}}|\varphi(t)|\,\mathrm{d}t<\infty\), then \(\mu\) is absolutely continuous with a bounded continuous density

\[
f(x)
=
\frac{1}{2\pi}
\int_{\mathbb{R}} e^{-itx}\,\varphi(t)\,\mathrm{d}t.
\]

    \item
    \emph{Atoms.} For every \(a\in\mathbb{R}\),

\[
\mu(\{a\})
=
\lim_{T\to\infty}
\frac{1}{2T}
\int_{-T}^{T} e^{-ita}\,\varphi(t)\,\mathrm{d}t.
\]

    \item
    \emph{Integer lattice.} If \(\mu(\mathbb{Z})=1\), then \(\varphi\) is \(2\pi\)-periodic and

\[
\mu(\{k\})
=
\frac{1}{2\pi}
\int_{-\pi}^{\pi} e^{-itk}\,\varphi(t)\,\mathrm{d}t,
\qquad
k\in\mathbb{Z}.
\]
\end{enumerate}

Items~(i)--(ii) are \citep[Theorems~3.3.4 and~3.3.5]{durrett2010}; (iii)--(iv) are the corresponding Fourier formulae on a compact interval, as in \citep[Chapter~6]{chung1968}. The same density formula reappears as \cref{def:chf-inversion}. Riemann--Lebesgue: if \(\mu\) is absolutely continuous then \(\varphi(t)\to 0\) as \(|t|\to\infty\); the contrapositive detects atoms and lattice laws.
\end{theorem}

\begin{example}[Poisson, both directions]
\label{ex:chf-poisson}
Let \(X\sim\mathrm{Poisson}(\lambda)\) with \(\lambda>0\), so \(\Pbb(X=k)=e^{-\lambda}\lambda^{k}/k!\) for \(k=0,1,2,\dots\).

\emph{Forward.} By \cref{prop:chf-compute}(i),

\[
\varphi_{X}(t)
=
\sum_{k=0}^{\infty} e^{itk}\,e^{-\lambda}\frac{\lambda^{k}}{k!}
=
e^{-\lambda}
\sum_{k=0}^{\infty} \frac{(\lambda e^{it})^{k}}{k!}
=
e^{-\lambda}\exp(\lambda e^{it})
=
\exp\bigl(\lambda(e^{it}-1)\bigr).
\]

The rearrangement is absolute convergence of the exponential series.

\emph{Reverse.} Start from \(\varphi(t)=\exp(\lambda(e^{it}-1))\) and apply the lattice formula \cref{thm:chf-inversion}(iv). Expand the exponential:

\[
\varphi(t)
=
e^{-\lambda}\exp(\lambda e^{it})
=
e^{-\lambda}
\sum_{m=0}^{\infty} \frac{\lambda^{m}}{m!}\,e^{itm}.
\]

The characters \(t\mapsto e^{itm}\) are orthogonal on \([-\pi,\pi]\):

\[
\frac{1}{2\pi}
\int_{-\pi}^{\pi} e^{-itk}\,e^{itm}\,\mathrm{d}t
=
\mathbf{1}_{\{m=k\}}.
\]

Termwise integration (justified by uniform convergence on the compact interval, or by positivity of the Poisson masses) therefore yields

\[
\Pbb(X=k)
=
\frac{1}{2\pi}
\int_{-\pi}^{\pi} e^{-itk}\,\varphi(t)\,\mathrm{d}t
=
e^{-\lambda}\frac{\lambda^{k}}{k!}.
\]

Thus the inverse Fourier series of \(\exp(\lambda(e^{it}-1))\) is exactly the Poisson pmf. Note that \(|\varphi(t)|=\exp(\lambda(\cos t-1))\) is \(2\pi\)-periodic and does not tend to \(0\) as \(|t|\to\infty\), so Riemann--Lebesgue already forbids an absolutely continuous law.
\end{example}

\begin{example}[Point mass, both directions]
\label{ex:chf-dirac}
Let \(X\sim\delta_{a}\), i.e.\ \(\Pbb(X=a)=1\).

\emph{Forward.} The sum in \cref{prop:chf-compute}(i) has a single term:

\[
\varphi_{X}(t)
=
e^{ita}.
\]

In particular \(|\varphi_{X}|\equiv 1\), so again \(\varphi_{X}(t)\not\to 0\).

\emph{Reverse.} Start from \(\varphi(t)=e^{ita}\) and recover the atom by \cref{thm:chf-inversion}(iii). For any \(x\in\mathbb{R}\),

\[
\frac{1}{2T}
\int_{-T}^{T} e^{-itx}\,e^{ita}\,\mathrm{d}t
=
\frac{1}{2T}
\int_{-T}^{T} e^{it(a-x)}\,\mathrm{d}t
=
\begin{cases}
1, & x=a,\\[0.35em]
\dfrac{\sin\bigl(T(a-x)\bigr)}{T(a-x)}, & x\neq a.
\end{cases}
\]

The sinc quotient tends to \(0\) as \(T\to\infty\) whenever \(x\neq a\). Hence

\[
\lim_{T\to\infty}
\frac{1}{2T}
\int_{-T}^{T} e^{-itx}\,\varphi(t)\,\mathrm{d}t
=
\mathbf{1}_{\{x=a\}},
\]

which is \(\delta_{a}(\{x\})\). The same conclusion is L\'evy's formula \cref{thm:chf-inversion}(i): the right-hand side equals \(1\) if \(a\in(c,d)\), equals \(1/2\) if \(a\in\{c,d\}\), and equals \(0\) otherwise, which is the law \(\delta_{a}\). Uniqueness is the zero-effort variant: \(e^{ita}\) is already the first row of \cref{note:chf-table}.
\end{example}

\begin{example}[Gamma, both directions]
\label{ex:chf-gamma}
Let \(X\sim\mathrm{Gamma}(\alpha,\beta)\) in the shape-rate parameterisation of \cref{note:chf-table}: density \(f(x)=\beta^{\alpha}x^{\alpha-1}e^{-\beta x}/\Gamma(\alpha)\) on \((0,\infty)\), with \(\alpha>0\) and \(\beta>0\).

\emph{Forward.} By \cref{prop:chf-compute}(ii),

\[
\varphi_{X}(t)
=
\frac{\beta^{\alpha}}{\Gamma(\alpha)}
\int_{0}^{\infty} x^{\alpha-1} \exp\bigl(-(\beta-it)x\bigr)\,\mathrm{d}x.
\]

The Gamma integral \(\int_{0}^{\infty} x^{\alpha-1}e^{-zx}\,\mathrm{d}x=\Gamma(\alpha)\,z^{-\alpha}\) holds for every \(z\in\mathbb{C}\) with \(\operatorname{Re} z>0\), the power being \(z^{-\alpha}=\exp(-\alpha\operatorname{Log} z)\) with the principal logarithm (argument in \((-\pi/2,\pi/2)\), since \(\operatorname{Re} z>0\)). Taking \(z=\beta-it\) (which has real part \(\beta>0\)) gives

\[
\varphi_{X}(t)
=
\beta^{\alpha}(\beta-it)^{-\alpha}
=
\Bigl(1-\frac{it}{\beta}\Bigr)^{-\alpha}.
\]

The same computation at \(\alpha=1\) is the exponential row of the table; the \(\chi^{2}(k)\) row is the case \(\alpha=k/2\), \(\beta=1/2\).

\emph{Reverse.} Three routes, in increasing computational cost.

\begin{enumerate}[label=\textup{(\alph*)}]
    \item
    \emph{Uniqueness.} The forward computation plus \cref{thm:chf-inversion}(i) says that any random variable whose characteristic function is \((1-it/\beta)^{-\alpha}\) is \(\mathrm{Gamma}(\alpha,\beta)\). This is the exam device: match the closed form to \cref{note:chf-table}.

    \item
    \emph{Convolution, integer shape.} If \(\alpha=n\in\mathbb{N}\), then \(\varphi_{X}(t)=\bigl(\beta/(\beta-it)\bigr)^{n}=\bigl(\varphi_{E}(t)\bigr)^{n}\) where \(E\sim\mathrm{Exp}(\beta)\). Independence in \cref{prop:chf-compute}(iii) therefore yields \(X\stackrel{(d)}{=}E_{1}+\cdots+E_{n}\) with \(E_{j}\) i.i.d.\ exponential, which is the Erlang law \(\mathrm{Gamma}(n,\beta)\).

    \item
    \emph{Density inversion, \(\alpha>1\).} One has \(|\varphi(t)|\sim |t/\beta|^{-\alpha}\) as \(|t|\to\infty\), so \(\int|\varphi|<\infty\) if and only if \(\alpha>1\). In that range \cref{thm:chf-inversion}(ii) applies. The exponential law \(\alpha=1\) is the boundary case where the density formula is \emph{not} available (and indeed \(|\varphi(t)|=\beta/\sqrt{\beta^{2}+t^{2}}\) fails to be integrable). As a complete residue computation take \(\alpha=2\), so \(\varphi(t)=\bigl(\beta/(\beta-it)\bigr)^{2}\) and \(\int|\varphi|<\infty\). For \(x>0\) the integrand \(e^{-itx}(\beta/(\beta-it))^{2}\) decays in the lower half-plane (\(|e^{-itx}|=e^{x\operatorname{Im} t}\)), and the only singularity is a pole of order \(2\) at \(t=-i\beta\). The substitution \(z=t+i\beta\) gives \(\beta-it=-iz\) and \(e^{-itx}=e^{-\beta x}e^{-izx}\), hence the integrand is \(-\beta^{2} e^{-\beta x}\,e^{-izx}z^{-2}\). Expanding \(e^{-izx}=1-ixz+O(z^{2})\) yields residue \(ix\beta^{2}e^{-\beta x}\) at \(z=0\). The contour is clockwise, so

\[
\int_{\mathbb{R}} e^{-itx}\Bigl(\frac{\beta}{\beta-it}\Bigr)^{2}\,\mathrm{d}t
=
-2\pi i\cdot(ix\beta^{2}e^{-\beta x})
=
2\pi\,\beta^{2} x e^{-\beta x}.
\]

Dividing by \(2\pi\) recovers the \(\mathrm{Gamma}(2,\beta)\) density \(\beta^{2} x e^{-\beta x}\) on \((0,\infty)\). For \(x<0\) one closes in the upper half-plane, where there are no poles, and the density vanishes. For non-integer \(\alpha>1\) the same inversion integral is legitimate, but the singularity at \(t=-i\beta\) is a branch point rather than a pole; uniqueness~(a) is then the cheaper identification.
\end{enumerate}
\end{example}

\begin{remark}[Forward versus reverse, on an exam]
\label{rem:chf-exam}
The three examples are the two computational directions. Forward is always \cref{prop:chf-compute}: a series for Poisson and \(\delta_{a}\), a Gamma integral for \(\mathrm{Gamma}(\alpha,\beta)\). Reverse is \cref{thm:chf-inversion}: a Fourier series on \([-\pi,\pi]\) for a lattice law, a Ces\`aro mean for an atom, and either uniqueness or (when \(\int|\varphi|<\infty\)) a Fourier integral for a density. In Problem~2 one uses neither the series nor the inversion integral. One computes \(\varphi_{S_{n}/b_{n}}(t)\) from independence, passes to the limit, and invokes L\'evy continuity \cref{thm:levy-continuity} to identify \(\mathcal{N}(0,1)\). That is uniqueness applied to a \emph{limit} of characteristic functions, which is the typical QE pattern; cf.\ \cref{study:pdf-methods}.
\end{remark}

\begin{definition}[Convergence in distribution]
\label{def:conv-d}
A sequence of real random variables \(X_{n}\) \emph{converges in distribution} (equivalently: \emph{converges weakly}) to \(X\), written \(X_{n} \xrightarrow{(d)} X\) or \(X_{n} \Rightarrow X\), if

\[
F_{n}(x) \coloneqq \Pbb(X_{n} \le x)
\to
F(x) \coloneqq \Pbb(X \le x)
\]

at every continuity point \(x\) of \(F\). Equivalently, \(\E[g(X_{n})] \to \E[g(X)]\) for every bounded continuous \(g \colon \mathbb{R} \to \mathbb{R}\). Equivalently, writing \(\mu_{n} = \mathrm{Law}(X_{n})\) and \(\mu = \mathrm{Law}(X)\): \(\liminf_{n} \mu_{n}(G) \ge \mu(G)\) for every open \(G\); \(\limsup_{n} \mu_{n}(K) \le \mu(K)\) for every closed \(K\); and \(\mu_{n}(A) \to \mu(A)\) whenever \(\mu(\partial A) = 0\). See \citep[Theorems~3.2.3 and~3.2.5]{durrett2010}.
\end{definition}

\begin{theorem}[L\'evy continuity theorem]
\label{thm:levy-continuity}
Let \(\{\mu_{n}\}_{1 \le n \le \infty}\) be probability measures on \(\mathbb{R}\) with characteristic functions \(\{\varphi_{n}\}_{1 \le n \le \infty}\).

\begin{enumerate}[label=\textup{(\roman*)}]
    \item
    If \(\mu_{n} \Rightarrow \mu_{\infty}\), then \(\varphi_{n}(t) \to \varphi_{\infty}(t)\) for every \(t \in \mathbb{R}\). The convergence is uniform on every compact interval, and \(\{\varphi_{n}\}\) is equicontinuous on \(\mathbb{R}\).

    \item
    Conversely, if \(\varphi_{n}(t) \to \varphi(t)\) for every \(t \in \mathbb{R}\) and the pointwise limit \(\varphi\) is continuous at \(t = 0\), then \(\{\mu_{n}\}\) is tight, \(\mu_{n} \Rightarrow \mu\) for a probability measure \(\mu\), and \(\varphi\) is the characteristic function of \(\mu\).

    \item
    In particular, if \(\mu\) is already a probability measure with characteristic function \(\varphi\), then

\[
\mu_{n} \Rightarrow \mu
\qquad\Longleftrightarrow\qquad
\varphi_{n}(t) \to \varphi(t)
\quad \text{for all } t \in \mathbb{R}.
\]

    A characteristic function is everywhere continuous, so the continuity-at-\(0\) hypothesis in (ii) is automatic in this case. Applied to Problem~2: \(S_{n}/b_{n} \xrightarrow{(d)} \mathcal{N}(0,1)\) if and only if

\[
\E\bigl[\exp\bigl(it S_{n}/b_{n}\bigr)\bigr]
\to
\exp(-t^{2}/2)
\qquad \text{for every } t \in \mathbb{R}.
\]
\end{enumerate}

Continuity of the pointwise limit at \(0\) cannot be omitted. If \(\mu_{n} = \mathcal{N}(0,n)\), then \(\varphi_{n}(t) = \exp(-nt^{2}/2) \to 0\) for \(t \neq 0\) while \(\varphi_{n}(0) = 1\), so the pointwise limit is discontinuous at \(0\) and \(\{\mu_{n}\}\) is not tight. See \citep[Theorem~3.3.6]{durrett2010} and \citep[Theorems~6.3.1 and~6.3.2]{chung1968}; the compact uniform upgrade is \citep[Exercise~3.3.16(ii)]{durrett2010}.
\end{theorem}

\begin{remark}[References]
\label{rem:levy-refs}
\Cref{def:chf} follows \citep[Theorems~3.3.1 and~3.3.2]{durrett2010} and \citep[Chapter~6]{chung1968}. \Cref{prop:chf-compute,thm:chf-inversion} are the computational unpacking of those theorems: forward formulae in \citep[Theorems~3.3.1 and~3.3.2]{durrett2010}, L\'evy inversion and density inversion in \citep[Theorems~3.3.4 and~3.3.5]{durrett2010}, with the atom and lattice formulae as in \citep[Chapter~6]{chung1968}. \Cref{def:conv-d} is the Portmanteau package in \citep[Theorems~3.2.3 and~3.2.5]{durrett2010}. \Cref{thm:levy-continuity} is \citep[Theorem~3.3.6, \S3.3.2, pp.~97--98]{durrett2010} together with \citep[Theorems~6.3.1 and~6.3.2, \S6.3]{chung1968}. Chung writes \(\xrightarrow{v}\) for vague convergence; once the limit is a probability measure this is ordinary weak convergence, i.e.\ \(\xrightarrow{(d)}\) for the corresponding random variables. When both sides are already probability measures, Chung records \(\mu_{n} \xrightarrow{v} \mu \iff \varphi_{n} \xrightarrow{u} \varphi\) as display~(6) in that section. The \(\mathcal{N}(0,n)\) counterexample is the remark after \citep[Theorem~3.3.6]{durrett2010}.
\end{remark}

\subsubsection{Solution of Problem 2}

\paragraph{Answer.}

Take \(b_{n}=\sqrt{n(n+1)(2n+1)/18}\). In particular \(b_{n}\sim n^{3/2}/3\).

Fix \(t\in\mathbb{R}\) and write
\[
v_{n}
\coloneqq
\sum_{j=1}^{n}j^{2}
=
\frac{n(n+1)(2n+1)}{6},
\qquad
\sum_{j=1}^{n}j^{4}
=
O(n^{5}).
\]
Two scales will appear:

\[
b_{n}
=
\sqrt{\frac{v_{n}}{3}}
=
\sqrt{\frac{n(n+1)(2n+1)}{18}},
\qquad
\tilde{b}_{n}
=
\sqrt{\frac{2v_{n}}{3}}
=
\frac13\sqrt{n(n+1)(2n+1)}.
\]

Thus \(\tilde{b}_{n}=\sqrt{2}\,b_{n}\) exactly. The first is the truncated-variance scale used in the proof; the second is the manuscript scale \(\sqrt{\mathrm{Var}(S_{n})}\).

\begin{proof}[Correct proof]
\textbf{Step 0 (two pieces of the variance).}
Each \(X_{j}\) is symmetric, so \(\E[X_{j}]=0\). The second moment splits by atoms:

\[
\begin{aligned}
\E\bigl[X_{j}^{2}\mathbf{1}_{\{|X_{j}|=j\}}\bigr]
&=
2\cdot j^{2}\cdot\frac16
=
\frac{j^{2}}{3},
\\
\E\bigl[X_{j}^{2}\mathbf{1}_{\{|X_{j}|=j^{2}\}}\bigr]
&=
2\cdot j^{4}\cdot\frac{1}{6j^{2}}
=
\frac{j^{2}}{3},
\\
\mathrm{Var}(X_{j})
&=
\frac{2j^{2}}{3}.
\end{aligned}
\]

Hence

\[
\sum_{j=1}^{n}\E\bigl[X_{j}^{2}\mathbf{1}_{\{|X_{j}|=j\}}\bigr]
=
\frac{v_{n}}{3}
=
b_{n}^{2},
\qquad
\mathrm{Var}(S_{n})
=
\frac{2v_{n}}{3}
=
\tilde{b}_{n}^{2}.
\]

The two pieces are equal as numbers. They are not equal as contributions to the Lévy limit: the second piece is carried by atoms of size \(j^{2}\), and \(\sum j^{2}\) is dominated by \(j\asymp n\).

\textbf{Step 1 (characteristic functions).}
The five-point law gives

\[
\varphi_{X_{j}}(t)
=
\frac{1}{3j^{2}}\cos(tj^{2})
+
\frac13\cos(tj)
+
\frac23
-
\frac{1}{3j^{2}}.
\]

Independence and \cref{def:chf} give \(\varphi_{S_{n}/b_{n}}(t)=\prod_{j=1}^{n}\varphi_{X_{j}}(t/b_{n})\). Rewrite each factor as

\[
\varphi_{X_{j}}(t/b_{n})
=
1+a_{n,j}+c_{n,j},
\]

where

\[
a_{n,j}
=
\frac13\Bigl(\cos\bigl(tj/b_{n}\bigr)-1\Bigr),
\qquad
c_{n,j}
=
\frac{1}{3j^{2}}\Bigl(\cos\bigl(tj^{2}/b_{n}\bigr)-1\Bigr).
\]

The term \(a_{n,j}\) is the \(\pm j\) contribution; \(c_{n,j}\) is the \(\pm j^{2}\) contribution. Use \(|\cos u-1|\le\min(2,u^{2}/2)\) throughout.

\textbf{Step 2 (Taylor for \(a_{n,j}\) is uniform in \(j\le n\)).}
Let \(\theta_{n,j}=tj/b_{n}\). Then \(|\theta_{n,j}|\le |t|n/b_{n}\sim 3|t|n^{-1/2}\to 0\), uniformly in \(j\le n\). The Taylor remainder \(|\cos\theta-1+\theta^{2}/2|\le\theta^{4}/24\) therefore applies to every factor, and

\[
a_{n,j}
=
-\frac{t^{2}j^{2}}{6b_{n}^{2}}
+
r_{n,j},
\qquad
|r_{n,j}|
\le
\frac{t^{4}j^{4}}{72\,b_{n}^{4}}.
\]

Summing and using \(b_{n}^{2}=v_{n}/3\),

\[
\sum_{j=1}^{n}a_{n,j}
=
-\frac{t^{2}v_{n}}{6b_{n}^{2}}
+
O\Bigl(\frac{n^{5}}{b_{n}^{4}}\Bigr)
=
-\frac{t^{2}}{2}
+
O(n^{-1}).
\]

\textbf{Step 3 (Taylor for \(c_{n,j}\) is not uniform; the sum still vanishes).}
The naive expansion \(\cos u=1-u^{2}/2+O(u^{4})\) with \(u=tj^{2}/b_{n}\) would give \(c_{n,j}\approx -t^{2}j^{2}/(6b_{n}^{2})\), hence \(\sum c_{n,j}\approx -t^{2}/2\), doubling the exponent. This is illegal: \(n^{2}/b_{n}\sim 3\sqrt{n}\not\to 0\), so \(u_{n,n}\) is not small.

To estimate the true sum, split at \(J_{n}=\lfloor n^{2/3}\rfloor\). Then \(J_{n}^{2}/b_{n}\sim n^{4/3}/n^{3/2}=n^{-1/6}\to 0\), so on \(\{j\le J_{n}\}\) the argument \(tj^{2}/b_{n}\) \emph{is} uniformly small and the quadratic bound is legitimate:

\[
\sum_{j\le J_{n}}|c_{n,j}|
\le
\sum_{j\le J_{n}}\frac{t^{2}j^{2}}{6b_{n}^{2}}
=
O\Bigl(\frac{J_{n}^{3}}{b_{n}^{2}}\Bigr)
=
O(n^{-1}).
\]

On the complementary range the cosine is merely bounded:

\[
\sum_{j>J_{n}}|c_{n,j}|
\le
\sum_{j>J_{n}}\frac{2}{3j^{2}}
\le
\int_{J_{n}}^{\infty}\frac{2}{3x^{2}}\,\mathrm{d}x
=
O(n^{-2/3}).
\]

Therefore \(\sum_{j=1}^{n}c_{n,j}\to 0\). In particular the \(\pm j^{2}\) atoms contribute nothing to \(\log\varphi_{S_{n}/b_{n}}(t)\).

The same split explains why ``half the variance'' is invisible. The rare-jump variance \(\sum_{j}j^{2}/3\) is carried by \(j\asymp n\), which is exactly the range \(j>J_{n}\) where Taylor fails. On the range where Taylor is valid, the rare-jump contribution is \(O(J_{n}^{3}/b_{n}^{2})=O(n^{-1})\).

\begin{remark}[The split \(J_{n}\) is not a parameter of the limit]
\label{rem:p2-cutoff-window}
Two different ratios appear, and they must not be fused.

The ratio \(n/b_{n}\) is forced by the scaling. With the correct choice \(b_{n}\sim n^{3/2}/3\) one has \(n/b_{n}\sim 3n^{-1/2}\to 0\), so Taylor for every \(a_{n,j}\) is automatic. If one took \(b_{n}\) only of order \(n\), this ratio would not tend to \(0\), the expansion of \(a_{n,j}\) would fail, and \(S_{n}/b_{n}\) would typically fail to converge to \(\mathcal{N}(0,1)\) (the truncated variance is of order \(n^{3}\)). If one took \(b_{n}\gg n^{3/2}\), Taylor would be even better, but \(S_{n}/b_{n}\to 0\) in probability: the limit exists and is degenerate. Thus ``too slow / too fast shrinking'' of \(j/b_{n}\) is a statement about the choice of \(b_{n}\), not about a free cutoff.

The integer \(J_{n}\) in Step~3 is only a proof device. It never appears in the definition of \(S_{n}/b_{n}\). Any sequence with
\[
J_{n}\to\infty
\qquad\text{and}\qquad
J_{n}=o(n^{3/4})
\]
works: the first forces \(\sum_{j>J_{n}}j^{-2}\to 0\), the second forces \(J_{n}^{2}/b_{n}\to 0\) so that Taylor is legal on \(\{j\le J_{n}\}\). The choice \(J_{n}=n^{2/3}\) sits in the interior of this window. If \(J_{n}\) stays bounded, the crude tail bound \(O(J_{n}^{-1})\) need not vanish, so \emph{this estimate} fails. If \(J_{n}\) is as large as \(n\), then \(J_{n}^{2}/b_{n}\not\to 0\), which is exactly the manuscript's illegal Taylor. In both cases the object \(\sum_{j}c_{n,j}\) is unchanged: a bad split does not create or destroy the distributional limit, it only makes the present comparison fail.

Oscillation of \(\cos(tj^{2}/b_{n})\) does not threaten existence of the limit of the sum. The bound \(\lvert\sum c_{n,j}\rvert\le\sum\lvert c_{n,j}\rvert\to 0\) yields \(\sum c_{n,j}\to 0\) by squeezing, whether or not individual cosines settle.

Finally, Lévy's theorem does not run backwards. \Cref{thm:levy-continuity} says: if \(\varphi_{S_{n}/b_{n}}(t)\) converges pointwise to a function continuous at \(0\), then that function is a characteristic function and \(S_{n}/b_{n}\) converges in distribution. One must still compute the pointwise limit (Steps~2--4). One cannot argue that ``Taylor fails, therefore the limit of \(\varphi_{n}\) might fail to exist, therefore it cannot fail, therefore \(S_{n}/b_{n}\) converges''. Failure of Taylor on the \(\pm j^{2}\) atoms is compatible with \(\varphi_{n}(t)\to e^{-t^{2}/2}\). Conversely, pointwise convergence of characteristic functions to a discontinuous function (e.g.\ \(\mathcal{N}(0,n)\)) yields no tight limit, which is the caveat already recorded after \cref{thm:levy-continuity}.
\end{remark}

\textbf{Step 4 (logarithm).}
From the bounds already obtained,
\[
\max_{j\le n}|a_{n,j}|
=
O\Bigl(\frac{n^{2}}{b_{n}^{2}}\Bigr)
=
O(n^{-1}),
\qquad
\max_{j\le n}|c_{n,j}|
\le
\max_{j\le n}\min\Bigl(\frac{2}{3j^{2}},\frac{t^{2}j^{2}}{6b_{n}^{2}}\Bigr).
\]
The minimum is largest when \(2/(3j^{2})\asymp t^{2}j^{2}/b_{n}^{2}\), i.e.\ at \(j\asymp n^{3/4}\), and that value is \(O(n^{-3/2})\). Thus \(\max_{j}|a_{n,j}+c_{n,j}|\to 0\). For large \(n\) every factor lies in a disk about \(1\) on which the principal logarithm is analytic and \(\log(1+z)=z+O(z^{2})\) uniformly, so

\[
\log\varphi_{S_{n}/b_{n}}(t)
=
\sum_{j=1}^{n}(a_{n,j}+c_{n,j})
+
O\Biggl(\sum_{j=1}^{n}|a_{n,j}+c_{n,j}|^{2}\Biggr).
\]

The quadratic remainder vanishes:

\[
\sum_{j=1}^{n}a_{n,j}^{2}
=
O\Bigl(\frac{n^{5}}{b_{n}^{4}}\Bigr)
=
O(n^{-1}),
\qquad
\sum_{j=1}^{n}c_{n,j}^{2}
\le
\bigl(\max_{j}|c_{n,j}|\bigr)\sum_{j}|c_{n,j}|
\to 0.
\]

Combining Steps 2--4,

\[
\log\varphi_{S_{n}/b_{n}}(t)
=
-\frac{t^{2}}{2}+o(1),
\qquad
\varphi_{S_{n}/b_{n}}(t)
\to
\exp(-t^{2}/2).
\]

\Cref{thm:levy-continuity} yields \(S_{n}/b_{n}\xrightarrow{(d)}\mathcal{N}(0,1)\).
\end{proof}

\subsubsection{Manuscript proof idea}

The calculation below follows the handwritten notes, in the same order and with the same expansions. \mserr{Red} marks a false claim or a missing justification.

Independence gives

\[
f_{S_{n}/b_{n}}(t)
=
\E\bigl[\exp(it S_{n}/b_{n})\bigr]
=
\E\Biggl[\exp\Biggl(it\sum_{j=1}^{n}X_{j}/b_{n}\Biggr)\Biggr]
=
\prod_{j=1}^{n}
f_{X_{j}}(t/b_{n}).
\]

This step is correct. Expanding the five-point law,

\[
\begin{aligned}
f_{X_{j}}(t)
&=
e^{itj^{2}}\cdot\frac{1}{6j^{2}}
+
e^{-itj^{2}}\cdot\frac{1}{6j^{2}}
+
\frac16\bigl(e^{itj}+e^{-itj}\bigr)
+
1\cdot\Bigl(\frac23-\frac{1}{3j^{2}}\Bigr)
\\
&=
\frac{1}{3j^{2}}\cos(tj^{2})
+
\frac13\cos(tj)
+
\frac23
-
\frac{1}{3j^{2}}.
\end{aligned}
\]

This identity is correct. Substituting \(t\mapsto t/b_{n}\),

\[
f_{S_{n}/b_{n}}(t)
=
\prod_{j=1}^{n}
\Biggl(
\frac23
+
\frac13\cos\Bigl(\frac{tj}{b_{n}}\Bigr)
+
\frac{1}{3j^{2}}\cos\Bigl(\frac{tj^{2}}{b_{n}}\Bigr)
-
\frac{1}{3j^{2}}
\Biggr).
\]

The notes then write: only when

\[
\mserr{\frac{j^{2}}{b_{n}}\to 0 \qquad \forall j}
\]

does the display above tend to a Gaussian characteristic function. \mserr{False: if \(b_{n}\) is of order \(n^{3/2}\), the worst index \(j=n\) gives \(n^{2}/b_{n}\sim 3\sqrt{n}\not\to 0\). The cosine Taylor expansion is therefore not available uniformly in \(j\). It is valid for the ordinary atoms \(\pm j\), since \(j/b_{n}\le n/b_{n}\to 0\), but not for the rare atoms \(\pm j^{2}\).}

\errpanel{%
\centering
\textbf{\mserr{The error, in one box.}}
\smallskip

\noindent
\mserr{Do not Taylor-expand} \(\displaystyle\cos(tj^{2}/b_{n})\). For \(j\sim n\) and \(b_{n}\asymp n^{3/2}\) one has \(tj^{2}/b_{n}\asymp\sqrt{n}\), so the cosine is oscillating in \([-1,1]\). Replacing it by \(1-\frac12(tj^{2}/b_{n})^{2}\) replaces a bounded oscillation by a parabola of size \(\asymp n\).

\smallskip
\noindent
\textbf{That does not make the characteristic-function term large.} The term is still
\[
c_{n,j}
=
\frac{1}{3j^{2}}\bigl(\cos(tj^{2}/b_{n})-1\bigr).
\]
Since \(\lvert\cos u-1\rvert\le 2\),
\[
\lvert c_{n,j}\rvert
\le
\frac{2}{3j^{2}}.
\]
For \(j\sim n\) this is \(O(n^{-2})\): a \emph{small} oscillation, amplitude \(1/j^{2}\). It is not an \(O(1)\) constant, and the oscillation does not survive as a random phase in the Lévy limit.

\smallskip
\noindent
\textbf{What to do instead.} Abandon the Taylor expansion of this cosine. Keep \(c_{n,j}\) as written and bound the sum. Split at \(J_{n}=n^{2/3}\):
\begin{itemize}
\item on \(\{j\le J_{n}\}\) one \emph{may} Taylor, but \(\sum_{j\le J_{n}}\lvert c_{n,j}\rvert=O(n^{-1})\);
\item on \(\{j>J_{n}\}\) use only \(\lvert c_{n,j}\rvert\le 2/(3j^{2})\), hence \(\sum_{j>J_{n}}\lvert c_{n,j}\rvert=O(n^{-2/3})\).
\end{itemize}
Therefore \(\sum_{j=1}^{n}c_{n,j}\to 0\): the whole \(\pm j^{2}\) family is an infinitesimal in \(\log\varphi_{S_{n}/b_{n}}(t)\). Continue the argument with the \(\pm j\) terms \(a_{n,j}\) only.

\smallskip
\noindent
The illegal Taylor is dangerous because it \emph{inflates} a small term. For \(j\sim n\) it pretends \(\lvert c_{n,j}\rvert\asymp t^{2}j^{2}/b_{n}^{2}\asymp n^{-1}\). There are \(\asymp n\) such indices, so the fake errors sum to \(O(1)\). The true terms sum to \(O(n^{-1})\). That spurious \(O(1)\) is exactly the extra \(-t^{2}/2\) which led to \(\tilde{b}_{n}\) instead of \(b_{n}\).
}

\begin{warning}[If the \(\pm j^{2}\) Taylor cannot be performed]
\label{warn:p2-no-taylor}
The cosine \(\cos(tj^{2}/b_{n})\) is oscillating, but after multiplication by the rare-atom weight \(1/(3j^{2})\) the contribution \(c_{n,j}\) is small, and \(\sum_{j}c_{n,j}=o(1)\). Oscillation is not a licence to treat the term as a constant, nor as a large remainder one may ignore or Taylor anyway; it is the reason the quadratic approximation is false, after which the bound \(\lvert\cos u-1\rvert\le 2\) is enough to kill the sum. See the box above and Step~3 of the correct proof.
\end{warning}

Using \(\cos x=1-x^{2}/2+\cdots\) on both cosines inside the logarithm, as in the notes,

\[
f_{S_{n}/b_{n}}(t)
=
\exp\Biggl[
\sum_{j=1}^{n}
\log\Biggl(
\frac23
+
\frac13\Biggl(1-\frac12\Bigl(\frac{tj}{b_{n}}\Bigr)^{2}+\cdots\Biggr)
+
\mserr{\frac{1}{3j^{2}}\Biggl(1-\frac12\Bigl(\frac{tj^{2}}{b_{n}}\Bigr)^{2}+\cdots\Biggr)}
-
\frac{1}{3j^{2}}
\Biggr)
\Biggr].
\]

\mserr{The red cosine expansion is illegitimate: it requires \(tj^{2}/b_{n}\to 0\).} The constants cancel and one obtains

\[
\exp\Biggl[
\sum_{j=1}^{n}
\log\Biggl(
1
-
\frac{t^{2}j^{2}}{6b_{n}^{2}}
-
\mserr{\frac{t^{2}j^{2}}{6b_{n}^{2}}}
+
\mserr{O\Bigl(\frac{n^{4}}{b_{n}^{4}}\Bigr)}
\Biggr)
\Biggr].
\]

The second quadratic term is the \(\pm j^{2}\) contribution; keeping it double-counts the Gaussian variance. \mserr{Gap: the remainder was written as \(O(n^{4}/b_{n}^{4})\), independently of \(j\). The cosine remainder for the \(\pm j\) atoms is \(O(j^{4}/b_{n}^{4})\); for the \(\pm j^{2}\) atoms it is not even of that order when \(j^{2}/b_{n}\) is large. Also, one must check that each factor stays near \(1\) so that the principal logarithm is defined: this needs \(\max_{j\le n}|f_{X_{j}}(t/b_{n})-1|\to 0\), which was not proved.}

Then \(\mserr{\log(1+x)\approx x}\) produces

\[
\exp\Biggl(
\sum_{j=1}^{n}
\Biggl(
-\frac{t^{2}j^{2}}{3b_{n}^{2}}
+
O\Bigl(\frac{j^{4}}{3b_{n}^{4}}\Bigr)
\Biggr)
\Biggr)
=
\exp\Biggl(
-\frac{t^{2}}{3b_{n}^{2}}\sum_{j=1}^{n}j^{2}
+
\sum_{j=1}^{n}O\Bigl(\frac{j^{4}}{b_{n}^{4}}\Bigr)
\Biggr).
\]

\mserr{Gap: \(\log(1+x)=x+O(x^{2})\) requires a uniform smallness of \(x\). The remainder cannot be discarded without summing it: \(\sum_{j=1}^{n}j^{4}/b_{n}^{4}=O(n^{5}/b_{n}^{4})\). For the (wrong) scale \(b_{n}^{2}\sim n^{3}/9\) this happens to be \(O(n^{-1})\), but that was not checked.}

The sum of squares \(\sum_{j=1}^{n}j^{2}=n(n+1)(2n+1)/6\) is correct. Matching \(\exp(-t^{2}/2)\) as in the notes,

\[
\frac{n(n+1)(2n+1)}{18b_{n}^{2}}
=
\frac12
\qquad\Longrightarrow\qquad
\mserr{b_{n}
=
\frac13\sqrt{n(n+1)(2n+1)}}.
\]

\mserr{False conclusion: this matching is only valid under the illegal expansion. The next subsection substitutes \(\tilde{b}_{n}\) into the legitimate expansion of the correct proof.}

\subsubsection{Why the manuscript matching \(\exp(-t^{2}/2)\) does not yield \(\mathcal{N}(0,1)\)}

The two coefficients \(\frac13\) and \(\frac16\) below are not interchangeable, and neither is the \(\frac13\) in \(b_{n}^{2}=v_{n}/3\). Track them separately.

From Step~2, the \(\pm j\) contribution is
\[
a_{n,j}
=
\frac13\Bigl(\cos\bigl(tj/\beta_{n}\bigr)-1\Bigr)
=
\frac13\Bigl(-\frac12\Bigl(\frac{tj}{\beta_{n}}\Bigr)^{2}+O\Bigl(\frac{j^{4}}{\beta_{n}^{4}}\Bigr)\Bigr)
=
-\frac{t^{2}j^{2}}{6\beta_{n}^{2}}
+
O\Bigl(\frac{j^{4}}{\beta_{n}^{4}}\Bigr).
\]
The \(\frac16\) is \(\frac13\cdot\frac12\): the atom probability \(\frac13\) times the cosine Taylor \(\frac12\). Summing over \(j\) gives the true leading term
\[
\sum_{j=1}^{n}a_{n,j}
=
-\frac{t^{2}v_{n}}{6\beta_{n}^{2}}
+
O\Bigl(\frac{n^{5}}{\beta_{n}^{4}}\Bigr).
\]
(The remainder is \(O(n^{-1})\) for either scale \(b_{n}\) or \(\tilde{b}_{n}\).) Equivalently,
\[
-\frac{t^{2}v_{n}}{6\beta_{n}^{2}}
=
-\frac{t^{2}}{2}\cdot\frac{v_{n}/3}{\beta_{n}^{2}}
=
-\frac{t^{2}}{2}\cdot\frac{\text{truncated variance}}{\beta_{n}^{2}}.
\]
The identity \(b_{n}^{2}=v_{n}/3\) is this truncated variance, not a Taylor coefficient.

The manuscript instead expands \emph{both} cosines. The illegal \(\pm j^{2}\) Taylor produces another copy of the same quadratic,
\[
c_{n,j}
\;\stackrel{\mathrm{formal}}{\approx}\;
-\frac{t^{2}j^{2}}{6\beta_{n}^{2}},
\]
and adding the two copies replaces \(\frac16\) by \(\frac13\):
\[
\log f_{S_{n}/\beta_{n}}(t)
\;\stackrel{\mathrm{formal}}{=}\;
-\frac{t^{2}v_{n}}{3\beta_{n}^{2}}
+
\text{(unchecked remainders)}
=
-\frac{t^{2}}{2}\cdot\frac{2v_{n}/3}{\beta_{n}^{2}}
=
-\frac{t^{2}}{2}\cdot\frac{\mathrm{Var}(S_{n})}{\beta_{n}^{2}}.
\]
Choosing \(\beta_{n}=\tilde{b}_{n}\) (i.e.\ \(\beta_{n}^{2}=2v_{n}/3=\mathrm{Var}(S_{n})\)) makes this formal right-hand side equal to \(-t^{2}/2\). That is why the calculation appears to match the standard Gaussian characteristic function. The formal identity is false: \(\sum c_{n,j}\to 0\), so the true logarithm keeps only the \(\frac16\) term. For a general scale \(\beta_{n}\asymp n^{3/2}\),

\[
\log\varphi_{S_{n}/\beta_{n}}(t)
=
-\frac{t^{2}v_{n}}{6\beta_{n}^{2}}
+
o(1).
\]

Plug in the two scales. For the correct choice \(\beta_{n}=b_{n}\), one has \(v_{n}/b_{n}^{2}=3\), hence \(-t^{2}v_{n}/(6b_{n}^{2})=-t^{2}/2\). For the manuscript choice \(\beta_{n}=\tilde{b}_{n}\), one has \(v_{n}/\tilde{b}_{n}^{2}=3/2\), hence

\[
\log\varphi_{S_{n}/\tilde{b}_{n}}(t)
=
-\frac{t^{2}}{4}
+
o(1).
\]

Since the characteristic function of \(\mathcal{N}(0,\sigma^{2})\) is \(\exp(-\sigma^{2}t^{2}/2)\), the limit \(\exp(-t^{2}/4)\) is the law \(\mathcal{N}(0,1/2)\). Equivalently, the correct proof already gives \(S_{n}/b_{n}\xrightarrow{(d)}\mathcal{N}(0,1)\) and \(\tilde{b}_{n}=\sqrt{2}\,b_{n}\), so

\[
\frac{S_{n}}{\tilde{b}_{n}}
=
\frac{1}{\sqrt{2}}\cdot\frac{S_{n}}{b_{n}}
\xrightarrow{(d)}
\mathcal{N}\bigl(0,1/2\bigr).
\]

The factor \(\sqrt{2}\) is exactly the overcount: the manuscript treated \(\sum c_{n,j}\) as another copy of \(\sum a_{n,j}\). Size by size, that overcount is produced by the large indices, not by a single outrageous term. Compare, for \(\beta_{n}\asymp n^{3/2}\) and \(j\asymp n\),

\[
\begin{aligned}
\text{naive Taylor:}\quad
|c_{n,j}|
&\;\stackrel{\mathrm{formal}}{\approx}\;
\frac{t^{2}j^{2}}{6\beta_{n}^{2}}
\asymp
\frac{1}{n},
\\
\text{truth:}\quad
|c_{n,j}|
&\;\le\;
\frac{2}{3j^{2}}
\asymp
\frac{1}{n^{2}}.
\end{aligned}
\]

There are \(\asymp n\) such indices, so the naive errors add to order \(1\), while the true sum over \(j\asymp n\) is \(O(1/n)\). In other words, Taylor pretends \(|\cos(tj^{2}/\beta_{n})-1|\approx (tj^{2}/\beta_{n})^{2}/2\asymp n\) when \(j\sim n\), but a cosine is at most \(2\); multiplying the garbage size \(n\) by the rare-atom weight \(1/(3j^{2})\asymp n^{-2}\) still leaves \(O(1/n)\) per term, which sums to a spurious constant.

The same comparison shows that one cannot rescue the manuscript by taking \(\beta_{n}\) large enough that \(n^{2}/\beta_{n}\to 0\). That would force \(\beta_{n}\gg n^{2}\), hence \(v_{n}/\beta_{n}^{2}\to 0\), and \(S_{n}/\beta_{n}\to 0\) in probability. The Gaussian window \(\beta_{n}\asymp n^{3/2}\) and the uniform Taylor window \(\beta_{n}\gg n^{2}\) are incompatible.

\begin{remark}[Lindeberg]
\label{rem:problem-2-trap}
The same obstruction in the triangular array \(X_{j}/\tilde{b}_{n}\) is Lindeberg's condition, cf.\ \citep[Theorem~3.4.5]{durrett2010}: for \(\varepsilon>0\),

\[
\frac{1}{\tilde{b}_{n}^{2}}
\sum_{j=1}^{n}
\E\bigl[X_{j}^{2}\mathbf{1}_{\{|X_{j}|>\varepsilon\tilde{b}_{n}\}}\bigr]
\]
does not tend to \(0\), because the events \(\{|X_{j}|=j^{2}\}\) with \(j^{2}\gtrsim \tilde{b}_{n}\) retain a positive fraction of \(\mathrm{Var}(S_{n})\) in the limit. The characteristic-function split \(a_{n,j}+c_{n,j}\) is the Fourier counterpart of truncating those jumps.
\end{remark}

\subsection{Problem 3. [10 pts.]}
We have 100 noodles in a bowl. Each noodle has two free ends, so initially there are 200 free ends in total. We are blindfolded and repeatedly perform the following operation: choose uniformly at random two free ends among all available free ends and connect them together. We continue until there are no free ends left. At the end of this procedure, the noodles form a random collection of disjoint loops. Compute the expectation of the number of loops.

\setcounter{section}{3}
\setcounter{definition}{0}

\subsubsection{Solution of Problem 3}

\paragraph{Answer.}

For \(n\) noodles the expected number of loops is

\[
E_{n}
=
\sum_{k=1}^{n}\frac{1}{2k-1}
=
H_{2n}-\frac12 H_{n},
\]

where \(H_{m}=\sum_{j=1}^{m}1/j\). In particular

\[
E_{100}
=
\sum_{k=1}^{100}\frac{1}{2k-1}
=
H_{200}-\frac12 H_{100}.
\]

\begin{proof}[Correct proof]
Write \(n\) for the number of noodles (the exam is the case \(n=100\)). At every moment the uneaten pieces are finitely many \emph{strands}: each strand is a path of noodles with exactly two free ends. Initially there are \(n\) strands of length one. A step consists of choosing two free ends uniformly and joining them. There are two possibilities, and both decrease the number of strands by \(1\):
\begin{itemize}
\item
the two ends belong to the same strand: that strand closes into a loop and leaves the pool;
\item
the two ends belong to distinct strands: those two strands concatenate into one longer strand.
\end{itemize}

\begin{center}
\begin{tikzpicture}[
  scale=0.92,
  >=Stealth,
  end/.style={circle, fill=black, inner sep=1.35pt},
  strand/.style={line width=1.15pt, line cap=round, black!75},
  pair/.style={red!75!black, line width=1.2pt, densely dashed},
  lab/.style={font=\scriptsize, text=black!70}
]
  \begin{scope}[shift={(0,0)}]
    \node[lab, anchor=south] at (1.35,1.55) {$k=3$ strands, $6$ free ends};
    \draw[strand] (0,1.05) -- (2.7,1.05);
    \node[end] at (0,1.05) {};
    \node[end] at (2.7,1.05) {};
    \node[lab, left] at (-0.12,1.05) {$S_{1}$};
    \draw[strand] (0,0.45) -- (2.7,0.45);
    \node[end] at (0,0.45) {};
    \node[end] at (2.7,0.45) {};
    \node[lab, left] at (-0.12,0.45) {$S_{2}$};
    \draw[strand] (0,-0.15) -- (2.7,-0.15);
    \node[end] at (0,-0.15) {};
    \node[end] at (2.7,-0.15) {};
    \node[lab, left] at (-0.12,-0.15) {$S_{3}$};
  \end{scope}

  \draw[->, thick, gray!60!black] (3.15,0.45) -- (3.85,0.45);

  \begin{scope}[shift={(4.15,0.85)}]
    \node[lab, anchor=south, text=red!70!black] at (1.2,1.05) {same strand, \(\mathbb{P}=\frac{1}{2k-1}\)};
    \draw[strand] (0,0.35) -- (2.4,0.35);
    \node[end] at (0,0.35) {};
    \node[end] at (2.4,0.35) {};
    \draw[pair] (0,0.35) to[out=110,in=70] (2.4,0.35);
    \node[lab, below] at (1.2,-0.05) {one loop; \(k-1\) strands left};
    \draw[strand, black!40] (0,-0.55) -- (2.4,-0.55);
    \node[end, fill=black!40] at (0,-0.55) {};
    \node[end, fill=black!40] at (2.4,-0.55) {};
    \draw[strand, black!40] (0,-1.05) -- (2.4,-1.05);
    \node[end, fill=black!40] at (0,-1.05) {};
    \node[end, fill=black!40] at (2.4,-1.05) {};
  \end{scope}

  \begin{scope}[shift={(8.05,0.85)}]
    \node[lab, anchor=south] at (1.55,1.05) {distinct strands, \(\mathbb{P}=\frac{2k-2}{2k-1}\)};
    \draw[strand] (0,0.35) -- (1.35,0.35);
    \node[end] at (0,0.35) {};
    \node[end] at (1.35,0.35) {};
    \draw[strand] (1.85,0.35) -- (3.2,0.35);
    \node[end] at (1.85,0.35) {};
    \node[end] at (3.2,0.35) {};
    \draw[pair] (1.35,0.35) -- (1.85,0.35);
    \node[lab, below] at (1.6,-0.05) {no loop; still \(k-1\) strands};
    \draw[strand, black!40] (0,-0.55) -- (3.2,-0.55);
    \node[end, fill=black!40] at (0,-0.55) {};
    \node[end, fill=black!40] at (3.2,-0.55) {};
  \end{scope}
\end{tikzpicture}
\end{center}

Both branches leave a uniformly random pairing on the remaining \(2(k-1)\) ends, i.e.\ a copy of the original problem of size \(k-1\). The process therefore visits the states ``\(k\) remaining strands'' for \(k=n,n-1,\dots,1\), and then stops.

\begin{center}
\begin{tikzpicture}[
  >=Stealth,
  box/.style={draw=gray!60, rounded corners=2pt, align=center, inner sep=5pt, font=\small, fill=white},
  pbox/.style={font=\scriptsize, text=red!70!black, align=center}
]
  \node[box] (En) {\(k\) strands};
  \node[box, fill=red!8] (loop) at (4.3,1.15) {loop \(+\) \(k-1\) strands};
  \node[box] (noloop) at (4.3,-1.15) {\(k-1\) strands, no new loop};
  \draw[->] (En.east) to[out=25,in=180] node[pbox, above, pos=0.45] {\(\mathbb{P}(I_{k}=1)=\dfrac{1}{2k-1}\)} (loop.west);
  \draw[->] (En.east) to[out=-25,in=180] node[pbox, below, pos=0.45] {\(\mathbb{P}(I_{k}=0)=\dfrac{2k-2}{2k-1}\)} (noloop.west);
\end{tikzpicture}
\end{center}

Let \(I_{k}\) be the indicator that a loop is formed at the step with \(k\) strands remaining. Every loop is created at exactly one such step, so the total number of loops is \(L_{n}=\sum_{k=1}^{n}I_{k}\). Among the \(\binom{2k}{2}=k(2k-1)\) pairs of free ends, exactly \(k\) pairs close a strand. Hence

\[
\Pbb(I_{k}=1)
=
\frac{k}{k(2k-1)}
=
\frac{1}{2k-1}.
\]

Linearity gives
\[
E_{n}
=
\E[L_{n}]
=
\sum_{k=1}^{n}\frac{1}{2k-1}.
\]
The same identity is the recurrence \(E_{n}=E_{n-1}+1/(2n-1)\) with \(E_{1}=1\): the first step contributes a loop with probability \(1/(2n-1)\), and in either case one is left with \(n-1\) strands whose remaining pairing is a copy of the original problem of size \(n-1\).

The harmonic evaluation is
\[
H_{2n}
=
\sum_{k=1}^{n}\frac{1}{2k-1}
+
\sum_{k=1}^{n}\frac{1}{2k}
=
E_{n}
+
\frac12 H_{n}.
\]

\begin{remark}[Asymptotics, after Qingshu]
\label{rem:p3-qingshu-mvt}
The exam asks for the exact expectation \(E_{n}=H_{2n}-\tfrac12 H_{n}\). The size of the sum is the harmonic asymptotic from Qingshu's 11th contest notes \citep{qingshu11}. The package taught there is: existence of Euler's constant (Example~7.5), the next term by Stolz (Example~7.10), and the odd-harmonic evaluation (Example~8.24). Two sibling expansions are recorded below: Stirling for \(n!\), and Euler--Maclaurin (Qingshu's EM formula) for the same \(1/(2n)\) correction that Stolz produces. Lagrange's mean value theorem recovers only the leading logarithm.

\emph{Existence of \(\gamma\).} Write \(H_{n}=\sum_{k=1}^{n}1/k\). Taylor's formula with Peano remainder (Qingshu, Theorem~7.1) gives \(\ln(1-x)=-x+O(x^{2})\) as \(x\to 0\), hence \(x+\ln(1-x)=O(x^{2})\). Telescoping \(\ln n=\sum_{k=2}^{n}\ln(k/(k-1))\) produces

\[
H_{n}-\ln n
=
1+\sum_{k=2}^{n}\Bigl(\frac{1}{k}-\ln\frac{k}{k-1}\Bigr)
=
1+\sum_{k=2}^{n}\Bigl(\frac{1}{k}+\ln\bigl(1-\tfrac{1}{k}\bigr)\Bigr).
\]

For \(k\ge 2\) one has \(1/k\in(0,1/2]\), so the general term is \(O(1/k^{2})\). The comparison \(\sum_{k=2}^{n}k^{-2}\le\int_{1}^{n}x^{-2}\,\mathrm{d}x<\infty\) (the \(p=2\) case of Qingshu's Lemma~7.1 / Proposition~7.4) therefore yields absolute convergence of the series of remainders, and

\[
\gamma
\coloneqq
\lim_{n\to\infty}(H_{n}-\ln n)
\]

exists in \(\mathbb{R}\). Equivalently, the complete leading statement is

\[
H_{n}
=
\ln n+\gamma+o(1).
\]

\emph{The next term.} Example~7.10 upgrades the remainder by Stolz. Since \(H_{n}-\ln n-\gamma\to 0\),

\[
\lim_{n\to\infty}n(H_{n}-\ln n-\gamma)
=
\lim_{n\to\infty}
\frac{H_{n}-\ln n-\gamma}{1/n}
=
\lim_{n\to\infty}
\frac{\frac{1}{n+1}-\ln(n+1)+\ln n}{\frac{1}{n+1}-\frac{1}{n}}.
\]

The numerator is \(\frac{1}{n+1}-\ln(1+1/n)\). Expanding \(\frac{1}{n+1}=\frac{1}{n}-\frac{1}{n^{2}}+O(n^{-3})\) and \(\ln(1+1/n)=\frac{1}{n}-\frac{1}{2n^{2}}+O(n^{-3})\) gives numerator \(-\frac{1}{2n^{2}}+O(n^{-3})\) and denominator \(-\frac{1}{n(n+1)}=-\frac{1}{n^{2}}+O(n^{-3})\), so the limit equals \(\frac12\). Thus

\[
H_{n}
=
\ln n+\gamma+\frac{1}{2n}+o\bigl(n^{-1}\bigr).
\]

\emph{Euler--Maclaurin (Qingshu's EM).} The same \(1/(2n)\) is the first endpoint correction in the Euler--Maclaurin formula. For \(f\in C^{2}[1,\infty)\),
\[
\sum_{k=1}^{n}f(k)
=
\int_{1}^{n}f(x)\,\mathrm{d}x
+
\frac{f(1)+f(n)}{2}
+
\frac{B_{2}}{2!}\bigl(f'(n)-f'(1)\bigr)
+
O\bigl(\|f''\|_{[n,\infty)}\bigr),
\]
with \(B_{2}=1/6\). (Higher even Bernoulli numbers continue the expansion.) Taking \(f(x)=1/x\) gives \(\int_{1}^{n}x^{-1}\,\mathrm{d}x=\ln n\) and \(f'(x)=-x^{-2}\), hence
\[
H_{n}
=
\ln n
+
C
+
\frac{1}{2n}
-
\frac{1}{12n^{2}}
+
O(n^{-4}),
\]
where the constant \(C\) is identified with \(\gamma\) by Example~7.5. Truncating after the trapezoid term already yields Qingshu's Stolz conclusion \(H_{n}=\ln n+\gamma+1/(2n)+o(n^{-1})\); the \(B_{2}\) term is optional on the exam.

\emph{Stirling.} Euler--Maclaurin applied to \(f(x)=\ln x\) is Stirling's formula. The form one narrates is
\[
n!
\sim
\sqrt{2\pi n}\,\Bigl(\frac{n}{e}\Bigr)^{n},
\qquad\text{i.e.}\qquad
\lim_{n\to\infty}
\frac{n!}{e^{-n}n^{n+1/2}}
=
\sqrt{2\pi},
\]
or the two-sided bound \(e(n/e)^{n}\le n!\le e\,n\,(n/e)^{n}\). Equivalently \(\ln n!=n\ln n-n+\tfrac12\ln(2\pi n)+o(1)\). This expands \(\sum_{k=1}^{n}\ln k\), not \(H_{n}=\sum 1/k\); it is recorded here because it is the same EM machine on a different summand, and because it is the standard factorial asymptotic in the QE corpus (e.g.\ the 2024 Fall hint). It is not needed to evaluate \(E_{n}\).

\emph{Odd harmonics, and \(E_{n}\).} Splitting even and odd terms as in the proof, or equivalently Example~8.24,

\[
\sum_{k=1}^{n}\frac{2}{2k-1}
=
2H_{2n}-H_{n}
=
2\bigl(\ln(2n)+\gamma+o(1)\bigr)-\bigl(\ln n+\gamma+o(1)\bigr)
=
\ln n+2\ln 2+\gamma+o(1).
\]

Dividing by \(2\) is the expansion of the exam sum:

\[
E_{n}
=
H_{2n}-\tfrac12 H_{n}
=
\tfrac12\ln n+\ln 2+\tfrac12\gamma+o(1).
\]

The \(\frac{1}{2n}\) terms cancel: \(\frac{1}{4n}-\frac12\cdot\frac{1}{2n}=0\), so the same display holds with \(o(n^{-1})\) if one uses the refined expansion.

\emph{What MVT alone gives.} The slogan ``Lagrange's mean value theorem preserves the order'' (Chapter~7, Example~7.17) only recovers \(E_{n}\sim\tfrac12\ln n\). On each interval \([2k-1,2k+1]\),

\[
\ln(2k+1)-\ln(2k-1)=\frac{2}{\theta_{k}},
\qquad
\theta_{k}\in(2k-1,2k+1).
\]

Since \(\theta_{k}>2k-1\), one has \(\frac{2}{\theta_{k}}<\frac{2}{2k-1}\). Summing the left-hand side telescopes to \(\ln(2n+1)\), hence \(\frac12\ln(2n+1)<E_{n}\). Shifting the interval to \([2k-3,2k-1]\) for \(k\ge 2\) likewise gives \(E_{n}<1+\frac12\ln(2n-1)\). The two bounds yield \(E_{n}\sim\tfrac12\ln n\); the constant \(\gamma\) is not produced by that comparison.
\end{remark}
\end{proof}

\subsubsection{Manuscript proof idea}

The calculation below follows the handwritten notes. \mserr{Red} marks a false claim or a missing justification. The recurrence and the odd-harmonic sum are correct; the unfinished integral does not compute \(E_{n}\).

Let \(E_{n}\) be the expected number of loops from \(n\) noodles, and condition on the partner of a given free end (equivalently, on the first uniformly chosen pair). The notes write

\[
E_{n}
=
\frac{1}{2n-1}\bigl(1+E_{n-1}\bigr)
+
\frac{2n-2}{2n-1}E_{n-1}
=
E_{n-1}+\frac{1}{2n-1},
\]

with \(E_{1}=1\), hence \(E_{n}=\sum_{k=1}^{n}1/(2k-1)\) and \(E_{100}=\sum_{k=1}^{100}1/(2k-1)\). The algebra is correct, and this is the answer.

\mserr{Gap: after two distinct noodles are joined, the system is \(n-1\) \emph{strands}, not \(n-1\) original noodles. The remaining \(2(n-1)\) ends are still uniformly paired, so the conditional law equals the original problem of size \(n-1\). That identification should be stated; without it the notation \(E_{n-1}\) on the merged state is not justified. The same gap appears if one starts from a distinguished noodle rather than from a distinguished end: the probability that a \emph{specified} noodle is the first pair is \(1/\binom{2n}{2}\), not \(1/(2n-1)\). The factor \(1/(2n-1)\) is the probability that the partner of a given end is the other end of the same strand, or equivalently that a uniformly chosen pair of ends closes some strand.}

The notes then introduce
\[
f(t)
=
\sum_{i=1}^{2n}\frac{t^{i}}{i},
\qquad
f'(t)
=
\frac{1-t^{2n}}{1-t}
\quad(t\neq 1),
\]
and
\[
f(t)-f(-t)
=
\sum_{\substack{1\le i\le 2n\\ i\text{ odd}}}\frac{2t^{i}}{i}.
\]
Evaluating at \(t=1\) (i.e.\ \(f(1)-f(-1)\)) does give \(2E_{n}\). This identity is correct. \mserr{Gap: \(f(1)=H_{2n}\) and \(f(-1)=\sum_{i=1}^{2n}(-1)^{i}/i\) should be named; the representation \(E_{n}=H_{2n}-\frac12 H_{n}\) is then immediate and is the natural closed form.}

The notes next write
\[
f(1)
=
\int_{0}^{1}\frac{1-s^{2n}}{1-s}\,\mathrm{d}s
\]
and substitute \(s=\sin^{2}\theta\). \mserr{This integral equals \(H_{2n}\), not \(E_{n}\) and not \(2E_{n}\). The substitution, even if completed, computes the wrong object. The notes abandon the calculation (it is crossed out); it is not needed once the odd-harmonic sum is known.}

\subsection{Problem 4. [10 pts.]}
Let \( \{B_t\}_{t \ge 0} \) be a one-dimensional standard Brownian motion started from the origin. Find and prove the limit (in distribution) of

\[
\left[ \int_0^{T^2} \exp(B_t)\, dt \right]^{1/T} \qquad \text{as } T \to \infty.
\]

Write your answer as a function of the standard normal distribution.

\setcounter{section}{4}
\setcounter{definition}{0}

\subsubsection{Definitions and warmup calculations for Problem 4}

The notes below do not solve Problem~4. They record the Wiener-process calculus from \citep{brzezniak1999} and \citep[Chapter~8]{durrett2010}. The two-time Gaussian and independence rules used in the examples are collected in \cref{prop:bm-calculus}; the examples then only specialise those identities.

\begin{definition}[Standard Brownian motion / Wiener process]
\label{def:brownian-motion}
A real process \(\{W(t)\}_{t\ge 0}\) is a \emph{Wiener process} (one-dimensional standard Brownian motion) if
\begin{enumerate}[label=\textup{(\roman*)}]
    \item
    \(W(0)=0\) almost surely;
    \item
    \(t\mapsto W(t)\) is almost surely continuous;
    \item
    for \(0<t_1<\cdots<t_n\) and Borel sets \(A_1,\dots,A_n\subset\mathbb{R}\),

\[
\begin{aligned}
&\Pbb\bigl(W(t_1)\in A_1,\dots,W(t_n)\in A_n\bigr)\\
&\qquad=
\int_{A_1}\cdots\int_{A_n}
p(t_1,0,x_1)\,p(t_2-t_1,x_1,x_2)\cdots p(t_n-t_{n-1},x_{n-1},x_n)\,dx_1\cdots dx_n,
\end{aligned}
\]

where the Gaussian transition density is

\[
p(t,x,y)
=
\frac{1}{\sqrt{2\pi t}}
\exp\Bigl(-\frac{(x-y)^2}{2t}\Bigr),
\qquad
t>0,\ x,y\in\mathbb{R}.
\]
\end{enumerate}
Equivalently: \(W\) has stationary independent increments and \(W(t)-W(s)\sim\mathcal{N}(0,t-s)\) for \(0\le s<t\). In particular \(W(t)\sim\mathcal{N}(0,t)\). See \citep[Definition~6.9, Proposition~6.2, Theorem~6.3]{brzezniak1999} and \citep[Chapter~8]{durrett2010}.
\end{definition}

\begin{remark}[\(W\) versus \(B\)]
\label{rem:W-vs-B}
There is only one process here. The letter \(W\) is the SUMS-book convention (Wiener process, \citep[Definition~6.9]{brzezniak1999}); the letter \(B\) is the exam's convention (Brownian motion, also Durrett's). Both names refer to the same law: a continuous centered Gaussian process with covariance \(\min\{s,t\}\), started at \(0\). In these notes we identify them by

\[
B_t
\coloneqq
W(t),
\qquad
t\ge 0.
\]

So \(\{B_t\}_{t\ge 0}\) in Problem~4 is exactly the process \(\{W(t)\}_{t\ge 0}\) of \cref{def:brownian-motion}. The two spellings \(W(t)\) and \(B_t\) are likewise interchangeable; nothing in the arguments depends on which letter is used. After this remark the notes follow the exam and write \(B_t\).
\end{remark}

\begin{proposition}[Brownian two-time calculus]
\label{prop:bm-calculus}
Let \(\{B_t\}_{t\ge 0}\) be standard Brownian motion as in \cref{def:brownian-motion}, and write \(\mathcal{F}_t\coloneqq\sigma(B_u:0\le u\le t)\). The following hold for all \(s,t\ge 0\) and \(a,b,\lambda\in\mathbb{R}\).
\begin{enumerate}[label=\textup{(\roman*)}, ref=\textup{(\roman*)}]
    \item
    \label{itm:bm-gaussian}
    \emph{Gaussian process.} For \(0\le t_1<\cdots<t_n\), the vector \((B_{t_1},\dots,B_{t_n})\) is centered multivariate Gaussian. In particular \(B_t\sim\mathcal{N}(0,t)\).
    \item
    \label{itm:bm-indep}
    \emph{Increments versus values.} If \(s<t\), then \(B_t-B_s\sim\mathcal{N}(0,t-s)\) and \(B_t-B_s\perp\mathcal{F}_s\). In particular \(B_s\perp(B_t-B_s)\), but \(B_t\not\perp(B_t-B_s)\). By contrast the \emph{values} \(B_s\) and \(B_t\) are \emph{not} independent whenever \(s,t>0\): they share the common past up to \(\min\{s,t\}\).
    \item
    \label{itm:bm-cov}
    \emph{Covariance.} \(\E[B_s B_t]=\min\{s,t\}\), and \(\E[(B_t-B_s)^2]=|t-s|\). Equivalently, if \(s\le t\) then \((B_s,B_t)\) is centered Gaussian with covariance \(\bigl(\begin{smallmatrix}s&s\\s&t\end{smallmatrix}\bigr)\), so \(\Corr(B_s,B_t)=\sqrt{s/t}\).
    \item
    \label{itm:bm-lincomb}
    \emph{Linear combinations.} \(a B_s+b B_t\sim\mathcal{N}\bigl(0,\,a^2 s+b^2 t+2ab\min\{s,t\}\bigr)\). In particular

\[
B_s+B_t
\sim
\mathcal{N}\bigl(0,\,s+t+2\min\{s,t\}\bigr).
\]

If \(s\le t\) the variance is \(t+3s\).
    \item
    \label{itm:bm-mg}
    \emph{Martingale.} If \(s\le t\), then \(\E[B_t\mid\mathcal{F}_s]=B_s\).
    \item
    \label{itm:bm-exp}
    \emph{Exponential moments.} \(\E[e^{\lambda B_t}]=\exp(\lambda^2 t/2)\). Conditionally, if \(s\le t\) then \(\E[e^{\lambda B_t}\mid\mathcal{F}_s]=e^{\lambda B_s}\exp\bigl(\lambda^2(t-s)/2\bigr)\), so \(M_t^\lambda\coloneqq\exp(\lambda B_t-\lambda^2 t/2)\) is a martingale. In particular \(\E[e^{B_s+B_t}]=\exp\bigl(\tfrac12(s+t+2\min\{s,t\})\bigr)\).
    \item
    \label{itm:bm-sym}
    \emph{Symmetry.} \(\{-B_t\}_{t\ge 0}\) is again standard Brownian motion.
    \item
    \label{itm:bm-shift}
    \emph{Markov shift.} For each fixed \(T\ge 0\), the process \(V(u)\coloneqq B_{T+u}-B_T\), \(u\ge 0\), is a standard Brownian motion independent of \(\mathcal{F}_T\).
\end{enumerate}
\end{proposition}

\begin{proof}
Independent Gaussian increments from \cref{def:brownian-motion} are the only input; cf.\ \citep[Theorem~6.3, Exercises~6.19--6.27]{brzezniak1999} and \citep[Chapter~8]{durrett2010}.
\begin{enumerate}[label=\textup{(\roman*)}]
    \item
    Fix \(0=t_0<t_1<\cdots<t_n\). The increments \(\Delta_k\coloneqq B_{t_k}-B_{t_{k-1}}\) are independent \(\mathcal{N}(0,t_k-t_{k-1})\). Then \(B_{t_k}=\Delta_1+\cdots+\Delta_k\), a linear image of a Gaussian vector, hence jointly Gaussian and centered.
    \item
    The increment law and \(B_t-B_s\perp\mathcal{F}_s\) are the definition. Non-independence of the \emph{values}: if \(0<s\le t\), then \(\E[B_s B_t]=\min\{s,t\}=s>0\) by \ref{itm:bm-cov} below, while independence of centered random variables would force the covariance to vanish. Non-independence of the later coordinate with its own increment: \(B_t=B_s+(B_t-B_s)\) still contains \(B_t-B_s\), and \(\E[B_t(B_t-B_s)]=\E[(B_t-B_s)^2]=t-s\neq 0\).
    \item
    Assume \(s\le t\) (otherwise swap names). Split the later time, \(B_t=B_s+(B_t-B_s)\), and pull out the \(\mathcal{F}_s\)-measurable factor:

\[
\E[B_s B_t]
=
\E[B_s^2]
+
\E\bigl[B_s(B_t-B_s)\bigr]
=
s
+
\E[B_s]\,\E[B_t-B_s]
=
s,
\]

the cross term factoring because \(B_s\perp(B_t-B_s)\). The increment second moment is \(\E[(B_t-B_s)^2]=\Var(B_t-B_s)=t-s\). The covariance matrix of \((B_s,B_t)\) is then \(\bigl(\begin{smallmatrix}s&s\\s&t\end{smallmatrix}\bigr)\), and \(\Corr(B_s,B_t)=s/\sqrt{st}=\sqrt{s/t}\).
    \item
    Assume \(s\le t\). Rewrite as an independent-increment split:

\[
a B_s+b B_t
=
(a+b)B_s
+
b(B_t-B_s).
\]

The two summands are independent centered Gaussians of variances \((a+b)^2 s\) and \(b^2(t-s)\), so the sum is Gaussian of variance

\[
(a+b)^2 s+b^2(t-s)
=
a^2 s+2ab s+b^2 s+b^2 t-b^2 s
=
a^2 s+b^2 t+2ab s,
\]

and \(s=\min\{s,t\}\). Taking \(a=b=1\) gives \(B_s+B_t\sim\mathcal{N}(0,s+t+2\min\{s,t\})\); if \(s\le t\) this is \(t+3s\).
    \item
    If \(s\le t\), then \(\E[B_t\mid\mathcal{F}_s]=\E[B_s+(B_t-B_s)\mid\mathcal{F}_s]=B_s+\E[B_t-B_s]=B_s\), using that the increment is independent of \(\mathcal{F}_s\) and centered.
    \item
    The law \(B_t\sim\mathcal{N}(0,t)\) has moment generating function \(\exp(\lambda^2 t/2)\) (complete the square in the Gaussian integral, or analytically continue the characteristic function in \cref{note:chf-table}). Conditionally: \(\E[e^{\lambda B_t}\mid\mathcal{F}_s]=e^{\lambda B_s}\E[e^{\lambda(B_t-B_s)}]=e^{\lambda B_s}\exp(\lambda^2(t-s)/2)\). Multiplying by the deterministic factor \(e^{-\lambda^2 t/2}\) yields the martingale property of \(M^\lambda\). For two times, \(B_s+B_t\) is centered Gaussian by \ref{itm:bm-lincomb}, so \(\E[e^{B_s+B_t}]=\exp\bigl(\tfrac12\Var(B_s+B_t)\bigr)=\exp\bigl(\tfrac12(s+t+2\min\{s,t\})\bigr)\).
    \item
    The process \(-B\) starts at \(0\), has continuous paths, and has independent increments \((-B_t)-(-B_s)=-(B_t-B_s)\sim\mathcal{N}(0,t-s)\).
    \item
    This is \citep[Exercise~6.27]{brzezniak1999}. Continuity and \(V(0)=0\) are immediate. For \(0\le u<v\), the increment \(V(v)-V(u)=B_{T+v}-B_{T+u}\) is \(\mathcal{N}(0,v-u)\) and independent of \(\mathcal{F}_{T+u}\), hence of \(\sigma(V(r):r\le u)\). The whole process \(V\) is assembled from increments after time \(T\), so \(V\perp\mathcal{F}_T\).
\end{enumerate}
\end{proof}

\begin{remark}[What is independent of what]
\label{rem:bm-indep-slogan}
The usable slogan is: \emph{future increments} after time \(s\) are independent of \(\mathcal{F}_{s}\); the \emph{values} \(B_{s}\) and \(B_{t}\) (both strictly positive) are not. Scaling is \cref{prop:bm-scaling}; the strong Markov property (the same statement with the deterministic time \(s\) replaced by a stopping time) is \cref{thm:bm-strong-markov}; the reflection principle is \cref{thm:bm-reflection}. Time inversion (\(t\mapsto t B_{1/t}\), with value \(0\) at \(t=0\)) is again Brownian motion \citep[Chapter~8]{durrett2010}, but is not used below. Lévy's characterisation (\(B\) and \(B_{t}^{2}-t\) are continuous martingales iff \(B\) is Brownian motion) is \citep[Theorem~6.4]{brzezniak1999}.
\end{remark}

\begin{proposition}[Brownian scaling]
\label{prop:bm-scaling}
If \(\{B_t\}_{t\ge 0}\) is standard Brownian motion and \(c>0\), then

\[
V(t)
\coloneqq
\frac{1}{c}\,B_{c^{2}t},
\qquad
t\ge 0,
\]

is again standard Brownian motion. In particular, as processes on \([0,1]\),

\[
\bigl(T^{-1}B_{uT^{2}}\bigr)_{u\in[0,1]}
\stackrel{(d)}{=}
\bigl(B_u\bigr)_{u\in[0,1]}.
\]
\end{proposition}

\begin{proof}
This is \citep[Exercise~6.28]{brzezniak1999}. Continuity and \(V(0)=0\) are immediate. Independent Gaussian increments are preserved because

\[
V(t)-V(s)
=
c^{-1}\bigl(B_{c^{2}t}-B_{c^{2}s}\bigr)
\sim
\mathcal{N}\bigl(0,c^{-2}(c^{2}t-c^{2}s)\bigr)
=
\mathcal{N}(0,t-s).
\]
\end{proof}

\begin{example}[Mean, variance, covariance]
\label{ex:bm-moments}
From \cref{itm:bm-gaussian}, \(B_t\sim\mathcal{N}(0,t)\), so \(\E[B_t]=0\) and \(\Var(B_t)=t\). The covariance \(\E[B_s B_t]=\min\{s,t\}\) is \cref{itm:bm-cov}, which is \citep[Exercises~6.19--6.21]{brzezniak1999}. The only point that is easy to get backwards is \emph{which} coordinate one may factor, and that is \cref{itm:bm-indep,itm:bm-mg}. Take \(0\le s\le t\). Then \(B_s\) is \(\mathcal{F}_s\)-measurable, while \(B_t-B_s\perp\mathcal{F}_s\), so one splits the \emph{later} time and pulls out the \emph{earlier} coordinate:

\[
\E[B_s B_t]
=
\E\bigl[B_s\bigl(B_s+(B_t-B_s)\bigr)\bigr]
=
\E[B_s^2]
+
\E[B_s]\,\E[B_t-B_s]
=
s.
\]

Equivalently \(\E[B_t\mid\mathcal{F}_s]=B_s\), hence \(\E[B_s B_t]=\E[B_s\,\E[B_t\mid\mathcal{F}_s]]=\E[B_s^2]\). Factoring the later coordinate does not work: \(B_t\not\perp(B_t-B_s)\) by \cref{itm:bm-indep}, and \(\E[B_t(B_t-B_s)]=t-s\neq 0\).
\end{example}

\begin{example}[Exponential moments]
\label{ex:bm-exp-moment}
The one-time formula \(\E[e^{\lambda B_t}]=\exp(\lambda^2 t/2)\) is \cref{itm:bm-exp}; it is also the moment generating function of \(\mathcal{N}(0,t)\) in \cref{note:chf-table}, or the Gaussian integral after completing the square:

\[
\begin{aligned}
\E\bigl[e^{\lambda B_t}\bigr]
&=
\int_{\mathbb{R}}
e^{\lambda x}
\frac{1}{\sqrt{2\pi t}}
\exp\Bigl(-\frac{x^2}{2t}\Bigr)\,dx
=
e^{\lambda^2 t/2}
\int_{\mathbb{R}}
\frac{1}{\sqrt{2\pi t}}
\exp\Bigl(-\frac{(x-\lambda t)^2}{2t}\Bigr)\,dx
=
e^{\lambda^2 t/2},
\end{aligned}
\]

since the last integrand is the \(\mathcal{N}(\lambda t,t)\) density. This is \citep[Exercise~6.22]{brzezniak1999}. The two-time formula is the same computation applied to the Gaussian law in \cref{itm:bm-lincomb}: \(B_s+B_t\sim\mathcal{N}(0,s+t+2\min\{s,t\})\), hence by \cref{itm:bm-exp}

\[
\E\bigl[e^{B_s+B_t}\bigr]
=
\exp\bigl(\tfrac12(s+t+2\min\{s,t\})\bigr).
\]
\end{example}

\begin{example}[Fubini for \(\int e^{B_t}\,dt\)]
\label{ex:bm-fubini-exp}
The random variable in Problem~4 is a power of an integral of \(e^{B_t}\). Before any large-\(T\) limit, the \emph{expectation} of the integral is elementary. The map \((t,\omega)\mapsto e^{B_t(\omega)}\) is nonnegative, so Fubini--Tonelli applies with no integrability check:

\[
\E\Bigl[\int_0^T e^{B_t}\,dt\Bigr]
=
\int_0^T \E\bigl[e^{B_t}\bigr]\,dt
=
\int_0^T e^{t/2}\,dt
=
2\bigl(e^{T/2}-1\bigr),
\]

where the inner expectation is \cref{itm:bm-exp} at \(\lambda=1\) (cf.\ \cref{ex:bm-exp-moment}). The same device computes a variance if needed: \(\E[(\int_0^T e^{B_t}\,dt)^2]=\int_0^T\int_0^T \E[e^{B_s+B_t}]\,ds\,dt\), and the integrand is the two-time formula in \cref{itm:bm-exp}. (A close cousin is the 2025 fall QE geometric-Brownian-motion integral, with \(e^{B_t}\) replaced by \(e^{B_t-t/2}\).)

This computes an expectation. Problem~4 asks for a limit \emph{in distribution} of a nonlinear functional of the whole path, so one still has to pass from \(L^1\) information to the law of the path. Scaling \cref{prop:bm-scaling} is the usual way to put the \(T^2\) horizon onto the unit interval.
\end{example}

\begin{theorem}[\(L^{p}\to L^{\infty}\) sandwich; Laplace]
\label{thm:lp-to-linfty}
Let \(\mu\) be a finite measure on a space \(E\), and let \(f\colon E\to\mathbb{R}\) be measurable and bounded above, with \(M\coloneqq\mathrm{ess\,sup}\,f<\infty\). Write

\[
\|e^{f}\|_{p}
\coloneqq
\Bigl(\int_{E} e^{p f}\,\mathrm{d}\mu\Bigr)^{1/p},
\qquad
0<p<\infty.
\]

This is the \(L^{p}(\mu)\) norm of the function \(e^{f}\). For every \(\varepsilon>0\),

\[
e^{M-\varepsilon}\,
\mu\bigl(\{f>M-\varepsilon\}\bigr)^{1/p}
\le
\|e^{f}\|_{p}
\le
e^{M}\,
\mu(E)^{1/p}.
\]

If \(\mu(\{f>M-\varepsilon\})>0\) for every \(\varepsilon>0\) (automatic if \(f\) is continuous on a compact interval, hence attains its maximum), then \(\mu(E)^{1/p}\to 1\) and \(\mu(\{f>M-\varepsilon\})^{1/p}\to 1\) as \(p\to\infty\), so

\[
\lim_{p\to\infty}
\|e^{f}\|_{p}
=
e^{M}
=
\|e^{f}\|_{\infty}.
\]

Equivalently, in logarithmic form (the elementary compact case of Varadhan's lemma),

\[
\lim_{p\to\infty}
\frac1p\log\int_{E} e^{p f}\,\mathrm{d}\mu
=
M.
\]
\end{theorem}

\begin{proof}
The upper bound is \(e^{f}\le e^{M}\) pointwise \(\mu\)-a.e., so \(\int e^{pf}\,\mathrm{d}\mu\le e^{pM}\mu(E)\). For the lower bound, restrict the integral to \(\{f>M-\varepsilon\}\):

\[
\int_{E} e^{p f}\,\mathrm{d}\mu
\ge
e^{p(M-\varepsilon)}\,
\mu\bigl(\{f>M-\varepsilon\}\bigr).
\]

Raise both sides to the power \(1/p\). The limit statement follows because \(a^{1/p}\to 1\) for every \(a\in(0,\infty)\).
\end{proof}

The same comparison in finite dimensions is the reason the inequalities exist at all. For \(a_1,\dots,a_n\ge 0\) and \(p>0\),

\[
\max_{1\le i\le n}a_i
\le
\Bigl(\sum_{i=1}^{n} a_i^{p}\Bigr)^{1/p}
\le
n^{1/p}\max_{1\le i\le n}a_i.
\]

The left side is ``keep only the largest term''; the right side is ``at most \(n\) copies of the largest term''. Sending \(p\to\infty\) kills the combinatorial factor \(n^{1/p}\to 1\), so the \(\ell^{p}\) sum is asymptotically the max. The integral in \cref{thm:lp-to-linfty} is the continuum version of that sum, with \(a_i\) replaced by \(e^{f}\).

\begin{remark}[What this is not]
\label{rem:not-jensen}
Jensen's inequality for the exponential, \(\exp(\int f\,\mathrm{d}\nu)\le\int e^{f}\,\mathrm{d}\nu\) on a probability \(\nu\), exchanges \(\exp\) with an \emph{average}. That is a different comparison: it produces an \(L^{1}\) lower bound of order \(\exp(\text{mean of }f)\), which is typically much smaller than \(\exp(\sup f)\). The object in Problem~4 is \(\bigl(\int e^{Tf}\bigr)^{1/T}=\|e^{f}\|_{T}\), i.e.\ an \(L^{T}\) norm, not an exponential of a mean. The two operations (power outside the integral, versus exponential of the integral) do not commute, and the sandwich above is the one that matches the \(1/T\) power.
\end{remark}

\begin{example}[A deterministic check]
\label{ex:laplace-two-values}
On \([0,1]\) let \(f=\mathbf{1}_{[0,1/2]}\log 2\) (so \(M=\log 2\), \(e^{f}\) equals \(2\) on a set of measure \(1/2\) and equals \(1\) elsewhere). Then

\[
\Bigl(\int_{0}^{1} e^{p f}\,\mathrm{d}u\Bigr)^{1/p}
=
\bigl(\tfrac12\cdot 2^{p}+\tfrac12\cdot 1^{p}\bigr)^{1/p}
=
\bigl(2^{p-1}+\tfrac12\bigr)^{1/p},
\]

which is squeezed between \(2\cdot(1/2)^{1/p}\) and \(2\cdot 1^{1/p}\), and tends to \(2=e^{M}\). The mass of the maximising set survives only as the harmless factor \((1/2)^{1/p}\to 1\).
\end{example}

\begin{definition}[Stopping time]
\label{def:stopping-time}
A random time \(\tau:\Omega\to[0,\infty]\) is a \emph{stopping time} for \(\{\mathcal{F}_t\}_{t\ge 0}\) if \(\{\tau\le t\}\in\mathcal{F}_t\) for every \(t\ge 0\). The stopped \(\sigma\)-algebra is

\[
\mathcal{F}_{\tau}
\coloneqq
\bigl\{A\in\mathcal{F}_{\infty}: A\cap\{\tau\le t\}\in\mathcal{F}_{t}\text{ for all }t\ge 0\bigr\},
\]

where \(\mathcal{F}_{\infty}\coloneqq\sigma\bigl(\bigcup_{t\ge 0}\mathcal{F}_{t}\bigr)\): events in \(\mathcal{F}_{\tau}\) are those decided by time \(\tau\). First hitting times of closed sets are stopping times for continuous processes. In particular, if \(\tau_{a}\coloneqq\inf\{s\ge 0:B_{s}=a\}\) (with \(\inf\emptyset=\infty\)), then \(\{\tau_{a}\le t\}=\{M_{t}\ge a\}\in\mathcal{F}_{t}\), where \(M_{t}\coloneqq\sup_{0\le s\le t}B_{s}\).
\end{definition}

\begin{theorem}[Strong Markov property]
\label{thm:bm-strong-markov}
Let \(\tau\) be a stopping time for the natural filtration \(\mathcal{F}_{t}=\sigma(B_{u}:u\le t)\). On \(\{\tau<\infty\}\) define

\[
V^{\tau}(u)
\coloneqq
B_{\tau+u}-B_{\tau},
\qquad
u\ge 0.
\]

Then \(V^{\tau}\) is a standard Brownian motion, and \(V^{\tau}\) is independent of \(\mathcal{F}_{\tau}\). In particular, path continuity forces \(B_{\tau_{a}}=a\) almost surely on \(\{\tau_{a}<\infty\}\), so \(\bigl(B_{\tau_{a}+u}-a\bigr)_{u\ge 0}\) is a standard Brownian motion independent of \(\mathcal{F}_{\tau_{a}}\). This is the stopping-time upgrade of the deterministic Markov shift \cref{itm:bm-shift}; see \citep[Chapter~8]{durrett2010}.
\end{theorem}

\begin{proof}
The only input beyond \cref{itm:bm-shift} is path continuity (so discrete stopping times can be passed to the limit).
\begin{enumerate}[label=\textup{(\roman*)}]
    \item
    If \(\tau\equiv T\in[0,\infty)\) is deterministic, this is exactly \cref{itm:bm-shift}.
    \item
    Suppose \(\tau\) takes values in a finite set \(\{t_{1},\dots,t_{n}\}\subset[0,\infty)\). Fix times \(0\le u_{1}<\cdots<u_{m}\) and a bounded Borel \(g:\mathbb{R}^{m}\to\mathbb{R}\). For \(A\in\mathcal{F}_{\tau}\) the atom \(A\cap\{\tau=t_{k}\}\) lies in \(\mathcal{F}_{t_{k}}\), and on that atom \(V^{\tau}(u_{i})=B_{t_{k}+u_{i}}-B_{t_{k}}\). The deterministic shift \cref{itm:bm-shift} at time \(t_{k}\) therefore factors:

\[
\begin{aligned}
\E\bigl[\mathbf{1}_{A}\,g\bigl(V^{\tau}(u_{1}),\dots,V^{\tau}(u_{m})\bigr)\bigr]
&=
\sum_{k}
\E\bigl[\mathbf{1}_{A\cap\{\tau=t_{k}\}}\bigr]
\,
\E\bigl[g(B_{u_{1}},\dots,B_{u_{m}})\bigr]\\
&=
\Pbb(A)\,
\E\bigl[g(B_{u_{1}},\dots,B_{u_{m}})\bigr].
\end{aligned}
\]

Thus the finite-dimensional distributions of \(V^{\tau}\) are those of Brownian motion, and \(V^{\tau}\perp\mathcal{F}_{\tau}\). Continuity of paths is inherited from \(B\).
    \item
    For a general stopping time with \(\tau<\infty\) a.s., dyadically overshoot: on \(\{(k-1)2^{-n}<\tau\le k2^{-n}\}\) set \(\tau_{n}\coloneqq k2^{-n}\), and set \(\tau_{n}\coloneqq\infty\) on \(\{\tau=\infty\}\). Then \(\tau\le\tau_{n}\le\tau+2^{-n}\) on \(\{\tau<\infty\}\), and \(\{\tau_{n}\le k2^{-n}\}=\{\tau\le k2^{-n}\}\in\mathcal{F}_{k2^{-n}}\), so each \(\tau_{n}\) is a discrete stopping time. Path continuity gives \(V^{\tau_{n}}\to V^{\tau}\) uniformly on compact time intervals, hence in finite-dimensional distributions. For \(A\in\mathcal{F}_{\tau}\) one has \(A\in\mathcal{F}_{\tau_{n}}\) (because \(\tau\le\tau_{n}\)), so the identity in (ii) for \(\tau_{n}\) passes to the limit and yields the same identity for \(V^{\tau}\).
\end{enumerate}
If \(\Pbb(\tau=\infty)>0\), restrict the construction to \(\{\tau<\infty\}\); the process \(V^{\tau}\) is defined only there.
\end{proof}

\begin{theorem}[Reflection: the running maximum]
\label{thm:bm-reflection}
Let \(M_t\coloneqq\sup_{0\le s\le t}B_s\), and write \(M\coloneqq M_1\). For every \(a>0\) and \(t>0\),

\[
\Pbb(M_t\ge a)
=
2\Pbb(B_t\ge a).
\]

In particular, for \(x>0\),

\[
\Pbb(M\ge x)
=
2\Pbb(B_1\ge x)
=
\Pbb\bigl(|B_1|\ge x\bigr),
\]

the second equality by symmetry of \(B_1\sim\mathcal{N}(0,1)\). Thus \(M\stackrel{(d)}{=}|B_1|\): the running maximum is nonnegative, hence not Gaussian. See \citep[Chapter~8]{durrett2010}.
\end{theorem}

\begin{proof}
Fix \(a>0\) and \(t>0\). Let \(\tau_{a}\coloneqq\inf\{s\ge 0: B_{s}=a\}\), with \(\inf\emptyset=\infty\). Path continuity and \(B_{0}=0\) give \(\{\tau_{a}\le t\}=\{M_{t}\ge a\}\). The intermediate-value theorem also yields \(\{B_{t}\ge a\}\subset\{\tau_{a}\le t\}\), so \(\{B_{t}\ge a\}=\{\tau_{a}\le t,\, B_{t}\ge a\}\). Gaussian tails give \(\Pbb(B_{t}=a)=0\).

On \(\{\tau_{a}<\infty\}\) reflect the path after the hitting time:

\[
B^{*}_{s}
\coloneqq
\begin{cases}
B_{s} & \text{if }s\le\tau_{a},\\
2a-B_{s} & \text{if }s>\tau_{a}.
\end{cases}
\]

The time \(\tau_{a}\) is a stopping time by \cref{def:stopping-time}. Path continuity forces \(B_{\tau_{a}}=a\) on \(\{\tau_{a}<\infty\}\). \Cref{thm:bm-strong-markov} therefore says that the post-hitting process \(\bigl(B_{\tau_{a}+u}-a\bigr)_{u\ge 0}\) is a standard Brownian motion independent of \(\mathcal{F}_{\tau_{a}}\), hence equal in law to its negative by \cref{itm:bm-sym}. Reflecting the path after \(\tau_{a}\) is exactly replacing that Brownian motion by its negative, so \(B^{*}\) is again a standard Brownian motion. On \(\{\tau_{a}\le t\}\) one has \(B_{t}\ge a\) if and only if \(B^{*}_{t}\le a\). Equating the two laws of the endpoint,

\[
\Pbb(\tau_{a}\le t,\, B_{t}\ge a)
=
\Pbb(\tau_{a}\le t,\, B_{t}\le a).
\]

The left side is \(\Pbb(B_{t}\ge a)\). Adding the two events (their intersection is the null set \(\{B_{t}=a\}\)) recovers \(\{\tau_{a}\le t\}\), so

\[
\Pbb(M_{t}\ge a)
=
\Pbb(\tau_{a}\le t)
=
2\Pbb(B_{t}\ge a).
\]

Specialising to \(t=1\) and using symmetry of \(B_{1}\sim\mathcal{N}(0,1)\) gives, for \(x>0\),

\[
\Pbb(M\ge x)
=
2\Pbb(B_{1}\ge x)
=
\Pbb(B_{1}\ge x)+\Pbb(B_{1}\le -x)
=
\Pbb\bigl(|B_{1}|\ge x\bigr).
\]

Both \(M\) and \(|B_{1}|\) are supported on \([0,\infty)\) and have the same survival function, hence \(M\stackrel{(d)}{=}|B_{1}|\).
\end{proof}

\begin{remark}[How to use the toolkit on Problem~4]
\label{rem:p4-toolkit}
The three computations above are the Gaussian and Fubini steps. The remaining analytic ingredients, which are left for the solution, are: (i) Brownian scaling to rewrite the \(T^{2}\)-horizon integral in terms of a path on \([0,1]\); (ii) the elementary factor \(T^{2/T}\to 1\); (iii) identifying the large-\(T\) behaviour of \(\bigl(\int_{0}^{1} \exp(T\,B_{u})\,\mathrm{d}u\bigr)^{1/T}\) by the sandwich in \cref{thm:lp-to-linfty} (the integral is dominated by the path maximum, not by an \(L^{1}\) mean; cf.\ \cref{rem:not-jensen}); (iv) translating that maximum into a function of a standard normal via \cref{thm:bm-reflection}. The maximum is \(\stackrel{(d)}{=}|B_1|\), not \(\stackrel{(d)}{=}B_1\).
\end{remark}

\begin{remark}[Hints, in order]
\label{rem:p4-hints}
Do not start from \(\E\int e^{B_t}\,dt\). That is the wrong scale: the \(1/T\) power turns an \(L^{1}\) mean into an \(L^{T}\) norm, cf.\ \cref{rem:not-jensen,thm:lp-to-linfty}.
\begin{enumerate}[label=\textup{(\arabic*)}]
    \item
    Picture: \(e^{B_t}\) is largest where the path \(B\) is highest. For large \(T\), the integral \(\int_0^{T^{2}}e^{B_t}\,dt\) should feel like ``height of the highest peak, times a slowly varying width''.
    \item
    Change of time \(t=uT^{2}\). Write the integral as a constant (depending only on \(T\)) times \(\int_0^{1}\exp(B_{uT^{2}})\,\mathrm{d}u\). What does Brownian scaling \cref{prop:bm-scaling} do to the process \(u\mapsto B_{uT^{2}}\)?
    \item
    After scaling you have something of the shape \(c(T)\bigl(\int_0^{1}e^{T B_u}\,\mathrm{d}u\bigr)^{1/T}\). Compute \(\lim_{T\to\infty}c(T)\) by taking logarithms.
    \item
    The remaining factor is exactly \(\|e^{B}\|_{L^{T}([0,1])}\). Apply the sandwich of \cref{thm:lp-to-linfty} pathwise (Brownian paths are continuous on \([0,1]\), so the essential supremum is an ordinary maximum). What random variable appears?
    \item
    The exam asks for a function of a standard normal. Convert the running maximum on \([0,1]\) by \cref{thm:bm-reflection}.
\end{enumerate}
\end{remark}

\subsubsection{Solution of Problem 4}

\paragraph{Answer.}

Let \(G\sim\mathcal{N}(0,1)\). Then

\[
\Bigl(\int_0^{T^2} e^{B_t}\,\mathrm{d}t\Bigr)^{1/T}
\xrightarrow{(d)}
e^{|G|}
\qquad\text{as }T\to\infty.
\]

\begin{proof}[Exam proof]
Write

\[
I_T
\coloneqq
\Bigl(\int_0^{T^2} e^{B_t}\,\mathrm{d}t\Bigr)^{1/T},
\qquad
T>0.
\]

The integrand is continuous and strictly positive, so the integral is a well-defined element of \((0,\infty)\) almost surely, and the power \(1/T\) is unambiguous. The argument has four steps: a change of time putting the horizon onto \([0,1]\); Brownian scaling, which replaces \(B_{uT^2}\) by \(T\) times a standard Brownian motion on \([0,1]\); an elementary prefactor \(T^{2/T}\to 1\); and a pathwise Laplace sandwich identifying the remaining \(L^T\) norm with the exponential of the running maximum. The reflection principle then rewrites that maximum as \(|G|\).

\textbf{Step 1 (time change).}
Substitute \(t=uT^2\), so \(\mathrm{d}t=T^2\,\mathrm{d}u\) and \(u\) runs over \([0,1]\). Then

\[
\int_0^{T^2} e^{B_t}\,\mathrm{d}t
=
T^2\int_0^1 e^{B_{uT^2}}\,\mathrm{d}u,
\]

and therefore

\[
I_T
=
T^{2/T}
\Bigl(\int_0^1 e^{B_{uT^2}}\,\mathrm{d}u\Bigr)^{1/T}.
\]

\textbf{Step 2 (Brownian scaling).}
The process \(u\mapsto T^{-1}B_{uT^2}\) on \([0,1]\) depends on \(T\), so it cannot be the object of a pathwise \(T\to\infty\) limit. What scaling gives is an equality of laws, for each fixed \(T\). By \cref{prop:bm-scaling} with \(c=T\), as random elements of \(C[0,1]\),

\[
\bigl(T^{-1}B_{uT^2}\bigr)_{u\in[0,1]}
\stackrel{(d)}{=}
\bigl(B_u\bigr)_{u\in[0,1]}.
\]

The map \(\omega\mapsto\bigl(\int_0^1 e^{T\omega(u)}\,\mathrm{d}u\bigr)^{1/T}\) is continuous on \(C[0,1]\) (the exponential is continuous and \([0,1]\) is compact, so the integrand is bounded, and the integral is continuous for the uniform topology). Hence, writing

\[
J_T
\coloneqq
T^{2/T}
\Bigl(\int_0^1 e^{T B_u}\,\mathrm{d}u\Bigr)^{1/T}
\]

for the same functional of the original \(B\) restricted to \([0,1]\), one has \(I_T\stackrel{(d)}{=}J_T\) for every \(T>0\). (Any other standard Brownian motion on \([0,1]\) would do; only the law of the path matters.) We now prove that this \(T\)-independent process yields \(J_T\to e^{M}\) almost surely, where
\[
M
\coloneqq
\max_{u\in[0,1]} B_u.
\]
This is the time-\(1\) case of \(M_t\) in \cref{thm:bm-reflection}. Equality in law for each \(T\) then gives \(I_T\xrightarrow{(d)}e^{M}\). Step~5 identifies the law of this same \(M\); it is not Gaussian.

\textbf{Step 3 (the prefactor).}
For \(T>1\),

\[
\log\bigl(T^{2/T}\bigr)
=
\frac{2}{T}\log T
\to 0
\]

because \(\log T=o(T)\). Exponentiating, \(T^{2/T}\to 1\).

\textbf{Step 4 (Laplace sandwich, pathwise).}
The random variable \(M=\max_{u\in[0,1]}B_u\) was defined in Step~2. Fix a continuous path \(u\mapsto B_u\) on the compact interval \([0,1]\). The maximum is attained. For every \(\varepsilon>0\) the set \(\{u\in[0,1]: B_u>M-\varepsilon\}\) is relatively open in \([0,1]\) and nonempty, hence has strictly positive Lebesgue measure; write \(\lambda_{\varepsilon}\) for that measure, so \(\lambda_{\varepsilon}\in(0,1]\). Apply \cref{thm:lp-to-linfty} on \(([0,1],\mathrm{d}u)\) to the bounded continuous function \(B\), with \(p=T\):

\[
e^{M-\varepsilon}\,\lambda_{\varepsilon}^{1/T}
\le
\Bigl(\int_0^1 e^{T B_u}\,\mathrm{d}u\Bigr)^{1/T}
\le
e^{M}\cdot 1^{1/T}
=
e^{M}.
\]

As \(T\to\infty\) one has \(\lambda_{\varepsilon}^{1/T}\to 1\), so the liminf of the middle term is at least \(e^{M-\varepsilon}\). The upper bound is \(e^{M}\) for every \(T\). Since \(\varepsilon>0\) is arbitrary,

\[
\Bigl(\int_0^1 e^{T B_u}\,\mathrm{d}u\Bigr)^{1/T}
\to
e^{M}
=
\exp\Bigl(\max_{u\in[0,1]} B_u\Bigr).
\]

This last display is real analysis: it holds for every deterministic continuous function on \([0,1]\). It is not yet a statement about random variables. Write \(\Omega_c\) for the event that \(u\mapsto B_u\) is continuous on \([0,1]\). By \cref{def:brownian-motion}(ii) one has \(\Pbb(\Omega_c)=1\). Restricting to a single \(\omega\in\Omega_c\), the path \(B(\cdot,\omega)\) is just such a continuous function, so the sandwich applies at that \(\omega\): writing
\[
X_T(\omega)
\coloneqq
\Bigl(\int_0^1 e^{T B_u(\omega)}\,\mathrm{d}u\Bigr)^{1/T},
\]
the real sequence \(X_T(\omega)\) converges to \(e^{M(\omega)}\). Equivalently,
\[
\Pbb\bigl(X_T\to e^{M}\bigr)
=
\Pbb(\Omega_c)
=
1.
\]
That is the definition of almost-sure convergence of the random variables \(X_T\): pointwise convergence off a null set. This is \citep[Definition~5.4.1]{ren2024}, and is the language of \citep[\S2.3]{durrett2010}. It is not a weaker mode that still needs to be upgraded. Hogg's Chapter~5 is the statistical packaging of the same circle of ideas, but it never defines a.s.\ convergence: it works only with \(\xrightarrow{\Pbb}\) and \(\xrightarrow{(d)}\). The implications used below are
\[
\text{a.s.}
\implies
\text{in probability}
\implies
\text{in distribution},
\]
which are \citep[Proposition~5.4.3]{ren2024} (and \citep[\S2.3]{durrett2010}) and \citep[Proposition~9.2.1]{ren2024} respectively. The second arrow is also \citep[Theorem~5.2.1]{hogg2019}, and is the Portmanteau package already recorded as \cref{def:conv-d} from \citep[Theorems~3.2.3 and~3.2.5]{durrett2010}. The lattice around those two arrows is \cref{rem:conv-modes}: the exam target is the weakest mode, and \(I_T\stackrel{(d)}{=}J_T\) does not copy the almost-sure statement from \(J_T\) onto \(I_T\).

\begin{remark}[Modes of convergence]
\label{rem:conv-modes}
Four modes appear in this exam (Problems~2, 4, and~8). The first three compare random variables on a common space; in distribution compares laws only. Solid arrows in the diagram are unconditional; dashed arrows need an extra hypothesis. Sources: \citep[\S2.3, Theorems~3.2.3 and~3.2.5]{durrett2010}, \citep[\S5.1--5.2]{hogg2019}, \citep[Definitions~5.4.1--5.4.2, Propositions~5.4.3, 9.2.1, 9.2.2]{ren2024}.

{\footnotesize
\setlength{\tabcolsep}{3pt}
\renewcommand{\arraystretch}{1.22}
\noindent
\begin{tabularx}{\linewidth}{|
  >{\raggedright\arraybackslash}p{2.05cm}|
  >{\raggedright\arraybackslash}p{2.35cm}|
  >{\raggedright\arraybackslash}X|
  >{\raggedright\arraybackslash}p{2.15cm}|
  >{\raggedright\arraybackslash}X|
}
\hline
\textbf{Mode} & \textbf{Notation} & \textbf{Definition} & \textbf{Same \(\Omega\)?} & \textbf{Always implies; converse} \\
\hline
Almost sure
  & \(X_n\xrightarrow{\mathrm{a.s.}}X\)
  & \(\Pbb(X_n(\omega)\to X(\omega))=1\)
  & Yes: pathwise
  & \(\Rightarrow\xrightarrow{\Pbb}\). Converse fails (typewriter / moving bump). Here: \(X_T\to e^{M}\) on \(\Omega_c\). \\
\hline
In probability
  & \(X_n\xrightarrow{\Pbb}X\)
  & \(\Pbb(|X_n-X|\ge\varepsilon)\to 0\) for every \(\varepsilon>0\)
  & Yes
  & \(\Rightarrow\xrightarrow{(d)}\). Subsequence \(\Rightarrow\) a.s.\ \citep[\S2.3]{durrett2010}. Definition in \cref{def:conv-p}. \\
\hline
\(L^{p}\) (\(p\ge 1\))
  & \(X_n\xrightarrow{L^{p}}X\)
  & \(\E[|X_n-X|^{p}]\to 0\)
  & Yes
  & \(\Rightarrow\xrightarrow{\Pbb}\) by Markov. If \(q>p\) then \(L^{q}\Rightarrow L^{p}\) on a probability space. Converse: need uniform integrability (Vitali). Incomparable with a.s.\ in general. \\
\hline
In distribution
  & \(X_n\xrightarrow{(d)}X\)
  & \(F_n(x)\to F(x)\) at continuity points of \(F\); equivalently \(\E[g(X_n)]\to\E[g(X)]\) for bounded continuous \(g\)
  & No: laws only
  & Weakest of the four. Portmanteau: \cref{def:conv-d}. L\'evy: \cref{thm:levy-continuity}. If the limit is a constant \(c\), then \(\xrightarrow{(d)}c\iff\xrightarrow{\Pbb}c\) \citep[Proposition~9.2.2]{ren2024}. \\
\hline
Continuous mapping
  & \(g\) continuous
  & \(X_n\xrightarrow{(d)}X\) \(\Rightarrow\) \(g(X_n)\xrightarrow{(d)}g(X)\)
  & Laws
  & Used in Step~2 (the integral functional on \(C[0,1]\)) and Step~5 (\(y\mapsto e^{y}\)). \\
\hline
Slutsky
  & \(A_n+B_n X_n\)
  & \(X_n\xrightarrow{(d)}X\), \(A_n\xrightarrow{\Pbb}a\), \(B_n\xrightarrow{\Pbb}b\)
  & Mixed
  & Then \(A_n+B_n X_n\xrightarrow{(d)}a+bX\). Recorded as \cref{lem:slutsky}. Not needed for the pathwise product \(c_T X_T\) on \(\Omega_c\). \\
\hline
\end{tabularx}
}

\vspace{0.55em}
\noindent
\begin{center}
\begin{tikzpicture}[
  >=Stealth,
  x=1cm, y=1cm,
  box/.style={
    draw=#1!60!black, fill=#1!9, rounded corners=3pt,
    align=center, font=\small, inner xsep=6pt, inner ysep=5pt,
    minimum width=3.35cm, minimum height=1.18cm
  },
  arr/.style={->, very thick, draw=black!52},
  darr/.style={->, thick, dashed, draw=black!42},
  elab/.style={font=\scriptsize, inner sep=1pt, fill=white, text=black!62}
]
  \node[box=blue]  (lp) at (4.85, 2.25)
    {\textbf{\(L^{q}\Rightarrow L^{p}\)}\\[-1pt]{\scriptsize \(q>p\ge 1\)}};
  \node[box=teal]  (as) at (0, 0)
    {\textbf{almost sure}\\[-1pt]{\scriptsize \(\xrightarrow{\mathrm{a.s.}}\)}};
  \node[box=orange](pr) at (4.85, 0)
    {\textbf{in probability}\\[-1pt]{\scriptsize \(\xrightarrow{\Pbb}\)}};
  \node[box=purple](di) at (9.7, 0)
    {\textbf{in distribution}\\[-1pt]{\scriptsize \(\xrightarrow{(d)}\ /\) weak}};

  \draw[arr] (as) -- node[elab, above] {always} (pr);
  \draw[arr] (pr) -- node[elab, above] {always} (di);
  \draw[arr] (lp) -- node[elab, right] {Markov} (pr);

  \draw[darr] (di.south) to[out=-120, in=-60, looseness=1.15]
    node[elab, below] {limit a constant}
    (pr.south);
  \draw[darr] (pr.south) to[out=-120, in=-60, looseness=1.15]
    node[elab, below] {subsequence}
    (as.south);
  \draw[darr] (pr.west) to[out=125, in=-125]
    node[elab, left] {+~UI}
    (lp.west);
\end{tikzpicture}
\end{center}
{\scriptsize
solid = unconditional; dashed = extra hypothesis.
This step: \(X_T\xrightarrow{\mathrm{a.s.}}e^{M}\) and \(c_T\to 1\), hence \(J_T\xrightarrow{\mathrm{a.s.}}e^{M}\); then \(I_T\stackrel{(d)}{=}J_T\) gives only \(I_T\xrightarrow{(d)}e^{M}\).
}
\end{remark}

The prefactor \(c_T\coloneqq T^{2/T}\) is deterministic, and \(c_T\to 1\) by Step~3. A deterministic real sequence is a degenerate random sequence, so \(c_T\xrightarrow{\Pbb}1\) \citep[Exercise~5.1.1]{hogg2019}; equivalently, on a common space, weak convergence to a constant is convergence in probability \citep[Proposition~9.2.2]{ren2024}. On \(\Omega_c\) one is multiplying two convergent real sequences, hence \(c_T X_T\to e^{M}\) almost surely, with no appeal to Slutsky: this is again \citep[Definition~5.4.1]{ren2024}. Slutsky is the algebraic tool if one only wants the weaker \(\xrightarrow{(d)}\) conclusion from \(X_T\xrightarrow{(d)}e^{M}\) and \(c_T\xrightarrow{\Pbb}1\); the argument below records it because the exam's target mode is in distribution.

\begin{lemma}[Slutsky]
\label{lem:slutsky}
Let \(X_n\), \(X\), \(A_n\), and \(B_n\) be random variables, and let \(a\) and \(b\) be constants. If \(X_n\xrightarrow{(d)}X\), \(A_n\xrightarrow{\Pbb}a\), and \(B_n\xrightarrow{\Pbb}b\), then
\[
A_n+B_n X_n
\xrightarrow{(d)}
a+bX.
\]
This is the QE-probability form: the converging-together lemma in \citep[\S3.2]{durrett2010} (continuous mapping of \((x,y)\mapsto a+by\) after the second coordinate converges to a constant). The statistical form is \citep[Theorem~5.2.5]{hogg2019} (Hogg writes \(\stackrel{D}{\to}\) for \(\xrightarrow{(d)}\)). Ren does not isolate the name; the ingredients are weak convergence of constants \citep[Proposition~9.2.2]{ren2024} and the bounded-continuous characterisation already used in the proof of \citep[Proposition~9.2.1]{ren2024}. In particular, taking \(A_n\equiv 0\) and \(B_n=c_n\) deterministic with \(c_n\to c\), one has \(c_n X_n\xrightarrow{(d)}cX\).
\end{lemma}

Applied here with \(A_T\equiv 0\), \(B_T=T^{2/T}\xrightarrow{\Pbb}1\), and \(X_T\) as above (which converges almost surely, hence in distribution), \cref{lem:slutsky} gives \(c_T X_T\xrightarrow{(d)}e^{M}\). The same product also converges almost surely, by the pathwise argument of the previous paragraph (replace the continuous parameter \(T\to\infty\) by an arbitrary sequence \(T_k\to\infty\) if one wants the sequential statements in Durrett, Hogg, and Ren literally). Thus \(J_T=c_T X_T\) satisfies

\[
J_T
\xrightarrow{\mathrm{a.s.}}
e^{M}.
\]

By Step~2, \(I_T\stackrel{(d)}{=}J_T\) for each \(T\), so \(I_T\xrightarrow{(d)}e^{M}\). This is still a random variable: \(M=\max_{u\in[0,1]}B_u\). It remains to write \(M\) as a function of a single normal.

\textbf{Step 5 (reflection).}
The claim is \(M\stackrel{(d)}{=}|B_1|\), not \(M\stackrel{(d)}{=}B_1\). Indeed \(M\ge B_0=0\) almost surely, so \(M\) cannot be Gaussian. \Cref{thm:bm-reflection} supplies the tail identity \(\Pbb(M\ge x)=2\Pbb(B_1\ge x)=\Pbb(|B_1|\ge x)\) for \(x>0\), hence \(M\stackrel{(d)}{=}|B_1|\). Writing \(G\coloneqq B_1\sim\mathcal{N}(0,1)\), the continuous map \(y\mapsto e^{y}\) yields

\[
e^{M}
\stackrel{(d)}{=}
e^{|G|}.
\]

Combining with the conclusion of Step~4, \(I_T\xrightarrow{(d)}e^{|G|}\).
\end{proof}

\subsection{Problem 5. [10 pts.]}
A Galton--Watson branching process is a stochastic process \( \{Z_n\}_{n \in \mathbb{N} \cup \{0\}} \) which evolves according to the recurrence formula

\[
Z_0 := 1, \qquad Z_{n+1} := \sum_{j=1}^{Z_n} X(j, n),
\]

where \( \{X(j, n) : j \in \mathbb{N},\, n \in \mathbb{N} \cup \{0\}\} \) is a set of i.i.d.\ \( \mathbb{N} \cup \{0\} \)-valued random variables with mean \( m := \mathbb{E}[X(1, 1)] < \infty \). Let \( p_k := \mathbb{P}[X(1, 1) = k] \), \( k \in \mathbb{N} \cup \{0\} \), and we assume \( p_0 > 0 \), \( p_0 + p_1 < 1 \). Let \( f(s) := \mathbb{E}[s^{X(1,1)}] \) for \( s \in [0, 1] \) (with the convention \( 0^0 = 1 \)). Let

\[
\mathcal{F}_n := \sigma(Z_0, \dots, Z_n), \qquad n \in \mathbb{N} \cup \{0\}.
\]

\begin{enumerate}[label=(\alph*)]
    \item Compute

    \[
    \mathbb{E}[Z_{n+1} \mid \mathcal{F}_n],\quad
    \mathbb{E}[Z_n],\quad
    \mathbb{E}[s^{Z_{n+1}} \mid \mathcal{F}_n],\quad
    \mathbb{E}[s^{Z_n}], \qquad n \in \mathbb{N} \cup \{0\},\; s \in [0, 1].
    \]
    \item Determine the number of solutions of \( f(s) = s \) on \( [0, 1] \).
    \item Let \( \tau := \inf\{n \in \mathbb{N} : Z_n = 0\} \) (with \( \inf \emptyset := \infty \)). Call the event \( \{\tau < \infty\} \) extinction. Compute \( \mathbb{P}[\tau < \infty] \). Write your answer as a function of \( q \), where \( q \) is the smallest solution of \( f(s) = s \) on \( [0, 1] \).
\end{enumerate}

\setcounter{section}{5}
\setcounter{definition}{0}

\subsubsection{Definitions and warmup calculations for Problem 5}

The notes below first separate the general branching mechanism from the Galton--Watson process that the exam actually uses, then unfold the generating-function calculus. Parts~(b)--(c) are left as hints. Sources: \citep[Exercise~5.11, Proposition~5.6, Remark~5.2]{brzezniak1999} (Markov chain on \(\{0,1,2,\dots\}\)) and \citep[\S5.3.4]{durrett2010} (generating-function calculus). In both texts ``branching process'' means Galton--Watson; the names are not interchangeable in the wider literature.

\paragraph{Learning path (Problems~5--6 as one package).}
Read in this order. Each step has a motivation, then a definition later on this page or in Problem~6. Do not skip to the tree or to percolation before the generating function is automatic.

\begin{enumerate}[label=\textup{Step~\arabic*.}]
    \item
    \emph{Independent lineages.} Motivation: a population model is usable only if disjoint families do not interact. That is the branching property \cref{def:branching-property}. Density dependence (epidemics after takeoff, logistic growth) is the typical way the property fails.
    \item
    \emph{Discrete generations, one type, count heads.} Restrict the clock to ``year \(n\)'', give everyone the same offspring law \((p_k)\), and record only how many people are alive. That is Galton--Watson \cref{def:galton-watson}, the exam's \(\{Z_n\}\). The tree \(T\) of Problem~6 is the same object with the parent pointers kept.
    \item
    \emph{Markov chain on \(\mathbb{N}_0\).} Motivation: given the current headcount, the next generation is an i.i.d.\ sum, so the past is irrelevant. State \(0\) is absorbing. This is why Brzeźniak puts the model in Chapter~5 rather than in a branching chapter.
    \item
    \emph{Probability generating functions.} Motivation: independence turns sums into products, \(s^{Y_1+\cdots+Y_k}=s^{Y_1}\cdots s^{Y_k}\). One number \(f(s)=\E[s^X]\) therefore encodes the whole offspring law, and a random number of summands becomes the composition \(f_n=f^{\circ n}\). Without this, extinction probabilities are an infinite hierarchy of convolutions.
    \item
    \emph{Mean \(m\) versus almost-sure extinction.} Motivation: \(\E[Z_n]=m^n\) only tracks the average. A lottery with rare huge families can have \(m>1\) while still dying out with positive probability. The trichotomy subcritical / critical / supercritical \cref{def:criticality} is about \(m\); the extinction probability is about the smallest root of \(f(s)=s\). These two numbers agree only in the subcritical and critical cases.
    \item
    \emph{From counts to the tree} (Problem~6). Motivation: percolation and inherited properties need edges, not just \(Z_n\). The recursive split into child subtrees is the branching property written on the graph.
    \item
    \emph{Inherited events, then percolation} (Problem~6). Motivation: on a supercritical tree the infinite core looks the same from every vertex along a surviving lineage, so many properties cannot take intermediate probabilities. Bernoulli percolation is the first nontrivial inherited property the exam asks you to name: ``open cluster of the root is finite''. Thinning reduces that question, annealed, back to Step~5.
\end{enumerate}

\paragraph{The general mechanism.}
A branching process is a population model whose defining feature is independence of disjoint lineages. Once an individual is born, the subprocess of its descendants is an independent copy of the original process, and the whole population is the superposition of those copies. There is no single canonical state space: individuals may carry a type, a spatial location, or a continuous-time lifetime. What is required is the branching property, not the generation-count recursion of the exam.

\begin{definition}[Branching property]
\label{def:branching-property}
Let a population be started from one individual, and write \(T\) for the (random) descendant structure of that individual. The process has the \emph{branching property} if, whenever the ancestor produces \(k\) children, the \(k\) descendant structures rooted at those children are i.i.d.\ copies of \(T\), independent of one another and of the ancestor's reproduction. The empty collection of children gives the empty population, which stays empty.
\end{definition}

The picture is a two-step sampling, not a global interaction. First the ancestor samples a number of children \(k\); then one attaches \(k\) independent fresh copies of the original law. Those copies never look at one another, and they never look back at \(k\).

\begin{center}
\begin{tikzpicture}[
    person/.style={circle, draw, inner sep=1.15pt, font=\scriptsize, minimum size=4.2mm},
    lab/.style={draw=none, font=\scriptsize, inner sep=1pt},
    edge from parent/.style={draw, -{Stealth[length=1.4mm]}},
    level distance=9.5mm,
    sibling distance=13mm
]
    \node[person] (o) {\(o\)}
        child { node[person] (v1) {\(v_1\)}
            child { node[person] (v11) {} }
            child { node[person] (v12) {} }
        }
        child { node[person] (v2) {\(v_2\)}
            child { node[person] (v21) {} }
        }
        child { node[person] (v3) {\(v_3\)} };
    \node[lab, left] at ($(o)+(-0.55,0)$) {ancestor};
    \draw[dashed, rounded corners, black!55]
        ($(v1)+(-0.72,0.38)$) rectangle ($(v12)+(0.55,-0.42)$);
    \draw[dashed, rounded corners, black!55]
        ($(v2)+(-0.48,0.38)$) rectangle ($(v21)+(0.48,-0.42)$);
    \draw[dashed, rounded corners, black!55]
        ($(v3)+(-0.42,0.38)$) rectangle ($(v3)+(0.42,-0.42)$);
    \node[lab] at ($(v1)+(0,-1.55)$) {\(T_1\cong T\)};
    \node[lab] at ($(v2)+(0,-1.15)$) {\(T_2\cong T\)};
    \node[lab, right] at ($(v3)+(0.55,0)$) {\(T_3\cong T\)};
    \node[lab, align=left, anchor=west] at (3.35,0.15) {\(k=3\) children first,};
    \node[lab, align=left, anchor=west] at (3.35,-0.35) {then \(T_1,T_2,T_3\) i.i.d.};
    \node[lab, align=left, anchor=west] at (3.35,-0.85) {\(T_i\perp T_j\) and \(T_i\perp k\).};
\end{tikzpicture}
\end{center}

The Galton--Watson process of \cref{def:galton-watson} is the discrete-generation, single-type special case in which one records only the generation sizes. Continuous-time, multitype, and spatial variants keep \cref{def:branching-property} and drop one or more of those restrictions. Reproduction that depends on the total population size (density dependence) destroys the property, because lineages then interact: a birth in one family changes the environment seen by another.

\begin{center}
\begin{tikzpicture}[
    person/.style={circle, draw, inner sep=1.05pt, font=\scriptsize, minimum size=3.8mm},
    lab/.style={draw=none, font=\scriptsize, inner sep=0.6pt},
    >=Stealth
]
    % left: independent lineages
    \begin{scope}[shift={(0,0)}]
        \node[lab] at (1.15,1.55) {branching: lineages ignore each other};
        \node[person] (a) at (0.35,0.95) {};
        \node[person] (a1) at (0.0,0.15) {};
        \node[person] (a2) at (0.7,0.15) {};
        \draw[->] (a) -- (a1);
        \draw[->] (a) -- (a2);
        \node[person] (b) at (1.95,0.95) {};
        \node[person] (b1) at (1.6,0.15) {};
        \node[person] (b2) at (2.3,0.15) {};
        \draw[->] (b) -- (b1);
        \draw[->] (b) -- (b2);
        \node[lab] at (1.15,-0.35) {\(T_{\mathrm{left}}\perp T_{\mathrm{right}}\)};
    \end{scope}
    % right: interaction
    \begin{scope}[shift={(5.15,0)}]
        \node[lab] at (1.35,1.55) {not branching: shared finite pool};
        \draw[dashed, rounded corners, black!50, fill=black!4]
            (-0.25,-0.15) rectangle (2.95,1.35);
        \node[person] (c) at (0.45,0.95) {};
        \node[person] (c1) at (0.1,0.2) {};
        \node[person] (c2) at (0.8,0.2) {};
        \draw[->] (c) -- (c1);
        \draw[->] (c) -- (c2);
        \node[person] (d) at (2.25,0.95) {};
        \node[person] (d1) at (1.9,0.2) {};
        \node[person, draw=black!40, text=black!40] (d2) at (2.6,0.2) {};
        \draw[->] (d) -- (d1);
        \draw[->, densely dashed, black!45] (d) -- (d2);
        \draw[line width=0.7pt, red!70!black] (2.48,0.32) -- (2.72,0.08);
        \draw[line width=0.7pt, red!70!black] (2.48,0.08) -- (2.72,0.32);
        \node[lab] at (1.35,-0.35) {one birth uses up a susceptible};
    \end{scope}
\end{tikzpicture}
\end{center}

\begin{example}[Four models that share the branching property]
\label{ex:branching-models}
All four keep independent lineages; they differ in the clock and in what is recorded.
\begin{enumerate}[label=\textup{(\arabic*)}]
    \item
    \emph{Galton--Watson} (this exam, \cref{def:galton-watson}): discrete generations, one type, i.i.d.\ offspring, state \(Z_n=\) generation size. Durrett and Brze\'zniak simply call this a branching process.
    \item
    \emph{Continuous-time Markov branching}: each living individual independently waits an exponential time, then is replaced by a random number of children. Binary fission (always two children) is the Yule process.
    \item
    \emph{Multitype Galton--Watson}: finitely many types; an individual's offspring is a random vector. The scalar mean \(m\) becomes a mean matrix, and criticality is spectral radius \(1\).
    \item
    \emph{Branching random walk}: each child is independently displaced from its parent. Problem~6's percolation on a Galton--Watson tree is a close cousin: keep or delete edges independently, then ask whether the open cluster of the root is infinite.
\end{enumerate}
\end{example}

The four clocks and state spaces, drawn on the same small family. Only the first one is the exam's \(\{Z_n\}\); the other three keep \cref{def:branching-property} and change what is recorded.

\begin{center}
\begin{tikzpicture}[
    person/.style={circle, draw, inner sep=0.9pt, font=\tiny, minimum size=3.4mm},
    lab/.style={draw=none, font=\scriptsize, inner sep=0.4pt},
    ttl/.style={draw=none, font=\small, inner sep=0.4pt},
    >=Stealth
]
    % (1) Galton--Watson: discrete generations
    \begin{scope}[shift={(0,0)}]
        \fill[black!6] (-0.55,1.05) rectangle (2.55,1.55);
        \fill[black!4] (-0.55,0.35) rectangle (2.55,0.85);
        \fill[black!6] (-0.55,-0.45) rectangle (2.55,0.05);
        \node[person] (g0) at (1.0,1.3) {\(o\)};
        \node[person] (g11) at (0.25,0.6) {};
        \node[person] (g12) at (1.0,0.6) {};
        \node[person] (g13) at (1.75,0.6) {};
        \node[person] (g21) at (0.0,-0.2) {};
        \node[person] (g22) at (0.5,-0.2) {};
        \node[person] (g23) at (1.0,-0.2) {};
        \draw[->] (g0) -- (g11);
        \draw[->] (g0) -- (g12);
        \draw[->] (g0) -- (g13);
        \draw[->] (g11) -- (g21);
        \draw[->] (g11) -- (g22);
        \draw[->] (g12) -- (g23);
        \node[lab, left] at (-0.55,1.3) {\(n{=}0\)};
        \node[lab, left] at (-0.55,0.6) {\(n{=}1\)};
        \node[lab, left] at (-0.55,-0.2) {\(n{=}2\)};
        \node[ttl] at (1.0,-0.95) {(1) Galton--Watson};
        \node[lab] at (1.0,-1.3) {one year \(=\) one generation};
    \end{scope}

    % (2) continuous-time Markov branching / Yule
    \begin{scope}[shift={(5.55,0)}]
        \draw[->] (-0.15,-0.35) -- (3.55,-0.35) node[right] {\(t\)};
        \node[person] (c0) at (0.0,1.15) {};
        \draw[thick] (0.0,1.15) -- (1.15,1.15);
        \node[person] (csplit) at (1.15,1.15) {};
        \draw[thick] (1.15,1.15) -- (1.15,1.65) -- (3.15,1.65);
        \node[person] at (3.15,1.65) {};
        \draw[thick] (1.15,1.15) -- (1.15,0.45) -- (2.05,0.45);
        \node[person] (c2) at (2.05,0.45) {};
        \draw[thick] (2.05,0.45) -- (2.05,0.85) -- (3.15,0.85);
        \draw[thick] (2.05,0.45) -- (2.05,0.05) -- (3.15,0.05);
        \node[person] at (3.15,0.85) {};
        \node[person] at (3.15,0.05) {};
        \node[lab] at (1.15,2.0) {split};
        \node[ttl] at (1.7,-0.95) {(2) continuous time};
        \node[lab] at (1.7,-1.3) {overlapping lifetimes};
    \end{scope}

    % (3) multitype
    \begin{scope}[shift={(0,-3.55)}]
        \node[person, fill=black!18] (m0) at (1.1,1.25) {\(A\)};
        \node[person, fill=black!18] (m1) at (0.25,0.35) {\(A\)};
        \node[person] (m2) at (1.1,0.35) {\(B\)};
        \node[person, fill=black!18] (m3) at (1.95,0.35) {\(A\)};
        \node[person] (m4) at (0.7,-0.45) {\(B\)};
        \node[person, fill=black!18] (m5) at (1.5,-0.45) {\(A\)};
        \draw[->] (m0) -- (m1);
        \draw[->] (m0) -- (m2);
        \draw[->] (m0) -- (m3);
        \draw[->] (m2) -- (m4);
        \draw[->] (m2) -- (m5);
        \node[ttl] at (1.1,-0.95) {(3) multitype};
        \node[lab] at (1.1,-1.3) {shaded \(=A\), open \(=B\)};
    \end{scope}

    % (4) branching random walk
    \begin{scope}[shift={(5.55,-3.55)}]
        \draw[->] (-0.35,-0.55) -- (3.45,-0.55) node[right] {\(x\)};
        \foreach \x in {0,1,2,3}
            \draw (\x,-0.5) -- (\x,-0.6);
        \node[person] (r0) at (1.0,1.25) {};
        \node[person] (r1) at (0.0,0.4) {};
        \node[person] (r2) at (2.2,0.4) {};
        \node[person] (r3) at (0.55,-0.15) {};
        \node[lab, above] at (1.0,1.25) {\(0\)};
        \node[lab, left] at (0.0,0.4) {\({-}1\)};
        \node[lab, right] at (2.2,0.4) {\(2\)};
        \node[lab, below] at (0.55,-0.15) {\(0\)};
        \draw[->] (r0) -- (r1);
        \draw[->] (r0) -- (r2);
        \draw[->] (r1) -- (r3);
        \node[ttl] at (1.55,-0.95) {(4) branching random walk};
        \node[lab] at (1.55,-1.3) {child \(=\) parent \(+\) jump};
    \end{scope}
\end{tikzpicture}
\end{center}

\begin{remark}[What is Galton--Watson, and how the four models differ]
\label{rem:branching-vs-gw}
The recursive picture ``the root hangs \(K\) independent copies of itself'' is \cref{def:branching-property}, not a Galton--Watson exclusive. It does not yet specify a clock, marks on individuals, or which coordinate is recorded. Write \(T\) for the descendant structure of one individual. After that individual produces \(k\) children,
\[
T
\;=\;
\bigl(\text{\(k\) children}\bigr)
\;+\;
(T_1,\dots,T_k),
\]
where \(T_1,\dots,T_k\) are i.i.d.\ copies of \(T\), independent of one another and of \(k\). If \(k=0\) the population is empty and stays empty. That is the whole first principle: lineages, once split, never look at one another.

Galton--Watson is that mechanism plus three extra restrictions, which produce the exam's \(\{Z_n\}\) in \cref{def:galton-watson}.
\begin{enumerate}[label=\textup{(\roman*)}]
    \item
    \emph{Discrete generations, in lockstep.} Generation \(n\) reproduces simultaneously and disappears. There are no overlapping lifetimes. The time index is the generation number \(n=0,1,2,\dots\).
    \item
    \emph{Single type, one offspring law \((p_k)\).} Children carry no type and no location. Every individual draws offspring from the same law.
    \item
    \emph{Record only the headcount \(Z_n\).} Parent--child edges may be discarded. The tree \(T\) of Problem~6 is the same reproduction with the edges kept; \(\{Z_n\}\) is the same reproduction with the edges dropped.
\end{enumerate}
Those three cuts turn the hanging of subtrees into
\[
Z_{n+1}
=
\sum_{j=1}^{Z_n} X(j,n).
\]
The composition \(f_n=f^{\circ n}\), the mean \(m^n\), and criticality by the scalar \(m\) are consequences of the cuts, not of \cref{def:branching-property} alone.

The four models of \cref{ex:branching-models} all keep independent hanging of subtrees. They differ in the clock, in the marks on individuals, and in the recorded state.

{\footnotesize
\setlength{\tabcolsep}{3.5pt}
\renewcommand{\arraystretch}{1.22}
\noindent
\begin{tabularx}{\linewidth}{@{}
  >{\raggedright\arraybackslash}p{2.45cm}
  >{\raggedright\arraybackslash}Y
  >{\raggedright\arraybackslash}Y
  >{\raggedright\arraybackslash}Y
  >{\raggedright\arraybackslash}Y
  @{}}
\toprule
\textbf{Model}
  & \textbf{Clock}
  & \textbf{Marks}
  & \textbf{Recorded state}
  & \textbf{Criticality} \\
\midrule
(1) Galton--Watson
  & discrete generations, lockstep
  & none
  & \(Z_n=\) generation size
  & scalar \(m=\E[X]\) \\
(2) CT Markov branching
  & each individual waits \(\mathrm{Exp}\), then is replaced
  & none
  & \(Z(t)=\) living at time \(t\)
  & still scalar; the Malthusian parameter includes the waiting rate \\
(3) Multitype GW
  & still discrete generations
  & finitely many types
  & vector of type counts
  & spectral radius of the mean matrix \\
(4) Branching random walk
  & typically discrete generations
  & spatial location
  & a cloud of positions
  & front versus population growth \\
\bottomrule
\end{tabularx}
}

Read on the same three-child family drawn above. Model~(1) asks only how many people are alive in each generation. Model~(2) lets the three lives run on independent exponential clocks, so lifetimes overlap; the Yule process is the special case of binary fission, with randomness only in the waiting times. Model~(3) equips children with types, so the mean is a matrix and criticality is spectral radius \(1\). Model~(4) equips children with spatial displacements. Problem~6's percolation is a cousin of~(4): independently keep or delete edges on a Galton--Watson tree, rather than writing \(\mathbb{R}^d\)-coordinates on vertices, and ask whether the open cluster of the root is infinite.
\end{remark}

\begin{example}[Applications of those models]
\label{ex:branching-applications}
\begin{enumerate}[label=\textup{(\arabic*)}]
    \item
    \emph{Surnames.} Galton and Watson asked whether a family name dies out. Only sons transmit the name; their number is the offspring law. This is the historical Galton--Watson problem, and is \cref{ex:gw-binomial}(2) below.
    \item
    \emph{Chain reactions.} A neutron produces a random number of further neutrons on collision. Extinction of the branching process is the event that the reaction dies out; \(m=1\) is criticality.
    \item
    \emph{Early epidemics.} While infectives are rare, each produces an approximately i.i.d.\ number of secondary cases (mean \(R_0\)). The outbreak dies out with the extinction probability of that Galton--Watson process. The approximation fails once the susceptible pool is depleted, because lineages then compete: density dependence, so not a branching process any more.
    \item
    \emph{\(M/G/1\) busy period.} Customers arriving during one service are the ``children'' of the customer in service. The busy period is finite if and only if that Galton--Watson process goes extinct.
\end{enumerate}
\end{example}

\begin{center}
\begin{tikzpicture}[
    person/.style={circle, draw, inner sep=0.85pt, font=\tiny, minimum size=3.3mm},
    lab/.style={draw=none, font=\scriptsize, inner sep=0.4pt},
    ttl/.style={draw=none, font=\small, inner sep=0.4pt},
    >=Stealth
]
    % surnames
    \begin{scope}[shift={(0,0)}]
        \node[person, fill=black!18] (s0) at (1.05,1.2) {};
        \node[person, fill=black!18] (s1) at (0.15,0.4) {};
        \node[person] (s2) at (1.05,0.4) {};
        \node[person, fill=black!18] (s3) at (1.95,0.4) {};
        \node[person, fill=black!18] (s11) at (0.15,-0.35) {};
        \draw[->] (s0) -- (s1);
        \draw[->] (s0) -- (s2);
        \draw[->] (s0) -- (s3);
        \draw[->] (s1) -- (s11);
        \node[lab, below] at (1.05,0.28) {name stops};
        \node[ttl] at (1.05,-0.85) {(1) surnames};
        \node[lab] at (1.05,-1.2) {shaded \(=\) sons only};
    \end{scope}

    % chain reaction
    \begin{scope}[shift={(4.35,0)}]
        \node[person] (n0) at (1.0,1.2) {\(n\)};
        \node[person] (n1) at (0.2,0.35) {\(n\)};
        \node[person] (n2) at (1.0,0.35) {\(n\)};
        \node[person] (n3) at (1.8,0.35) {\(n\)};
        \node[person, draw=black!40, text=black!40] (n4) at (1.0,-0.4) {};
        \draw[->] (n0) -- (n1);
        \draw[->] (n0) -- (n2);
        \draw[->] (n0) -- (n3);
        \draw[->, densely dashed, black!45] (n2) -- (n4);
        \draw[line width=0.65pt, red!70!black] (0.88,-0.28) -- (1.12,-0.52);
        \draw[line width=0.65pt, red!70!black] (0.88,-0.52) -- (1.12,-0.28);
        \node[lab] at (1.0,-0.05) {};
        \node[ttl] at (1.0,-0.85) {(2) chain reaction};
        \node[lab] at (1.0,-1.2) {absorbed \(=\) zero offspring};
    \end{scope}

    % epidemic
    \begin{scope}[shift={(8.55,0)}]
        \node[person, fill=red!20] (e0) at (0.9,1.2) {};
        \node[person, fill=red!20] (e1) at (0.25,0.45) {};
        \node[person, fill=red!20] (e2) at (1.55,0.45) {};
        \node[person, fill=red!20] (e11) at (0.0,-0.25) {};
        \node[person, fill=red!20] (e12) at (0.5,-0.25) {};
        \node[person, draw=black!40] (e21) at (1.3,-0.25) {};
        \node[person, draw=black!40] (e22) at (1.8,-0.25) {};
        \draw[->] (e0) -- (e1);
        \draw[->] (e0) -- (e2);
        \draw[->] (e1) -- (e11);
        \draw[->] (e1) -- (e12);
        \draw[->, densely dashed, black!45] (e2) -- (e21);
        \draw[->, densely dashed, black!45] (e2) -- (e22);
        \node[lab] at (0.35,0.9) {early};
        \node[lab] at (1.85,0.9) {later};
        \node[ttl] at (0.95,-0.85) {(3) epidemic};
        \node[lab] at (0.95,-1.2) {dashed: pool depleted};
    \end{scope}

    % M/G/1 --- placed below surnames+chain to avoid overflow
\end{tikzpicture}
\end{center}

\begin{center}
\begin{tikzpicture}[
    lab/.style={draw=none, font=\scriptsize, inner sep=0.4pt},
    ttl/.style={draw=none, font=\small, inner sep=0.4pt},
    person/.style={circle, draw, inner sep=0.85pt, font=\tiny, minimum size=3.3mm},
    >=Stealth
]
    \draw[->] (0,0) -- (8.4,0) node[right] {\(t\)};
    \draw[line width=2.4pt, black!75] (0.4,0) -- (3.6,0);
    \node[lab] at (2.0,0.55) {service of customer \(0\)};
    \foreach \x/\labtxt in {1.1/1, 2.0/2, 3.1/3}
    {
        \fill (\x,0) circle (1.35pt);
        \node[lab] at (\x, -0.35) {\labtxt};
    }
    \node[person] (q0) at (9.35,1.35) {\(0\)};
    \node[person] (q1) at (8.55,0.45) {\(1\)};
    \node[person] (q2) at (9.35,0.45) {\(2\)};
    \node[person] (q3) at (10.15,0.45) {\(3\)};
    \draw[->] (q0) -- (q1);
    \draw[->] (q0) -- (q2);
    \draw[->] (q0) -- (q3);
    \draw[->, thick, black!50] (3.85,0.35) to[out=25,in=160] (8.15,1.15);
    \node[ttl] at (5.2,-0.95) {(4) \(M/G/1\) busy period};
    \node[lab] at (5.2,-1.3) {arrivals during one service \(=\) children};
\end{tikzpicture}
\end{center}

\paragraph{What the exam's process is.}
One individual is alive in year \(0\). Every individual lives exactly one year and, independently of everyone else, produces \(k\) children with probability \(p_k\), then dies. The children form the next generation, and the story repeats. The exam's \(Z_n\) is the number of individuals alive in year \(n\). This is exactly the Vugiel population \(\xi_n\) of \citep[Exercise~5.11]{brzezniak1999}. The picture records both the family tree and the generation counts: horizontal bands are years, and the labels \(X(j,n)\) are the offspring numbers in the exam's sum.

\begin{center}
\begin{tikzpicture}[
    person/.style={circle, draw, inner sep=1.1pt, font=\scriptsize, minimum size=4.4mm},
    lab/.style={draw=none, font=\scriptsize, inner sep=0.5pt},
    >=Stealth
]
    \fill[black!5] (-2.15,0.85) rectangle (4.55,1.55);
    \fill[black!3] (-2.15,-0.05) rectangle (4.55,0.65);
    \fill[black!5] (-2.15,-1.15) rectangle (4.55,-0.25);
    \node[person] (o) at (1.2,1.2) {\(o\)};
    \node[person] (a) at (0.0,0.3) {\(1\)};
    \node[person] (b) at (1.2,0.3) {\(2\)};
    \node[person] (c) at (2.4,0.3) {\(3\)};
    \node[person] (a1) at (-0.45,-0.7) {};
    \node[person] (a2) at (0.45,-0.7) {};
    \node[person] (b1) at (1.2,-0.7) {};
    \draw[->] (o) -- (a);
    \draw[->] (o) -- (b);
    \draw[->] (o) -- (c);
    \draw[->] (a) -- (a1);
    \draw[->] (a) -- (a2);
    \draw[->] (b) -- (b1);
    \node[lab, left] at (-2.15,1.2) {year \(0\)};
    \node[lab, left] at (-2.15,0.3) {year \(1\)};
    \node[lab, left] at (-2.15,-0.7) {year \(2\)};
    \node[lab, right] at (4.55,1.2) {\(Z_0=1\)};
    \node[lab, right] at (4.55,0.3) {\(Z_1=3\)};
    \node[lab, right] at (4.55,-0.7) {\(Z_2=3\)};
    \node[lab] at (0.0,-0.05) {\(X(1,1)=2\)};
    \node[lab] at (1.2,-0.05) {\(X(2,1)=1\)};
    \node[lab] at (2.4,-0.05) {\(X(3,1)=0\)};
    \node[lab] at (1.2,1.55) {\(X(1,0)=3\)};
\end{tikzpicture}
\end{center}

In the picture, the three individuals of generation \(1\) have \(2\), \(1\), and \(0\) children respectively. The exam writes those three numbers as \(X(1,1)\), \(X(2,1)\), \(X(3,1)\), so
\[
Z_2
=
X(1,1)+X(2,1)+X(3,1)
=
2+1+0
=
3.
\]
If some generation is empty, there is nobody left to reproduce, and every later generation stays empty. That is extinction.

\begin{definition}[Galton--Watson process]
\label{def:galton-watson}
Let \(\{X(j,n): j\ge 1,\ n\ge 0\}\) be i.i.d.\ nonnegative integer-valued random variables (the \emph{offspring} of the \(j\)th individual in generation \(n\)), with
\[
p_k
\coloneqq
\Pbb\bigl(X(1,1)=k\bigr),
\qquad
k=0,1,2,\dots,
\qquad
m
\coloneqq
\E\bigl[X(1,1)\bigr]
\in(0,\infty).
\]
The process started at one individual is defined by \(Z_0=1\) and
\[
Z_{n+1}
\coloneqq
\sum_{j=1}^{Z_n} X(j,n),
\]
with the empty sum equal to \(0\). Equivalently, in the notation of \citep[Exercise~5.11]{brzezniak1999}: if generation \(n\) has size \(i\), then generation \(n+1\) is the sum of \(i\) independent copies of the offspring law \((p_k)\). See also \citep[(5.3.2)]{durrett2010}.
\end{definition}

Relative to \cref{def:branching-property}, this adds the three restrictions of \cref{rem:branching-vs-gw}: time is the generation index (no overlapping lifetimes); every individual in every generation has the same offspring law \((p_k)\); the recorded state is only the generation size \(Z_n\), not ages, types, or locations. The exam's standing assumptions \(p_0>0\) and \(p_0+p_1<1\) are extra nondegeneracy conditions on that offspring law, not part of the branching property: they mean, respectively, that one individual can have no children, and that the offspring law is not supported on \(\{0,1\}\) (so two or more children occur with positive probability).

\begin{definition}[Criticality]
\label{def:criticality}
The process is \emph{subcritical} if \(m<1\), \emph{critical} if \(m=1\), and \emph{supercritical} if \(m>1\). Equivalently, these are the three possible values of the slope \(f'(1-)\) at the right endpoint of the generating-function graph.
\end{definition}

\begin{remark}[Why the mean is not the extinction probability]
\label{rem:mean-vs-extinction}
The recursion \(\E[Z_{n+1}\mid\mathcal{F}_n]=mZ_n\) iterates to \(\E[Z_n]=m^n\). If \(m<1\) this already forces \(Z_n\to 0\) in probability, hence almost-sure extinction (Markov plus integer values). If \(m>1\) the mean explodes, but that only says a small-probability set of explosive families pulls the average up; it does not forbid a positive probability of dying out. Generating functions are the device that sees that leftover probability. The assumptions \(p_0>0\) and \(p_0+p_1<1\) exclude the two degenerate pictures where this distinction collapses: \(p_0=0\) means the process can never hit \(0\); \(p_0+p_1=1\) makes \(f\) linear, so the graph of \(f\) cannot bend below the diagonal.
\end{remark}

\begin{definition}[Markov chain, for comparison]
\label{def:markov-chain-brz}
An \(S\)-valued sequence \(\{\xi_n\}_{n\ge 0}\) on a countable set \(S\) is a \emph{Markov chain} if
\[
\Pbb(\xi_{n+1}=s\mid \xi_0,\dots,\xi_n)
=
\Pbb(\xi_{n+1}=s\mid \xi_n)
\]
for every \(n\) and every \(s\in S\). This is \citep[Definition~5.1, (5.10)]{brzezniak1999}. If the right-hand side does not depend on \(n\), the chain is time-homogeneous, and one writes \(p(j\mid i)\coloneqq\Pbb(\xi_1=j\mid \xi_0=i)\) for the transition probabilities.
\end{definition}

The process \(\{Z_n\}\) is a time-homogeneous Markov chain on \(S=\{0,1,2,\dots\}\), because the law of \(Z_{n+1}\) given the whole past depends on the past only through the current population size \(Z_n\). The transition probabilities are the laws of i.i.d.\ sums: if \(X_1,X_2,\dots\) are i.i.d.\ copies of \(X(1,1)\), then
\[
p(j\mid i)
=
\Pbb(X_1+\cdots+X_i=j),
\qquad
p(j\mid 0)
=
\mathbf{1}_{\{j=0\}}.
\]
State \(0\) is absorbing: once the population is extinct it stays extinct. Brzeźniak computes two special cases in Exercise~5.12: binomial offspring gives another binomial,
\[
p(j\mid i)
=
\binom{Ni}{j}p^j(1-p)^{Ni-j},
\]
and Poisson(\(\lambda\)) offspring gives Poisson(\(i\lambda\)),
\[
p(j\mid i)
=
e^{-i\lambda}\frac{(i\lambda)^j}{j!}.
\]
The exam does not fix a parametric family; it works with a general \((p_k)\).

\begin{remark}[Notation dictionary]
\label{rem:gw-notation}
\begin{center}
\setlength{\tabcolsep}{4pt}
\renewcommand{\arraystretch}{1.2}
\begin{tabularx}{\linewidth}{@{}Y l Y@{}}
\toprule
object & exam / these notes & Brze\'zniak / Durrett \\
\midrule
population in generation \(n\) & \(Z_n\) & \(\xi_n\) \citep{brzezniak1999}; \(Z_n\) \citep{durrett2010} \\
offspring of one individual & \(X(j,n)\) & \(X_j\) \citep{brzezniak1999}; \(\xi_i^{n+1}\) \citep{durrett2010} \\
offspring mean & \(m\) & \(\lambda\) in Prop.~5.6; \(\mu\) in Durrett \\
offspring pgf & \(f(s)\) & \(F(x)\) in (5.33) \\
extinction probability from \(1\) individual & \(\Pbb(\tau<\infty)\) & \(\phi(1)\) in (5.29) \\
\bottomrule
\end{tabularx}
\end{center}
The filtration in Durrett is the full offspring history \(\mathcal{F}_n=\sigma(X(j,k): j\ge 1,\ 0\le k<n)\). The smaller \(\sigma(Z_0,\dots,Z_n)\) is enough for every computation in (a): given \(Z_n=k\), the next generation is simply a sum of \(k\) fresh copies of \(X(1,1)\).
\end{remark}

\begin{remark}[Why generating functions]
\label{rem:why-pgf}
Ordinary generating functions are the exponential generating functions of a nonnegative integer. Three facts make them the correct transform here, and no other.
\begin{enumerate}[label=\textup{(\roman*)}]
    \item
    Evaluation at \(0\) and \(1\): \(f(0)=p_0=\Pbb(X=0)\) and \(f(1)=1\). Extinction in one generation is a point on the graph.
    \item
    Independent sums become products: \cref{ex:pgf-iid-sum}. A random number of i.i.d.\ summands therefore becomes composition, which is why \(\E[s^{Z_n}]\) is an iterate of \(f\) rather than a convolution written \(n\) times.
    \item
    The derivative at \(1\) is the mean: \(f'(1-)=m\). Convexity of \(f\) is the second-moment structure \(f''\ge 0\). The whole criticality picture of \cref{def:criticality} is visible in one sketch of \(y=f(s)\) against \(y=s\).
\end{enumerate}
Characteristic functions do the same job for laws on \(\mathbb{R}\); they are the wrong tool for a random sum of a random number of nonnegative integers, because one needs the composition identity \(f_n=f^{\circ n}\), which is special to \(s\mapsto s^k\).
\end{remark}

\begin{definition}[Offspring probability generating function]
\label{def:offspring-pgf}
The \emph{probability generating function} of one individual's offspring is
\[
f(s)
\coloneqq
\E\bigl[s^{X(1,1)}\bigr]
=
\sum_{k=0}^{\infty} p_k s^k,
\qquad
s\in[0,1].
\]
Write \(f_n(s)\coloneqq\E[s^{Z_n}]\) for the generating function of generation \(n\), and \(f^{\circ n}\) for the \(n\)-fold composition \(f\circ\cdots\circ f\).
\end{definition}

\begin{proposition}[{Shape of \(f\) on the unit interval}]
\label{prop:pgf-shape}
The function \(f\) in \cref{def:offspring-pgf} satisfies:
\begin{enumerate}[label=\textup{(\roman*)}]
    \item
    \(f(0)=p_0>0\) and \(f(1)=1\);
    \item
    \(f\) is nondecreasing and convex on \([0,1]\), with
    \[
    f'(s)
    =
    \sum_{k\ge 1} k p_k s^{k-1}
    \ge 0,
    \qquad
    f''(s)
    =
    \sum_{k\ge 2} k(k-1)p_k s^{k-2}
    \ge 0
    \]
    for \(s\in[0,1)\) (termwise differentiation is justified by monotone convergence, or by \citep[Theorem~A.5.2]{durrett2010});
    \item
    \(f'(1-)=m\);
    \item
    \(p_0+p_1<1\) forces \(p_k>0\) for some \(k\ge 2\), hence \(f''>0\) on \((0,1)\) and \(f\) is \emph{strictly} convex.
\end{enumerate}
In particular \(f\) is not the identity map, and the graph of \(y=f(s)\) can cross the diagonal \(y=s\) at most twice on \([0,1]\).
\end{proposition}

\begin{proof}
(i) Immediate from \cref{def:offspring-pgf} and the standing assumption \(p_0>0\): \(f(0)=p_0\) and \(f(1)=\sum_k p_k=1\).

(ii) The power series of \(f\) has nonnegative coefficients and converges at \(s=1\), hence on the whole of \([0,1]\). Termwise differentiation on \([0,1)\) is \citep[Theorem~A.5.2]{durrett2010}, or monotone convergence of difference quotients. Nonnegative coefficients give \(f'\ge 0\) and \(f''\ge 0\) on \([0,1)\), so \(f\) is nondecreasing and convex there; continuity extends both properties to the compact interval \([0,1]\).

(iii) For \(s\in[0,1)\) one has \(f'(s)=\sum_{k\ge 1} k p_k s^{k-1}=\E\bigl[X(1,1)\,s^{X(1,1)-1}\bigr]\). The integrand increases to \(X(1,1)\) as \(s\uparrow 1\), so monotone convergence yields \(f'(1-)=\E[X(1,1)]=m\).

(iv) The assumption \(p_0+p_1<1\) means \(p_k>0\) for some \(k\ge 2\). That term contributes \(k(k-1)p_k s^{k-2}>0\) for every \(s\in(0,1)\), hence \(f''>0\) on \((0,1)\). A \(C^2\) function with strictly positive second derivative on the open interval is strictly convex on the closed interval.

In particular, \(f(s)=s\) for all \(s\) would force \(p_1=1\), which is forbidden. The difference \(g(s)=f(s)-s\) is strictly convex, so \(g'\) is strictly increasing and has at most one zero; therefore \(g\) has at most two zeros on \([0,1]\). Equivalently, a strictly convex graph meets a straight line at most twice.
\end{proof}

\begin{example}[A random sum, with \(N\) non-random]
\label{ex:pgf-iid-sum}
This is the only algebraic identity needed for (a). If \(Y_1,\dots,Y_k\) are i.i.d.\ copies of \(X(1,1)\) and \(k\) is a deterministic integer, then
\[
\E\bigl[s^{Y_1+\cdots+Y_k}\bigr]
=
\E[s^{Y_1}]\cdots\E[s^{Y_k}]
=
f(s)^k,
\]
and likewise \(\E[Y_1+\cdots+Y_k]=km\). The empty sum (\(k=0\)) has generating function \(s^0=1=f(s)^0\) and mean \(0\). Given \(Z_n\), the next generation is exactly this random sum with \(k=Z_n\).
\end{example}

\begin{proposition}[The four identities in (a)]
\label{prop:gw-four}
Let \(\mathcal{F}_n=\sigma(Z_0,\dots,Z_n)\). Then, for every \(n\ge 0\) and every \(s\in[0,1]\),
\[
\E[Z_{n+1}\mid\mathcal{F}_n]
=
m Z_n,
\qquad
\E[Z_n]
=
m^n,
\]
\[
\E\bigl[s^{Z_{n+1}}\mid\mathcal{F}_n\bigr]
=
f(s)^{Z_n},
\qquad
\E\bigl[s^{Z_n}\bigr]
=
f^{\circ n}(s),
\]
where \(f^{\circ 0}(s)\coloneqq s\) and \(f^{\circ 1}=f\).
\end{proposition}

\begin{proof}
On \(\{Z_n=k\}\) the sum in \cref{def:galton-watson} is a sum of \(k\) i.i.d.\ copies of \(X(1,1)\), independent of \(\mathcal{F}_n\). \cref{ex:pgf-iid-sum} therefore gives
\[
\E[Z_{n+1}\mid\mathcal{F}_n]
=
m Z_n,
\qquad
\E\bigl[s^{Z_{n+1}}\mid\mathcal{F}_n\bigr]
=
f(s)^{Z_n}
\]
almost surely (both sides are \(0\) and \(1\) respectively when \(k=0\)). Taking expectations in the first display and iterating from \(Z_0=1\) yields \(\E[Z_n]=m^n\). Taking expectations in the second display yields
\[
\E\bigl[s^{Z_{n+1}}\bigr]
=
\E\bigl[f(s)^{Z_n}\bigr]
=
f_n\bigl(f(s)\bigr),
\]
so \(f_{n+1}=f_n\circ f\). Starting from \(f_0(s)=s\) (since \(Z_0=1\)) one obtains \(f_n=f^{\circ n}\). This is the iteration written as \(F^{(n)}\) in \citep[(5.34)]{brzezniak1999}.
\end{proof}

\begin{example}[Binary fission, and Galton--Watson's original sex-of-children law]
\label{ex:gw-binomial}
Two concrete offspring laws, both satisfying \(p_0>0\) and \(p_0+p_1<1\).
\begin{enumerate}[label=\textup{(\arabic*)}]
    \item
    \emph{Binary:} \(p_0=p_2=\tfrac12\). Then \(m=1\) and \(f(s)=\tfrac12+\tfrac12 s^2\). The equation \(f(s)=s\) rearranges as \((s-1)^2=0\), so the only root in \([0,1]\) is \(s=1\).
    \item
    Durrett's Exercise~5.3.13: each family has three children, independently male or female with equal probability, and only sons keep the name. Then
    \[
    f(s)
    =
    \Bigl(\frac{1+s}{2}\Bigr)^3
    =
    \frac18+\frac38 s+\frac38 s^2+\frac18 s^3,
    \]
    with \(m=f'(1)=\tfrac32>1\). This is the generating function drawn in \citep[Figure~5.3]{durrett2010}.
    \item
    Brzeźniak's Poisson case, Exercise~5.12(2) and Proposition~5.6: \(p_k=e^{-\lambda}\lambda^k/k!\). Then
    \[
    f(s)
    =
    \exp\bigl(\lambda(s-1)\bigr),
    \qquad
    m=\lambda,
    \]
    and \(f(s)=s\) is the equation \(s=e^{\lambda(s-1)}\) written as (5.28) in the book. If \(\lambda\le 1\) the only root in \([0,1]\) is \(s=1\); if \(\lambda>1\) there is a second root \(\hat r\in(0,1)\). Remark~5.2 of the same book says the same dichotomy holds for a general offspring law, with \(\lambda\) replaced by the mean \(m\).
\end{enumerate}
\end{example}

The same three laws as small family trees. Binary is critical: typical realisations die, a rare fat tree keeps the mean at \(1\). The surname law is supercritical: many realisations still die (no sons), but a surviving name tends to explode. Poisson interpolates by the value of \(\lambda\).

\begin{center}
\begin{tikzpicture}[
    person/.style={circle, draw, inner sep=0.9pt, font=\tiny, minimum size=3.3mm},
    lab/.style={draw=none, font=\scriptsize, inner sep=0.4pt},
    ttl/.style={draw=none, font=\small, inner sep=0.4pt},
    >=Stealth
]
    % binary: two typical outcomes
    \begin{scope}[shift={(0,0.2)}]
        \node[person] (b0) at (0.7,1.15) {};
        \node[lab] at (0.7,0.55) {\(p_0=\tfrac12\): dies};
        \node[person] (b1) at (2.55,1.15) {};
        \node[person] (b11) at (2.15,0.45) {};
        \node[person] (b12) at (2.95,0.45) {};
        \draw[->] (b1) -- (b11);
        \draw[->] (b1) -- (b12);
        \node[lab] at (2.55,-0.05) {\(p_2=\tfrac12\): two children};
        \node[ttl] at (1.6,-0.55) {binary, \(m=1\)};
    \end{scope}

    % three children, sons/daughters
    \begin{scope}[shift={(5.4,0.2)}]
        \node[person, fill=black!18] (t0) at (1.2,1.15) {};
        \node[person, fill=black!18] (t1) at (0.3,0.45) {};
        \node[person] (t2) at (1.2,0.45) {};
        \node[person, fill=black!18] (t3) at (2.1,0.45) {};
        \draw[->] (t0) -- (t1);
        \draw[->] (t0) -- (t2);
        \draw[->] (t0) -- (t3);
        \node[lab] at (1.2,-0.05) {always 3 children; open \(=\) daughter};
        \node[ttl] at (1.2,-0.55) {surnames, \(m=\tfrac32\)};
    \end{scope}
\end{tikzpicture}
\end{center}

\begin{center}
\begin{tikzpicture}[scale=4.2, >=Stealth]
    \draw[->] (0,0) -- (1.12,0) node[right, font=\small] {\(s\)};
    \draw[->] (0,0) -- (0,1.12) node[above, font=\small] {\(f(s)\)};
    \draw[densely dotted] (0,1) -- (1,1);
    \draw[densely dotted] (1,0) -- (1,1);
    \draw[line width=0.6pt] (0,0) -- (1,1) node[right, font=\small] {\(y=s\)};
    \draw[line width=1pt, domain=0:1, samples=80, smooth]
        plot (\x, {0.18 + 0.12*\x + 0.28*\x*\x + 0.42*\x*\x*\x});
    \fill (0,0.18) circle (0.7pt) node[left, font=\small] {\(p_0\)};
    \fill (1,1) circle (0.7pt) node[above right, font=\small] {\((1,1)\)};
    \fill (0.22,0.22) circle (0.7pt) node[below right, font=\small] {\(q\)};
    \node[font=\small] at (0.62,0.38) {\(y=f(s)\)};
\end{tikzpicture}
\end{center}

The sketch is the supercritical picture \(m=f'(1-)>1\): the graph leaves \((1,1)\) with slope larger than \(1\), so it sits below the diagonal on a left neighbourhood of \(1\), and must cross \(y=s\) once more in \([0,1)\) because \(f(0)=p_0>0\). If instead \(m\le 1\), strict convexity plus \(f(1)=1\) keeps the graph strictly above the diagonal on \([0,1)\), and the only crossing is at \(s=1\).

\begin{remark}[How to use the toolkit on Problem~5]
\label{rem:p5-toolkit}
Part~(a) is \cref{prop:gw-four}. Parts~(b)--(c) are the geometry of \cref{prop:pgf-shape} together with the extinction probabilities \(\theta_n\coloneqq\Pbb(Z_n=0)\). Brzeźniak computes this for Poisson offspring in Proposition~5.6, by showing \(\theta_n=f^{\circ n}(0)\) and \(\theta_n\uparrow\hat r\), the smallest root of \(f(s)=s\); Remark~5.2 says the argument does not use the Poisson form. Durrett's Theorem~5.3.9 is the same statement. The exam's \(\tau\) is the first hitting time of \(0\), so \(\{\tau<\infty\}=\{Z_n=0\text{ for some }n\}\); the identification \(\Pbb(\tau<\infty)=\lim_n\theta_n\) is \cref{lem:extinction-nested}.
\end{remark}

\begin{remark}[Hints for (b)--(c), in order]
\label{rem:p5-hints}
\begin{enumerate}[label=\textup{(\arabic*)}]
    \item
    For (b): \(s=1\) is always a root. Use \cref{prop:pgf-shape} to decide whether there is another root in \([0,1)\). The answer depends on whether \(m\le 1\) or \(m>1\); it is not a single integer that is independent of \(m\). The three examples above are the three regimes.
    \item
    For (c): let \(\theta_n=\Pbb(Z_n=0)=f^{\circ n}(0)\). Conditioning on the first generation gives \(\theta_{n+1}=f(\theta_n)\) with \(\theta_0=0\) (or \(\theta_1=p_0\)). The sequence is increasing and bounded, hence converges to a fixed point of \(f\). Which fixed point, and why not the other one? Compare with \(q\). This is the argument of \citep[Proposition~5.6]{brzezniak1999} with \(F\) renamed \(f\).
\end{enumerate}
\end{remark}

\begin{lemma}[Continuity of probability]
\label{lem:prob-continuity}
Let \((A_n)_{n\ge 0}\) be events on a probability space.
\begin{enumerate}[label=\textup{(\roman*)}]
    \item
    If \(A_n\subset A_{n+1}\) for every \(n\), then
    \[
    \Pbb\Bigl(\bigcup_{n\ge 0} A_n\Bigr)
    =
    \lim_{n\to\infty}\Pbb(A_n).
    \]
    \item
    If \(A_n\supset A_{n+1}\) for every \(n\), then
    \[
    \Pbb\Bigl(\bigcap_{n\ge 0} A_n\Bigr)
    =
    \lim_{n\to\infty}\Pbb(A_n).
    \]
\end{enumerate}
This is countable additivity after a disjoint decomposition, not Boole's inequality \(\Pbb(\bigcup A_n)\le\sum\Pbb(A_n)\). The latter is the wrong bound here: the events \(\{Z_n=0\}\) overlap massively, so the sum diverges even when the union has probability strictly less than \(1\).
\end{lemma}

\begin{proof}
For (i), disjointify: set \(B_0\coloneqq A_0\) and \(B_k\coloneqq A_k\setminus A_{k-1}\) for \(k\ge 1\). Then the \(B_k\) are pairwise disjoint, \(A_n=\bigsqcup_{k=0}^n B_k\), and \(\bigcup_n A_n=\bigsqcup_{k\ge 0} B_k\). Countable additivity gives
\[
\Pbb\Bigl(\bigcup_n A_n\Bigr)
=
\sum_{k=0}^\infty \Pbb(B_k)
=
\lim_{n\to\infty}\sum_{k=0}^n \Pbb(B_k)
=
\lim_{n\to\infty}\Pbb(A_n).
\]
Item~(ii) is (i) applied to the complements: \(A_n^c\uparrow\bigl(\bigcap_n A_n\bigr)^c\), and \(\Pbb(A_n^c)=1-\Pbb(A_n)\).
\end{proof}

\begin{lemma}[Extinction is a nested union]
\label{lem:extinction-nested}
State \(0\) is absorbing, so \(\{Z_n=0\}\subset\{Z_{n+1}=0\}\) for every \(n\). The extinction time satisfies
\[
\{\tau<\infty\}
=
\bigcup_{n\ge 0}\{Z_n=0\}
=
\bigcup_{n\ge 0}\{\tau\le n\}.
\]
Consequently \cref{lem:prob-continuity}(i) yields \(\Pbb(\tau<\infty)=\lim_n\theta_n\), where \(\theta_n=\Pbb(Z_n=0)\).
\end{lemma}

\begin{proof}
If \(Z_n=0\), the empty-sum convention in \cref{def:galton-watson} forces \(Z_{n+1}=0\), hence \(Z_m=0\) for all \(m\ge n\). Thus \(\{Z_n=0\}=\{\tau\le n\}\), the events increase, and \(\tau<\infty\) if and only if \(Z_n=0\) for some (equivalently: all large) \(n\).
\end{proof}

\subsubsection{Solution of Problem 5}

\paragraph{Answer.}
∆
(a) For every \(n\ge 0\) and \(s\in[0,1]\),
\[
\E[Z_{n+1}\mid\mathcal{F}_n]=mZ_n,
\qquad
\E[Z_n]=m^n,
\]
\[
\E\bigl[s^{Z_{n+1}}\mid\mathcal{F}_n\bigr]=f(s)^{Z_n},
\qquad
\E\bigl[s^{Z_n}\bigr]=f^{\circ n}(s),
\]
with \(\mathcal{F}_n=\sigma(Z_0,\dots,Z_n)\), as in \cref{prop:gw-four}.

(b) The equation \(f(s)=s\) always has \(s=1\) as a solution. If \(m\le 1\), it is the \emph{only} solution in \([0,1]\). If \(m>1\), there is exactly one further solution \(q\in(0,1)\) (the smallest root). The answer is therefore not a fixed integer independent of \(m\): one solution when \(m\le 1\), two when \(m>1\).

(c) Writing \(q\) for the smallest root of \(f(s)=s\) on \([0,1]\),
\[
\Pbb(\tau<\infty)
=
\begin{cases}
1, & m\le 1,\\
q, & m>1.
\end{cases}
\]

\begin{proof}[Exam proof]
\textbf{(a).}
This is \cref{prop:gw-four}; the four identities are the same computation written twice, once for means and once for generating functions. By \cref{def:galton-watson},
\[
Z_{n+1}
=
\sum_{j=1}^{Z_n} X(j,n),
\]
with the empty sum equal to \(0\). The summand \(X(j,n)\) is the number of children of the \(j\)th individual in generation \(n\). These variables are i.i.d.\ copies of \(X(1,1)\), and the whole array \(\{X(j,n): j\ge 1\}\) is independent of \(\mathcal{F}_n=\sigma(Z_0,\dots,Z_n)\). The only difficulty is that the number of summands is itself random.

Fix an integer \(k\ge 0\). On the event \(\{Z_n=k\}\) the upper limit is deterministic, so
\[
Z_{n+1}
=
X(1,n)+\cdots+X(k,n)
\qquad\text{on }\{Z_n=k\}.
\]
This is the fixed-length i.i.d.\ sum of \cref{ex:pgf-iid-sum}. The mean of \(k\) copies is \(km\), hence \(\E[Z_{n+1}\mid Z_n=k]=km\). Since this holds for every \(k\),
\[
\E[Z_{n+1}\mid\mathcal{F}_n]
=
m Z_n
\]
almost surely (both sides are \(0\) on \(\{Z_n=0\}\)). Taking expectations, \(\E[Z_{n+1}]=m\,\E[Z_n]\). Starting from \(Z_0=1\) and iterating gives \(\E[Z_n]=m^n\).

For the generating function, independence turns a sum into a product. On \(\{Z_n=k\}\),
\[
\E\bigl[s^{Z_{n+1}}\bigm|Z_n=k\bigr]
=
\E\bigl[s^{X(1,n)+\cdots+X(k,n)}\bigr]
=
\E[s^{X(1,1)}]^k
=
f(s)^k,
\]
with the convention \(f(s)^0=1\) when \(k=0\) (the empty population contributes \(s^0=1\)). Therefore
\[
\E\bigl[s^{Z_{n+1}}\mid\mathcal{F}_n\bigr]
=
f(s)^{Z_n}
\]
almost surely. Taking expectations replaces the random exponent \(Z_n\) by an ordinary expectation:
\[
\E\bigl[s^{Z_{n+1}}\bigr]
=
\E\bigl[f(s)^{Z_n}\bigr].
\]
The right-hand side is the generating function of \(Z_n\), evaluated at the number \(f(s)\) rather than at \(s\). Writing \(g_n(u)\coloneqq\E[u^{Z_n}]\) and \(f^{\circ 0}(s)\coloneqq s\), \(f^{\circ (n+1)}\coloneqq f^{\circ n}\circ f\), the previous display is the composition identity
\[
g_{n+1}(s)
=
g_n\bigl(f(s)\bigr),
\qquad n\ge 0.
\]
The claim \(g_n=f^{\circ n}\) is now ordinary induction on \(n\).

\emph{Base step.} \(Z_0=1\), so \(g_0(s)=\E[s^1]=s=f^{\circ 0}(s)\).

\emph{Inductive step.} Assume \(g_n=f^{\circ n}\). Then
\[
g_{n+1}(s)
=
g_n\bigl(f(s)\bigr)
=
f^{\circ n}\bigl(f(s)\bigr)
=
\bigl(f^{\circ n}\circ f\bigr)(s)
=
f^{\circ(n+1)}(s).
\]
The first two iterates are worth writing out: \(g_1(s)=g_0(f(s))=f(s)\), and \(g_2(s)=g_1(f(s))=f(f(s))\). After \(n\) steps the argument of \(g_0\) has been replaced by \(f\) exactly \(n\) times, which is \(f^{\circ n}(s)\). Therefore, for every \(n\ge 0\) and every \(s\in[0,1]\),
\[
\E\bigl[s^{Z_n}\bigr]
=
g_n(s)
=
f^{\circ n}(s).
\]
Collecting the four identities: \(\E[Z_{n+1}\mid\mathcal{F}_n]=mZ_n\), \(\E[Z_n]=m^n\), \(\E[s^{Z_{n+1}}\mid\mathcal{F}_n]=f(s)^{Z_n}\), and \(\E[s^{Z_n}]=f^{\circ n}(s)\). That is (a).

\textbf{(b).}
Always \(f(1)=\sum_k p_k=1\), so \(s=1\) is a root. By \cref{prop:pgf-shape}, \(f\) is strictly convex on \([0,1]\) and \(f'(1-)=m\). The number of roots is the number of crossings of \(y=f(s)\) with the diagonal \(y=s\), and the slope at the common point \((1,1)\) decides the picture.

\begin{center}
\begin{tikzpicture}[scale=3.2, >=Stealth]
    % left: m \le 1
    \begin{scope}[shift={(0,0)}]
        \draw[->] (0,0) -- (1.18,0) node[right, font=\small] {\(s\)};
        \draw[->] (0,0) -- (0,1.18) node[above, font=\small] {\(f(s)\)};
        \draw[densely dotted] (0,1) -- (1,1);
        \draw[densely dotted] (1,0) -- (1,1);
        \draw[line width=0.55pt] (0,0) -- (1,1);
        \node[font=\scriptsize, rotate=45, anchor=south west, inner sep=1pt] at (0.42,0.42) {\(y=s\)};
        \draw[line width=1.05pt, domain=0:1, samples=80, smooth]
            plot (\x, {0.40 + 0.35*\x + 0.20*\x*\x + 0.05*\x*\x*\x});
        \fill (0,0.40) circle (0.85pt) node[left, font=\small] {\(p_0\)};
        \fill (1,1) circle (0.85pt) node[above right, font=\small] {\((1,1)\)};
        \draw[line width=0.65pt, densely dashed]
            (0.70, {1-0.85*0.30}) -- (1.14, {1+0.85*0.14});
        \node[font=\scriptsize, anchor=west] at (1.02,0.78) {\(m\le 1\)};
        \node[font=\small] at (0.48,0.18) {\(y=f(s)\)};
        \node[font=\small] at (0.50,-0.40) {\(m\le 1\): one root};
        \node[font=\scriptsize] at (0.50,-0.64) {\(f(s)>s\) on \([0,1)\)};
    \end{scope}
    % right: m > 1
    \begin{scope}[shift={(2.55,0)}]
        \draw[->] (0,0) -- (1.18,0) node[right, font=\small] {\(s\)};
        \draw[->] (0,0) -- (0,1.18) node[above, font=\small] {\(f(s)\)};
        \draw[densely dotted] (0,1) -- (1,1);
        \draw[densely dotted] (1,0) -- (1,1);
        \draw[line width=0.55pt] (0,0) -- (1,1);
        \node[font=\scriptsize, rotate=45, anchor=south west, inner sep=1pt] at (0.58,0.58) {\(y=s\)};
        \draw[line width=1.05pt, domain=0:1, samples=80, smooth]
            plot (\x, {0.18 + 0.12*\x + 0.28*\x*\x + 0.42*\x*\x*\x});
        \fill (0,0.18) circle (0.85pt) node[left, font=\small] {\(p_0\)};
        \fill (1,1) circle (0.85pt) node[above right, font=\small] {\((1,1)\)};
        \fill (0.22,0.22) circle (0.85pt) node[below right, font=\small] {\(q\)};
        \draw[densely dotted] (0.22,0) -- (0.22,0.22) -- (0,0.22);
        \draw[line width=0.65pt, densely dashed]
            (0.78, {1-1.95*0.22}) -- (1.10, {1+1.95*0.10});
        \node[font=\scriptsize, anchor=west] at (1.02,1.16) {\(m>1\)};
        \node[font=\small] at (0.62,0.34) {\(y=f(s)\)};
        \node[font=\small] at (0.50,-0.40) {\(m>1\): two roots};
        \node[font=\scriptsize] at (0.50,-0.64) {\(q\in(0,1)\) and \(s=1\)};
    \end{scope}
\end{tikzpicture}
\end{center}

Left: the graph leaves \((1,1)\) no steeper than the diagonal, so strict convexity keeps it strictly above \(y=s\) on \([0,1)\). (The critical case \(m=1\) is the same picture with the dashed slope coinciding with the diagonal at \((1,1)\).) Right: the graph leaves \((1,1)\) steeper than the diagonal, so it dips below \(y=s\) just to the left of \(1\); together with \(f(0)=p_0>0\) this forces a second crossing \(q\).

If \(m\le 1\), then for every \(s\in[0,1)\) strict convexity and the endpoint values \(f(0)=p_0>0\), \(f(1)=1\) force the graph of \(f\) to lie strictly above the diagonal \(y=s\). (Equivalently, for \(s<1\) one has \(f(s)>s\): otherwise the secant from \((0,p_0)\) to \((1,1)\) would contradict \(f'(1-)\le 1\).) Hence no \(s\in[0,1)\) satisfies \(f(s)=s\).

If \(m>1\), then \(f'(1-)>1\), so \(f(s)<s\) for \(s\) in a left neighbourhood of \(1\). Since \(f(0)=p_0>0\), the intermediate-value theorem gives at least one \(q\in(0,1)\) with \(f(q)=q\). Strict convexity allows at most two crossings of \(y=s\), and \(1\) is already one of them, so there is exactly one additional root \(q\in(0,1)\), necessarily the smallest.

\textbf{(c).}
Let \(\theta_n\coloneqq\Pbb(Z_n=0)\). Because \(Z_0=1\), one has \(\theta_0=0\). For \(n\ge 0\),
\[
\theta_{n+1}
=
\Pbb(Z_{n+1}=0)
=
\E\bigl[\Pbb(Z_{n+1}=0\mid Z_n)\bigr]
=
\E\bigl[p_0^{Z_n}\bigr]
=
f(\theta_n),
\]
using \(\Pbb(\text{all offspring }0\mid Z_n=k)=p_0^k\) and \(\E[s^{Z_n}]=f^{\circ n}(s)\) at \(s=p_0\), or directly \(f(\theta_n)=\sum_k \Pbb(Z_n=k)p_0^k\). Thus \(\theta_n=f^{\circ n}(0)\).

One extra step is needed before passing from \(\theta_n\) to \(\Pbb(\tau<\infty)\). State \(0\) is absorbing: if \(Z_n=0\), the empty-sum convention in \cref{def:galton-watson} forces \(Z_{n+1}=0\), hence \(Z_m=0\) for all \(m\ge n\). Therefore \(\{Z_n=0\}\subset\{Z_{n+1}=0\}\), and
\[
\{\tau<\infty\}
=
\bigcup_{n\ge 0}\{Z_n=0\}
=
\bigcup_{n\ge 0}\{\tau\le n\}.
\]
The events on the right are nested and increasing, so \cref{lem:prob-continuity}(i) (equivalently \cref{lem:extinction-nested}) yields \(\Pbb(\tau<\infty)=\lim_n\theta_n\). In particular \((\theta_n)\) is nondecreasing and bounded by \(1\), hence \(\theta_n\uparrow \theta_\infty\) for some \(\theta_\infty\in[0,1]\), and \(\Pbb(\tau<\infty)=\theta_\infty\). Continuity of \(f\) on \([0,1]\) gives \(\theta_\infty=f(\theta_\infty)\), so \(\theta_\infty\) is a fixed point of \(f\) in \([0,1]\).

If \(m\le 1\), part~(b) shows that the only fixed point is \(1\). Therefore \(\theta_\infty=1\), and \(\Pbb(\tau<\infty)=\lim_n\theta_n=1\).

If \(m>1\), part~(b) supplies exactly two roots of \(f(s)=s\) in \([0,1]\): the smallest root \(q\in(0,1)\), and \(s=1\). Thus \(\theta_\infty\in\{q,1\}\). It remains to show that the iteration started at \(0\) selects \(q\), not \(1\).

By the right-hand picture in (b), and strict convexity,
\[
f(s)>s
\quad\text{on }[0,q),
\qquad
f(q)=q,
\qquad
f(s)<s
\quad\text{on }(q,1).
\]
Also \(f\) is strictly increasing on \([0,1]\) (\cref{prop:pgf-shape}: \(f''>0\) on \((0,1)\) forces \(f'\) strictly increasing, hence \(f'>0\) on \((0,1]\)).

The claim is \(0=\theta_0<\theta_1<\theta_2<\cdots<q\). Indeed \(\theta_0=0<q\). If \(\theta_n\in[0,q)\), then \(\theta_{n+1}=f(\theta_n)>\theta_n\) because \(f(s)>s\) on \([0,q)\), while \(\theta_{n+1}=f(\theta_n)<f(q)=q\) because \(f\) is strictly increasing. Starting from \(\theta_0=0\) one therefore obtains an increasing sequence trapped in \([0,q)\), so \(\theta_\infty\le q\). The only fixed point in \([0,q]\) is \(q\), hence \(\theta_\infty=q\). In particular the orbit never enters \((q,1]\), so it cannot converge to the other fixed point \(1\). (The inequality \(f(s)<s\) on \((q,1)\) is the companion fact that \(1\) is repulsive from the left: any seed in \((q,1)\) is driven down towards \(q\), not up to \(1\).)

\begin{center}
\begin{tikzpicture}[scale=3.35, >=Stealth]
    \draw[->] (0,0) -- (1.18,0) node[right, font=\small] {\(s\)};
    \draw[->] (0,0) -- (0,1.18) node[above, font=\small] {\(f(s)\)};
    \draw[line width=0.5pt] (0,0) -- (1,1);
    \node[font=\scriptsize, rotate=45, anchor=south west, inner sep=1pt] at (0.58,0.58) {\(y=s\)};
    \draw[line width=1.05pt, domain=0:1, samples=80, smooth]
        plot (\x, {0.18 + 0.12*\x + 0.28*\x*\x + 0.42*\x*\x*\x});
    \node[font=\small] at (0.72,0.38) {\(y=f(s)\)};
    % cobweb: 0 -> 0.18 -> 0.213 -> 0.222 \approx q
    \draw[line width=0.7pt, red!70!black]
        (0,0) -- (0,0.18) -- (0.18,0.18) -- (0.18,0.213)
        -- (0.213,0.213) -- (0.213,0.222) -- (0.222,0.222);
    \fill (0,0) circle (0.8pt) node[below left, font=\small] {\(\theta_0=0\)};
    \fill (0.18,0.18) circle (0.7pt) node[below right, font=\scriptsize] {\(\theta_1\)};
    \fill (0.22,0.22) circle (0.85pt) node[above left, font=\small] {\(q\)};
    \fill (1,1) circle (0.85pt) node[above right, font=\small] {\(1\)};
    \draw[densely dotted] (0.22,0) -- (0.22,0.22) -- (0,0.22);
    \node[font=\small] at (0.50,-0.42) {iteration from \(0\) is trapped below \(q\)};
\end{tikzpicture}
\end{center}

Therefore \(\Pbb(\tau<\infty)=\theta_\infty=q\).
\end{proof}

\begin{remark}[Sanity checks]
\label{rem:p5-checks}
Part~(a) is pure conditioning plus \cref{ex:pgf-iid-sum}. Part~(b) is a picture of strict convexity against \(y=s\); the mean \(m=f'(1-)\) decides whether the graph enters the region below the diagonal near \(1\). Part~(c) identifies extinction with the attractor of the iteration \(s\mapsto f(s)\) started at \(0\); in the supercritical case that attractor is \(q\), not \(1\), even though \(1\) is also a fixed point.
\end{remark}

\subsection{Problem 6. [15 pts.]}
In the above definition of Galton--Watson branching process, \( Z_n \) can be thought of as the number of descendants in the \( n \)-th generation, and \( X(j, n) \) can be thought of as the number of children of the \( j \)-th of these descendants. By tracking genealogical relationships, we obtain a tree \( T \) rooted at the single individual in generation 0 with a vertex for each individual in the progeny and an edge for each parent-child relationship. \( T \) is called a Galton--Watson tree.

A property of rooted trees is said to be inherited if this property holds for all finite trees, and whenever it holds for a tree, it also holds for all subtrees rooted at the children of the root.
\begin{enumerate}[label=(\alph*)]
    \item Prove that for a Galton--Watson tree \( T \), conditioned on non-extinction, an inherited property \( A \) has a probability of either 0 or 1. (Hint: One may use \( s \le f(s) \) implies \( s \in [0, q] \cup \{1\} \).)
    \item Let \( T \) be the Galton--Watson tree for an offspring distribution with mean \( m > 1 \). Perform percolation on \( T \) with density \( p \) (i.e., each edge of \( T \) is open with probability \( p \) and closed with probability \( 1 - p \), independently of all other edges). Let \( \mathcal{C}_0 \) be the maximal subtree of the root formed by open edges in \( T \). Define the critical probability

    \[
    p_c(T) := \sup\{p \in [0, 1] : \mathbb{P}_p[|\mathcal{C}_0| = \infty \mid T] = 0\},
    \]

    where \( |\mathcal{C}_0| \) is the number of vertices in \( \mathcal{C}_0 \). Conditioned on non-extinction of \( T \), compute \( p_c(T) \).
\end{enumerate}

\setcounter{section}{6}
\setcounter{definition}{0}

\subsubsection{Definitions and background for Problem 6}

The notes below do not solve Problem~6. They add the extra language the exam uses on top of Problem~5: the \emph{tree} (not just generation sizes), \emph{inherited} properties, and Bernoulli \emph{percolation}. Neither \citep{brzezniak1999} nor \citep[\S5.3.4]{durrett2010} treats the family tree; both stop at \(\{Z_n\}\). The nearest previous Qiuzhen item is 2023 fall regular-tree percolation, where the tree is a deterministic \(d\)-regular graph and \(p_c\) is a number; here the tree itself is random, and \(p_c(T)\) is a random variable.

Assume throughout the standing hypotheses of Problem~5, and \(m>1\) whenever the exam conditions on non-extinction. Write \(q\) for the smallest root of \(f(s)=s\) on \([0,1]\), as in Problem~5(c).

\paragraph{Learning path (continue from Problem~5).}
Problem~6 is not a new process. It is the same Galton--Watson mechanism, with the parent pointers kept, plus two extra languages. Read in this order.

\begin{enumerate}[label=\textup{Step~\arabic*.}, start=8]
    \item
    \emph{Keep the edges.} Motivation: \(Z_n\) forgets who begat whom. Percolation is an operation on edges, and ``subtree rooted at a child'' is an operation on vertices. Both require the family tree \(T\) of \cref{def:gw-tree}. Finite \(T\) is exactly extinction.
    \item
    \emph{Recursive split.} Motivation: the branching property \cref{def:branching-property} on the tree is \cref{prop:gw-tree-split}: given \(k\) children, the \(k\) hanging subtrees are i.i.d.\ copies of \(T\). Every generating-function identity in Problem~5 is this split, evaluated on the functional \(\Phi(T)=s^{|T\cap\text{generation }1|}\). Problem~6(a) evaluates the same split on \(\Phi=\mathbf{1}_A\).
    \item
    \emph{Inherited properties.} Motivation: one wants a 0--1 law on the infinite core. The class has to contain every finite tree (otherwise extinction contaminates the conditional probability) and has to pass from a tree to its child subtrees (otherwise the recursive split does not close). That is \cref{def:inherited}. The two extreme examples---all trees, and only finite trees---already achieve conditional probabilities \(1\) and \(0\) given non-extinction. The theorem is that nothing else occurs.
    \item
    \emph{Why condition on non-extinction.} On extinction, \(T\) is finite, every inherited property holds, and \(p_c(T)=1\). Conditioning discards that trivial atom of mass \(q\) and asks what the infinite tree looks like. This is the same reason one studies \(Z_n/m^n\) only in the supercritical case.
    \item
    \emph{Bernoulli percolation as a second branching process.} Motivation: independently deleting edges with probability \(1-p\) is, from the root's point of view, independently discarding children. The annealed open cluster is therefore again Galton--Watson, with thinned offspring law \(f_p(s)=f(1-p+ps)\) and mean \(mp\) \cref{def:thinning}. Criticality of that thinned process is the comparison \(p\) versus \(1/m\). The quenched threshold \(p_c(T)\) is a function of the frozen tree; the exam asks for its conditional law. The event ``\(p\)-percolation from the root dies'' is inherited, so Step~10 applies pathwise in \(p\).
\end{enumerate}

\paragraph{Process versus tree.}
The process \(\{Z_n\}\) only remembers how many people are alive in each generation. The \emph{Galton--Watson tree} \(T\) remembers who is whose child. The two objects determine each other: \(Z_n\) is the number of vertices at graph distance \(n\) from the root, and \(T\) is finite if and only if \(\{Z_n\}\) becomes extinct.

\begin{center}
\begin{tikzpicture}[
    every node/.style={circle, draw, inner sep=1.3pt, font=\scriptsize},
    level distance=11mm,
    sibling distance=16mm,
    edge from parent/.style={draw, -{Stealth[length=1.6mm]}}
]
    \node (root) {\(o\)}
        child { node {\(v_1\)}
            child { node {} }
            child { node {} }
        }
        child { node {\(v_2\)}
            child { node {} }
        }
        child { node {\(v_3\)} };
    \node[draw=none, font=\small, anchor=west] at (2.9,0) {\(T\): one vertex per individual};
    \node[draw=none, font=\small, anchor=west] at (2.9,-1.1) {generation \(1\): \(Z_1=3\) children of the root};
    \node[draw=none, font=\small, anchor=west] at (2.9,-2.2) {dashed boxes: subtrees \(T_1,T_2,T_3\)};
    \draw[dashed, rounded corners, black!50]
        (-1.55,-0.55) rectangle (0.05,-2.55);
    \draw[dashed, rounded corners, black!50]
        (0.35,-0.55) rectangle (1.35,-2.55);
    \node[draw=none, font=\scriptsize] at (-0.75,-2.75) {\(T_1\)};
    \node[draw=none, font=\scriptsize] at (0.85,-2.75) {\(T_2\)};
    \node[draw=none, font=\scriptsize] at (1.85,-0.55) {\(T_3=\{\text{leaf}\}\)};
\end{tikzpicture}
\end{center}

\begin{definition}[Galton--Watson tree]
\label{def:gw-tree}
Let \(\{Z_n\}\) be the Galton--Watson process of \cref{def:galton-watson}. The associated \emph{Galton--Watson tree} \(T\) is the rooted family tree whose vertices in generation \(n\) are the \(Z_n\) individuals, with an edge from each individual to each of its children. Write \(o\) for the root, and \(T_1,\dots,T_{Z_1}\) for the subtrees rooted at the children of \(o\) (an empty list if \(Z_1=0\)). Extinction of \(\{Z_n\}\) is exactly the event that \(T\) is a finite tree.
\end{definition}

The one structural fact used throughout both (a) and (b) is the recursive decomposition, which is the tree-level restatement of ``given \(Z_n=k\), the next generation is \(k\) independent copies of the offspring law''.

\begin{proposition}[Recursive structure]
\label{prop:gw-tree-split}
Conditionally on \(Z_1=k\), the subtrees \(T_1,\dots,T_k\) are i.i.d.\ copies of \(T\), independent of everything else. In particular, for any bounded measurable functional \(\Phi\) of rooted trees,
\[
\E\bigl[\Phi(T)\bigr]
=
\E\Bigl[\,
\Psi_k\bigl(\Phi(T_1),\dots,\Phi(T_k)\bigr)
\Bigm| Z_1=k
\Bigr]
\]
reduces to an average against the offspring law \((p_k)\). Taking \(\Phi=\mathbf{1}_A\) for an event \(A\) that depends on \(T\) only through whether the child subtrees lie in \(A\) produces a generating-function identity of the shape \(\E[s^{Z_1}]=f(s)\), as in Problem~5.
\end{proposition}

\paragraph{Inherited properties.}
The exam defines this in the statement of the problem. Restated:

\begin{definition}[Inherited property]
\label{def:inherited}
A property \(A\) of rooted trees is \emph{inherited} if
\begin{enumerate}[label=\textup{(\roman*)}]
    \item
    every finite tree has \(A\);
    \item
    whenever a tree has \(A\), every subtree rooted at a child of the root also has \(A\).
\end{enumerate}
\end{definition}

Item (i) forces extinction \(\subset A\), hence \(\Pbb(T\in A)\ge q\). Item (ii) is a downward closure: \(A\) cannot hold at the root unless it holds at every child. Equivalently, the complementary event is \emph{upward-closed}: if some child subtree fails \(A\), then the whole tree fails \(A\).

Two examples that satisfy the definition, and one that does not:
\begin{enumerate}[label=\textup{(\arabic*)}]
    \item
    \(A=\{\text{all rooted trees}\}\). Then \(\Pbb(T\in A\mid\text{non-extinction})=1\).
    \item
    \(A=\{\text{finite trees}\}\). Then \(\Pbb(T\in A\mid\text{non-extinction})=0\). This is exactly extinction.
    \item
    Non-example: ``the root has an even number of children''. Finite trees need not have this, so (i) fails; and even if the root has even degree, a child subtree need not.
\end{enumerate}

The exam's (a) is the statement that, on a Galton--Watson tree, there are no other possibilities after conditioning on non-extinction: an inherited \(A\) behaves like one of the two examples above. The algebraic input is the same random-sum identity as in Problem~5.

\begin{lemma}[Generating-function bound]
\label{lem:inherited-pgf}
If \(A\) is inherited and \(\eta\coloneqq\Pbb(T\in A)\), then
\[
\eta
\le
f(\eta).
\]
\end{lemma}

\begin{proof}
By \cref{def:inherited}(ii), \(\{T\in A\}\subset\{T_1\in A,\dots,T_{Z_1}\in A\}\) (vacuous if \(Z_1=0\)). The \(T_i\) are i.i.d.\ copies of \(T\) given \(Z_1\) by \cref{prop:gw-tree-split}, so
\[
\Pbb(T_1\in A,\dots,T_{Z_1}\in A)
=
\E\bigl[\eta^{Z_1}\bigr]
=
f(\eta).
\]
\end{proof}

Together with \(\eta\ge q\) and the geometry of \(f\) on \([q,1]\) already drawn for Problem~5, this is the only input (a) needs. The remaining step is to pass from the unconditional probability \(\eta\) to the conditional probability given non-extinction.

\paragraph{Bernoulli percolation.}
The 2023 fall QE ran the same experiment on a \emph{fixed} infinite \(d\)-regular tree: each edge is kept independently with probability \(p\). Here the underlying graph is the random tree \(T\).

\begin{center}
\begin{tikzpicture}[
    every node/.style={circle, draw, inner sep=1.4pt, font=\scriptsize},
    level distance=11mm,
    sibling distance=15mm
]
    \node[fill=black!12] (o) {\(o\)}
        child { node[fill=black!12] {\(v_1\)}
            child { node[fill=black!12] {} }
            child { node {} }
            edge from parent[line width=1.1pt]
        }
        child { node {\(v_2\)}
            child { node {} }
            edge from parent[dashed, thick]
        }
        child { node[fill=black!12] {\(v_3\)}
            edge from parent[line width=1.1pt]
        };
    \node[draw=none, font=\small, anchor=west] at (2.7,0) {solid = open (prob.\ \(p\))};
    \node[draw=none, font=\small, anchor=west] at (2.7,-0.9) {dashed = closed (prob.\ \(1-p\))};
    \node[draw=none, font=\small, anchor=west] at (2.7,-1.8) {shaded = \(\mathcal{C}_0\), the open cluster of \(o\)};
\end{tikzpicture}
\end{center}

\begin{definition}[Bernoulli percolation on \(T\)]
\label{def:tree-percolation}
Fix \(p\in[0,1]\). Independently declare each edge of \(T\) \emph{open} with probability \(p\) and \emph{closed} with probability \(1-p\). Two vertices are in the same \emph{open cluster} if they are joined by a path of open edges. Write \(\mathcal{C}_0\) for the open cluster containing the root, and \(\Pbb_p(\,\cdot\mid T)\) for this coin-flip law given the tree (the \emph{quenched} law: \(T\) is frozen, only the edges are random). The \emph{percolation critical value} of that tree is
\[
p_c(T)
\coloneqq
\sup\bigl\{p\in[0,1]:\Pbb_p\bigl(|\mathcal{C}_0|=\infty\mid T\bigr)=0\bigr\}.
\]
\end{definition}

Three remarks on the definition, before any computation.
\begin{enumerate}[label=\textup{(\arabic*)}]
    \item
    \(|\mathcal{C}_0|=\infty\) means there is an infinite open ray from the root. If \(T\) itself is finite this is impossible for every \(p\), so \(p_c(T)=1\) on the extinction event. The exam asks for the law of \(p_c(T)\) after conditioning on non-extinction.
    \item
    For each frozen tree \(t\), the map \(p\mapsto\Pbb_p(|\mathcal{C}_0|=\infty\mid t)\) is nondecreasing, so the set in the supremum is an initial interval of \([0,1]\) and \(p_c(t)\) is well-defined in \([0,1]\). On a deterministic \(d\)-regular tree the 2023 fall exam asked for a single number \(p_c\); here \(p_c(T)\) may a priori depend on the realisation of \(T\).
    \item
    There are two \(\sigma\)-algebras in play. The \emph{quenched} probability \(\Pbb_p(\,\cdot\mid T)\) treats \(T\) as given. The \emph{annealed} probability \(\Pbb_p\) averages over both the tree and the edge coins. They are related by
    \[
    \Pbb_p\bigl(|\mathcal{C}_0|=\infty\bigr)
    =
    \E\bigl[\Pbb_p\bigl(|\mathcal{C}_0|=\infty\mid T\bigr)\bigr].
    \]
    In particular, if the annealed probability is zero then the quenched probability is zero for almost every tree; if the annealed probability is strictly positive then the event \(\{\Pbb_p(|\mathcal{C}_0|=\infty\mid T)=0\}\) cannot have probability one.
\end{enumerate}

\paragraph{Thinning.}
Unconditionally, the individuals that remain connected to the root by an open path form a new branching process. Each person still has \(X\sim(p_k)\) children in \(T\), but the edge to a given child is open with probability \(p\), independently. The number of \emph{kept} children is therefore a \(p\)-thinning of \(X\):

\begin{definition}[\(p\)-thinning]
\label{def:thinning}
If \(X\) is a nonnegative integer-valued random variable and, given \(X=k\), one keeps each of \(k\) items independently with probability \(p\), the number kept is a \(p\)-thinning of \(X\). Its generating function is
\[
f_p(s)
\coloneqq
\E\bigl[(1-p+ps)^{X}\bigr]
=
f(1-p+ps),
\]
with mean \(m p\). (Check: differentiate at \(s=1\), or condition on \(X\).)
\end{definition}

Thus, under the annealed law, \(|\mathcal{C}_0|=\infty\) if and only if this thinned Galton--Watson process survives. The extinction criterion is the one already set up in Problem~5, applied to \(f_p\) rather than to \(f\). That is the only input (b) needs from the branching-process calculus; the rest is to pass from the annealed statement to a statement about the random variable \(p_c(T)\), using (a).

One inherited property that appears naturally in (b): for each fixed \(p\), write
\[
A_p
\coloneqq
\bigl\{\text{rooted trees }t:\Pbb_p\bigl(|\mathcal{C}_0|=\infty\mid t\bigr)=0\bigr\}
\]
for the collection of trees on which \(p\)-percolation from the root dies out almost surely. Finite trees lie in \(A_p\). If a child subtree \(t_i\) failed to lie in \(A_p\), the open cluster of that child inside \(t_i\) would be infinite with positive quenched probability; the edge from the root to that child is open with probability \(p\), and on that event the child's cluster sits inside \(\mathcal{C}_0\). So \(A_p\) is inherited in the sense of \cref{def:inherited} (the case \(p=0\) is vacuous: \(A_0\) contains every tree). Part~(a), once proved, therefore applies to each \(A_p\).

\begin{remark}[What to have in hand]
\label{rem:p6-toolkit}
From Problem~5: the offspring generating function \(f\), the mean \(m\), the smallest root \(q\) of \(f(s)=s\), and the geometry of \(f\) on \([q,1]\). From the notes above: the recursive split of \(T\) into i.i.d.\ child subtrees; the definition of inherited; the bound \(\eta\le f(\eta)\); quenched versus annealed percolation; and the identification of the annealed cluster with a thinned branching process of mean \(mp\).
\end{remark}

\subsubsection{Solution of Problem 6}

\paragraph{Answer.}

Assume \(m>1\) and write \(q\in(0,1)\) for the smallest root of \(f(s)=s\), as in Problem~5(c).

(a) For every inherited property \(A\),
\[
\Pbb(T\in A\mid T\text{ is infinite})
\in\{0,1\}.
\]
(Equivalently, conditioning on non-extinction \(\{Z_n>0\ \forall n\}\), which has probability \(1-q\), the conditional probability is \(0\) or \(1\).)

(b) Almost surely on that non-extinction event,
\[
p_c(T)=\frac{1}{m}.
\]

\begin{proof}[Exam proof]
\textbf{(a).}
Let \(\eta\coloneqq\Pbb(T\in A)\). \Cref{def:inherited}(i) says every finite tree has \(A\). Extinction means \(T\) is finite, and occurs with probability \(q\) by Problem~5(c), so \(\eta\ge q\). \Cref{lem:inherited-pgf} gives \(\eta\le f(\eta)\).

Since \(m>1\), Problem~5(b) gives a smallest fixed point \(q\in(0,1)\) with \(f(q)=q\), and strict convexity (\cref{prop:pgf-shape}) yields
\[
f(s)<s
\qquad\text{for every }s\in(q,1),
\qquad
f(s)>s
\qquad\text{for every }s\in[0,q).
\]
If \(\eta\in(q,1)\), then \(f(\eta)<\eta\), contradicting \(\eta\le f(\eta)\). Hence \(\eta\notin(q,1)\). Because \(\eta\ge q\), the only possibilities are \(\eta=q\) or \(\eta=1\).

Write \(N\) for the non-extinction event. Then \(\Pbb(N)=1-q\) and \(\{T\in A\}\supset\{T\text{ finite}\}\), so
\[
\Pbb(T\in A, N)
=
\Pbb(T\in A)-\Pbb(T\in A,\ T\text{ finite})
=
\eta-q.
\]
Therefore
\[
\Pbb(T\in A\mid N)
=
\frac{\eta-q}{1-q}
=
\begin{cases}
0, & \eta=q,\\
1, & \eta=1.
\end{cases}
\]

\textbf{(b).}
Fix \(p\in[0,1]\) and define
\[
A_p
\coloneqq
\bigl\{t:\ \Pbb_p\bigl(|\mathcal{C}_0|=\infty\mid t\bigr)=0\bigr\},
\]
the set of trees on which \(p\)-percolation from the root dies out almost surely. Finite trees lie in \(A_p\). If some child subtree \(t_i\notin A_p\), then with positive quenched probability that child has an infinite open cluster; when the root edge to that child is open (probability \(p\)), that cluster lies inside \(\mathcal{C}_0\), so \(t\notin A_p\). Thus \(A_p\) is inherited.

By~(a), \(\Pbb(T\in A_p\mid N)\in\{0,1\}\) for each fixed \(p\).

\emph{Annealed comparison.}
Under the annealed law, keep each child independently with probability \(p\). The number of kept children is a \(p\)-thinning of the offspring variable, with generating function \(f_p(s)=f(1-p+ps)\) and mean \(mp\) (\cref{def:thinning}). The open cluster from the root is infinite if and only if this thinned Galton--Watson process survives. By Problem~5(c) applied to \(f_p\), survival has positive probability exactly when \(mp>1\), i.e.\ when \(p>1/m\).

If \(p<1/m\), then \(\Pbb_p(|\mathcal{C}_0|=\infty)=0\), so \(\Pbb_p(T\in A_p^c)=0\) and hence \(\Pbb(T\in A_p\mid N)=1\). Thus \(p_c(T)\ge 1/m\) almost surely on \(N\).

If \(p>1/m\), then \(\Pbb_p(|\mathcal{C}_0|=\infty)>0\), so \(\Pbb(T\in A_p\mid N)=0\). For each fixed tree \(t\), the map \(r\mapsto \Pbb_r(|\mathcal{C}_0|=\infty\mid t)\) is nondecreasing, so the set \(\{r:\Pbb_r(|\mathcal{C}_0|=\infty\mid t)=0\}\) is an initial interval of \([0,1]\) with supremum \(p_c(t)\). Having \(\Pbb(T\in A_p\mid N)=0\) for all \(p>1/m\) forces \(p_c(T)\le 1/m\) almost surely on \(N\).

Combining the two inequalities, \(p_c(T)=1/m\) almost surely conditioned on non-extinction.
\end{proof}

\begin{remark}[Sanity checks]
\label{rem:p6-checks}
Part~(a) is the same fixed-point squeeze as Problem~5(c), with \(\eta\le f(\eta)\) in place of \(\theta_{n+1}=f(\theta_n)\). Part~(b) separates the quenched threshold \(p_c(T)\) from the annealed critical value \(1/m\): thinning reduces percolation to another Galton--Watson extinction problem, and the 0--1 law from~(a) upgrades the annealed dichotomy to a single tree-by-tree threshold. On extinction, \(p_c(T)=1\) because percolation cannot be infinite on a finite tree; that trivial case is exactly why the exam conditions on non-extinction.
\end{remark}

\subsection{Problem 7. [10 pts.]}
Let \( X_1, X_2, \dots, X_n \) be i.i.d.\ observations with pdf

\[
f(x; \theta) = \frac{1}{2\theta\sqrt{x}} e^{-\sqrt{x}/\theta} I_{\{x>0\}},
\]

where \( \theta > 0 \) is the unknown parameter.
\begin{enumerate}[label=(\alph*)]
    \item Show that the family \( f(x_1, \dots, x_n; \theta) \) (\( \theta > 0 \)) has a monotone likelihood ratio in a certain statistic \( T \).
    \item Find a size \( \alpha \) (\( 0 < \alpha < 1 \)) uniformly most powerful test for \( H_0 : \theta \ge 1 \) versus \( H_1 : \theta < 1/2 \). Make your rejection region expressed as explicitly as possible.
\end{enumerate}

\setcounter{section}{7}
\setcounter{definition}{0}

\subsubsection{Definitions from Hogg for Problem 7}

Testing vocabulary from \citep[\S8.1--\S8.2]{hogg2019}. The exam-level solution is assembled after the hints.

Write \(\mathbf{X}=(X_1,\dots,X_n)\) for the sample, \(\mathbf{x}\) for a realised point, and
\[
L(\theta,\mathbf{x})
\coloneqq
\prod_{i=1}^{n} f(x_i;\theta)
\]
for the likelihood. A (nonrandomised) test is identified with its critical region \(C\): reject \(H_0\) if \(\mathbf{X}\in C\). The power function is \(\gamma_C(\theta)\coloneqq\Pbb_{\theta}(\mathbf{X}\in C)\).

\begin{definition}[Size / significance level]
\label{def:test-size}
Let \(H_0:\theta\in\omega_0\) versus \(H_1:\theta\in\omega_1\), with \(\omega_0\cap\omega_1=\emptyset\). The \emph{size} (also called the \emph{significance level}) of the test with critical region \(C\) is
\[
\alpha
\coloneqq
\max_{\theta\in\omega_0}\Pbb_{\theta}(\mathbf{X}\in C)
=
\max_{\theta\in\omega_0}\gamma_C(\theta).
\]
This is \citep[(8.1.3)]{hogg2019}: the size is the worst Type~I error over the null. When \(\omega_0=\{\theta'\}\) is a singleton the maximum is just \(\Pbb_{\theta'}(\mathbf{X}\in C)\).
\end{definition}

\begin{remark}[Type~I, Type~II, and the power function]
\label{rem:type-i-ii-power}
The same critical region \(C\) produces two kinds of mistake, recorded as Hogg's \(2\times 2\) decision table \citep[Table~4.5.1]{hogg2019}. A Type~I error is a false rejection: \(H_0\) is true and \(\mathbf{X}\in C\). A Type~II error is a missed detection: \(H_1\) is true and \(\mathbf{X}\notin C\).

\begin{center}
\renewcommand{\arraystretch}{1.28}
\begin{tabular}{@{}lcc@{}}
\toprule
 & \multicolumn{2}{c}{True state of nature} \\
\cmidrule(lr){2-3}
Decision & \(H_0\) true \((\theta\in\omega_0)\) & \(H_1\) true \((\theta\in\omega_1)\) \\
\midrule
Reject \(H_0\) \((\mathbf{X}\in C)\)
  & Type~I error
  & Correct (power) \\
Accept \(H_0\) \((\mathbf{X}\notin C)\)
  & Correct
  & Type~II error \\
\bottomrule
\end{tabular}
\end{center}

The power function already written above, \(\gamma_C(\theta)=\Pbb_{\theta}(\mathbf{X}\in C)\), is the rejection probability as a function of the true parameter. Hogg writes the same formula on the alternative only \citep[(4.5.5), (8.1.4)]{hogg2019}; extending the domain to all of \(\omega_0\cup\omega_1\) is the useful convention, because then size is just the worst value of \(\gamma_C\) on the null. At a point of \(H_1\) one also writes \(\beta_C(\theta)\coloneqq 1-\gamma_C(\theta)\) for the Type~II probability. The four numbers are then

\begin{center}
\renewcommand{\arraystretch}{1.32}
\begin{tabular}{@{}llll@{}}
\toprule
Object & Formula & Meaning & Want \\
\midrule
\(\gamma_C(\theta)\)
  & \(\Pbb_{\theta}(\mathbf{X}\in C)\)
  & reject probability at that \(\theta\)
  & small on \(\omega_0\), large on \(\omega_1\) \\
\(\alpha\)
  & \(\max_{\theta\in\omega_0}\gamma_C(\theta)\)
  & size / significance level / worst Type~I
  & small; we \emph{fix} it \\
\(\beta_C(\theta)\)
  & \(1-\gamma_C(\theta)\), \(\theta\in\omega_1\)
  & Type~II error at that alternative
  & small \\
\(1-\beta_C(\theta)\)
  & \(\gamma_C(\theta)\), \(\theta\in\omega_1\)
  & power at that alternative
  & large \\
\bottomrule
\end{tabular}
\end{center}

Thus \(\alpha\) has several names for the same number \citep[after (4.5.8)]{hogg2019}: size, significance level, maximum Type~I error, and maximum of the power function on \(H_0\). Power at an alternative is the probability the test actually detects that alternative; equivalently, it is one minus Type~II error. Minimising Type~II is the same problem as maximising \(\gamma_C\) on \(\omega_1\).

The two errors cannot be minimised at once: they sit on a seesaw \citep[\S4.5]{hogg2019}. The empty critical region never rejects, so \(\alpha=0\) and \(\beta\equiv 1\); the whole sample space always rejects, so \(\alpha=1\) and \(\beta\equiv 0\). Enlarging \(C\) raises every value of \(\gamma_C\), which simultaneously inflates Type~I and inflates power. The Neyman--Pearson programme therefore treats Type~I as the error one is willing to cap, and then hunts for power: first restrict to tests of size \(\alpha\), then maximise \(\gamma_C\) on \(H_1\). That is the content of \cref{def:best-cr} for a simple pair, and of \cref{def:ump-size-alpha} once \(H_1\) is composite. Type~I is capped because a false rejection is regarded as the worse of the two mistakes; the usual numerical cap is \(\alpha=0.05\), but the theory only needs a fixed \(\alpha\in(0,1)\).
\end{remark}

\begin{definition}[Best critical region of size \(\alpha\)]
\label{def:best-cr}
Let \(H_0:\theta=\theta'\) and \(H_1:\theta=\theta''\) be simple. A set \(C\) is a \emph{best critical region of size \(\alpha\)} if
\begin{enumerate}[label=\textup{(\alph*)}]
    \item
    \(\Pbb_{\theta'}(\mathbf{X}\in C)=\alpha\);
    \item
    whenever \(A\) is another set with \(\Pbb_{\theta'}(\mathbf{X}\in A)=\alpha\), one has \(\Pbb_{\theta''}(\mathbf{X}\in C)\ge\Pbb_{\theta''}(\mathbf{X}\in A)\).
\end{enumerate}
That is: among all tests of exact size \(\alpha\), \(C\) maximises power at the single alternative \(\theta''\). See \citep[Definition~8.1.1]{hogg2019}. The Neyman--Pearson lemma \citep[Theorem~8.1.1]{hogg2019} realises such a \(C\) as a likelihood-ratio region
\[
\frac{L(\theta',\mathbf{x})}{L(\theta'',\mathbf{x})}\le k,
\]
with \(k\) chosen so that the size is \(\alpha\).
\end{definition}

\begin{definition}[UMP test of size \(\alpha\)]
\label{def:ump-size-alpha}
The critical region \(C\) is a \emph{uniformly most powerful (UMP) critical region of size \(\alpha\)} for testing a simple hypothesis \(H_0\) against a composite alternative \(H_1\) if \(C\) is a best critical region of size \(\alpha\) for testing \(H_0\) against \emph{each} simple hypothesis in \(H_1\). The test defined by \(C\) is then a UMP test of significance level \(\alpha\). See \citep[Definition~8.2.1]{hogg2019}.
\end{definition}

In words: one and the same size-\(\alpha\) rejection region must be Neyman--Pearson most powerful at every point of \(H_1\). If the most powerful region for \(\theta=\theta''\) changes when \(\theta''\) moves inside \(H_1\), no UMP test exists. That is the two-sided point-null situation of \cref{prop:ump-point-null}; the textbook instance is the normal mean \citep[Example~8.2.3]{hogg2019}.

When the null itself is composite, as in Problem~7, size is still \cref{def:test-size}. Hogg's one-sided package is
\[
H_0:\theta\le\theta'
\quad\text{versus}\quad
H_1:\theta>\theta',
\]
or the opposite one-sided pair \(H_0:\theta\ge\theta'\) versus \(H_1:\theta<\theta'\), which he records as completely analogous \citep[(8.2.1), (8.2.6)]{hogg2019}. The level is \(\max_{\theta\in\omega_0}\gamma_C(\theta)\); for a monotone family this maximum is attained at the boundary point \(\theta'\).

\begin{definition}[Monotone likelihood ratio]
\label{def:mlr}
The likelihood \(L(\theta,\mathbf{x})\) has \emph{monotone likelihood ratio (mlr)} in the statistic \(y=u(\mathbf{x})\) if, for every \(\theta_1<\theta_2\), the ratio
\[
\frac{L(\theta_1,\mathbf{x})}{L(\theta_2,\mathbf{x})}
\]
is a monotone function of \(y=u(\mathbf{x})\). See \citep[Definition~8.2.2]{hogg2019}.
\end{definition}

The definition does not fix the direction. If the ratio is a \emph{decreasing} function \(g(y)\) of \(y\), then large \(y\) favours large \(\theta\); if it is increasing, reverse every inequality below. Hogg works out the decreasing case in detail, then says the increasing case follows by changing the sense of the inequalities \citep[after Definition~8.2.2]{hogg2019}. Part~(a) of the exam is this definition: exhibit \(T=u(\mathbf{X})\) and the direction. Existence of a UMP test is the next theorem, not the definition.

\begin{theorem}[Karlin--Rubin]
\label{thm:karlin-rubin}
Assume \(L(\theta,\mathbf{x})\) has monotone decreasing likelihood ratio in \(Y=u(\mathbf{X})\), in the sense of \cref{def:mlr}. Fix \(\theta'\) and \(\alpha\in(0,1)\), and choose \(c_Y\) so that \(\Pbb_{\theta'}(Y\ge c_Y)=\alpha\).
\begin{enumerate}[label=\textup{(\roman*)}]
    \item
    For \(H_0:\theta\le\theta'\) versus \(H_1:\theta>\theta'\), the test
    \[
    \text{reject }H_0\text{ if }Y\ge c_Y
    \]
    is UMP of level \(\alpha\). This is the claim Hogg records as \citep[(8.2.3)]{hogg2019} immediately after Definition~8.2.2.
    \item
    For \(H_0:\theta\ge\theta'\) versus \(H_1:\theta<\theta'\), the test
    \[
    \text{reject }H_0\text{ if }Y\le c_Y',
    \]
    with \(c_Y'\) determined by \(\Pbb_{\theta'}(Y\le c_Y')=\alpha\), is UMP of level \(\alpha\). This is the analogue recorded at \citep[(8.2.6)]{hogg2019}.
\end{enumerate}
In both cases \(\gamma(\theta)\) is monotone (nondecreasing in (i), nonincreasing in (ii)), so the size over the composite null equals the Type~I error at the boundary \(\theta'\). The same critical region is Neyman--Pearson most powerful against every simple alternative on the indicated side, hence UMP in the sense of \cref{def:ump-size-alpha}.
\end{theorem}

Hogg does not number this as a theorem and does not use the name Karlin--Rubin; he writes ``we claim that the following test is UMP'' after Definition~8.2.2. The QE syllabus lists it under that name. The argument is the one he sketches there: for any fixed \(\theta''>\theta'\), Neyman--Pearson \citep[Theorem~8.1.1]{hogg2019} rejects for \(L(\theta',\mathbf{x})/L(\theta'',\mathbf{x})\le k\); decreasing mlr rewrites this as \(Y\ge g^{-1}(k)\), a threshold that does not depend on the particular \(\theta''>\theta'\). The opposite one-sided pair is the same argument with the inequalities reversed. For a regular exponential family with increasing \(p(\theta)\), the mlr statistic is \(Y=\sum_i K(X_i)\); that is the unnumbered argument after \citep[Example~8.2.5]{hogg2019}, including \citep[(8.2.5)--(8.2.6)]{hogg2019}.

\begin{remark}[What \(\alpha\) is, and which tail to reject]
\label{rem:ump-mlr-picture}
The letter \(\alpha\) in \cref{thm:karlin-rubin} is the size from \cref{def:test-size}, a probability in \((0,1)\). It is not a cutoff on \(Y\): the test is never ``reject if \(Y\ge\alpha\)''. The number compared with \(Y\) is an \(\alpha\)-quantile \(c_Y\) (or \(c_Y'\)) of the boundary law of \(Y\) at \(\theta=\theta'\). Which tail that quantile cuts is read off the likelihood ratio on each one-sided pair, as follows.

Assume decreasing mlr as in \cref{thm:karlin-rubin}: for every \(\theta_1<\theta_2\),

\[
\frac{L(\theta_1,\mathbf{x})}{L(\theta_2,\mathbf{x})}
=
g_{\theta_1,\theta_2}(y)
\]

is a decreasing function of \(y=u(\mathbf{x})\). Thus as \(y\) grows, the smaller parameter becomes relatively less likely than the larger one.

For \(H_0:\theta\le\theta'\) versus \(H_1:\theta>\theta'\), fix any simple alternative \(\theta''>\theta'\). Neyman--Pearson rejects when

\[
\frac{L(\theta',\mathbf{x})}{L(\theta'',\mathbf{x})}
\le
k.
\]

Here \(\theta'<\theta''\), so the left-hand side is decreasing in \(Y\), and the inequality is a right-tail event \(Y\ge c\). As \(Y\) increases, \(L(\theta')/L(\theta'')\) falls: under the observed sample the boundary null \(\theta'\) becomes relatively less likely than the larger alternative, which is precisely the event on which \(H_0:\theta\le\theta'\) fails. The constraint \(\Pbb_{\theta'}(Y\ge c_Y)=\alpha\) only chooses how far into that right tail one goes.

For \(H_0:\theta\ge\theta'\) versus \(H_1:\theta<\theta'\), fix \(\theta''<\theta'\). Neyman--Pearson still rejects when \(L(\theta')/L(\theta'')\le k\), but now \(\theta''<\theta'\), so decreasing mlr applies to \(L(\theta'')/L(\theta')\) and \(L(\theta')/L(\theta'')\) is increasing in \(Y\). The inequality is therefore a left-tail event \(Y\le c'\). As \(Y\) increases, \(L(\theta')/L(\theta'')\) rises, so the larger parameter \(\theta'\) becomes relatively more likely than the smaller alternative: large \(Y\) supports \(H_0\). The null \(\theta\ge\theta'\) becomes relatively less likely only when \(Y\) is small. The size constraint is now \(\Pbb_{\theta'}(Y\le c_Y')=\alpha\).

Problem~7 is this second pair with \(\theta'=1\) and \(Y=T=\sum_i\sqrt{X_i}\). For any \(\theta''<1\),

\[
\frac{L(1,\mathbf{x})}{L(\theta'',\mathbf{x})}
=
(\theta'')^{n}
\exp\bigl(T(\mathbf{x})(\tfrac{1}{\theta''}-1)\bigr),
\]

and \(\tfrac{1}{\theta''}-1>0\), so the ratio is increasing in \(T\). Neyman--Pearson therefore rejects for small \(T\), which is \cref{thm:karlin-rubin}(ii).

Those two tails are different regions. That is why a two-sided alternative has no UMP; see \cref{prop:ump-point-null}.
\end{remark}

\begin{proposition}[Point nulls: one-sided versus two-sided]
\label{prop:ump-point-null}
Assume the same decreasing mlr as in \cref{thm:karlin-rubin}, and fix \(\theta'\) and \(\alpha\in(0,1)\). Write \(C_+\) for the right-tail region \(\{Y\ge c_Y\}\) of \cref{thm:karlin-rubin}(i) and \(C_-\) for the left-tail region \(\{Y\le c_Y'\}\) of \cref{thm:karlin-rubin}(ii).
\begin{enumerate}[label=\textup{(\roman*)}]
    \item
    For the simple null \(H_0:\theta=\theta'\) versus the one-sided alternative \(H_1:\theta>\theta'\), the test with critical region \(C_+\) is UMP of size \(\alpha\). Versus \(H_1:\theta<\theta'\), the test with critical region \(C_-\) is UMP of size \(\alpha\). The null is now a singleton, so size is just \(\Pbb_{\theta'}(C)=\alpha\); the composite-null strengthening in \cref{thm:karlin-rubin} is not needed.
    \item
    For \(H_0:\theta=\theta'\) versus the two-sided alternative \(H_1:\theta\neq\theta'\), no UMP test of size \(\alpha\) exists.
\end{enumerate}
\end{proposition}

\begin{proof}
Part~(i) is the Neyman--Pearson half of the argument for \cref{thm:karlin-rubin}: against any fixed \(\theta''>\theta'\), decreasing mlr rewrites \(L(\theta')/L(\theta'')\le k\) as the same right tail \(C_+\), independent of the particular \(\theta''>\theta'\). Against any \(\theta''<\theta'\) the same computation produces the left tail \(C_-\). Either region is therefore best at every simple alternative on its side, hence UMP in the sense of \cref{def:ump-size-alpha}. Hogg records the right-sided instance as \citep[Example~8.2.2]{hogg2019}, and notes after \citep[Example~8.2.3]{hogg2019} that a one-sided alternative would restore a UMP test.

For part~(ii), suppose some region \(C\) were UMP of size \(\alpha\) for \(\theta=\theta'\) versus \(\theta\neq\theta'\). Then \(C\) would be a best critical region of size \(\alpha\) for every simple pair \((\theta',\theta'')\) with \(\theta''\neq\theta'\). Taking \(\theta''>\theta'\) forces \(C\) to be a right tail of \(Y\) (up to a \(\Pbb_{\theta'}\)-null set). Taking \(\theta''<\theta'\) forces \(C\) to be a left tail. For \(\alpha\in(0,1)\) these two sets are not equal \(\Pbb_{\theta'}\)-almost surely. Contradiction. This is \citep[Example~8.2.3]{hogg2019}, written for \(N(\theta,1)\) but using only that the Neyman--Pearson inequality flips with the sign of \(\theta''-\theta'\).
\end{proof}

\begin{remark}[What to write instead of a two-sided UMP]
\label{rem:two-sided-no-ump}
A point null is not itself the obstruction. The obstruction is asking one region to be most powerful on both sides of \(\theta'\) at once. The same one-sided test that is UMP against \(\{\theta>\theta'\}\) has monotone power, so on the opposite side its power falls below \(\alpha\); it is not even unbiased for \(H_1:\theta\neq\theta'\), in the sense of \citep[Definition~8.1.2]{hogg2019}. Restricting to unbiased tests restores a best procedure: the usual equal-tailed two-sided test is UMP unbiased. Hogg does not develop that theory; he points to Lehmann and notes that the likelihood-ratio test coincides with the UMPU test for a normal mean with known variance \citep[\S8.3, after (8.3.1)]{hogg2019}. On an exam, the expected move is: exhibit the two opposite Neyman--Pearson tails, conclude that no UMP exists, and if a two-sided procedure is still asked for, write the equal-tailed region or the LRT. The 2025 Spring Gamma item \(H_0:\theta=1\) versus \(H_1:\theta\neq 1\) is exactly this question.
\end{remark}

\begin{remark}[How to use the toolkit on Problem~7]
\label{rem:p7-toolkit}
Part~(a) is \cref{def:mlr}: exhibit a statistic \(T=u(\mathbf{X})\) such that \(L(\theta_1,\mathbf{x})/L(\theta_2,\mathbf{x})\) is monotone in \(T\) whenever \(\theta_1<\theta_2\). The family is a one-parameter exponential family in \(\sqrt{x}\), so a natural candidate is a sum of \(\sqrt{X_i}\). Record the direction (increasing or decreasing).

Part~(b) is \cref{def:ump-size-alpha} together with \cref{thm:karlin-rubin}(ii). The exam writes \(H_0:\theta\ge 1\) versus \(H_1:\theta<1/2\). These two sets are not complementary: the interval \((1/2,1)\) is in neither hypothesis. Size is still computed on \(H_0\), so the boundary point is \(\theta'=1\), not \(\theta=1/2\). Because a UMP region for the larger alternative \(\{\theta<1\}\) is most powerful at every \(\theta''<1\), it remains UMP when the alternative is restricted to the subset \(\{\theta<1/2\}\). The critical value needs the law of \(T\): that is \cref{thm:jacobian} applied to \(Y_i=\sqrt{X_i}\).
\end{remark}

\begin{remark}[Hints, in order]
\label{rem:p7-hints}
\begin{enumerate}[label=\textup{(\arabic*)}]
    \item
    Write the joint density. The \(\theta\)-dependent piece should be of the form \(\theta^{-n}\exp(-T/\theta)\) times a factor free of \(\theta\). What is \(T\)?
    \item
    For \(\theta_1<\theta_2\), form \(L(\theta_1)/L(\theta_2)\) and decide whether it is decreasing or increasing in \(T\). That is \cref{def:mlr}.
    \item
    The hypotheses are the left-sided pair of \cref{thm:karlin-rubin}(ii) with \(\theta'=1\). Reject for small \(T\) or for large \(T\), according to the direction in (a).
    \item
    Size \(\alpha\) is \cref{def:test-size} over \(\{\theta\ge 1\}\). If the power function is monotone, the maximum Type~I error is at \(\theta=1\). The null law of \(T\) is obtained from \(Y_i=\sqrt{X_i}\) by the Jacobian formula \cref{thm:jacobian}, then recognised as Gamma / \(\chi^{2}\). Solve \(\Pbb_{\theta=1}(T\in C)=\alpha\).
    \item
    The gap \((1/2,1)\) does not change the critical region. It does mean you never evaluate Type~I error at \(\theta=1/2\).
\end{enumerate}
\end{remark}

The size calculation needs the law of \(T=\sum_i\sqrt{X_i}\). That is a change of variables: the exam density is written in \(x\), the sufficient statistic is written in \(y=\sqrt{x}\). A density is not a probability. For a continuous random variable the infinitesimal statement is
\[
\Pbb(X\in\mathrm{d}x)
=
f_X(x)\,|\mathrm{d}x|.
\]
If \(Y=g(X)\) is one-to-one, the events \(\{Y\in\mathrm{d}y\}\) and \(\{X\in\mathrm{d}x\}\) are the same event, with \(x=g^{-1}(y)\). Probability is conserved, so
\[
f_Y(y)\,|\mathrm{d}y|
=
f_X(x)\,|\mathrm{d}x|.
\]
A nonlinear \(g\) stretches or compresses intervals: the factor \(|\mathrm{d}x/\mathrm{d}y|\) is how much \(x\)-length sits over a unit of \(y\)-length. That factor is the Jacobian of the inverse map. In several dimensions the same conservation holds with volume in place of length, and the stretching factor is \(|\det D\mathbf{w}|\), where \(\mathbf{x}=\mathbf{w}(\mathbf{y})\) is the inverse. The absolute value is there because probability cannot be negative, and because a decreasing map reverses orientation.

\begin{theorem}[Jacobian change of variables]
\label{thm:jacobian}
\begin{enumerate}[label=\textup{(\roman*)}]
    \item
    \emph{One dimension.} Let \(X\) be continuous with density \(f_X\) and support \(S_X\). Let \(Y=g(X)\), where \(g\) is \(C^{1}\) and one-to-one on \(S_X\), with inverse \(x=g^{-1}(y)\) and \(g^{-1}\) differentiable on \(S_Y=g(S_X)\). Then, for \(y\in S_Y\),
    \[
    f_Y(y)
    =
    f_X\bigl(g^{-1}(y)\bigr)
    \biggl|\frac{\mathrm{d}x}{\mathrm{d}y}\biggr|
    =
    f_X\bigl(g^{-1}(y)\bigr)
    \bigl|(g^{-1})'(y)\bigr|.
    \]
    This is \citep[Theorem~1.7.1]{hogg2019}. Hogg calls \(\mathrm{d}x/\mathrm{d}y\) the Jacobian of the transformation.
    \item
    \emph{\(n\) dimensions.} Let \(\mathbf{X}\) be continuous on an open set \(S\subset\mathbb{R}^n\), with joint density \(f_{\mathbf{X}}\). Let \(\mathbf{Y}=T(\mathbf{X})\), where \(T:S\to\mathcal{T}:=T(S)\) is a \(C^{1}\) diffeomorphism, with inverse \(\mathbf{x}=\mathbf{w}(\mathbf{y})\) and Jacobian
    \[
    J(\mathbf{y})
    \coloneqq
    \det
    \biggl(
    \frac{\partial w_i}{\partial y_j}(\mathbf{y})
    \biggr)_{1\le i,j\le n}
    \neq 0
    \qquad\text{on }\mathcal{T}.
    \]
    Then, for \(\mathbf{y}\in\mathcal{T}\),
    \[
    f_{\mathbf{Y}}(\mathbf{y})
    =
    f_{\mathbf{X}}\bigl(\mathbf{w}(\mathbf{y})\bigr)\,
    \bigl|J(\mathbf{y})\bigr|.
    \]
    The two-dimensional case is \citep[(2.2.3), \S2.2]{hogg2019}; the same formula in \(\mathbb{R}^n\) is the density of \(\eta=\varphi(\xi)\) in \citep[\S4.8, after Theorem~4.8.2]{ren2024}.
\end{enumerate}
\end{theorem}

\begin{proof}
\emph{(i), increasing case.} Assume \(g\) is strictly increasing, so \(g^{-1}\) is strictly increasing and \(\mathrm{d}x/\mathrm{d}y>0\). For \(y\in S_Y\),
\[
F_Y(y)
=
\Pbb(Y\le y)
=
\Pbb\bigl(g(X)\le y\bigr)
=
\Pbb\bigl(X\le g^{-1}(y)\bigr)
=
F_X\bigl(g^{-1}(y)\bigr).
\]
Differentiate by the chain rule:
\[
f_Y(y)
=
F_Y'(y)
=
f_X\bigl(g^{-1}(y)\bigr)\,
\frac{\mathrm{d}x}{\mathrm{d}y}
=
f_X\bigl(g^{-1}(y)\bigr)
\biggl|\frac{\mathrm{d}x}{\mathrm{d}y}\biggr|.
\]
This is Hogg's identity \citep[(1.7.12)--(1.7.13)]{hogg2019}.

\emph{(i), decreasing case.} Now \(g\) is strictly decreasing, so \(\mathrm{d}x/\mathrm{d}y<0\) and the inequality reverses:
\[
F_Y(y)
=
\Pbb\bigl(g(X)\le y\bigr)
=
\Pbb\bigl(X\ge g^{-1}(y)\bigr)
=
1-F_X\bigl(g^{-1}(y)\bigr).
\]
Differentiate:
\[
f_Y(y)
=
f_X\bigl(g^{-1}(y)\bigr)
\biggl(-\frac{\mathrm{d}x}{\mathrm{d}y}\biggr)
=
f_X\bigl(g^{-1}(y)\bigr)
\biggl|\frac{\mathrm{d}x}{\mathrm{d}y}\biggr|.
\]
The two cases are therefore the same formula. That is why the Jacobian enters in absolute value.

\emph{(ii).} Let \(B\subset\mathcal{T}\) be a Borel set of finite volume, and write \(A=T^{-1}(B)\). Because \(T\) is one-to-one,
\[
\Pbb(\mathbf{X}\in A)
=
\Pbb(\mathbf{Y}\in B).
\]
The left-hand side is \(\int_A f_{\mathbf{X}}(\mathbf{x})\,\mathrm{d}\mathbf{x}\). The calculus change of variables \(\mathbf{x}=\mathbf{w}(\mathbf{y})\) sends \(A\) to \(B\) and multiplies the volume element by \(|J(\mathbf{y})|\), so
\[
\int_A f_{\mathbf{X}}(\mathbf{x})\,\mathrm{d}\mathbf{x}
=
\int_B f_{\mathbf{X}}\bigl(\mathbf{w}(\mathbf{y})\bigr)\,
\bigl|J(\mathbf{y})\bigr|\,\mathrm{d}\mathbf{y}.
\]
The right-hand side is \(\Pbb(\mathbf{Y}\in B)\) for every such \(B\), hence the integrand is \(f_{\mathbf{Y}}\). This is the argument recorded at \citep[\S2.2]{hogg2019} and, with \(\varphi^{-1}\) in place of \(\mathbf{w}\), at \citep[\S4.8]{ren2024}.
\end{proof}

Hogg's computational order, after Theorem~1.7.1: (1)~support of \(Y\); (2)~solve \(x=g^{-1}(y)\); (3)~compute \(\mathrm{d}x/\mathrm{d}y\); (4)~multiply \(f_X(g^{-1}(y))\) by \(|\mathrm{d}x/\mathrm{d}y|\). The same four steps in several variables, with \(|J|\) in place of \(|\mathrm{d}x/\mathrm{d}y|\). If only a function of some coordinates is needed, adjoin dummy coordinates so the map is \(\mathbb{R}^n\to\mathbb{R}^n\), then integrate the dummies out.

\begin{example}[The map in Problem~7]
\label{ex:sqrt-jacobian}
On \(\{x>0\}\) the exam density is
\[
f(x;\theta)
=
\frac{1}{2\theta\sqrt{x}}\,e^{-\sqrt{x}/\theta}.
\]
Set \(Y=\sqrt{X}\). Then \(g(x)=\sqrt{x}\) is \(C^{1}\) and strictly increasing on \((0,\infty)\), with inverse \(x=y^{2}\) and Jacobian
\[
\frac{\mathrm{d}x}{\mathrm{d}y}
=
2y.
\]
The support of \(Y\) is \((0,\infty)\). For \(y>0\), \cref{thm:jacobian}(i) gives
\[
f_Y(y;\theta)
=
f(y^{2};\theta)\cdot |2y|
=
\frac{1}{2\theta y}\,e^{-y/\theta}\cdot 2y
=
\theta^{-1}e^{-y/\theta}.
\]
Thus \(Y\sim\mathrm{Exp}(1/\theta)\) in the rate parameterisation (equivalently \(\mathrm{Gamma}(1,1/\theta)\)). For the sample, the coordinatewise map \(\mathbf{x}\mapsto(\sqrt{x_1},\dots,\sqrt{x_n})\) is a diffeomorphism of \((0,\infty)^n\) with diagonal Jacobian, so
\[
|J(\mathbf{y})|
=
\prod_{i=1}^{n}(2y_i).
\]
Independence is preserved, and \(T=\sum_{i=1}^{n}Y_i\) is the sum of \(n\) i.i.d.\ exponentials.
\end{example}

\subsubsection{Solution of Problem 7}

\paragraph{Answer.}

Write \(T\coloneqq\sum_{i=1}^{n}\sqrt{X_i}\). The family has monotone decreasing likelihood ratio in \(T\), in the sense of \cref{def:mlr} \citep[Definition~8.2.2]{hogg2019}. Let \(\chi^{2}_{2n}(\alpha)\) denote the lower \(\alpha\)-quantile of the \(\chi^{2}(2n)\) law, i.e.\ \(\Pbb\bigl(\chi^{2}(2n)\le\chi^{2}_{2n}(\alpha)\bigr)=\alpha\). By \cref{thm:karlin-rubin}(ii) \citep[(8.2.6)]{hogg2019}, the test

\[
\text{reject }H_0\text{ if }T\le\tfrac12\chi^{2}_{2n}(\alpha)
\]

is UMP of size \(\alpha\) for \(H_0:\theta\ge 1\) versus \(H_1:\theta<1/2\). Equivalently, reject \(H_0\) if \(2T\le\chi^{2}_{2n}(\alpha)\).

\begin{proof}[Exam proof]
Write \(\mathbf{X}=(X_1,\dots,X_n)\) and \(L(\theta,\mathbf{x})\) as in the notes above. Throughout, \(\theta>0\) and the observations are supported on \((0,\infty)\), independently of \(\theta\).

\textbf{(a).}
On \(\{x_i>0\text{ for all }i\}\),

\[
L(\theta,\mathbf{x})
=
\prod_{i=1}^{n}
\frac{1}{2\theta\sqrt{x_i}}
\exp\bigl(-\sqrt{x_i}/\theta\bigr)
=
2^{-n}\theta^{-n}
\Bigl(\prod_{i=1}^{n}x_i\Bigr)^{-1/2}
\exp\bigl(-T(\mathbf{x})/\theta\bigr),
\]

where \(T(\mathbf{x})\coloneqq\sum_{i=1}^{n}\sqrt{x_i}\). For \(0<\theta_1<\theta_2\),

\[
\frac{L(\theta_1,\mathbf{x})}{L(\theta_2,\mathbf{x})}
=
\Bigl(\frac{\theta_2}{\theta_1}\Bigr)^{n}
\exp\Bigl(T(\mathbf{x})\bigl(\tfrac{1}{\theta_2}-\tfrac{1}{\theta_1}\bigr)\Bigr).
\]

The prefactor does not depend on the data. The coefficient \(\tfrac{1}{\theta_2}-\tfrac{1}{\theta_1}\) is strictly negative, so the ratio is a strictly decreasing function of \(T(\mathbf{x})\). This is monotone likelihood ratio in the statistic \(T=T(\mathbf{X})\), in the decreasing direction of \cref{def:mlr}: large \(T\) favours large \(\theta\).

\textbf{(b).}
The hypotheses are the left-sided pair of \cref{thm:karlin-rubin}(ii), except that the exam's alternative \(\{\theta<1/2\}\) is a proper subset of \(\{\theta<1\}\). Size is computed on \(\omega_0=\{\theta\ge 1\}\) by \cref{def:test-size}, so the boundary point is \(\theta'=1\), not \(\theta=1/2\). The interval \((1/2,1)\) is in neither hypothesis and never enters the size calculation.

\emph{Law of \(T\).}
Let \(Y_i\coloneqq\sqrt{X_i}\). By \cref{thm:jacobian}(i) and \cref{ex:sqrt-jacobian}, each \(Y_i\) has density \(\theta^{-1}e^{-y/\theta}\) on \((0,\infty)\), i.e.\ \(\mathrm{Exp}(1/\theta)\) in the rate parameterisation of \cref{note:chf-table}, equivalently \(\mathrm{Gamma}(1,1/\theta)\). The sum of \(n\) i.i.d.\ copies is \(T\sim\mathrm{Gamma}(n,1/\theta)\). Equivalently \(T\stackrel{(d)}{=}\theta S\) with \(S\sim\mathrm{Gamma}(n,1)\). Since \(\chi^{2}(2n)=\mathrm{Gamma}(n,1/2)\), one has \(2S\sim\chi^{2}(2n)\), and therefore

\[
\frac{2T}{\theta}\sim\chi^{2}(2n).
\]

\emph{Neyman--Pearson at the boundary.}
Fix any simple alternative \(\theta''<1\). By the Neyman--Pearson lemma recorded in \cref{def:best-cr}, a best critical region of size \(\alpha\) for \(H_0:\theta=1\) versus \(H_1:\theta=\theta''\) is

\[
\frac{L(1,\mathbf{x})}{L(\theta'',\mathbf{x})}
\le
k,
\]

with \(k\) chosen to give size \(\alpha\). Substituting the likelihood,

\[
\frac{L(1,\mathbf{x})}{L(\theta'',\mathbf{x})}
=
(\theta'')^{n}
\exp\bigl(T(\mathbf{x})(1/\theta''-1)\bigr).
\]

Here \(1/\theta''-1>0\), so the left-hand side is strictly increasing in \(T(\mathbf{x})\). The inequality is therefore equivalent to \(T(\mathbf{x})\le c\), for a threshold \(c\) that may depend on \(\alpha\) and \(n\) but not on the particular \(\theta''<1\). The same rejection region is thus best against every simple alternative \(\theta''<1\), hence UMP of size \(\alpha\) for the simple null \(\theta=1\) against the composite \(\{\theta<1\}\), in the sense of \cref{def:ump-size-alpha}. This is the left-sided case of \cref{thm:karlin-rubin}(ii). Restricting the alternative to the exam's subset \(\{\theta<1/2\}\) does not change the critical region: a region that is most powerful at every \(\theta''<1\) remains most powerful at every \(\theta''<1/2\).

\emph{Size over the composite null.}
For the test \(C=\{T\le c\}\) the power function is

\[
\gamma_C(\theta)
=
\Pbb_{\theta}(T\le c)
=
\Pbb(S\le c/\theta),
\]

which is nonincreasing in \(\theta\). Therefore

\[
\max_{\theta\ge 1}\gamma_C(\theta)
=
\gamma_C(1)
=
\Pbb_{\theta=1}(T\le c).
\]

Choosing \(c\) so that \(\gamma_C(1)=\alpha\) makes the test of exact size \(\alpha\) for \(H_0:\theta\ge 1\), in the sense of \cref{def:test-size}. Combined with the previous paragraph, \(C\) is UMP of size \(\alpha\) for the exam's pair of hypotheses.

\emph{Critical value.}
Under \(\theta=1\) one has \(2T\sim\chi^{2}(2n)\). The equation \(\Pbb_{\theta=1}(T\le c)=\alpha\) becomes \(\Pbb\bigl(\chi^{2}(2n)\le 2c\bigr)=\alpha\), hence \(2c=\chi^{2}_{2n}(\alpha)\). The law of \(T\) is continuous, so no randomisation is needed. This is the test recorded in the answer.
\end{proof}

\begin{remark}[Sanity checks]
\label{rem:p7-checks}
Small \(T\) means the \(\sqrt{X_i}\) are typically small, which is the small-\(\theta\) exponential law, so the test rejects toward \(H_1:\theta<1/2\) as it should. Calibrating the cutoff at \(\theta=1/2\) would produce a strictly smaller critical value, hence a test of size strictly less than \(\alpha\) on \(\{\theta\ge 1\}\), which cannot be UMP of size \(\alpha\). The quantile \(\chi^{2}_{2n}(\alpha)\) is the left-tail point, not the upper \(\alpha\) critical value used in two-sided chi-square tests.
\end{remark}

\subsection{Problem 8. [15 pts.]}
Let \( X_1, X_2, \dots, X_n \) be i.i.d.\ Bernoulli random variables with unknown success probability \( 0 < \theta < 1 \). Given \( 0 < p < 1 \), let \( \hat{\xi}_p \) be the sample quantile based on the \( n \) observations, i.e., \( \hat{\xi}_p = \inf\{x : F_n(x) \ge p\} \), where \( F_n \) denotes the empirical CDF based on the data.
\begin{enumerate}[label=(\alph*)]
    \item Does \( \hat{\xi}_p \) converge in probability to the population quantile \( \xi_p \) for each \( 0 < p < 1 \) as \( n \to \infty \)? Prove your answer.
    \item Is it possible that for a non-degenerate discrete distribution \( F \) we have \( \hat{\xi}_p \) converges in probability to \( \xi_p \) for all \( 0 < p < 1 \)? Prove your answer.
\end{enumerate}

\setcounter{section}{8}
\setcounter{definition}{0}

\subsubsection{Hogg examples near Problem 8 (not the exam item)}

Problem~8 is a discrete-quantile consistency question. The same objects appear in Hogg as the quantile definition \citep[Definition~1.7.2]{hogg2019}, order statistics \citep[\S4.4]{hogg2019}, convergence in probability \citep[\S5.1]{hogg2019}, and the empirical-cdf plug-in \citep[Remark~10.1.1, Exercise~10.1.3]{hogg2019}. The four modes of convergence used on this exam, and which implications are unconditional, are the table in \cref{rem:conv-modes}; the exam writes \(\xrightarrow{\Pbb}\), which is the in-probability row of that table, recorded locally as \cref{def:conv-p}. The exercises below are cousins, not clones. The exam-level solution is assembled after the hints.

\begin{definition}[Quantile of order \(p\)]
\label{def:quantile}
Let \(0<p<1\). A \emph{quantile of order \(p\)} of \(X\) is a number \(\xi_p\) such that
\[
\Pbb(X<\xi_p)\le p
\qquad\text{and}\qquad
\Pbb(X\le\xi_p)\ge p.
\]
See \citep[Definition~1.7.2]{hogg2019}. Equivalently, writing \(F\) for the cdf,
\[
F(\xi_p-)\le p\le F(\xi_p).
\]
If \(F\) is continuous and strictly increasing at the relevant point, this collapses to \(\xi_p=F^{-1}(p)\) as in \citep[\S4.4.1]{hogg2019}. For a discrete law the two inequalities are the only safe definition: \(\xi_p\) need not be unique, and \(F(\xi_p)=p\) can fail.
\end{definition}

The exam's \(\hat{\xi}_p=\inf\{x:F_n(x)\ge p\}\) is the left-inverse of the empirical cdf, i.e.\ the smallest number satisfying the second inequality in \cref{def:quantile} for \(F_n\). Hogg's sample quantile in the continuous chapter is an order statistic \(Y_k\) with \(k=\lfloor p(n+1)\rfloor\) \citep[\S4.4.1]{hogg2019}; for large \(n\) the two conventions agree. On a discrete \(F_n\), Hogg switches back to \cref{def:quantile}: see \citep[Remark~10.1.1]{hogg2019}.

\begin{definition}[Empirical cdf]
\label{def:ecdf}
For a sample \(X_1,\dots,X_n\),
\[
F_n(x)
\coloneqq
n^{-1}\#\{i:X_i\le x\}
=
n^{-1}\sum_{i=1}^{n}\mathbf{1}_{\{X_i\le x\}}.
\]
This is \citep[(10.1.1)]{hogg2019}. It is a discrete cdf putting mass \(n^{-1}\) at each sample point.
\end{definition}

\begin{definition}[Convergence in probability]
\label{def:conv-p}
\(X_n\xrightarrow{\Pbb}X\) if \(\Pbb(|X_n-X|\ge\varepsilon)\to 0\) for every \(\varepsilon>0\). See \citep[Definition~5.1.1]{hogg2019}. If \(X\) is the constant \(a\), one writes \(X_n\xrightarrow{\Pbb}a\).
\end{definition}

\begin{lemma}[Chebyshev]
\label{lem:chebyshev}
If \(\E[X^2]<\infty\), then for every \(\varepsilon>0\),
\[
\Pbb\bigl(|X-\E[X]|\ge\varepsilon\bigr)
\le
\frac{\Var(X)}{\varepsilon^2}.
\]
This is Markov applied to \((X-\E[X])^2\). It is the inequality behind \citep[Theorem~5.1.1]{hogg2019}.
\end{lemma}

\begin{theorem}[Weak law of large numbers]
\label{thm:wlln}
Let \(X_1,\dots,X_n\) be i.i.d.\ with \(\E[X_1]=\mu\) and \(\Var(X_1)=\sigma^2<\infty\). Then \(\overline{X}_n\xrightarrow{\Pbb}\mu\). See \citep[Theorem~5.1.1]{hogg2019}. The proof is \cref{lem:chebyshev} on the sample mean:
\[
\Pbb\bigl(|\overline{X}_n-\mu|\ge\varepsilon\bigr)
\le
\frac{\sigma^2}{n\varepsilon^2}
\to 0.
\]
\end{theorem}

\begin{lemma}[{Pointwise consistency of \(F_n\); cf.\ \citep[Exercise~5.1.2]{hogg2019}}]
\label{lem:Fn-pointwise}
Fix \(x\in\mathbb{R}\). Then \(F_n(x)\xrightarrow{\Pbb}F(x)\).
\end{lemma}

\begin{proof}
The indicators \(\mathbf{1}_{\{X_i\le x\}}\) are i.i.d.\ Bernoulli with mean \(F(x)\). Their average is \(F_n(x)\). This is exactly Hogg's Exercise~5.1.2(a) with success probability \(p=F(x)\): if \(Y_n\sim\mathrm{Bin}(n,p)\), then \(Y_n/n\xrightarrow{\Pbb}p\) by \cref{thm:wlln}, or directly by \cref{lem:chebyshev},
\[
\Pbb\bigl(|Y_n/n-p|\ge\varepsilon\bigr)
\le
\frac{p(1-p)}{n\varepsilon^2}
\to 0.
\]
\end{proof}

So \(F_n\) is consistent at every fixed \(x\). Problem~8 asks something sharper: whether the \emph{inverse} \(F_n^{-1}(p)\) is consistent at every \(p\). Inverses fail at flats and jumps of \(F\), which is why discrete laws are the interesting case.

\begin{exercise}[{Hogg 1.7.9(a): a discrete median}]
\label{ex:hogg-179a}
Let \(X\sim\mathrm{Bin}(4,1/4)\), with pmf
\[
p(x)
=
\binom{4}{x}\Bigl(\frac14\Bigr)^x\Bigl(\frac34\Bigr)^{4-x},
\qquad
x=0,1,2,3,4.
\]
Find the median \(\xi_{1/2}\) using \cref{def:quantile}.
\end{exercise}

\begin{proof}[Solution]
The masses are \(81,108,54,12,1\) over the common denominator \(256\), so
\[
F(0)
=
\frac{81}{256}
\approx 0.316,
\qquad
F(1)
=
\frac{189}{256}
\approx 0.738.
\]
Check \(\xi=1\): \(\Pbb(X<1)=F(0)\le 1/2\) and \(\Pbb(X\le 1)=F(1)\ge 1/2\). Check \(\xi=0\): \(\Pbb(X\le 0)=F(0)<1/2\), so \(0\) is not a median. Thus \(\xi_{1/2}=1\). (Hogg's selected answers record only parts (b)--(c) of this exercise, \citep[Appendix~F]{hogg2019}.)
\end{proof}

The same two inequalities are what one writes for Bernoulli(\(\theta\)) in Problem~8(a). For a two-point law the cdf has a single jump of height \(1-\theta\) at \(0\), so \(\xi_p\) is \(0\) or \(1\) according as \(p\le 1-\theta\) or \(p>1-\theta\).

\begin{exercise}[{Hogg 1.7.11(d): quartiles of a discrete cdf}]
\label{ex:hogg-1711d}
Let
\[
F(x)
=
\sum_{j=1}^{x}\Bigl(\frac12\Bigr)^{j}
\]
on the positive integers (and \(F(x)=F(\lfloor x\rfloor)\) in between). Find the pmf, the first quartile \(\xi_{1/4}\), and the \(0.60\) quantile \(\xi_{0.60}\).
\end{exercise}

\begin{proof}[Solution]
The sum is a geometric series: \(p(j)=(1/2)^j\) for \(j=1,2,\dots\), and \(\sum_{j\ge 1}p(j)=1\). The cdf at the atoms is \(F(k)=1-2^{-k}\), so \(F(1)=1/2\), \(F(2)=3/4\), \(F(3)=7/8\), and so on.

For \(\xi_{1/4}\): the candidate \(1\) satisfies \(\Pbb(X<1)=0\le 1/4\) and \(\Pbb(X\le 1)=1/2\ge 1/4\). No smaller support point exists, so \(\xi_{1/4}=1\).

For \(\xi_{0.60}\): \(F(1)=0.5<0.60\le 0.75=F(2)\), and \(\Pbb(X<2)=1/2\le 0.60\le \Pbb(X\le 2)\). Thus \(\xi_{0.60}=2\).
\end{proof}

Notice that a whole interval of \(p\), namely \((1/2,3/4]\), shares the same quantile \(2\). That flat-in-\(p\) / jump-in-\(x\) geometry is the obstruction in Problem~8(b).

\begin{exercise}[{Hogg 4.4.6(a): order statistic versus population median}]
\label{ex:hogg-446a}
Let \(X_1,X_2,X_3\) be i.i.d.\ with pdf \(f(x)=2x\) on \((0,1)\). Compute the probability that the sample minimum exceeds the population median.
\end{exercise}

\begin{proof}[Solution]
Here \(F(x)=x^2\) on \((0,1)\), so the median solves \(F(\xi_{1/2})=1/2\) and is \(\xi_{1/2}=2^{-1/2}\). The minimum \(Y_1\) exceeds \(\xi_{1/2}\) if and only if all three observations do, hence
\[
\Pbb(Y_1>\xi_{1/2})
=
\bigl(1-F(\xi_{1/2})\bigr)^3
=
\bigl(\tfrac12\bigr)^3
=
\frac18,
\]
which is the value recorded in \citep[Appendix~F]{hogg2019}. This is a finite-\(n\) exact probability, not a consistency statement; the continuous strictly increasing cdf makes the median unique, which is the situation Problem~8 deliberately leaves.
\end{proof}

\begin{example}[{Hogg 5.1.2: sample maximum, uniform \((0,\theta)\)}]
\label{ex:hogg-512-uniform}
Let \(X_i\) be i.i.d.\ uniform on \((0,\theta)\) and \(Y_n=\max_i X_i\). The cdf of \(Y_n\) is \((t/\theta)^n\) for \(0<t\le\theta\) \citep[(5.1.1)]{hogg2019}. For \(0<\varepsilon<\theta\),
\[
\Pbb(|Y_n-\theta|\ge\varepsilon)
=
\Pbb(Y_n\le\theta-\varepsilon)
=
\Bigl(\frac{\theta-\varepsilon}{\theta}\Bigr)^n
\to 0,
\]
so \(Y_n\xrightarrow{\Pbb}\theta\). This is consistency of an \emph{extreme} order statistic, proved from its exact cdf. The same pattern---write \(\Pbb(Y_n\) is on the wrong side of the target\()\), then send \(n\to\infty\)---is what one uses for sample quantiles, except that the exact cdf of \(\hat{\xi}_p\) is binomial rather than a power.
\end{example}

\begin{exercise}[{Hogg 5.1.7: sample minimum of a shifted exponential}]
\label{ex:hogg-517}
Let \(X_i\) be i.i.d.\ with pdf \(e^{-(x-\theta)}\mathbf{1}_{\{x>\theta\}}\), and \(Y_n=\min_i X_i\). Prove \(Y_n\xrightarrow{\Pbb}\theta\) by first obtaining the cdf of \(Y_n\).
\end{exercise}

\begin{proof}[Solution]
Always \(Y_n\ge\theta\). For \(t>\theta\),
\[
\Pbb(Y_n>t)
=
\Pbb(X_1>t)^n
=
e^{-n(t-\theta)},
\]
so the cdf is \(F_{Y_n}(t)=1-e^{-n(t-\theta)}\) for \(t>\theta\). Thus for \(\varepsilon>0\),
\[
\Pbb(|Y_n-\theta|\ge\varepsilon)
=
\Pbb(Y_n\ge\theta+\varepsilon)
=
e^{-n\varepsilon}
\to 0.
\]
This is the order-statistic consistency item. Continuity of the cdf to the right of \(\theta\) is used in an essential way: \(Y_n\) cannot overshoot \(\theta\) to the left.
\end{proof}

\begin{exercise}[{Hogg 10.1.3: the sample median as \(T(F_n)\)}]
\label{ex:hogg-1013}
Let \(T(F)=F^{-1}(1/2)\) denote the median functional, using \cref{def:quantile} when \(F\) is discrete. Show that the mean of \(F_n\) is \(\overline{X}\), and that \(T(F_n)\) is the sample median \(Q_2\) of \citep[(4.4.4)]{hogg2019}.
\end{exercise}

\begin{proof}[Solution]
The discrete law \(F_n\) puts mass \(n^{-1}\) at each \(X_i\), so its mean is \(n^{-1}\sum_i X_i=\overline{X}\).

For the median: if \(n=2m+1\) is odd, the middle order statistic \(X_{(m+1)}\) satisfies
\[
F_n\bigl(X_{(m+1)}-\bigr)
=
\frac{m}{n}
=
\frac12-\frac{1}{2n}
\le
\frac12
\le
F_n\bigl(X_{(m+1)}\bigr),
\]
hence is a median of \(F_n\) in the sense of \cref{def:quantile}. If \(n=2m\) is even, \(F_n\) equals \(1/2\) throughout \([X_{(m)},X_{(m+1)})\) at the averaged sample median of (4.4.4), and every point of that interval is a median of \(F_n\). This is the argument written in \citep[Remark~10.1.1]{hogg2019}.
\end{proof}

So the exam's \(\hat{\xi}_p=\inf\{x:F_n(x)\ge p\}\) is precisely Hogg's induced estimator \(T(F_n)\) for the left-inverse functional \(T(F)=\inf\{x:F(x)\ge p\}\).

\begin{remark}[How these items sit next to Problem~8]
\label{rem:p8-hogg-map}
\begin{enumerate}[label=\textup{(\arabic*)}]
    \item
    Discrete quantile arithmetic is \cref{def:quantile}, \cref{ex:hogg-179a}, and \cref{ex:hogg-1711d}. For Bernoulli(\(\theta\)), first write \(\xi_p\) as a function of \(p\) and \(\theta\) by the same two inequalities.
    \item
    \Cref{lem:Fn-pointwise} is the only asymptotic input. On \(\{0,1\}\), \(F_n(0)=1-\overline{X}_n\) and \(F_n(1)=1\), so \(\hat{\xi}_p\in\{0,1\}\) and the event \(\{\hat{\xi}_p=0\}\) is \(\{F_n(0)\ge p\}=\{\overline{X}_n\le 1-p\}\).
    \item
    \Cref{ex:hogg-512-uniform} and \cref{ex:hogg-517} show how to prove \(Y_n\xrightarrow{\Pbb}c\) from an exact tail. For Problem~8 the tail is binomial, not exponential or a power; the CLT (or just the fact that \(\overline{X}_n\xrightarrow{\Pbb}\theta\)) decides whether \(\Pbb(\overline{X}_n\le 1-p)\) tends to \(0\), \(1\), or \(1/2\).
    \item
    The continuous examples have a unique \(\xi_p\) with \(F\) strictly increasing there. That hypothesis is exactly what a nondegenerate discrete \(F\) lacks at its atoms, which is the content of part~(b).
    \item
    Hogg's Exercise~4.4.7 (minimum of five dice rolls) is worth looking at for a different discrete warning: the continuous order-statistic densities of \S4.4 are not applicable, because ties have positive probability. The exam avoids densities and works only with \(F_n\).
\end{enumerate}
\end{remark}

\begin{remark}[Hints for Problem~8, in order]
\label{rem:p8-hints}
\begin{enumerate}[label=\textup{(\arabic*)}]
    \item
    Write \(F\) and \(F_n\) on \(\mathbb{R}\). Both are two-step functions, with the only random height equal to the proportion of zeros.
    \item
    Using \cref{def:quantile} (or the exam's \(\inf\)), list \(\xi_p\) as a function of \(p\). There is a critical \(p\) equal to a jump height.
    \item
    Translate \(\{\hat{\xi}_p=0\}\) into an event about \(\overline{X}_n\). Apply \cref{lem:Fn-pointwise} / Exercise~5.1.2. The answer to (a) is not a uniform yes.
    \item
    For (b): if a nondegenerate discrete law has an atom \(a\) that is not the rightmost point of the support, consider \(p=F(a)\in(0,1)\). Compare \(\xi_p\) with the fluctuation of \(F_n(a)\) around \(F(a)\).
\end{enumerate}
\end{remark}

\subsubsection{Solution of Problem 8}

\paragraph{Answer.}

Use the same left-inverse for the population as the exam uses for the sample:
\[
\xi_p
\coloneqq
\inf\{x:F(x)\ge p\}.
\]
(a)~No. For Bernoulli(\(\theta\)) one has \(\hat{\xi}_p\xrightarrow{\Pbb}\xi_p\) if and only if \(p\neq 1-\theta\). At \(p=1-\theta\), the sample quantile takes the values \(0\) and \(1\) with probabilities tending to \(1/2\), so it does not converge in probability to any constant.
(b)~No. The same jump obstruction occurs at \(p=F(a)\) for any non-maximal atom \(a\) of a non-degenerate discrete law.

\begin{proof}[Exam proof]
Write \(F_n\) for the empirical cdf of \cref{def:ecdf}, and write \(\overline{X}_n=n^{-1}\sum_{i=1}^{n}X_i\). The sample quantile is the left-inverse \(\hat{\xi}_p=\inf\{x:F_n(x)\ge p\}\) of \(F_n\). Population quantiles are taken with the same convention, which is a version of \cref{def:quantile}. Convergence in probability is \cref{def:conv-p}.

\textbf{(a).}
The Bernoulli cdf is the two-step function
\[
F(x)
=
\begin{cases}
0, & x<0,\\
1-\theta, & 0\le x<1,\\
1, & x\ge 1.
\end{cases}
\]
Thus \(\xi_p=0\) if \(0<p\le 1-\theta\), and \(\xi_p=1\) if \(1-\theta<p<1\). (When \(p=1-\theta\), every point of \([0,1]\) satisfies the two inequalities of \cref{def:quantile}; the left-inverse selects \(0\).) The empirical cdf is the same two-step function with random height \(F_n(0)=1-\overline{X}_n\), so \(\hat{\xi}_p\in\{0,1\}\) and
\[
\hat{\xi}_p=0
\quad\Longleftrightarrow\quad
F_n(0)\ge p
\quad\Longleftrightarrow\quad
\overline{X}_n\le 1-p.
\]
By \cref{lem:Fn-pointwise}, \(F_n(0)\xrightarrow{\Pbb}F(0)=1-\theta\).

If \(p<1-\theta\), then \(1-\theta\) lies a fixed distance above \(p\), so \(\Pbb\bigl(F_n(0)<p\bigr)\to 0\). Hence \(\Pbb(\hat{\xi}_p=1)\to 0\), and \(\hat{\xi}_p\xrightarrow{\Pbb}0=\xi_p\).

If \(p>1-\theta\), then \(1-\theta\) lies a fixed distance below \(p\), so \(\Pbb\bigl(F_n(0)\ge p\bigr)\to 0\). Hence \(\Pbb(\hat{\xi}_p=0)\to 0\), and \(\hat{\xi}_p\xrightarrow{\Pbb}1=\xi_p\).

If \(p=1-\theta\), the threshold sits on the jump height. Then \(\hat{\xi}_p=0\) if and only if \(\overline{X}_n\le\theta\). The classical CLT gives
\[
\sqrt{n}\,(\overline{X}_n-\theta)
\xrightarrow{(d)}
\mathcal{N}\bigl(0,\theta(1-\theta)\bigr),
\]
\citep[\S5.3]{hogg2019}, \citep[Theorem~3.4.1]{durrett2010}. The limiting cdf is continuous at \(0\), so \cref{def:conv-d} yields
\[
\Pbb(\overline{X}_n\le\theta)
=
\Pbb\bigl(\sqrt{n}\,(\overline{X}_n-\theta)\le 0\bigr)
\to
\tfrac12.
\]
Therefore
\[
\Pbb(\hat{\xi}_p=0)\to\tfrac12,
\qquad
\Pbb(\hat{\xi}_p=1)\to\tfrac12.
\]
In particular \(\Pbb\bigl(|\hat{\xi}_p-0|\ge 1/2\bigr)\to 1/2\), so \(\hat{\xi}_p\) does not converge in probability to \(\xi_p=0\). The same two limits show that \(\hat{\xi}_p\) does not converge in probability to \(1\), nor to any other constant. Thus the answer to (a) is negative: consistency fails at the jump height \(p=1-\theta\).

\textbf{(b).}
No. Let \(F\) be non-degenerate and discrete, with support \(S\). Then \(S\) contains at least two points. Let \(a<a'\) be consecutive points of \(S\), so \(a,a'\in S\) and \((a,a')\cap S=\emptyset\). Set \(p\coloneqq F(a)\). The atom at \(a\) and the atom at \(a'\) give \(0<p<1\). For \(x<a\) one has \(F(x)\le F(a-)<F(a)=p\), hence \(\xi_p=a\).

\begin{center}
\begin{tikzpicture}[x=1.55cm, y=3.25cm, >=Stealth]
  \draw[->] (-0.12,0) -- (4.12,0) node[right, font=\small] {\(x\)};
  \draw[->] (0,-0.05) -- (0,1.22) node[above, font=\small] {\(F(x)\)};
  \draw[densely dotted] (0,1) -- (3.82,1);
  \draw[densely dotted] (0,0.58) -- (3.82,0.58);
  \node[font=\small, left] at (0,1) {\(1\)};
  \node[font=\small, left] at (0,0.58) {\(p=F(a)\)};
  \node[font=\small, left] at (0,0) {\(0\)};
  \foreach \xx/\lab in {0.85/\(s_1\), 1.85/\(a\), 3.05/\(a'\)} {
    \draw (\xx,0.018) -- (\xx,-0.018);
    \fill (\xx,0) circle (1.15pt);
    \node[font=\small, below] at (\xx,-0.065) {\lab};
  }
  \node[font=\scriptsize, below] at (1.85,-0.26) {\(\xi_p=a\)};
  \draw[line width=1.05pt]
    (-0.05,0) -- (0.85,0)
    -- (0.85,0.28) -- (1.85,0.28)
    -- (1.85,0.58);
  \draw[line width=1.28pt, blue!62!black]
    (1.85,0.58) -- (3.05,0.58);
  \draw[line width=1.05pt]
    (3.05,0.58) -- (3.05,1) -- (3.82,1);
  \fill (0.85,0.28) circle (1.05pt);
  \fill (1.85,0.58) circle (1.05pt);
  \fill (3.05,1) circle (1.05pt);
  \draw[fill=white, line width=0.6pt] (0.85,0) circle (1.05pt);
  \draw[fill=white, line width=0.6pt] (1.85,0.28) circle (1.05pt);
  \draw[fill=white, line width=0.6pt] (3.05,0.58) circle (1.05pt);
  \node[font=\scriptsize, text=blue!62!black] at (2.45,0.72) {\((a,a')\cap S=\emptyset\)};
  \node[font=\small] at (1.95,-0.50) {\(S=\{s_1,a,a'\}\)};
\end{tikzpicture}
\end{center}

The only atoms are the ticks on the axis. Consecutive means the blue plateau has no further jump: \(F\) is constant on \([a,a')\). Write \(A_n=\{F_n(a)<p\}\) and \(B_n=\{|\hat{\xi}_p-\xi_p|\ge a'-a\}\). On \(A_n\) the left-inverse of \(F_n\) cannot stop at \(a\) or to its left, and there is no support in \((a,a')\), so \(\hat{\xi}_p\ge a'\). Thus \(A_n\subset B_n\), and
\[
\Pbb(B_n)
=
\Pbb\bigl(|\hat{\xi}_p-\xi_p|\ge a'-a\bigr)
\ge
\Pbb(A_n)
=
\Pbb\bigl(F_n(a)<F(a)\bigr).
\]
The indicators \(\mathbf{1}_{\{X_i\le a\}}\) are i.i.d.\ Bernoulli with mean \(F(a)\in(0,1)\), and \(F_n(a)\) is their sample mean. The same CLT as in (a), now with variance \(F(a)\bigl(1-F(a)\bigr)>0\), gives \(\Pbb(A_n)\to 1/2\). Hence \(\liminf_n\Pbb(B_n)\ge 1/2\), so \(\Pbb(B_n)\) cannot tend to \(0\). (Unlike (a), the events \(A_n\) and \(B_n\) need not coincide: extra atoms to the left of \(a\) can make \(\hat{\xi}_p<a\) on \(A_n^c\). The CLT therefore yields \(1/2\) for \(\Pbb(A_n)\), not for \(\Pbb(B_n)\).) Thus \(\hat{\xi}_p\not\xrightarrow{\Pbb}\xi_p\) at this \(p\). Therefore no non-degenerate discrete \(F\) can have sample-quantile consistency at every \(p\in(0,1)\).
\end{proof}

\begin{remark}[Sanity checks]
\label{rem:p8-checks}
Part~(a) is decided by where \(p\) sits relative to the single jump of \(F\). Weak law / \cref{lem:Fn-pointwise} is enough off the jump; the CLT is used only on the jump, where \(F_n(0)\) still concentrates at \(1-\theta\) but undershoots and overshoots with equal limiting probability. Calibrating \(\xi_p\) as \(1\) rather than \(0\) at \(p=1-\theta\) does not repair (a): both atoms of \(\hat{\xi}_p\) keep mass \(1/2\). Part~(b) is the same picture with the Bernoulli jump \(1-\theta\) replaced by a general interior jump height \(F(a)\). A Dirac law is consistent at every \(p\), which is why the question asks for a non-degenerate \(F\).
\end{remark}

\subsection{Problem 9. [10 pts.]}
Suppose a single observation \( Y \mid \mu \sim N(\mu, 1) \), and the prior on \( \mu \) is \( \pi(\mu) = 0.5 \cdot I(\mu = -1) + 0.5 \cdot I(\mu = 1) \) where \( I(\cdot) \) is the indicator function. That is, under this prior, \( \mu \) is supported on \( \{-1, 1\} \) and takes either value with a probability 0.5. For estimating \( \mu \), let the action space be \( \{-1, 1\} \).
\begin{enumerate}[label=(\alph*)]
    \item Write down the posterior \( \pi(\mu \mid Y = y) \) and derive the Bayes rule \( \delta \) under the absolute error loss \( L(\mu, a) = |\mu - a| \).
    \item Is \( \delta \) above a minimax rule under the absolute error loss? Please justify your answer.
\end{enumerate}

\setcounter{section}{9}
\setcounter{definition}{0}

\subsubsection{Notes for Problem 9}

The sampling model is ordinary. The unusual piece is the \emph{two-point prior}: before seeing \(Y\), \(\mu\) is \(\pm 1\) with equal probability. After seeing \(Y=y\), those two masses are reweighted by the two Gaussian likelihoods. Absolute-error Bayes asks for a posterior median, not the posterior mean. Part~(b) leaves the Bayesian game and asks a frequentist question about the same rule.

Sources: posterior and odds \citep[\S1.3, (1.5)]{gelman2013}; two-point update of the same shape as \citep[\S1.4]{gelman2013}; \(L_1\mapsto\) posterior median in \citep{gelmannotes}; constant-risk Bayes is minimax in \citep[\S7.1]{hogg2019}. The conjugate table in BDA Chapter~2 is the wrong drawer: this prior is two-point, not a normal--normal pair.

\begin{remark}[Exam recipe from BDA]
\label{rem:p9-bda-steps}
\begin{enumerate}[label=\textup{(\arabic*)}]
    \item
    Posterior, \citep[\S1.3]{gelman2013}:
    \[
    \pi(\mu\mid y)\propto f(y\mid\mu)\,\pi(\mu).
    \]
    Discrete prior: posterior odds \(=\) prior odds \(\times\) likelihood ratio \citep[(1.5)]{gelman2013}. The conjugate table in BDA Chapter~2 is the wrong drawer: this prior is two-point, not a normal--normal pair. Prior odds \(1\), then multiply by a likelihood ratio.
    \item
    Decision: minimise posterior expected loss \citep[Chapter~9]{gelman2013}. The maps are \(L_2\mapsto\) posterior mean, \(L_1\mapsto\) posterior median, \(0\)--\(1\mapsto\) posterior mode \citep{gelmannotes}. The exam wrote \(|\mu-a|\), so median, not the posterior mean \(2p-1\).
    \item
    Compute that median on \(\{-1,1\}\).
    \item
    Part~(b) is frequentist. Compute the risk \(R(\mu,\delta)\). Two write-ups are given: quote constant-risk Bayes \citep[\S7.1]{hogg2019} together with \cref{rem:p9-theta}, or unroll \cref{def:risk-minimax} directly from the two-point risks and the Bayes property in (a). Bayes and minimax are parallel criteria, not a hierarchy; see \cref{rem:risk-vs-bayes}.
\end{enumerate}
\end{remark}

\begin{definition}[Prior and posterior]
\label{def:prior-posterior}
The model is
\[
Y\mid\mu\sim f(y\mid\mu),
\qquad
\mu\sim\pi.
\]
The \emph{posterior} is the conditional law of \(\mu\) given \(Y=y\),
\[
\pi(\mu\mid y)
=
\frac{f(y\mid\mu)\,\pi(\mu)}{m(y)},
\qquad
m(y)
=
\int f(y\mid\mu')\,\pi(\mathrm{d}\mu').
\]
See \citep[\S1.3, (1.5)]{gelman2013}. If \(\pi\) is discrete, the integral is a sum. Equivalently, for any two parameter values,
\[
\frac{\pi(\mu_1\mid y)}{\pi(\mu_2\mid y)}
=
\frac{\pi(\mu_1)}{\pi(\mu_2)}
\cdot
\frac{f(y\mid\mu_1)}{f(y\mid\mu_2)}
\]
(posterior odds \(=\) prior odds \(\times\) likelihood ratio).
\end{definition}

\begin{definition}[Posterior expected loss, and Bayes rule]
\label{def:bayes-rule}
A \emph{Bayes rule} for loss \(L(\mu,a)\) is any \(a=\delta(y)\) minimising the posterior expected loss
\[
\rho(a\mid y)
=
\E\bigl[L(\mu,a)\mid Y=y\bigr]
=
\int L(\mu,a)\,\pi(\mathrm{d}\mu\mid y).
\]
See \citep[Chapter~9]{gelman2013} (posterior expected loss \(\rho(d\mid y)\)). Averaging \(\rho(\delta(Y)\mid Y)\) over the marginal of \(Y\) is the Bayes risk \(r(\pi,\delta)=\int R(\mu,\delta)\,\pi(\mathrm{d}\mu)\). A Bayes rule therefore also minimises the Bayes risk.
\end{definition}

\begin{proposition}[\(L_2\) \(\mapsto\) posterior mean, \(L_1\) \(\mapsto\) posterior median]
\label{prop:l1-median}
If \(L(\mu,a)=(\mu-a)^2\), a Bayes estimate is the posterior mean. If \(L(\mu,a)=|\mu-a|\), a Bayes estimate is any posterior median: a number \(m\) with
\[
\pi(\mu<m\mid y)\le\tfrac12
\qquad\text{and}\qquad
\pi(\mu\le m\mid y)\ge\tfrac12.
\]
See \citep{gelmannotes}: squared loss \(\mapsto\) posterior mean; absolute error \(\mapsto\) posterior median. The two inequalities are \cref{def:quantile} applied to the posterior.
\end{proposition}

The \(L_1\) claim is the only non-obvious input to (a). For this exam one does not need the general proof: the posterior lives on \(\{-1,1\}\), so \(\rho(a\mid y)\) is elementary.

\begin{lemma}[Two-point \(L_1\) Bayes]
\label{lem:twopoint-l1}
Let \(\pi(\mu=1\mid y)=p\) and \(\pi(\mu=-1\mid y)=1-p\). Under \(L(\mu,a)=|\mu-a|\),
\[
\delta(y)
=
\begin{cases}
1, & p>\tfrac12,\\
-1, & p<\tfrac12,\\
\text{any point of }[-1,1], & p=\tfrac12.
\end{cases}
\]
\end{lemma}

\begin{proof}
Write \(\rho(a)=(1-p)|a+1|+p|a-1|\). The action \(a\) is unrestricted in \(\mathbb{R}\). The kink points of the two absolute values split the line into three pieces, and the slope on each piece is constant:

\begin{center}
\begin{tikzpicture}[x=0.70cm, y=0.98cm, >=Stealth]
  \tikzset{
    leftseg/.style={line width=1.15pt, red!72!black},
    midseg/.style={line width=1.28pt, blue!62!black},
    rightseg/.style={line width=1.15pt, orange!75!black},
    ax/.style={->, gray!55!black},
    ticklab/.style={font=\scriptsize},
    paneltitle/.style={font=\small, anchor=north},
    slopelab/.style={font=\scriptsize}
  }
  % helper: one panel. p-value encoded by rho at the two kinks.
  % p=1/4: rho(-1)=1/2, rho(1)=3/2
  \begin{scope}[shift={(0,0)}]
    \draw[ax] (-2.25,0) -- (2.45,0) node[right, font=\small] {\(a\)};
    \draw[ax] (0,-0.12) -- (0,2.55) node[above, font=\small] {\(\rho(a)\)};
    \draw[densely dotted, gray!55] (-1,0) -- (-1,2.35);
    \draw[densely dotted, gray!55] (1,0) -- (1,2.35);
    \node[ticklab, below] at (-1,-0.02) {\(-1\)};
    \node[ticklab, below] at (1,-0.02) {\(1\)};
    \draw[leftseg] (-2.15,1.65) -- (-1,0.50);
    \draw[midseg] (-1,0.50) -- (1,1.50);
    \draw[rightseg] (1,1.50) -- (2.05,2.55);
    \fill[red!72!black] (-1,0.50) circle (1.35pt);
    \node[slopelab, red!72!black, rotate=-38, anchor=south] at (-1.62,1.18) {slope \(-1\)};
    \node[slopelab, blue!62!black, rotate=22, anchor=south] at (0.12,1.12) {\(1-2p>0\)};
    \node[slopelab, orange!75!black, rotate=38, anchor=south west] at (1.18,1.78) {slope \(+1\)};
    \node[paneltitle] at (0,-0.78) {\(p<\tfrac12\): min at \(-1\)};
  \end{scope}
  % p=1/2: rho(-1)=rho(1)=1
  \begin{scope}[shift={(5.55,0)}]
    \draw[ax] (-2.25,0) -- (2.45,0) node[right, font=\small] {\(a\)};
    \draw[ax] (0,-0.12) -- (0,2.55) node[above, font=\small] {\(\rho(a)\)};
    \draw[densely dotted, gray!55] (-1,0) -- (-1,2.35);
    \draw[densely dotted, gray!55] (1,0) -- (1,2.35);
    \node[ticklab, below] at (-1,-0.02) {\(-1\)};
    \node[ticklab, below] at (1,-0.02) {\(1\)};
    \draw[leftseg] (-2.15,2.15) -- (-1,1.00);
    \draw[midseg] (-1,1.00) -- (1,1.00);
    \draw[rightseg] (1,1.00) -- (2.15,2.15);
    \draw[line width=0.7pt, blue!62!black] (-1,-0.08) -- (1,-0.08);
    \fill[blue!62!black] (-1,1.00) circle (1.2pt);
    \fill[blue!62!black] (1,1.00) circle (1.2pt);
    \node[slopelab, red!72!black, rotate=-42, anchor=south] at (-1.68,1.68) {slope \(-1\)};
    \node[slopelab, blue!62!black, above] at (0,1.00) {\(\rho\equiv 1\)};
    \node[slopelab, orange!75!black, rotate=42, anchor=south west] at (1.22,1.42) {slope \(+1\)};
    \node[paneltitle] at (0,-0.78) {\(p=\tfrac12\): any point of \([-1,1]\)};
  \end{scope}
  % p=3/4: rho(-1)=3/2, rho(1)=1/2
  \begin{scope}[shift={(11.10,0)}]
    \draw[ax] (-2.25,0) -- (2.45,0) node[right, font=\small] {\(a\)};
    \draw[ax] (0,-0.12) -- (0,2.55) node[above, font=\small] {\(\rho(a)\)};
    \draw[densely dotted, gray!55] (-1,0) -- (-1,2.35);
    \draw[densely dotted, gray!55] (1,0) -- (1,2.35);
    \node[ticklab, below] at (-1,-0.02) {\(-1\)};
    \node[ticklab, below] at (1,-0.02) {\(1\)};
    \draw[leftseg] (-2.05,2.55) -- (-1,1.50);
    \draw[midseg] (-1,1.50) -- (1,0.50);
    \draw[rightseg] (1,0.50) -- (2.15,1.65);
    \fill[orange!75!black] (1,0.50) circle (1.35pt);
    \node[slopelab, red!72!black, rotate=-38, anchor=south east] at (-1.05,2.05) {slope \(-1\)};
    \node[slopelab, blue!62!black, rotate=-22, anchor=south] at (-0.05,1.12) {\(1-2p<0\)};
    \node[slopelab, orange!75!black, rotate=38, anchor=south west] at (1.18,0.88) {slope \(+1\)};
    \node[paneltitle] at (0,-0.78) {\(p>\tfrac12\): min at \(1\)};
  \end{scope}
\end{tikzpicture}
\end{center}

Red is \(a<-1\), blue is \([-1,1]\), orange is \(a>1\). Explicitly,
\[
\rho(a)
=
\begin{cases}
-a+2p-1, & a\le -1,\\
1+a(1-2p), & -1\le a\le 1,\\
a+1-2p, & a\ge 1.
\end{cases}
\]
On \(a<-1\) the slope is \(-1\), so \(\rho\) is strictly decreasing toward \(-1\), and every point to the left of \(-1\) is worse than \(a=-1\). On \(a>1\) the slope is \(+1\), so \(\rho\) is strictly increasing away from \(1\), and
\[
\rho(a)=a+1-2p\ge 2(1-p)=\rho(1).
\]
Thus a global minimiser may be sought on \([-1,1]\) alone; the exam never restricts the action space, the \(L_1\) geometry does. On that interval the slope is \(1-2p\): if \(p>\tfrac12\) the graph decreases, hence the unique min is \(a=1\); if \(p<\tfrac12\) it increases, hence the unique min is \(a=-1\); if \(p=\tfrac12\) it is constantly \(1\), so every point of \([-1,1]\) is optimal.
\end{proof}

On two-point support, \(L_1\) Bayes therefore coincides with \(0\)--\(1\) Bayes (pick the more probable point), except at the tie \(p=\tfrac12\), where \(L_1\) also allows the whole interval \([-1,1]\).

\begin{definition}[Risk and minimax]
\label{def:risk-minimax}
The \emph{risk} of a decision rule \(\delta\) is
\[
R(\mu,\delta)
=
\E\bigl[L(\mu,\delta(Y))\mid\mu\bigr].
\]
A rule \(\delta_*\) is \emph{minimax} (on a specified parameter space \(\Theta\)) if
\[
\sup_{\mu\in\Theta} R(\mu,\delta_*)
=
\inf_{\delta}\sup_{\mu\in\Theta} R(\mu,\delta).
\]
Hogg introduces the risk function and the minimax principle in \S7.1, as an alternative to MVUE when one does not force unbiasedness. Gelman et al.\ do not develop this frequentist criterion.
\end{definition}

\begin{remark}[Risk is not a rule; Bayes and minimax are parallel]
\label{rem:risk-vs-bayes}
A \emph{decision rule} \(\delta\) is a map from data to action. Risk, Bayes, and minimax are not synonyms for that map, and they are not the same object under three names.

The loss \(L(\mu,a)\) is a number once \(\mu\) and \(a\) are given. Two different averages of it are in play:
\[
R(\mu,\delta)
=
\E\bigl[L(\mu,\delta(Y))\mid\mu\bigr],
\qquad
\rho(a\mid y)
=
\E\bigl[L(\mu,a)\mid Y=y\bigr].
\]
The left side (\emph{risk}) fixes the unknown parameter and averages over the sampling experiment \(Y\mid\mu\); it is a function of \(\mu\), and no prior enters. The right side (\emph{posterior expected loss}, \cref{def:bayes-rule}) fixes the observed \(y\) and averages over \(\mu\mid y\). Averaging the left side against a prior, or the right side against the marginal of \(Y\), yields the same number, the Bayes risk \(r(\pi,\delta)=\int R(\mu,\delta)\,\pi(\mathrm{d}\mu)\).

A \emph{Bayes rule} is a decision rule that minimises \(\rho(\,\cdot\mid y)\) for each \(y\), equivalently minimises \(r(\pi,\cdot)\). A \emph{minimax} rule is a decision rule that minimises the worst-case height \(\sup_{\mu\in\Theta}R(\mu,\cdot)\). These are two parallel optimality criteria for choosing among decision rules: they collapse the risk function \(R(\cdot,\delta)\) to a number in two different ways (average versus supremum) and then minimise that number over \(\delta\). The definition of minimax does not mention a prior. It is not ``first be Bayes, then also be minimax''.

The link used on this exam is a theorem, not a definition. \Cref{thm:constant-risk-bayes} says that a Bayes rule whose risk happens to be constant (for a prior with full support) is automatically minimax. That is a sufficient method for \emph{finding} a minimax rule---which is why (a) is reused in (b)---not a prerequisite for the notion. The second write-up of (b) unrolls \cref{def:risk-minimax} directly; the first write-up quotes the sufficient criterion.
\end{remark}

\begin{theorem}[Constant-risk Bayes is minimax]
\label{thm:constant-risk-bayes}
Let \(\Theta\) be finite. If \(\delta_*\) is Bayes for a prior \(\pi\) with full support on \(\Theta\), and \(R(\mu,\delta_*)\) does not depend on \(\mu\in\Theta\), then \(\delta_*\) is minimax on \(\Theta\), and \(\pi\) is least favourable.
\end{theorem}

\begin{proof}
Write \(r\) for the constant risk. For every competing \(\delta\),
\[
r
=
\sum_{\mu\in\Theta} R(\mu,\delta_*)\,\pi(\mu)
=
r(\pi,\delta_*)
\le
r(\pi,\delta)
=
\sum_{\mu\in\Theta} R(\mu,\delta)\,\pi(\mu)
\le
\sup_{\mu\in\Theta} R(\mu,\delta).
\]
The first inequality is the definition of a Bayes rule. Hence \(\sup R(\cdot,\delta_*)=r\le\sup R(\cdot,\delta)\).
\end{proof}

\begin{remark}[What is needed from the books]
\label{rem:p9-books}
Posterior odds are \citep[\S1.3]{gelman2013}. The \(L_1\mapsto\) median map is \citep{gelmannotes}. Frequentist risk and minimax are \citep[\S7.1]{hogg2019}, not a BDA topic. Hogg Chapter~11 is a statistical cousin of the same Bayes update, not the primary reference for (b).
\end{remark}

\begin{remark}[The parameter space for (b)]
\label{rem:p9-theta}
Minimax is a statement about a specified parameter space \(\Theta\). The prior lives on \(\{-1,1\}\), and that is the intended \(\Theta\). If one instead took \(\Theta=\mathbb{R}\), then \(R(\mu,\delta)=|\mu-\operatorname{sign}(Y)|\) has \(\sup_{\mu}R(\mu,\delta)=\infty\), so this \(\delta\) is not minimax on \(\mathbb{R}\). State \(\Theta=\{-1,1\}\) in the answer.
\end{remark}

\subsubsection{Solution of Problem 9}

\begin{proof}[Solution of (a)]
The two likelihoods are
\[
f(y\mid\mu)
=
(2\pi)^{-1/2}\exp\bigl(-(y-\mu)^2/2\bigr),
\qquad
\mu=\pm 1.
\]
The prior odds are \(1\). By \citep[(1.5)]{gelman2013} the posterior odds are the likelihood ratio
\[
\frac{\pi(\mu=1\mid y)}{\pi(\mu=-1\mid y)}
=
\frac{f(y\mid 1)}{f(y\mid -1)}
=
\exp\Bigl(-\frac{(y-1)^2}{2}+\frac{(y+1)^2}{2}\Bigr)
=
e^{2y}.
\]
Normalising,
\[
\pi(\mu=1\mid y)
=
\frac{e^{2y}}{1+e^{2y}}
=
\frac{1}{1+e^{-2y}},
\qquad
\pi(\mu=-1\mid y)
=
\frac{1}{1+e^{2y}}.
\]
In particular \(\pi(\mu=1\mid y)>\tfrac12\) if and only if \(y>0\). By \cref{lem:twopoint-l1}, an \(L_1\) Bayes rule is
\[
\delta(y)
=
\begin{cases}
1, & y>0,\\
-1, & y<0,\\
\text{any point of }[-1,1], & y=0.
\end{cases}
\]
The value at the single point \(y=0\) is irrelevant for risk, since \(\Pbb(Y=0\mid\mu)=0\). One may take \(\delta(y)=\operatorname{sign}(y)\) with the convention \(\operatorname{sign}(0)=0\).
\end{proof}

\begin{proof}[Solution of (b), via \cref{thm:constant-risk-bayes}]
This route uses \cref{rem:p9-theta} for the parameter space and then quotes the constant-risk Bayes criterion. Work on \(\Theta=\{-1,1\}\), \cref{rem:p9-theta}. For the rule of (a),
\[
R(1,\delta)
=
\E\bigl[|1-\delta(Y)|\mid\mu=1\bigr]
=
2\,\Pbb(Y<0\mid\mu=1)
=
2\Phi(-1),
\]
because \(\delta(Y)=-1\) on \(\{Y<0\}\) and the loss is then \(2\). Likewise
\[
R(-1,\delta)
=
2\,\Pbb(Y>0\mid\mu=-1)
=
2\Phi(-1).
\]
The risk is constant. The rule is Bayes for the uniform prior on \(\Theta\), which has full support. \Cref{thm:constant-risk-bayes} gives that \(\delta\) is minimax.
\end{proof}

\begin{proof}[Solution of (b), from \cref{def:risk-minimax}]
This route does not use \cref{rem:p9-theta} or \cref{thm:constant-risk-bayes}. It unrolls the definition of minimax on the two-point parameter space that the exam actually wrote down: the sampling model and the prior are both supported on \(\Theta=\{-1,1\}\), so the frequentist comparison is
\[
\sup_{\mu\in\{-1,1\}} R(\mu,\delta)
\stackrel{?}{=}
\inf_{\delta'}\sup_{\mu\in\{-1,1\}} R(\mu,\delta').
\]
Take \(\delta\) as in (a), with the convention \(\delta(0)=0\). On \(\{Y>0\}\) one has \(\delta(Y)=1\) and the loss is \(0\) if \(\mu=1\), or \(2\) if \(\mu=-1\). On \(\{Y<0\}\) the losses swap. The event \(\{Y=0\}\) is null under both Gaussians. Hence
\[
R(1,\delta)
=
2\,\Pbb(Y<0\mid\mu=1)
=
2\Phi(-1),
\qquad
R(-1,\delta)
=
2\,\Pbb(Y>0\mid\mu=-1)
=
2\Phi(-1).
\]
In particular \(\sup_{\mu\in\{-1,1\}}R(\mu,\delta)=2\Phi(-1)\). It remains to show that no competitor can have a smaller maximum risk.

Let \(\delta'\) be an arbitrary decision rule, and write \(\pi\) for the exam's uniform prior on \(\{-1,1\}\). The Bayes risk of any rule is the average of its two frequentist risks:
\[
r(\pi,\delta')
=
\tfrac12 R(-1,\delta')
+
\tfrac12 R(1,\delta')
\le
\sup_{\mu\in\{-1,1\}} R(\mu,\delta').
\]
Part~(a) already produced \(\delta\) as a Bayes rule for this \(\pi\), so \(r(\pi,\delta)\le r(\pi,\delta')\). Combined with the constant-risk computation above,
\[
2\Phi(-1)
=
r(\pi,\delta)
\le
r(\pi,\delta')
\le
\sup_{\mu\in\{-1,1\}} R(\mu,\delta').
\]
Thus every competitor has maximum risk at least \(2\Phi(-1)\), which is exactly \(\sup R(\cdot,\delta)\). By \cref{def:risk-minimax}, \(\delta\) is minimax on \(\{-1,1\}\).
\end{proof}

\begin{remark}[Sanity checks]
\label{rem:p9-checks}
The posterior log-odds equal \(2y\): large positive \(y\) votes for \(\mu=1\), as it should. The threshold \(y=0\) is the midpoint of \(\{-1,1\}\), which is the equal-prior MLE / likelihood-ratio test of \(\mu=1\) vs.\ \(\mu=-1\). The common risk \(2\Phi(-1)\approx 0.317\) is twice the Gaussian tail past one standard deviation: each true \(\mu\) is misclassified with probability \(\Phi(-1)\), and a mistake costs \(2\).
\end{remark}

\subsection{Problem 10. [10 pts.]}
Let \( X_1, X_2, \dots, X_n \) be i.i.d.\ random variables following \( \mathrm{Poisson}(\sqrt{\lambda}) \) distribution where \( \lambda > 0 \).
\begin{enumerate}[label=(\alph*)]
    \item Find the BUE (UMVUE) \( \tilde{\lambda} \) for \( \lambda \).
    \item Let \( \hat{\lambda} \) denote the MLE for \( \lambda \) based on \( X_1, X_2, \dots, X_n \), show that \( \sqrt{n}(\tilde{\lambda} - \hat{\lambda}) \to 0 \) in probability as \( n \to \infty \).
\end{enumerate}

\setcounter{section}{10}
\setcounter{definition}{0}

\subsubsection{Notes for Problem 10}

This is not a solution of the exam item. The mean of each \(X_i\) is \(\sqrt{\lambda}\), and the target \(\lambda\) is the square of that mean. The usual UMVUE machine applies: complete sufficient statistic, then an unbiased function of it. Part~(b) asks how close that function sits to the MLE.

Hogg writes MVUE for the same object \citep[Definition~7.1.1]{hogg2019}. After a complete sufficient statistic \(T\) is in hand, there are two computational techniques \citep[\S7.6]{hogg2019}: inspection of \(\E[\varphi(T)]\), and Rao--Blackwell of a crude unbiased estimator. The worked items below are those two techniques.

\begin{definition}[MVUE / UMVUE]
\label{def:mvue}
An estimator \(Y=u(X_1,\dots,X_n)\) of \(g(\theta)\) is a \emph{minimum variance unbiased estimator} (MVUE; also UMVUE) if \(\E_\theta[Y]=g(\theta)\) for every \(\theta\), and \(\Var_\theta(Y)\le\Var_\theta(Y')\) for every other unbiased estimator \(Y'\). See \citep[Definition~7.1.1]{hogg2019}. ``Uniformly'' means the variance comparison holds for all \(\theta\) in the parameter space.
\end{definition}

One cannot check this by listing competitors. The usable criterion is: restrict to functions of a complete sufficient statistic.

\begin{definition}[Completeness]
\label{def:complete}
A family \(\{f_T(t;\theta):\theta\in\Omega\}\) is \emph{complete} if \(\E_\theta[u(T)]=0\) for all \(\theta\in\Omega\) forces \(u(T)=0\) almost surely (under every member of the family). See \citep[\S7.4]{hogg2019}.
\end{definition}

\begin{theorem}[Lehmann--Scheff\'e]
\label{thm:ls}
Let \(T\) be a sufficient statistic for \(\theta\) whose family of distributions is complete. If some function \(\varphi(T)\) is unbiased for \(g(\theta)\), then \(\varphi(T)\) is the unique MVUE of \(g(\theta)\). See \citep[Theorem~7.4.1]{hogg2019}.
\end{theorem}

\begin{theorem}[Regular exponential family]
\label{thm:expfam-complete}
If the observations are i.i.d.\ from a regular one-parameter exponential family
\[
f(x;\theta)
=
\exp\bigl(p(\theta)K(x)+q(\theta)+S(x)\bigr),
\]
Assume the support does not depend on \(\theta\), and that \(p\) is nonconstant and continuously differentiable. Then \(T=\sum_{i=1}^n K(X_i)\) is a complete sufficient statistic for \(\theta\). See \citep[Theorems~7.5.1--7.5.2]{hogg2019}. A one-to-one reparametrisation \(\theta\mapsto h(\theta)\) does not change \(T\).
\end{theorem}

Poisson is the model case: \(K(x)=x\), so \(T=\sum X_i\) is complete sufficient for the mean, whatever name one gives that mean.

\begin{fact}[Poisson additivity]
\label{fact:poisson-additive}
If \(X_1,\dots,X_k\) are independent and \(X_j\sim\mathrm{Poisson}(\lambda_j)\), then
\[
\sum_{j=1}^{k} X_j
\sim
\mathrm{Poisson}\Bigl(\sum_{j=1}^{k}\lambda_j\Bigr).
\]
In particular, for i.i.d.\ \(X_i\sim\mathrm{Poisson}(\theta)\), the sample sum \(T=\sum_{i=1}^{n}X_i\sim\mathrm{Poisson}(n\theta)\). For Problem~10 with \(\theta=\sqrt{\lambda}\), therefore \(T\sim\mathrm{Poisson}(n\sqrt{\lambda})\).
\end{fact}

\begin{remark}[Independent sums that stay in the same family]
\label{rem:sum-closure}
Several standard families are \emph{closed under independent addition}: the sum has the same type, with parameters combined in a fixed rule. No proofs here; these are exam-side facts.
\begin{enumerate}[label=\textup{(\arabic*)}]
    \item
    \textbf{Bernoulli / binomial.} If \(X\sim\mathrm{Bernoulli}(p)\) and \(Y\sim\mathrm{Bernoulli}(p)\) are independent, then \(X+Y\sim\mathrm{Bin}(2,p)\). More generally, independent \(\mathrm{Bin}(n_j,p)\) with the same \(p\) sum to \(\mathrm{Bin}(\sum_j n_j,p)\). A sum of \(n\) i.i.d.\ Bernoulli\((p)\) is \(\mathrm{Bin}(n,p)\).
    \item
    \textbf{Poisson.} Rates add; see \cref{fact:poisson-additive}.
    \item
    \textbf{Normal.} If \(X\sim N(\mu_1,\sigma_1^2)\) and \(Y\sim N(\mu_2,\sigma_2^2)\) are independent, then \(X+Y\sim N(\mu_1+\mu_2,\sigma_1^2+\sigma_2^2)\). Any linear combination \(aX+bY\) with \(X,Y\) independent normal is again normal (means and variances combine by the usual rules).
    \item
    \textbf{Gamma / exponential / chi-square.} If \(X\sim\Gamma(\alpha_1,\beta)\) and \(Y\sim\Gamma(\alpha_2,\beta)\) are independent with the same rate \(\beta>0\) (shape--rate parametrisation), then \(X+Y\sim\Gamma(\alpha_1+\alpha_2,\beta)\). Hence i.i.d.\ \(\mathrm{Exp}(\beta)\) sum to \(\Gamma(n,\beta)\). Independent \(\chi_k^2\) variables sum to \(\chi_{k_1+k_2}^2\).
    \item
    \textbf{Negative binomial.} If \(X\sim\mathrm{NB}(r_1,p)\) and \(Y\sim\mathrm{NB}(r_2,p)\) are independent with the same success probability \(p\in(0,1)\), then \(X+Y\sim\mathrm{NB}(r_1+r_2,p)\). Geometric\((p)\) is the case \(r=1\).
    \item
    \textbf{Cauchy (same scale).} If \(X\sim\mathrm{Cauchy}(\mu_1,\gamma)\) and \(Y\sim\mathrm{Cauchy}(\mu_2,\gamma)\) are independent with the same scale \(\gamma\), then \(X+Y\sim\mathrm{Cauchy}(\mu_1+\mu_2,2\gamma)\). Cauchy is the textbook example of a \emph{stable} law: an appropriately scaled sum of i.i.d.\ copies stays Cauchy.
    \item
    \textbf{Multinomial (merge categories).} If \((X_1,\dots,X_k)\sim\mathrm{Mult}(n;\pi_1,\dots,\pi_k)\), then any subvector obtained by adding coordinatewise over a partition of the categories is again multinomial, with the same \(n\) and merged probability vector.
\end{enumerate}
Poisson additivity is the one Problem~10 uses directly: \(T=\sum X_i\) is not only sufficient, its law is explicit. The normal and gamma items are the other closures one sees most often in QE problems (CLT limits, waiting-time sums, chi-square tests).
\end{remark}

\begin{remark}[Two computational techniques]
\label{rem:umvue-recipe}
After \(T\) is in hand, there are two standard ways to produce \(\varphi(T)\).
\begin{enumerate}[label=\textup{(\arabic*)}]
    \item
    \emph{Inspection.} Compute \(\E[\psi(T)]\) for a natural function \(\psi\) (often a power, a falling factorial, or a plug-in) and rescale to kill the bias. Typical algebra: if \(\E[T]=\mu\) then \(\E[T^{2}]=\Var(T)+\mu^{2}\), so \(T^{2}\) overshoots \(\mu^{2}\) by \(\Var(T)\); subtract that, or use a falling factorial \(\E[T(T-1)]=\mu^{2}\) when \(T\) is Poisson. This is Hogg's first technique in \S7.6, and the route in \cref{ex:hogg-761} and \cref{ex:yau-2011-p1p}.
    \item
    \emph{Rao--Blackwell.} Start from any unbiased estimator \(Z\) that is easy to write down (one observation, an indicator, \((-1)^{X_1}\)), then replace \(Z\) by \(\E[Z\mid T]\). Completeness upgrades the result from ``better than \(Z\)'' to ``the MVUE''. This is Hogg's second technique, and the route in \cref{ex:hogg-767}.
\end{enumerate}
The written recipe is the same in either case: complete sufficient statistic \(\to\) unbiased function of it \(\to\) Lehmann--Scheff\'e.
\end{remark}

\begin{exercise}[{Hogg Example~7.5.3: Poisson mean}]
\label{ex:hogg-753}
Let \(X_1,\dots,X_n\) be i.i.d.\ \(\mathrm{Poisson}(\theta)\), \(\theta>0\). Find the MVUE of \(\theta\).
\end{exercise}

\begin{proof}[Solution]
The pmf is \(\exp\bigl(x\log\theta-\theta-\log(x!)\bigr)\), a regular exponential family with \(K(x)=x\). By \cref{thm:expfam-complete}, \(T=\sum X_i\) is complete sufficient. \(\E[T]=n\theta\), so \(\overline{X}=T/n\) is unbiased and a function of \(T\). \Cref{thm:ls} gives that \(\overline{X}\) is the MVUE of \(\theta\).
\end{proof}

That is the mean. Problem~10 asks for the \emph{square} of the mean. The next item is the Gaussian analogue of that square.

\begin{exercise}[{Hogg 7.6.1: MVUE of a squared mean}]
\label{ex:hogg-761}
Let \(X_1,\dots,X_n\) be i.i.d.\ \(N(\theta,1)\). Find the MVUE of \(\theta^2\).
\end{exercise}

\begin{proof}[Solution]
\(\overline{X}\) is complete sufficient for \(\theta\) \citep[Example~7.5.2]{hogg2019}, and \(\overline{X}\sim N(\theta,1/n)\), so
\[
\E[\overline{X}^2]
=
\Var(\overline{X})+\bigl(\E\overline{X}\bigr)^2
=
\frac1n+\theta^2.
\]
Thus \(\overline{X}^2-1/n\) is unbiased for \(\theta^2\) and is a function of \(\overline{X}\). \Cref{thm:ls} gives that it is the MVUE. (Hogg's hint is exactly this computation of \(\E[\overline{X}^2]\).)
\end{proof}

The Poisson version of the same inspection uses a falling factorial instead of subtracting a constant, because \(\E[T(T-1)]=(\E T)^2\) for Poisson \(T\). That identity is the only extra Poisson input; the logic is \cref{ex:hogg-761}.

\begin{exercise}[{Hogg 7.6.7: Poisson MVUE versus MLE}]
\label{ex:hogg-767}
Let \(X_1,\dots,X_n\) be i.i.d.\ \(\mathrm{Poisson}(\theta)\), \(\theta>0\), and \(T=\sum X_i\).
\begin{enumerate}[label=\textup{(\alph*)}]
    \item
    Show \(\E[(-1)^{X_1}\mid T=t]=(1-2/n)^t\), and conclude that \((1-2/n)^T\) is the MVUE of \(e^{-2\theta}\).
    \item
    The MLE of \(e^{-2\theta}\) is \(e^{-2\overline{X}}\). Show that the two estimators are asymptotically equivalent, in the sense \((1-2/n)^{n\overline{X}}\approx e^{-2\overline{X}}\) for large \(n\).
\end{enumerate}
\end{exercise}

\begin{proof}[Solution]
Given \(T=t\), the vector \((X_1,\dots,X_n)\) is multinomial with \(t\) trials and equal cell probabilities \(1/n\), so \(X_1\mid T=t\sim\mathrm{Bin}(t,1/n)\). The binomial pgf at \(-1\) is
\[
\E\bigl[(-1)^{X_1}\mid T=t\bigr]
=
\bigl(1-\tfrac1n-\tfrac1n\bigr)^t
=
\bigl(1-\tfrac2n\bigr)^t.
\]
Unconditionally, \(\E[(-1)^{X_1}]=e^{-2\theta}\) by the Poisson pgf (Hogg's Remark~7.6.1). Rao--Blackwell plus \cref{thm:ls} gives the MVUE \((1-2/n)^T\).

The MLE of \(\theta\) is \(\overline{X}\), so the MLE of \(e^{-2\theta}\) is \(e^{-2\overline{X}}\) by invariance. Writing \(T=n\overline{X}\),
\[
\bigl(1-\tfrac2n\bigr)^{n\overline{X}}
=
\exp\bigl(n\overline{X}\log(1-2/n)\bigr),
\]
and \(n\log(1-2/n)\to-2\), so the MVUE tracks the MLE. This is the pattern of Problem~10(b): both estimators are functions of the same complete sufficient statistic, and their difference is of smaller order than \(n^{-1/2}\).
\end{proof}

Hogg's Remark~7.6.1 also records a warning: for \(n=1\), the MVUE of \(e^{-2\theta}\) is \((-1)^X\in\{\pm1\}\), which is a terrible practical estimate of a number in \((0,1)\). Unbiasedness is a principle, not a theorem that the estimator looks like the target.

\begin{exercise}[{Yau 2011 Probability, inspection for \(p(1-p)\)}]
\label{ex:yau-2011-p1p}
Let \(X_1,\dots,X_n\) be i.i.d.\ \(\mathrm{Bernoulli}(p)\), and \(T=\sum X_i\). Find the UMVUE of \(v(p)=p(1-p)\). This is \citep[Problem~5(b)]{yau2011}.
\end{exercise}

\begin{proof}[Solution]
Bernoulli is a regular exponential family, so \(T\sim\mathrm{Bin}(n,p)\) is complete sufficient. The plug-in \(\hat p(1-\hat p)\) with \(\hat p=T/n\) is biased:
\[
\E[\hat p(1-\hat p)]
=
p-\bigl(\Var(\hat p)+p^2\bigr)
=
p(1-p)-\frac{p(1-p)}{n}
=
\frac{n-1}{n}\,p(1-p).
\]
Rescaling gives the unbiased function \(T(n-T)/(n(n-1))\). \Cref{thm:ls} finishes. The same inspection, with a second-moment correction, is Hogg's Example~7.6.1 for the binomial variance \citep[\S7.6]{hogg2019}.
\end{proof}

\begin{remark}[Nearby items]
\label{rem:p10-sources}
\begin{enumerate}[label=\textup{(\arabic*)}]
    \item
    The Yau-contest template is \citep{yau2011,yau2021}. Besides \cref{ex:yau-2011-p1p}, the 2021 linear-model item estimates \(\mathbf{x}_1\boldsymbol{\beta}/\sigma\) from the complete sufficient pair \((X\hat{\boldsymbol{\beta}},S^2)\) by correcting the \(\chi^2\) bias of \(S^{-1}\). That is inspection in a two-parameter exponential family.
    \item
    Qiuzhen 2025 Fall Problem~8 is the continuous cousin of \cref{ex:hogg-767}: exponential sample, UMVUE of \(e^{-\lambda x}\) by Rao--Blackwell of \(\mathbf{1}_{\{X_1>x\}}\). 2025 Fall Problem~9(a) is a Bernoulli function of \(p\), same \(T\) as \cref{ex:yau-2011-p1p}.
\end{enumerate}
\end{remark}

\begin{remark}[Hints for Problem~10, in order]
\label{rem:p10-hints}
\begin{enumerate}[label=\textup{(\arabic*)}]
    \item
    Write the Poisson(\(\sqrt{\lambda}\)) pmf as an exponential family in the mean \(\sqrt{\lambda}\). \Cref{thm:expfam-complete} names the complete sufficient statistic; it is the same \(T\) as in \cref{ex:hogg-753}.
    \item
    The target is the square of the mean. Compute \(\E[\overline{X}^2]\) (or \(\E[T(T-1)]\)) as in \cref{ex:hogg-761}, and correct the bias. Do not stop at \(\overline{X}^2\): that is the MLE, not the UMVUE.
    \item
    The MLE of \(\sqrt{\lambda}\) is \(\overline{X}\) (or a truncated version if you worry about the boundary \(\lambda=0\), which has probability tending to zero). Invariance gives \(\hat\lambda\). Both \(\tilde\lambda\) and \(\hat\lambda\) are functions of \(T\). Their difference is of smaller order than \(n^{-1/2}\), which is the same comparison as \cref{ex:hogg-767}(b).
\end{enumerate}
\end{remark}

\subsubsection{Solution of Problem 10}

\paragraph{Answer.}

Let \(T=\sum_{i=1}^{n}X_i\) and \(\overline{X}=T/n\).

(a)
\[
\tilde{\lambda}
=
\frac{T(T-1)}{n^{2}}
=
\overline{X}^{2}-\frac{\overline{X}}{n}.
\]
This is the UMVUE of \(\lambda\).

(b) The MLE is \(\hat{\lambda}=\overline{X}^{2}\). As \(n\to\infty\),
\[
\sqrt{n}\,(\tilde{\lambda}-\hat{\lambda})
=
-\frac{T}{n^{3/2}}
\xrightarrow{\Pbb}
0.
\]

\begin{proof}[Exam proof]
The pmf of \(X_1\) is
\[
\exp\Bigl(x\log\sqrt{\lambda}-\sqrt{\lambda}-\log(x!)\Bigr),
\]
a regular one-parameter exponential family with natural statistic \(K(x)=x\) and parameter \(\theta=\sqrt{\lambda}\). By \cref{thm:expfam-complete}, \(T=\sum_{i=1}^{n}X_i\) is complete sufficient for \(\sqrt{\lambda}\). Under the model, \(T\sim\mathrm{Poisson}(n\sqrt{\lambda})\) by \cref{fact:poisson-additive}.

\textbf{(a).}
The target is \(\lambda=(\sqrt{\lambda})^{2}\), the square of the Poisson mean. For \(T\sim\mathrm{Poisson}(\mu)\),
\[
\E[T(T-1)]=\mu^{2}.
\]
Here \(\mu=n\sqrt{\lambda}\), so
\[
\E\Bigl[\frac{T(T-1)}{n^{2}}\Bigr]
=
\frac{(n\sqrt{\lambda})^{2}}{n^{2}}
=
\lambda.
\]
The estimator \(\tilde{\lambda}=T(T-1)/n^{2}\) is therefore unbiased for \(\lambda\) and is a function of the complete sufficient statistic \(T\). \Cref{thm:ls} gives that it is the UMVUE. The algebra
\[
\frac{T(T-1)}{n^{2}}
=
\Bigl(\frac{T}{n}\Bigr)^{2}-\frac{T}{n^{2}}
=
\overline{X}^{2}-\frac{\overline{X}}{n}
\]
matches the inspection pattern of \cref{ex:hogg-761}.

\textbf{(b).}
The likelihood depends on the data only through \(\overline{X}\), which is the MLE of \(\sqrt{\lambda}\) (same argument as \cref{ex:hogg-753} with mean \(\sqrt{\lambda}\)). By invariance, \(\hat{\lambda}=\overline{X}^{2}\). Both estimators are functions of \(T\), and
\[
\tilde{\lambda}-\hat{\lambda}
=
\frac{T(T-1)-T^{2}}{n^{2}}
=
-\frac{T}{n^{2}}.
\]
Hence
\[
\sqrt{n}\,(\tilde{\lambda}-\hat{\lambda})
=
-\frac{T}{n^{3/2}}.
\]
It remains to show that this tends to \(0\) in the sense of \cref{def:conv-p}. Under the model, \(T\sim\mathrm{Poisson}(n\sqrt{\lambda})\), so
\[
\E[T]=n\sqrt{\lambda},
\qquad
\Var(T)=n\sqrt{\lambda}.
\]
Two equivalent writings:

\emph{Chebyshev on \(T/n\).} The sample mean of the original i.i.d.\ Poisson(\(\sqrt{\lambda}\)) observations is \(T/n=\overline{X}\). \Cref{lem:chebyshev} gives, for every \(\varepsilon>0\),
\[
\Pbb\bigl(|T/n-\sqrt{\lambda}|\ge\varepsilon\bigr)
\le
\frac{\Var(T/n)}{\varepsilon^2}
=
\frac{n\sqrt{\lambda}}{n^{2}\varepsilon^2}
=
\frac{\sqrt{\lambda}}{n\varepsilon^2}
\to 0.
\]
Thus \(T/n\xrightarrow{\Pbb}\sqrt{\lambda}\). This is also \cref{thm:wlln} applied to \(X_1,\dots,X_n\). Scaling by the deterministic factor \(n^{-1/2}\to 0\),
\[
\frac{T}{n^{3/2}}
=
n^{-1/2}\cdot\frac{T}{n}.
\]
\Cref{lem:slutsky} with \(A_n\equiv 0\), \(B_n=n^{-1/2}\), and \(X_n=T/n\) yields \(T/n^{3/2}\xrightarrow{\Pbb}0\cdot\sqrt{\lambda}=0\).

\emph{Chebyshev directly on the displayed object.} One need not pass through the weak law. Write \(Z_n=T/n^{3/2}\). Then
\[
\E[Z_n]
=
\frac{\sqrt{\lambda}}{\sqrt{n}}\to 0,
\qquad
\Var(Z_n)
=
\frac{\sqrt{\lambda}}{n^{2}}\to 0,
\]
so \cref{lem:chebyshev} yields \(Z_n-\E[Z_n]\xrightarrow{\Pbb}0\). Adding the deterministic sequence \(\E[Z_n]\to 0\) gives \(Z_n\xrightarrow{\Pbb}0\) in the sense of \cref{def:conv-p}. Equivalently, for every \(\varepsilon>0\) and all large \(n\) one has \(|\E[Z_n]|<\varepsilon/2\), hence
\[
\Pbb(|Z_n|\ge\varepsilon)
\le
\Pbb\bigl(|Z_n-\E[Z_n]|\ge\varepsilon/2\bigr)
\le
\frac{4\Var(Z_n)}{\varepsilon^{2}}
\to 0.
\]
Either way, \(\sqrt{n}\,(\tilde{\lambda}-\hat{\lambda})\xrightarrow{\Pbb}0\).
\end{proof}

\begin{remark}[Sanity checks]
\label{rem:p10-checks}
The falling factorial \(T(T-1)\) is the Poisson analogue of subtracting \(1/n\) from \(\overline{X}^{2}\) in \cref{ex:hogg-761}: it removes the \(\Var(\overline{X})\) bias exactly. Part~(b) is sharper than asymptotic normality: the difference is \(O_{\Pbb}(n^{-1})\), not merely \(O_{\Pbb}(n^{-1/2})\), because both estimators are the same function of \(T\) up to a deterministic \(1/n\) correction.
\end{remark}

\subsection{Problem 11. [15 pts.]}
Let \( Y = \sigma\bigl(\rho |U| + \sqrt{1 - \rho^2}\, V\bigr) \) where \( U, V \) are independent \( N(0, 1) \) random variables and \( \rho \in (-1, 1) \), \( \sigma \in (0, \infty) \) are unknown parameters.
\begin{enumerate}[label=(\alph*)]
    \item Show that the pdf of \( Y \) is given by

    \[
    f(y \mid \rho, \sigma)
    =
    \frac{a}{\sqrt{2\pi\sigma^2}}
    \exp\left(-\frac{y^2}{2\sigma^2}\right)
    \Phi\left( \frac{\rho y}{b\sigma\sqrt{1-\rho^2}} \right),
    \qquad -\infty < y < \infty
    \]

    for some positive constants \( a, b \). Write down the values of \( a \) and \( b \). In the above expression, \( \Phi(z) \) denotes the CDF for a standard normal random variable.
    \item Consider \( n \) observations \( Y_1, \dots, Y_n \), modeled as \( Y_i \stackrel{\mathrm{iid}}{\sim} f(y_i \mid \rho, \sigma) \), \( \rho \in (-1, 1) \), \( \sigma \in (0, \infty) \). Does a maximum likelihood estimate of \( (\rho, \sigma) \in (-1, 1) \times (0, \infty) \) always exist? Please justify your answer.
    \item For the same model, suppose we are interested in testing \( H_0 : \rho = 0 \) versus \( H_1 : \rho \neq 0 \). What is the value of \( c \) such that the test that rejects \( H_0 \) if

    \[
    \frac{\sqrt{n}\, |\overline{Y}|}{s_Y} > c
    \]

    has size \( \alpha = 0.05 \)? Identify \( c \) as a specific quantile of a named distribution. In the above expression, \( \overline{Y} \) is the sample mean and \( s_Y \) is the sample standard deviation.
\end{enumerate}

\setcounter{section}{11}
\setcounter{definition}{0}

\subsubsection{How to compute a pdf (and what not to reach for first)}

The relation you remembered is real, and it is written in the 2024 notes \citep{ren2024}. It is not the working device for part~(a).

\begin{definition}[Density as the derivative of the cdf]
\label{def:pdf-from-cdf}
If \(F(x)=\int_{-\infty}^{x}p(t)\,\mathrm{d}t\) with \(p\ge 0\) Riemann integrable, then \(p\) is a density of \(F\), and at every continuity point of \(p\) one has \(p(x)=F'(x)\). This is \citep[Definition~4.3.2 and notes~2--4]{ren2024}; the same calculus step is Hogg's cdf technique \citep[\S1.7.2]{hogg2019}. Changing \(p\) on a null set does not change \(F\).
\end{definition}

\begin{definition}[Characteristic function, and the inversion that actually recovers a density]
\label{def:chf-inversion}
The characteristic function of a real random variable \(\xi\) is \(f(t)=\E[e^{it\xi}]\). It always exists, and it determines the law uniquely \citep[\S8.2--8.4]{ren2024}. If \(\int_{\mathbb{R}}|f(t)|\,\mathrm{d}t<\infty\), then the law is absolutely continuous with
\[
p(x)
=
\frac{1}{2\pi}
\int_{-\infty}^{\infty}
e^{-itx}f(t)\,\mathrm{d}t.
\]
If moreover \(\int_{\mathbb{R}}|t|^{n}|f(t)|\,\mathrm{d}t<\infty\), then \(p\) is \(C^{n}\) and
\[
p^{(n)}(x)
=
\frac{1}{2\pi}
\int_{-\infty}^{\infty}
(-it)^{n}e^{-itx}f(t)\,\mathrm{d}t.
\]
This is \citep[Theorem~8.4.4]{ren2024}, and is the density case \cref{thm:chf-inversion}(ii) already recorded for Problem~2. The derivatives of \(f\) at \(0\), when they exist, recover moments: \(f^{(n)}(0)=i^{n}\E[\xi^{n}]\) under the usual integrability. They do \emph{not} produce \(p\). Hogg records the same distinction in \citep[Remark~1.9.1]{hogg2019}.
\end{definition}

\begin{study}[A short menu for ``find the pdf of \(Y=g(X)\)'']
\label{study:pdf-methods}
In exam problems the density is almost never obtained by writing down a characteristic function and inverting it. The notes themselves compute laws from characteristic functions by the \emph{product formula plus uniqueness}, not by evaluating the inversion integral \citep[\S9.1]{ren2024}. A practical order, matching \citep[\S4.7--4.8]{ren2024} and \citep[\S1.7, \S2.2--2.3]{hogg2019}:
\begin{enumerate}[label=\textup{(\arabic*)}]
    \item
    \textbf{Recognize a named law}, or an affine image of one (location-scale of \(\mathcal{N}(0,1)\), gamma, etc.).
    \item
    \textbf{Condition} on a convenient coordinate, then mix. If \(Y\mid W=w\) has a density \(f_{Y\mid W}(\,\cdot\mid w)\) and \(W\) has density \(f_{W}\),
    \[
    f_{Y}(y)
    =
    \int f_{Y\mid W}(y\mid w)\,f_{W}(w)\,\mathrm{d}w.
    \]
    This is the continuous form of Hogg's \(f(y)=\int f(y\mid x)f(x)\,\mathrm{d}x\) \citep[\S2.3]{hogg2019}.
    \item
    \textbf{Cdf, then differentiate.} Write \(F_{Y}(y)=\Pbb(g(X)\le y)\) by inverting \(g\) on each monotonic piece, then use \cref{def:pdf-from-cdf}. This is \citep[\S4.7]{ren2024} and \citep[\S1.7.2]{hogg2019}.
    \item
    \textbf{One-to-one map plus Jacobian.} This is \cref{thm:jacobian}: if \(Y=\phi(X)\) is \(C^{1}\) and one-to-one,
    \[
    f_{Y}(y)
    =
    f_{X}\bigl(\phi^{-1}(y)\bigr)
    \bigl|(\phi^{-1})'(y)\bigr|.
    \]
    Multivariate: adjoin dummy coordinates so the map is \(\mathbb{R}^{n}\to\mathbb{R}^{n}\), multiply by \(|\det D\phi^{-1}|\), then integrate the dummies out \citep[Theorem~1.7.1, \S2.2]{hogg2019}, \citep[\S4.8]{ren2024}.
    \item
    \textbf{Convolution} for a sum of jointly continuous variables \citep[Proposition~4.8.3]{ren2024}, \citep[(2.2.5)]{hogg2019}:
    \[
    f_{\xi+\eta}(z)
    =
    \int_{\mathbb{R}} p_{\xi,\eta}(z-x,x)\,\mathrm{d}x.
    \]
    If \(\xi\perp\eta\), the integrand factors.
    \item
    \textbf{Characteristic functions}, only when the object is a sum of independents (or a Poisson mixture, etc.) and the product \(f_{Y}(t)=\prod f_{X_{j}}(t)\) is a named transform. Then uniqueness identifies the law \citep[\S9.1]{ren2024}. Inversion \cref{def:chf-inversion} is the existence theorem behind that uniqueness; it is a poor first computational step unless \(\int|f|<\infty\) and the Fourier integral is obvious.
\end{enumerate}
Part~(a) is item~(2), equivalently item~(5) after writing \(Y=\sigma\rho\,|U|+\sigma\sqrt{1-\rho^{2}}\,V\). The characteristic function of \(Y\) exists and equals
\[
\E\bigl[e^{it\sigma\rho|U|}\bigr]\,
\exp\bigl(-\tfrac12\sigma^{2}(1-\rho^{2})t^{2}\bigr),
\]
but inverting that product by hand is strictly harder than completing the square in the Gaussian integral below. Do not start there.
\end{study}

\subsubsection{Definitions used in Problem 11}

\begin{notation}[Standard normal density and cdf]
\label{not:phi-Phi}
Write
\[
\phi(z)
=
\frac{1}{\sqrt{2\pi}}e^{-z^{2}/2},
\qquad
\Phi(z)
=
\int_{-\infty}^{z}\phi(u)\,\mathrm{d}u.
\]
A location-scale copy: if \(Z\sim\mathcal{N}(0,1)\) then \(\mu+\tau Z\) has density \(\tau^{-1}\phi\bigl((x-\mu)/\tau\bigr)\) for \(\tau>0\). Half-normal: \(|U|\) for \(U\sim\mathcal{N}(0,1)\) has density \(\sqrt{2/\pi}\,e^{-w^{2}/2}\) on \((0,\infty)\). This is the \(\alpha=1\) case of \citep[\S4.7, Exercise~1(d)]{ren2024}.
\end{notation}

\begin{definition}[Sample standard deviation]
\label{def:sample-sd}
For a sample \(Y_{1},\dots,Y_{n}\),
\[
s_{Y}^{2}
=
\frac{1}{n-1}\sum_{i=1}^{n}(Y_{i}-\overline{Y})^{2}.
\]
This is Hogg's \(S^{2}\) \citep[\S4.2, Example~4.2.1]{hogg2019}. Under \(Y_{i}\stackrel{\mathrm{iid}}{\sim}\mathcal{N}(\mu,\sigma^{2})\),
\[
T
=
\frac{\overline{Y}-\mu}{s_{Y}/\sqrt{n}}
\sim
t_{n-1}
\]
by \citep[Theorem~3.6.1(d)]{hogg2019}.
\end{definition}

\subsubsection{Solution of Problem 11}

\paragraph{Answer.}

(a) \(a=2\) and \(b=1\). (b)~No: if every observation has the same strict sign, the likelihood increases as \(\rho\) tends to that sign and the supremum is not attained on the open set \((-1,1)\times(0,\infty)\). (c)~Under \cref{def:sample-sd}, \(c=t_{n-1,\,0.975}\), the \(0.975\)-quantile of Student's \(t\) with \(n-1\) degrees of freedom.

\begin{proof}[Correct proof]
\textbf{(a).}
Let \(W:=|U|\). Since \(U\perp V\) and \(W\) is a function of \(U\), also \(W\perp V\). For \(w>0\), symmetry of \(\mathcal{N}(0,1)\) gives
\[
F_{W}(w)
=
\Pbb(|U|\le w)
=
\Pbb(-w\le U\le w)
=
\Phi(w)-\Phi(-w)
=
2\Phi(w)-1.
\]
Differentiating as in \cref{def:pdf-from-cdf} (equivalently, \(f_{W}(w)=f_{U}(w)+f_{U}(-w)\) on \((0,\infty)\)),
\[
f_{W}(w)
=
2\phi(w)
=
\sqrt{\frac{2}{\pi}}\,e^{-w^{2}/2},
\qquad w>0,
\]
the half-normal density recorded in \cref{not:phi-Phi}. Conditionally on \(W=w\),
\[
Y
=
\sigma\rho w
+
\sigma\sqrt{1-\rho^{2}}\,V
\sim
\mathcal{N}\bigl(\sigma\rho w,\,\sigma^{2}(1-\rho^{2})\bigr).
\]
Mixing as in \cref{study:pdf-methods}(2),
\begin{align*}
f(y\mid\rho,\sigma)
&=
\int_{0}^{\infty}
\frac{1}{\sigma\sqrt{1-\rho^{2}}}\,
\phi\Bigl(\frac{y-\sigma\rho w}{\sigma\sqrt{1-\rho^{2}}}\Bigr)
\sqrt{\frac{2}{\pi}}
e^{-w^{2}/2}
\,\mathrm{d}w.
\end{align*}
The exponent in the integrand is
\[
-\frac{(y-\sigma\rho w)^{2}}{2\sigma^{2}(1-\rho^{2})}
-
\frac{w^{2}}{2}
=
-\frac{y^{2}}{2\sigma^{2}}
-
\frac{\bigl(w-\rho y/\sigma\bigr)^{2}}{2(1-\rho^{2})}.
\]
(The cross term completes the square; the leftover in \(y\) is exactly the \(\mathcal{N}(0,\sigma^{2})\) exponent.) Hence
\[
f(y\mid\rho,\sigma)
=
\frac{1}{\pi\sigma\sqrt{1-\rho^{2}}}
\exp\Bigl(-\frac{y^{2}}{2\sigma^{2}}\Bigr)
\int_{0}^{\infty}
\exp\Bigl(
-\frac{\bigl(w-\rho y/\sigma\bigr)^{2}}{2(1-\rho^{2})}
\Bigr)
\,\mathrm{d}w,
\]
where the prefactor used \(\sqrt{2/\pi}\,/\,\sqrt{2\pi}=1/\pi\). Substitute \(t=\bigl(w-\rho y/\sigma\bigr)/\sqrt{1-\rho^{2}}\). The lower limit becomes \(-\rho y/(\sigma\sqrt{1-\rho^{2}})\), and
\[
\int_{0}^{\infty}
\exp\Bigl(
-\frac{\bigl(w-\rho y/\sigma\bigr)^{2}}{2(1-\rho^{2})}
\Bigr)
\,\mathrm{d}w
=
\sqrt{1-\rho^{2}}\,\sqrt{2\pi}\,
\Phi\Bigl(\frac{\rho y}{\sigma\sqrt{1-\rho^{2}}}\Bigr).
\]
Therefore
\[
f(y\mid\rho,\sigma)
=
\frac{2}{\sqrt{2\pi\sigma^{2}}}
\exp\Bigl(-\frac{y^{2}}{2\sigma^{2}}\Bigr)
\Phi\Bigl(\frac{\rho y}{\sigma\sqrt{1-\rho^{2}}}\Bigr).
\]
This is the claimed form with \(a=2\) and \(b=1\).

Two checks. If \(\rho=0\), then \(\Phi(0)=1/2\) and \(f\) collapses to the \(\mathcal{N}(0,\sigma^{2})\) density, matching \(Y=\sigma V\). If \(y\) and \(\rho\) have opposite signs and \(|\rho|\to 1\), the argument of \(\Phi\) tends to \(-\infty\) and the density on that half-line vanishes, matching the limit \(Y\to\sigma|U|\) (folded normal).

\textbf{(b).}
The parameter space \(\Theta=(-1,1)\times(0,\infty)\) is open. Write \(\lambda:=\rho/\sqrt{1-\rho^{2}}\in\mathbb{R}\), so the argument of \(\Phi\) is \(\lambda y/\sigma\). For an i.i.d.\ sample the likelihood is
\[
L(\rho,\sigma)
=
\Bigl(\frac{2}{\sigma\sqrt{2\pi}}\Bigr)^{n}
\exp\Bigl(-\frac{1}{2\sigma^{2}}\sum_{i=1}^{n}Y_{i}^{2}\Bigr)
\prod_{i=1}^{n}
\Phi\Bigl(\frac{\lambda Y_{i}}{\sigma}\Bigr).
\]
Suppose every \(Y_{i}\) has the same strict sign, say \(Y_{i}>0\) for all \(i\). For each fixed \(\sigma>0\), the product of \(\Phi\) is strictly less than \(1\) at every finite \(\lambda\), and tends to \(1\) as \(\lambda\to+\infty\) (i.e.\ \(\rho\to 1^{-}\)). Thus
\[
\sup_{\Theta}L
=
\sup_{\sigma>0}
\Bigl(\frac{2}{\sigma\sqrt{2\pi}}\Bigr)^{n}
\exp\Bigl(-\frac{1}{2\sigma^{2}}\sum_{i=1}^{n}Y_{i}^{2}\Bigr),
\]
the half-normal profile, and the supremum is not attained at any \(\rho\in(-1,1)\). The same argument with \(\lambda\to-\infty\) applies if every observation is negative. Hence an MLE on \(\Theta\) does not always exist.

(If the sample takes both strict signs, then \(L\to 0\) as \(|\lambda|\to\infty\), and also as \(\sigma\to 0\) or \(\sigma\to\infty\) whenever \(\sum Y_{i}^{2}>0\). Existence on a mixed-sign sample is a separate, typically yes, question. The exam asked whether an MLE \emph{always} exists.)

\textbf{(c).}
Under \(H_{0}:\rho=0\) one has \(Y_{i}=\sigma V_{i}\stackrel{\mathrm{iid}}{\sim}\mathcal{N}(0,\sigma^{2})\). By \cref{def:sample-sd},
\[
T
=
\frac{\sqrt{n}\,\overline{Y}}{s_{Y}}
\sim
t_{n-1}.
\]
The test is two-sided of size \(\alpha=0.05\), so
\[
c
=
t_{n-1,\,0.975},
\]
the unique number with \(\Pbb(|t_{n-1}|>c)=0.05\). (If an exam convention takes \(s_{Y}^{2}\) with divisor \(n\) rather than \(n-1\), the exact null law is no longer \(t_{n-1}\); under Hogg's \(S^{2}\) it is.)
\end{proof}

\begin{remark}[Why the \(\Phi\) factor is there]
\label{rem:p11-skew}
The factor \(2\Phi(\lambda y/\sigma)\) tilts a centred Gaussian towards the sign of \(\lambda\). It is not an extra parameterisation trick: it is what remains after integrating a Gaussian likelihood in \(w>0\) against a half-normal mixing measure. The same calculation is the stochastic representation of a (location-zero) skew-normal law; the exam does not require the name.
\end{remark}

\begin{thebibliography}{99}

\bibitem[Durrett(2010)]{durrett2010}
Rick Durrett.
\newblock \emph{Probability: Theory and Examples}.
\newblock 4th~ed., Cambridge University Press, 2010.
\newblock Almost-sure convergence: \S2.3. Weak convergence / Portmanteau: Theorems~3.2.3, 3.2.5. Characteristic functions: Theorems~3.3.1, 3.3.2. Inversion: Theorems~3.3.4, 3.3.5. L\'evy continuity: Theorem~3.3.6. Continuous mapping and the converging-together lemma (Slutsky): \S3.2.

\bibitem[Chung(1968)]{chung1968}
Kai-Lai Chung.
\newblock \emph{A Course in Probability Theory}.
\newblock New York, 1968.
\newblock Characteristic functions and inversion (including atoms and lattice laws): Chapter~6.

\bibitem[Brze\'zniak and Zastawniak(1999)]{brzezniak1999}
Zdzis\l{}aw Brze\'zniak and Tomasz Zastawniak.
\newblock \emph{Basic Stochastic Processes}.
\newblock Springer Undergraduate Mathematics Series, Springer-Verlag, London, 1999.
\newblock Wiener process: Chapter~6, Definition~6.9.

\bibitem[Hogg, McKean and Craig(2019)]{hogg2019}
Robert~V. Hogg, Joseph~W. McKean, and Allen~T. Craig.
\newblock \emph{Introduction to Mathematical Statistics}.
\newblock 8th~ed., Pearson, 2019.
\newblock Decision table / Type~I--II / power: Table~4.5.1, \S4.5, (4.5.4)--(4.5.5). Size: \S8.1, (8.1.3); power function: (8.1.4). Best critical region: Definition~8.1.1. Neyman--Pearson: Theorem~8.1.1. Unbiased test: Definition~8.1.2. UMP of size \(\alpha\): Definition~8.2.1. Monotone likelihood ratio: Definition~8.2.2. Karlin--Rubin (unnumbered claim after Definition~8.2.2): (8.2.3), (8.2.6); exponential-family mlr: after Example~8.2.5. One-sided point null: Example~8.2.2. Two-sided point null, no UMP: Example~8.2.3. Likelihood-ratio / UMPU pointer: \S8.3. One-dimensional Jacobian: Theorem~1.7.1. Multivariate Jacobian: \S2.2, (2.2.3). Quantiles: Definition~1.7.2. Sample quantiles and order statistics: \S4.4. Convergence in probability: \S5.1, Definition~5.1.1; deterministic sequences: Exercise~5.1.1. Convergence in distribution: Definition~5.2.1. In probability implies in distribution: Theorem~5.2.1. Slutsky: Theorem~5.2.5. Empirical cdf: (10.1.1), Remark~10.1.1. Risk and minimax principle: \S7.1. MVUE: Definition~7.1.1. Completeness: \S7.4. Lehmann--Scheff\'e: Theorem~7.4.1. Regular exponential family: Theorems~7.5.1--7.5.2. Functions of a parameter: \S7.6. Prior and posterior: \S11.1.1, (11.1.5). Bayes estimator: \S11.1.2, (11.1.10); \(L_1\) \(\mapsto\) posterior median.
\bibitem[Gelman et al.(2013)]{gelman2013}
Andrew Gelman, John~B. Carlin, Hal~S. Stern, David~B. Dunson, Aki Vehtari, and Donald~B. Rubin.
\newblock \emph{Bayesian Data Analysis}.
\newblock 3rd~ed., CRC Press, 2013.
\newblock Bayes theorem and posterior: \S1.3. Posterior odds: (1.5). Discrete two-point update: \S1.4. Decision analysis / expected utility: Chapter~9.

\bibitem[BDA lecture notes(2026)]{gelmannotes}
Lecture notes for Gelman et al., \emph{Bayesian Data Analysis}.
\newblock Bayes theorem and posterior odds: Chapter~1. Posterior expected loss; \(L_2\mapsto\) mean, \(L_1\mapsto\) median, \(0\)--\(1\mapsto\) mode: Chapter~9.

\bibitem[Yau Stat Review(2025)]{yaustat2025}
\textit{Yau Competition in Mathematics, Statistics Track: Advanced Review}.
\newblock Personal final-prep notes, 25~June 2025, 55~pp.
\newblock Bayesian framework: Definition~5.1. Point estimators: Definition~5.4. Decision theory / \(L_1\) \(\mapsto\) posterior median: Theorem~5.5. Checklist: \S8.1.

\bibitem[Yau Contest(2011)]{yau2011}
S.~T.~Yau High School Mathematics Awards, Probability and Statistics, 2011, Problem~5.
\newblock Two-sample Bernoulli UMVUE of \(p_1-p_2\) and of \(p_1(1-p_1)\).

\bibitem[Yau Contest(2021)]{yau2021}
S.~T.~Yau High School Mathematics Awards, Probability and Statistics, 2021, linear-model item.
\newblock UMVUE of \(\mathbf{x}_1\boldsymbol{\beta}/\sigma\) in a normal linear model, by bias-correcting a plug-in of the complete sufficient pair \((X\hat{\boldsymbol{\beta}},S^2)\).

\bibitem[Qingshu(11th)]{qingshu11}
清疏.
\newblock \emph{大学生数学竞赛班第十一届讲义}.
\newblock Euler--Mascheroni constant: Example~7.5. Next term by Stolz: Example~7.10, equivalently the first Euler--Maclaurin correction. Odd harmonic sum: Example~8.24. Lagrange MVT preserves order: Example~7.17 / Chapter~7. Stirling is Euler--Maclaurin on \(\ln x\).

\bibitem[Ren and Liu(2024)]{ren2024}
任佳刚, 刘继成.
\newblock \emph{概率论教程}.
\newblock Lecture notes, compiled 10~September 2024, 300~pp.
\newblock Density as \(F'\): Definition~4.3.2. Transforms: \S4.7--4.8. Jacobian change of variables: \S4.8, after Theorem~4.8.2. Almost-sure convergence: Definition~5.4.1. In probability: Definition~5.4.2. Almost sure implies in probability: Proposition~5.4.3. Characteristic functions and inversion: Theorem~8.4.4. Computing laws from characteristic functions: \S9.1. In probability implies weak: Proposition~9.2.1. Weak convergence to a constant implies in probability: Proposition~9.2.2.
\end{thebibliography}

\end{document}
